\documentclass{openreader}
\title{Κατὰ Ἰωάννην}
\date{}
\setmainlanguage{english}
\setmainfont{Charis SIL}
\setotherlanguage{greek}
\newfontfamily\greekfont[Script=Greek]{SBL BibLit}
\begin{document}
\ORBselectlanguage{greek}
\maketitle
\raggedbottom 
\fontsize{16pt}{24pt}\selectfont

\ChapterHeader{1}{Chapter 1}

\MainTextVerseMark{1}{1}Ἐν ἀρχῇ ἦν\FnParse{a}{imperfect active indicative 3rd singular} ὁ λόγος, καὶ ὁ λόγος ἦν\FnParse{b}{imperfect active indicative 3rd singular} πρὸς τὸν θεόν, καὶ θεὸς ἦν\FnParse{c}{imperfect active indicative 3rd singular} ὁ λόγος. 
\MainTextVerseMark{1}{2}οὗτος ἦν\FnParse{a}{imperfect active indicative 3rd singular} ἐν ἀρχῇ πρὸς τὸν θεόν. 
\MainTextVerseMark{1}{3}πάντα δι’ αὐτοῦ ἐγένετο,\FnParse{a}{aorist middle indicative 3rd singular} καὶ χωρὶς αὐτοῦ ἐγένετο\FnParse{b}{aorist middle indicative 3rd singular} οὐδὲ ἕν. ὃ γέγονεν\FnParse{c}{perfect active indicative 3rd singular} 
\MainTextVerseMark{1}{4}ἐν αὐτῷ ζωὴ ἦν,\FnParse{a}{imperfect active indicative 3rd singular} καὶ ἡ ζωὴ ἦν\FnParse{b}{imperfect active indicative 3rd singular} τὸ φῶς τῶν ἀνθρώπων· 
\MainTextVerseMark{1}{5}καὶ τὸ φῶς ἐν τῇ σκοτίᾳ\FnParseFormGloss{a}{dative feminine singular}{σκοτία}{darkness} φαίνει,\FnParse{b}{present active indicative 3rd singular} καὶ ἡ σκοτία\FnParseFormGloss{c}{nominative feminine singular}{σκοτία}{darkness} αὐτὸ οὐ κατέλαβεν.\FnParseFormGloss{d}{aorist active indicative 3rd singular}{καταλαμβάνω}{to obtain} 
\MainTextVerseMark{1}{6}Ἐγένετο\FnParse{a}{aorist middle indicative 3rd singular} ἄνθρωπος ἀπεσταλμένος\FnParse{b}{perfect passive participle nominative masculine singular} παρὰ θεοῦ, ὄνομα αὐτῷ Ἰωάννης· 
\MainTextVerseMark{1}{7}οὗτος ἦλθεν\FnParse{a}{aorist active indicative 3rd singular} εἰς μαρτυρίαν, ἵνα μαρτυρήσῃ\FnParse{b}{aorist active subjunctive 3rd singular} περὶ τοῦ φωτός, ἵνα πάντες πιστεύσωσιν\FnParse{c}{aorist active subjunctive 3rd plural} δι’ αὐτοῦ. 
\MainTextVerseMark{1}{8}οὐκ ἦν\FnParse{a}{imperfect active indicative 3rd singular} ἐκεῖνος τὸ φῶς, ἀλλ’ ἵνα μαρτυρήσῃ\FnParse{b}{aorist active subjunctive 3rd singular} περὶ τοῦ φωτός. 
\MainTextVerseMark{1}{9}ἦν\FnParse{a}{imperfect active indicative 3rd singular} τὸ φῶς τὸ ἀληθινὸν\FnParseFormGloss{b}{nominative neuter singular}{ἀληθινός}{true} ὃ φωτίζει\FnParseFormGloss{c}{present active indicative 3rd singular}{φωτίζω}{to give light} πάντα ἄνθρωπον ἐρχόμενον\FnParse{d}{present middle participle accusative masculine singular} εἰς τὸν κόσμον. 
\MainTextVerseMark{1}{10}Ἐν τῷ κόσμῳ ἦν,\FnParse{a}{imperfect active indicative 3rd singular} καὶ ὁ κόσμος δι’ αὐτοῦ ἐγένετο,\FnParse{b}{aorist middle indicative 3rd singular} καὶ ὁ κόσμος αὐτὸν οὐκ ἔγνω.\FnParse{c}{aorist active indicative 3rd singular} 
\MainTextVerseMark{1}{11}εἰς τὰ ἴδια ἦλθεν,\FnParse{a}{aorist active indicative 3rd singular} καὶ οἱ ἴδιοι αὐτὸν οὐ παρέλαβον.\FnParse{b}{aorist active indicative 3rd plural} 
\MainTextVerseMark{1}{12}ὅσοι δὲ ἔλαβον\FnParse{a}{aorist active indicative 3rd plural} αὐτόν, ἔδωκεν\FnParse{b}{aorist active indicative 3rd singular} αὐτοῖς ἐξουσίαν τέκνα θεοῦ γενέσθαι,\FnParse{c}{aorist middle infinitive} τοῖς πιστεύουσιν\FnParse{d}{present active participle dative masculine plural} εἰς τὸ ὄνομα αὐτοῦ, 
\MainTextVerseMark{1}{13}οἳ οὐκ ἐξ αἱμάτων οὐδὲ ἐκ θελήματος σαρκὸς οὐδὲ ἐκ θελήματος ἀνδρὸς ἀλλ’ ἐκ θεοῦ ἐγεννήθησαν.\FnParse{a}{aorist passive indicative 3rd plural} 
\MainTextVerseMark{1}{14}Καὶ ὁ λόγος σὰρξ ἐγένετο\FnParse{a}{aorist middle indicative 3rd singular} καὶ ἐσκήνωσεν\FnParseFormGloss{b}{aorist active indicative 3rd singular}{σκηνόω}{to live} ἐν ἡμῖν, καὶ ἐθεασάμεθα\FnParseFormGloss{c}{aorist middle indicative 1st plural}{θεάομαι}{to see} τὴν δόξαν αὐτοῦ, δόξαν ὡς μονογενοῦς\FnParseFormGloss{d}{genitive masculine singular}{μονογενής}{one and only} παρὰ πατρός, πλήρης\FnParseFormGloss{e}{nominative masculine singular}{πλήρης}{full} χάριτος καὶ ἀληθείας· 
\MainTextVerseMark{1}{15}(Ἰωάννης μαρτυρεῖ\FnParse{a}{present active indicative 3rd singular} περὶ αὐτοῦ καὶ κέκραγεν\FnParse{b}{perfect active indicative 3rd singular} λέγων·\FnParse{c}{present active participle nominative masculine singular} Οὗτος ἦν\FnParse{d}{imperfect active indicative 3rd singular} ⸂ὃν εἶπον⸃·\FnParse{e}{aorist active indicative 1st singular} Ὁ ὀπίσω μου ἐρχόμενος\FnParse{f}{present middle participle nominative masculine singular} ἔμπροσθέν μου γέγονεν,\FnParse{g}{perfect active indicative 3rd singular} ὅτι πρῶτός μου ἦν·)\FnParse{h}{imperfect active indicative 3rd singular} 
\MainTextVerseMark{1}{16}⸀ὅτι ἐκ τοῦ πληρώματος\FnParseFormGloss{a}{genitive neuter singular}{πλήρωμα}{fullness} αὐτοῦ ἡμεῖς πάντες ἐλάβομεν,\FnParse{b}{aorist active indicative 1st plural} καὶ χάριν ἀντὶ\FnFormGloss{c}{ἀντί}{in exchange for} χάριτος· 
\MainTextVerseMark{1}{17}ὅτι ὁ νόμος διὰ Μωϋσέως ἐδόθη,\FnParse{a}{aorist passive indicative 3rd singular} ἡ χάρις καὶ ἡ ἀλήθεια διὰ Ἰησοῦ Χριστοῦ ἐγένετο.\FnParse{b}{aorist middle indicative 3rd singular} 
\MainTextVerseMark{1}{18}θεὸν οὐδεὶς ἑώρακεν\FnParse{a}{perfect active indicative 3rd singular} πώποτε·\FnFormGloss{b}{πώποτε}{ever} ⸂μονογενὴς\FnParseFormGloss{c}{nominative masculine singular}{μονογενής}{one and only} θεὸς⸃ ὁ ὢν\FnParse{d}{present active participle nominative masculine singular} εἰς τὸν κόλπον\FnParseFormGloss{e}{accusative masculine singular}{κόλπος}{side} τοῦ πατρὸς ἐκεῖνος ἐξηγήσατο.\FnParseFormGloss{f}{aorist middle indicative 3rd singular}{ἐξηγέομαι}{to tell} 
\MainTextVerseMark{1}{19}Καὶ αὕτη ἐστὶν\FnParse{a}{present active indicative 3rd singular} ἡ μαρτυρία τοῦ Ἰωάννου ὅτε ⸀ἀπέστειλαν\FnParse{b}{aorist active indicative 3rd plural} οἱ Ἰουδαῖοι ἐξ Ἱεροσολύμων ἱερεῖς καὶ Λευίτας\FnParseFormGloss{c}{accusative masculine plural}{Λευίτης}{Levite} ἵνα ἐρωτήσωσιν\FnParse{d}{aorist active subjunctive 3rd plural} αὐτόν· Σὺ τίς εἶ;\FnParse{e}{present active indicative 2nd singular} 
\MainTextVerseMark{1}{20}καὶ ὡμολόγησεν\FnParseFormGloss{a}{aorist active indicative 3rd singular}{ὁμολογέω}{to confess} καὶ οὐκ ἠρνήσατο,\FnParse{b}{aorist middle indicative 3rd singular} καὶ ὡμολόγησεν\FnParseFormGloss{c}{aorist active indicative 3rd singular}{ὁμολογέω}{to confess} ὅτι ⸂Ἐγὼ οὐκ εἰμὶ⸃\FnParse{d}{present active indicative 1st singular} ὁ χριστός. 
\MainTextVerseMark{1}{21}καὶ ἠρώτησαν\FnParse{a}{aorist active indicative 3rd plural} αὐτόν· Τί οὖν; ⸂σὺ Ἠλίας\FnParseFormGloss{b}{nominative masculine singular}{Ἠλίας}{Elijah} εἶ⸃;\FnParse{c}{present active indicative 2nd singular} καὶ λέγει·\FnParse{d}{present active indicative 3rd singular} Οὐκ εἰμί.\FnParse{e}{present active indicative 1st singular} Ὁ προφήτης εἶ\FnParse{f}{present active indicative 2nd singular} σύ; καὶ ἀπεκρίθη·\FnParse{g}{aorist passive indicative 3rd singular} Οὔ.\FnFormGloss{h}{οὔ}{no} 
\MainTextVerseMark{1}{22}εἶπαν\FnParse{a}{aorist active indicative 3rd plural} οὖν αὐτῷ· Τίς εἶ;\FnParse{b}{present active indicative 2nd singular} ἵνα ἀπόκρισιν\FnParseFormGloss{c}{accusative feminine singular}{ἀπόκρισις}{answer} δῶμεν\FnParse{d}{aorist active subjunctive 1st plural} τοῖς πέμψασιν\FnParse{e}{aorist active participle dative masculine plural} ἡμᾶς· τί λέγεις\FnParse{f}{present active indicative 2nd singular} περὶ σεαυτοῦ; 
\MainTextVerseMark{1}{23}ἔφη·\FnParse{a}{imperfect active indicative 3rd singular} Ἐγὼ φωνὴ βοῶντος\FnParseFormGloss{b}{present active participle genitive masculine singular}{βοάω}{to call} ἐν τῇ ἐρήμῳ· Εὐθύνατε\FnParseFormGloss{c}{aorist active imperative 2nd plural}{εὐθύνω}{to make straight} τὴν ὁδὸν κυρίου, καθὼς εἶπεν\FnParse{d}{aorist active indicative 3rd singular} Ἠσαΐας\FnParseFormGloss{e}{nominative masculine singular}{Ἠσαΐας}{Isaiah} ὁ προφήτης. 
\MainTextVerseMark{1}{24}⸀Καὶ ἀπεσταλμένοι\FnParse{a}{perfect passive participle nominative masculine plural} ἦσαν\FnParse{b}{imperfect active indicative 3rd plural} ἐκ τῶν Φαρισαίων. 
\MainTextVerseMark{1}{25}καὶ ἠρώτησαν\FnParse{a}{aorist active indicative 3rd plural} αὐτὸν καὶ εἶπαν\FnParse{b}{aorist active indicative 3rd plural} αὐτῷ· Τί οὖν βαπτίζεις\FnParse{c}{present active indicative 2nd singular} εἰ σὺ οὐκ εἶ\FnParse{d}{present active indicative 2nd singular} ὁ χριστὸς ⸂οὐδὲ Ἠλίας\FnParseFormGloss{e}{nominative masculine singular}{Ἠλίας}{Elijah} οὐδὲ⸃ ὁ προφήτης; 
\MainTextVerseMark{1}{26}ἀπεκρίθη\FnParse{a}{aorist passive indicative 3rd singular} αὐτοῖς ὁ Ἰωάννης λέγων·\FnParse{b}{present active participle nominative masculine singular} Ἐγὼ βαπτίζω\FnParse{c}{present active indicative 1st singular} ἐν ὕδατι· ⸀μέσος ὑμῶν ⸀ἕστηκεν\FnParse{d}{perfect active indicative 3rd singular} ὃν ὑμεῖς οὐκ οἴδατε,\FnParse{e}{perfect active indicative 2nd plural} 
\MainTextVerseMark{1}{27}⸀ὁ ὀπίσω μου ⸀ἐρχόμενος,\FnParse{a}{present middle participle nominative masculine singular} οὗ ⸂οὐκ εἰμὶ⸃\FnParse{b}{present active indicative 1st singular} ἄξιος ἵνα λύσω\FnParse{c}{aorist active subjunctive 1st singular} αὐτοῦ τὸν ἱμάντα\FnParseFormGloss{d}{accusative masculine singular}{ἱμάς}{thong} τοῦ ὑποδήματος.\FnParseFormGloss{e}{genitive neuter singular}{ὑπόδημα}{sandal} 
\MainTextVerseMark{1}{28}ταῦτα ἐν Βηθανίᾳ\FnParseFormGloss{a}{dative feminine singular}{Βηθανία}{Bethany} ἐγένετο\FnParse{b}{aorist middle indicative 3rd singular} πέραν\FnFormGloss{c}{πέραν}{on the other side} τοῦ Ἰορδάνου,\FnParseFormGloss{d}{genitive masculine singular}{Ἰορδάνης}{Jordan} ὅπου ἦν\FnParse{e}{imperfect active indicative 3rd singular} ⸀ὁ Ἰωάννης βαπτίζων.\FnParse{f}{present active participle nominative masculine singular} 
\MainTextVerseMark{1}{29}Τῇ ἐπαύριον\FnFormGloss{a}{ἐπαύριον}{the next day} βλέπει\FnParse{b}{present active indicative 3rd singular} τὸν Ἰησοῦν ἐρχόμενον\FnParse{c}{present middle participle accusative masculine singular} πρὸς αὐτόν, καὶ λέγει·\FnParse{d}{present active indicative 3rd singular} Ἴδε ὁ ἀμνὸς\FnParseFormGloss{e}{nominative masculine singular}{ἀμνός}{lamb} τοῦ θεοῦ ὁ αἴρων\FnParse{f}{present active participle nominative masculine singular} τὴν ἁμαρτίαν τοῦ κόσμου. 
\MainTextVerseMark{1}{30}οὗτός ἐστιν\FnParse{a}{present active indicative 3rd singular} ⸀ὑπὲρ οὗ ἐγὼ εἶπον·\FnParse{b}{aorist active indicative 1st singular} Ὀπίσω μου ἔρχεται\FnParse{c}{present middle indicative 3rd singular} ἀνὴρ ὃς ἔμπροσθέν μου γέγονεν,\FnParse{d}{perfect active indicative 3rd singular} ὅτι πρῶτός μου ἦν·\FnParse{e}{imperfect active indicative 3rd singular} 
\MainTextVerseMark{1}{31}κἀγὼ οὐκ ᾔδειν\FnParse{a}{pluperfect active indicative 1st singular} αὐτόν, ἀλλ’ ἵνα φανερωθῇ\FnParse{b}{aorist passive subjunctive 3rd singular} τῷ Ἰσραὴλ διὰ τοῦτο ἦλθον\FnParse{c}{aorist active indicative 1st singular} ἐγὼ ⸀ἐν ὕδατι βαπτίζων.\FnParse{d}{present active participle nominative masculine singular} 
\MainTextVerseMark{1}{32}καὶ ἐμαρτύρησεν\FnParse{a}{aorist active indicative 3rd singular} Ἰωάννης λέγων\FnParse{b}{present active participle nominative masculine singular} ὅτι Τεθέαμαι\FnParseFormGloss{c}{perfect middle indicative 1st singular}{θεάομαι}{to see} τὸ πνεῦμα καταβαῖνον\FnParse{d}{present active participle accusative neuter singular} ⸀ὡς περιστερὰν\FnParseFormGloss{e}{accusative feminine singular}{περιστερά}{dove} ἐξ οὐρανοῦ, καὶ ἔμεινεν\FnParse{f}{aorist active indicative 3rd singular} ἐπ’ αὐτόν· 
\MainTextVerseMark{1}{33}κἀγὼ οὐκ ᾔδειν\FnParse{a}{pluperfect active indicative 1st singular} αὐτόν, ἀλλ’ ὁ πέμψας\FnParse{b}{aorist active participle nominative masculine singular} με βαπτίζειν\FnParse{c}{present active infinitive} ἐν ὕδατι ἐκεῖνός μοι εἶπεν·\FnParse{d}{aorist active indicative 3rd singular} Ἐφ’ ὃν ἂν ἴδῃς\FnParse{e}{aorist active subjunctive 2nd singular} τὸ πνεῦμα καταβαῖνον\FnParse{f}{present active participle accusative neuter singular} καὶ μένον\FnParse{g}{present active participle accusative neuter singular} ἐπ’ αὐτόν, οὗτός ἐστιν\FnParse{h}{present active indicative 3rd singular} ὁ βαπτίζων\FnParse{i}{present active participle nominative masculine singular} ἐν πνεύματι ἁγίῳ· 
\MainTextVerseMark{1}{34}κἀγὼ ἑώρακα,\FnParse{a}{perfect active indicative 1st singular} καὶ μεμαρτύρηκα\FnParse{b}{perfect active indicative 1st singular} ὅτι οὗτός ἐστιν\FnParse{c}{present active indicative 3rd singular} ὁ ⸀ἐκλεκτὸς\FnParseFormGloss{d}{nominative masculine singular}{ἐκλεκτός}{elect} τοῦ θεοῦ. 
\MainTextVerseMark{1}{35}Τῇ ἐπαύριον\FnFormGloss{a}{ἐπαύριον}{the next day} πάλιν εἱστήκει\FnParse{b}{pluperfect active indicative 3rd singular} ⸀ὁ Ἰωάννης καὶ ἐκ τῶν μαθητῶν αὐτοῦ δύο, 
\MainTextVerseMark{1}{36}καὶ ἐμβλέψας\FnParseFormGloss{a}{aorist active participle nominative masculine singular}{ἐμβλέπω}{to look} τῷ Ἰησοῦ περιπατοῦντι\FnParse{b}{present active participle dative masculine singular} λέγει·\FnParse{c}{present active indicative 3rd singular} Ἴδε ὁ ἀμνὸς\FnParseFormGloss{d}{nominative masculine singular}{ἀμνός}{lamb} τοῦ θεοῦ. 
\MainTextVerseMark{1}{37}καὶ ἤκουσαν\FnParse{a}{aorist active indicative 3rd plural} ⸂οἱ δύο μαθηταὶ αὐτοῦ⸃ λαλοῦντος\FnParse{b}{present active participle genitive masculine singular} καὶ ἠκολούθησαν\FnParse{c}{aorist active indicative 3rd plural} τῷ Ἰησοῦ. 
\MainTextVerseMark{1}{38}στραφεὶς\FnParseFormGloss{a}{aorist passive participle nominative masculine singular}{στρέφω}{to turn} δὲ ὁ Ἰησοῦς καὶ θεασάμενος\FnParseFormGloss{b}{aorist middle participle nominative masculine singular}{θεάομαι}{to see} αὐτοὺς ἀκολουθοῦντας\FnParse{c}{present active participle accusative masculine plural} λέγει\FnParse{d}{present active indicative 3rd singular} αὐτοῖς· Τί ζητεῖτε;\FnParse{e}{present active indicative 2nd plural} οἱ δὲ εἶπαν\FnParse{f}{aorist active indicative 3rd plural} αὐτῷ· Ῥαββί\FnParseFormGloss{g}{masculine singular}{ῥαββί}{Rabbi} (ὃ λέγεται\FnParse{h}{present passive indicative 3rd singular} ⸀μεθερμηνευόμενον\FnParseFormGloss{i}{present passive participle nominative neuter singular}{μεθερμηνεύω}{to translate} Διδάσκαλε), ποῦ μένεις;\FnParse{j}{present active indicative 2nd singular} 
\MainTextVerseMark{1}{39}λέγει\FnParse{a}{present active indicative 3rd singular} αὐτοῖς· Ἔρχεσθε\FnParse{b}{present middle imperative 2nd plural} καὶ ⸂ὄψεσθε.\FnParse{c}{future middle indicative 2nd plural} ἦλθαν\FnParse{d}{aorist active indicative 3rd plural} οὖν⸃ καὶ εἶδαν\FnParse{e}{aorist active indicative 3rd plural} ποῦ μένει,\FnParse{f}{present active indicative 3rd singular} καὶ παρ’ αὐτῷ ἔμειναν\FnParse{g}{aorist active indicative 3rd plural} τὴν ἡμέραν ἐκείνην· ὥρα ἦν\FnParse{h}{imperfect active indicative 3rd singular} ὡς δεκάτη.\FnParseFormGloss{i}{nominative feminine singular}{δέκατος}{tenth} 
\MainTextVerseMark{1}{40}ἦν\FnParse{a}{imperfect active indicative 3rd singular} Ἀνδρέας\FnParseFormGloss{b}{nominative masculine singular}{Ἀνδρέας}{Andrew} ὁ ἀδελφὸς Σίμωνος Πέτρου εἷς ἐκ τῶν δύο τῶν ἀκουσάντων\FnParse{c}{aorist active participle genitive masculine plural} παρὰ Ἰωάννου καὶ ἀκολουθησάντων\FnParse{d}{aorist active participle genitive masculine plural} αὐτῷ· 
\MainTextVerseMark{1}{41}εὑρίσκει\FnParse{a}{present active indicative 3rd singular} οὗτος ⸀πρῶτον τὸν ἀδελφὸν τὸν ἴδιον Σίμωνα καὶ λέγει\FnParse{b}{present active indicative 3rd singular} αὐτῷ· Εὑρήκαμεν\FnParse{c}{perfect active indicative 1st plural} τὸν Μεσσίαν\FnParseFormGloss{d}{accusative masculine singular}{Μεσσίας}{Messiah} (ὅ ἐστιν\FnParse{e}{present active indicative 3rd singular} μεθερμηνευόμενον\FnParseFormGloss{f}{present passive participle nominative neuter singular}{μεθερμηνεύω}{to translate} χριστός). 
\MainTextVerseMark{1}{42}⸀ἤγαγεν\FnParse{a}{aorist active indicative 3rd singular} αὐτὸν πρὸς τὸν Ἰησοῦν. ἐμβλέψας\FnParseFormGloss{b}{aorist active participle nominative masculine singular}{ἐμβλέπω}{to look} αὐτῷ ὁ Ἰησοῦς εἶπεν·\FnParse{c}{aorist active indicative 3rd singular} Σὺ εἶ\FnParse{d}{present active indicative 2nd singular} Σίμων ὁ υἱὸς ⸀Ἰωάννου, σὺ κληθήσῃ\FnParse{e}{future passive indicative 2nd singular} Κηφᾶς\FnParseFormGloss{f}{nominative masculine singular}{Κηφᾶς}{Cephas} (ὃ ἑρμηνεύεται\FnParseFormGloss{g}{present passive indicative 3rd singular}{ἑρμηνεύω}{to translate} Πέτρος). 
\MainTextVerseMark{1}{43}Τῇ ἐπαύριον\FnFormGloss{a}{ἐπαύριον}{the next day} ἠθέλησεν\FnParse{b}{aorist active indicative 3rd singular} ἐξελθεῖν\FnParse{c}{aorist active infinitive} εἰς τὴν Γαλιλαίαν. καὶ εὑρίσκει\FnParse{d}{present active indicative 3rd singular} Φίλιππον καὶ λέγει\FnParse{e}{present active indicative 3rd singular} αὐτῷ ὁ Ἰησοῦς· Ἀκολούθει\FnParse{f}{present active imperative 2nd singular} μοι. 
\MainTextVerseMark{1}{44}ἦν\FnParse{a}{imperfect active indicative 3rd singular} δὲ ὁ Φίλιππος ἀπὸ Βηθσαϊδά,\FnParseFormGloss{b}{genitive feminine singular}{Βηθσαϊδά}{Bethsaida} ἐκ τῆς πόλεως Ἀνδρέου\FnParseFormGloss{c}{genitive masculine singular}{Ἀνδρέας}{Andrew} καὶ Πέτρου. 
\MainTextVerseMark{1}{45}εὑρίσκει\FnParse{a}{present active indicative 3rd singular} Φίλιππος τὸν Ναθαναὴλ\FnParseFormGloss{b}{accusative masculine singular}{Ναθαναήλ}{Nathanael} καὶ λέγει\FnParse{c}{present active indicative 3rd singular} αὐτῷ· Ὃν ἔγραψεν\FnParse{d}{aorist active indicative 3rd singular} Μωϋσῆς ἐν τῷ νόμῳ καὶ οἱ προφῆται εὑρήκαμεν,\FnParse{e}{perfect active indicative 1st plural} ⸀Ἰησοῦν υἱὸν τοῦ Ἰωσὴφ τὸν ἀπὸ Ναζαρέτ.\FnParseFormGloss{f}{genitive feminine singular}{Ναζαρέτ}{Nazareth} 
\MainTextVerseMark{1}{46}καὶ εἶπεν\FnParse{a}{aorist active indicative 3rd singular} αὐτῷ Ναθαναήλ·\FnParseFormGloss{b}{nominative masculine singular}{Ναθαναήλ}{Nathanael} Ἐκ Ναζαρὲτ\FnParseFormGloss{c}{genitive feminine singular}{Ναζαρέτ}{Nazareth} δύναταί\FnParse{d}{present middle indicative 3rd singular} τι ἀγαθὸν εἶναι;\FnParse{e}{present active infinitive} λέγει\FnParse{f}{present active indicative 3rd singular} αὐτῷ ⸀ὁ Φίλιππος· Ἔρχου\FnParse{g}{present middle imperative 2nd singular} καὶ ἴδε.\FnParse{h}{aorist active imperative 2nd singular} 
\MainTextVerseMark{1}{47}εἶδεν\FnParse{a}{aorist active indicative 3rd singular} ⸀ὁ Ἰησοῦς τὸν Ναθαναὴλ\FnParseFormGloss{b}{accusative masculine singular}{Ναθαναήλ}{Nathanael} ἐρχόμενον\FnParse{c}{present middle participle accusative masculine singular} πρὸς αὐτὸν καὶ λέγει\FnParse{d}{present active indicative 3rd singular} περὶ αὐτοῦ· Ἴδε ἀληθῶς\FnFormGloss{e}{ἀληθῶς}{truly} Ἰσραηλίτης\FnParseFormGloss{f}{nominative masculine singular}{Ἰσραηλίτης}{Israelite} ἐν ᾧ δόλος\FnParseFormGloss{g}{nominative masculine singular}{δόλος}{deceit} οὐκ ἔστιν.\FnParse{h}{present active indicative 3rd singular} 
\MainTextVerseMark{1}{48}λέγει\FnParse{a}{present active indicative 3rd singular} αὐτῷ Ναθαναήλ·\FnParseFormGloss{b}{nominative masculine singular}{Ναθαναήλ}{Nathanael} Πόθεν\FnFormGloss{c}{πόθεν}{from where} με γινώσκεις;\FnParse{d}{present active indicative 2nd singular} ἀπεκρίθη\FnParse{e}{aorist passive indicative 3rd singular} Ἰησοῦς καὶ εἶπεν\FnParse{f}{aorist active indicative 3rd singular} αὐτῷ· Πρὸ τοῦ σε Φίλιππον φωνῆσαι\FnParse{g}{aorist active infinitive} ὄντα\FnParse{h}{present active participle accusative masculine singular} ὑπὸ τὴν συκῆν\FnParseFormGloss{i}{accusative feminine singular}{συκῆ}{fig tree} εἶδόν\FnParse{j}{aorist active indicative 1st singular} σε. 
\MainTextVerseMark{1}{49}ἀπεκρίθη\FnParse{a}{aorist passive indicative 3rd singular} ⸂αὐτῷ Ναθαναήλ⸃·\FnParseFormGloss{b}{nominative masculine singular}{Ναθαναήλ}{Nathanael} Ῥαββί,\FnParseFormGloss{c}{masculine singular}{ῥαββί}{Rabbi} σὺ εἶ\FnParse{d}{present active indicative 2nd singular} ὁ υἱὸς τοῦ θεοῦ, σὺ ⸂βασιλεὺς εἶ⸃\FnParse{e}{present active indicative 2nd singular} τοῦ Ἰσραήλ. 
\MainTextVerseMark{1}{50}ἀπεκρίθη\FnParse{a}{aorist passive indicative 3rd singular} Ἰησοῦς καὶ εἶπεν\FnParse{b}{aorist active indicative 3rd singular} αὐτῷ· Ὅτι εἶπόν\FnParse{c}{aorist active indicative 1st singular} σοι ⸀ὅτι εἶδόν\FnParse{d}{aorist active indicative 1st singular} σε ὑποκάτω\FnFormGloss{e}{ὑποκάτω}{under} τῆς συκῆς\FnParseFormGloss{f}{genitive feminine singular}{συκῆ}{fig tree} πιστεύεις;\FnParse{g}{present active indicative 2nd singular} μείζω τούτων ⸀ὄψῃ.\FnParse{h}{future middle indicative 2nd singular} 
\MainTextVerseMark{1}{51}καὶ λέγει\FnParse{a}{present active indicative 3rd singular} αὐτῷ· Ἀμὴν ἀμὴν λέγω\FnParse{b}{present active indicative 1st singular} ⸀ὑμῖν, ὄψεσθε\FnParse{c}{future middle indicative 2nd plural} τὸν οὐρανὸν ἀνεῳγότα\FnParse{d}{perfect active participle accusative masculine singular} καὶ τοὺς ἀγγέλους τοῦ θεοῦ ἀναβαίνοντας\FnParse{e}{present active participle accusative masculine plural} καὶ καταβαίνοντας\FnParse{f}{present active participle accusative masculine plural} ἐπὶ τὸν υἱὸν τοῦ ἀνθρώπου. 
\ChapterHeader{2}{Chapter 2}

\MainTextVerseMark{2}{1}Καὶ τῇ ⸂ἡμέρᾳ τῇ τρίτῃ⸃ γάμος\FnParseFormGloss{a}{nominative masculine singular}{γάμος}{wedding banquet} ἐγένετο\FnParse{b}{aorist middle indicative 3rd singular} ἐν Κανὰ\FnParseFormGloss{c}{dative feminine singular}{Κανά}{Cana} τῆς Γαλιλαίας, καὶ ἦν\FnParse{d}{imperfect active indicative 3rd singular} ἡ μήτηρ τοῦ Ἰησοῦ ἐκεῖ· 
\MainTextVerseMark{2}{2}ἐκλήθη\FnParse{a}{aorist passive indicative 3rd singular} δὲ καὶ ὁ Ἰησοῦς καὶ οἱ μαθηταὶ αὐτοῦ εἰς τὸν γάμον.\FnParseFormGloss{b}{accusative masculine singular}{γάμος}{wedding banquet} 
\MainTextVerseMark{2}{3}καὶ ὑστερήσαντος\FnParseFormGloss{a}{aorist active participle genitive masculine singular}{ὑστερέω}{to lack} οἴνου λέγει\FnParse{b}{present active indicative 3rd singular} ἡ μήτηρ τοῦ Ἰησοῦ πρὸς αὐτόν· Οἶνον οὐκ ἔχουσιν.\FnParse{c}{present active indicative 3rd plural} 
\MainTextVerseMark{2}{4}⸀καὶ λέγει\FnParse{a}{present active indicative 3rd singular} αὐτῇ ὁ Ἰησοῦς· Τί ἐμοὶ καὶ σοί, γύναι; οὔπω\FnFormGloss{b}{οὔπω}{not yet} ἥκει\FnParseFormGloss{c}{present active indicative 3rd singular}{ἥκω}{to come} ἡ ὥρα μου. 
\MainTextVerseMark{2}{5}λέγει\FnParse{a}{present active indicative 3rd singular} ἡ μήτηρ αὐτοῦ τοῖς διακόνοις·\FnParseFormGloss{b}{dative masculine plural}{διάκονος}{servant} ⸂Ὅ τι⸃ ἂν λέγῃ\FnParse{c}{present active subjunctive 3rd singular} ὑμῖν ποιήσατε.\FnParse{d}{aorist active imperative 2nd plural} 
\MainTextVerseMark{2}{6}ἦσαν\FnParse{a}{imperfect active indicative 3rd plural} δὲ ἐκεῖ ⸂λίθιναι\FnParseFormGloss{b}{nominative feminine plural}{λίθινος}{made of stone} ὑδρίαι⸃\FnParseFormGloss{c}{nominative feminine plural}{ὑδρία}{water jar} ἓξ\FnParseFormGloss{d}{nominative feminine plural}{ἕξ}{six} ⸂κατὰ τὸν καθαρισμὸν\FnParseFormGloss{e}{accusative masculine singular}{καθαρισμός}{cleansing} τῶν Ἰουδαίων κείμεναι⸃,\FnParseFormGloss{f}{present middle participle nominative feminine plural}{κεῖμαι}{to lay} χωροῦσαι\FnParseFormGloss{g}{present active participle nominative feminine plural}{χωρέω}{to go} ἀνὰ\FnFormGloss{h}{ἀνά}{each} μετρητὰς\FnParseFormGloss{i}{accusative masculine plural}{μετρητής}{measure} δύο ἢ τρεῖς. 
\MainTextVerseMark{2}{7}λέγει\FnParse{a}{present active indicative 3rd singular} αὐτοῖς ὁ Ἰησοῦς· Γεμίσατε\FnParseFormGloss{b}{aorist active imperative 2nd plural}{γεμίζω}{to fill} τὰς ὑδρίας\FnParseFormGloss{c}{accusative feminine plural}{ὑδρία}{water jar} ὕδατος· καὶ ἐγέμισαν\FnParseFormGloss{d}{aorist active indicative 3rd plural}{γεμίζω}{to fill} αὐτὰς ἕως ἄνω.\FnFormGloss{e}{ἄνω}{above} 
\MainTextVerseMark{2}{8}καὶ λέγει\FnParse{a}{present active indicative 3rd singular} αὐτοῖς· Ἀντλήσατε\FnParseFormGloss{b}{aorist active imperative 2nd plural}{ἀντλέω}{to draw} νῦν καὶ φέρετε\FnParse{c}{present active imperative 2nd plural} τῷ ἀρχιτρικλίνῳ·\FnParseFormGloss{d}{dative masculine singular}{ἀρχιτρίκλινος}{master of the banquet} ⸂οἱ δὲ⸃ ἤνεγκαν.\FnParse{e}{aorist active indicative 3rd plural} 
\MainTextVerseMark{2}{9}ὡς δὲ ἐγεύσατο\FnParseFormGloss{a}{aorist middle indicative 3rd singular}{γεύομαι}{to taste} ὁ ἀρχιτρίκλινος\FnParseFormGloss{b}{nominative masculine singular}{ἀρχιτρίκλινος}{master of the banquet} τὸ ὕδωρ οἶνον γεγενημένον,\FnParse{c}{perfect passive participle accusative neuter singular} καὶ οὐκ ᾔδει\FnParse{d}{pluperfect active indicative 3rd singular} πόθεν\FnFormGloss{e}{πόθεν}{from where} ἐστίν,\FnParse{f}{present active indicative 3rd singular} οἱ δὲ διάκονοι\FnParseFormGloss{g}{nominative masculine plural}{διάκονος}{servant} ᾔδεισαν\FnParse{h}{pluperfect active indicative 3rd plural} οἱ ἠντληκότες\FnParseFormGloss{i}{perfect active participle nominative masculine plural}{ἀντλέω}{to draw} τὸ ὕδωρ, φωνεῖ\FnParse{j}{present active indicative 3rd singular} τὸν νυμφίον\FnParseFormGloss{k}{accusative masculine singular}{νυμφίος}{bridegroom} ὁ ἀρχιτρίκλινος\FnParseFormGloss{l}{nominative masculine singular}{ἀρχιτρίκλινος}{master of the banquet} 
\MainTextVerseMark{2}{10}καὶ λέγει\FnParse{a}{present active indicative 3rd singular} αὐτῷ· Πᾶς ἄνθρωπος πρῶτον τὸν καλὸν οἶνον τίθησιν,\FnParse{b}{present active indicative 3rd singular} καὶ ὅταν ⸀μεθυσθῶσιν\FnParseFormGloss{c}{aorist passive subjunctive 3rd plural}{μεθύω}{to get drunk} τὸν ἐλάσσω·\FnParseFormGloss{d}{accusative masculine singular}{ἐλάσσων}{lesser} σὺ τετήρηκας\FnParse{e}{perfect active indicative 2nd singular} τὸν καλὸν οἶνον ἕως ἄρτι. 
\MainTextVerseMark{2}{11}ταύτην ⸀ἐποίησεν\FnParse{a}{aorist active indicative 3rd singular} ἀρχὴν τῶν σημείων ὁ Ἰησοῦς ἐν Κανὰ\FnParseFormGloss{b}{dative feminine singular}{Κανά}{Cana} τῆς Γαλιλαίας καὶ ἐφανέρωσεν\FnParse{c}{aorist active indicative 3rd singular} τὴν δόξαν αὐτοῦ, καὶ ἐπίστευσαν\FnParse{d}{aorist active indicative 3rd plural} εἰς αὐτὸν οἱ μαθηταὶ αὐτοῦ. 
\MainTextVerseMark{2}{12}Μετὰ τοῦτο κατέβη\FnParse{a}{aorist active indicative 3rd singular} εἰς Καφαρναοὺμ\FnParseFormGloss{b}{accusative feminine singular}{Καφαρναούμ}{Capernaum} αὐτὸς καὶ ἡ μήτηρ αὐτοῦ καὶ οἱ ⸀ἀδελφοὶ καὶ οἱ μαθηταὶ αὐτοῦ, καὶ ἐκεῖ ἔμειναν\FnParse{c}{aorist active indicative 3rd plural} οὐ πολλὰς ἡμέρας. 
\MainTextVerseMark{2}{13}Καὶ ἐγγὺς ἦν\FnParse{a}{imperfect active indicative 3rd singular} τὸ πάσχα\FnParseFormGloss{b}{nominative neuter singular}{πάσχα}{Passover} τῶν Ἰουδαίων, καὶ ἀνέβη\FnParse{c}{aorist active indicative 3rd singular} εἰς Ἱεροσόλυμα ὁ Ἰησοῦς. 
\MainTextVerseMark{2}{14}καὶ εὗρεν\FnParse{a}{aorist active indicative 3rd singular} ἐν τῷ ἱερῷ τοὺς πωλοῦντας\FnParseFormGloss{b}{present active participle accusative masculine plural}{πωλέω}{to sell} βόας\FnParseFormGloss{c}{accusative masculine plural}{βοῦς}{ox} καὶ πρόβατα καὶ περιστερὰς\FnParseFormGloss{d}{accusative feminine plural}{περιστερά}{dove} καὶ τοὺς κερματιστὰς\FnParseFormGloss{e}{accusative masculine plural}{κερματιστής}{money exchanger} καθημένους,\FnParse{f}{present middle participle accusative masculine plural} 
\MainTextVerseMark{2}{15}καὶ ποιήσας\FnParse{a}{aorist active participle nominative masculine singular} φραγέλλιον\FnParseFormGloss{b}{accusative neuter singular}{φραγέλλιον}{whip} ἐκ σχοινίων\FnParseFormGloss{c}{genitive neuter plural}{σχοινίον}{cords} πάντας ἐξέβαλεν\FnParse{d}{aorist active indicative 3rd singular} ἐκ τοῦ ἱεροῦ τά τε πρόβατα καὶ τοὺς βόας,\FnParseFormGloss{e}{accusative masculine plural}{βοῦς}{ox} καὶ τῶν κολλυβιστῶν\FnParseFormGloss{f}{genitive masculine plural}{κολλυβιστής}{money exchanger} ἐξέχεεν\FnParseFormGloss{g}{aorist active indicative 3rd singular}{ἐκχέω}{to pour out} ⸂τὰ κέρματα⸃\FnParseFormGloss{h}{accusative neuter plural}{κέρμα}{coin} καὶ τὰς τραπέζας\FnParseFormGloss{i}{accusative feminine plural}{τράπεζα}{table} ⸀ἀνέστρεψεν,\FnParseFormGloss{j}{aorist active indicative 3rd singular}{ἀναστρέφω}{to conduct oneself} 
\MainTextVerseMark{2}{16}καὶ τοῖς τὰς περιστερὰς\FnParseFormGloss{a}{accusative feminine plural}{περιστερά}{dove} πωλοῦσιν\FnParseFormGloss{b}{present active participle dative masculine plural}{πωλέω}{to sell} εἶπεν·\FnParse{c}{aorist active indicative 3rd singular} Ἄρατε\FnParse{d}{aorist active imperative 2nd plural} ταῦτα ἐντεῦθεν,\FnFormGloss{e}{ἐντεῦθεν}{from here} μὴ ποιεῖτε\FnParse{f}{present active imperative 2nd plural} τὸν οἶκον τοῦ πατρός μου οἶκον ἐμπορίου.\FnParseFormGloss{g}{genitive neuter singular}{ἐμπόριον}{market} 
\MainTextVerseMark{2}{17}⸀ἐμνήσθησαν\FnParseFormGloss{a}{aorist passive indicative 3rd plural}{μιμνῄσκομαι}{to remember} οἱ μαθηταὶ αὐτοῦ ὅτι γεγραμμένον\FnParse{b}{perfect passive participle nominative neuter singular} ἐστίν·\FnParse{c}{present active indicative 3rd singular} Ὁ ζῆλος\FnParseFormGloss{d}{nominative masculine singular}{ζῆλος}{zeal} τοῦ οἴκου σου καταφάγεταί\FnParseFormGloss{e}{future middle indicative 3rd singular}{κατεσθίω}{to eat up} με. 
\MainTextVerseMark{2}{18}ἀπεκρίθησαν\FnParse{a}{aorist passive indicative 3rd plural} οὖν οἱ Ἰουδαῖοι καὶ εἶπαν\FnParse{b}{aorist active indicative 3rd plural} αὐτῷ· Τί σημεῖον δεικνύεις\FnParse{c}{present active indicative 2nd singular} ἡμῖν, ὅτι ταῦτα ποιεῖς;\FnParse{d}{present active indicative 2nd singular} 
\MainTextVerseMark{2}{19}ἀπεκρίθη\FnParse{a}{aorist passive indicative 3rd singular} Ἰησοῦς καὶ εἶπεν\FnParse{b}{aorist active indicative 3rd singular} αὐτοῖς· Λύσατε\FnParse{c}{aorist active imperative 2nd plural} τὸν ναὸν τοῦτον καὶ ἐν τρισὶν ἡμέραις ἐγερῶ\FnParse{d}{future active indicative 1st singular} αὐτόν. 
\MainTextVerseMark{2}{20}εἶπαν\FnParse{a}{aorist active indicative 3rd plural} οὖν οἱ Ἰουδαῖοι· Τεσσεράκοντα\FnParseFormGloss{b}{dative neuter plural}{τεσσεράκοντα}{forty} καὶ ἓξ\FnParseFormGloss{c}{dative neuter plural}{ἕξ}{six} ἔτεσιν οἰκοδομήθη\FnParse{d}{aorist passive indicative 3rd singular} ὁ ναὸς οὗτος, καὶ σὺ ἐν τρισὶν ἡμέραις ἐγερεῖς\FnParse{e}{future active indicative 2nd singular} αὐτόν; 
\MainTextVerseMark{2}{21}ἐκεῖνος δὲ ἔλεγεν\FnParse{a}{imperfect active indicative 3rd singular} περὶ τοῦ ναοῦ τοῦ σώματος αὐτοῦ. 
\MainTextVerseMark{2}{22}ὅτε οὖν ἠγέρθη\FnParse{a}{aorist passive indicative 3rd singular} ἐκ νεκρῶν, ἐμνήσθησαν\FnParseFormGloss{b}{aorist passive indicative 3rd plural}{μιμνῄσκομαι}{to remember} οἱ μαθηταὶ αὐτοῦ ὅτι τοῦτο ἔλεγεν,\FnParse{c}{imperfect active indicative 3rd singular} καὶ ἐπίστευσαν\FnParse{d}{aorist active indicative 3rd plural} τῇ γραφῇ καὶ τῷ λόγῳ ⸀ὃν εἶπεν\FnParse{e}{aorist active indicative 3rd singular} ὁ Ἰησοῦς. 
\MainTextVerseMark{2}{23}Ὡς δὲ ἦν\FnParse{a}{imperfect active indicative 3rd singular} ἐν τοῖς Ἱεροσολύμοις ἐν τῷ πάσχα\FnParseFormGloss{b}{dative neuter singular}{πάσχα}{Passover} ἐν τῇ ἑορτῇ,\FnParseFormGloss{c}{dative feminine singular}{ἑορτή}{feast} πολλοὶ ἐπίστευσαν\FnParse{d}{aorist active indicative 3rd plural} εἰς τὸ ὄνομα αὐτοῦ, θεωροῦντες\FnParse{e}{present active participle nominative masculine plural} αὐτοῦ τὰ σημεῖα ἃ ἐποίει·\FnParse{f}{imperfect active indicative 3rd singular} 
\MainTextVerseMark{2}{24}αὐτὸς ⸀δὲ Ἰησοῦς οὐκ ἐπίστευεν\FnParse{a}{imperfect active indicative 3rd singular} ⸀αὑτὸν αὐτοῖς διὰ τὸ αὐτὸν γινώσκειν\FnParse{b}{present active infinitive} πάντας 
\MainTextVerseMark{2}{25}καὶ ὅτι οὐ χρείαν εἶχεν\FnParse{a}{imperfect active indicative 3rd singular} ἵνα τις μαρτυρήσῃ\FnParse{b}{aorist active subjunctive 3rd singular} περὶ τοῦ ἀνθρώπου, αὐτὸς γὰρ ἐγίνωσκεν\FnParse{c}{imperfect active indicative 3rd singular} τί ἦν\FnParse{d}{imperfect active indicative 3rd singular} ἐν τῷ ἀνθρώπῳ. 
\ChapterHeader{3}{Chapter 3}

\MainTextVerseMark{3}{1}Ἦν\FnParse{a}{imperfect active indicative 3rd singular} δὲ ἄνθρωπος ἐκ τῶν Φαρισαίων, Νικόδημος\FnParseFormGloss{b}{nominative masculine singular}{Νικόδημος}{Nicodemus} ὄνομα αὐτῷ, ἄρχων τῶν Ἰουδαίων· 
\MainTextVerseMark{3}{2}οὗτος ἦλθεν\FnParse{a}{aorist active indicative 3rd singular} πρὸς αὐτὸν νυκτὸς καὶ εἶπεν\FnParse{b}{aorist active indicative 3rd singular} αὐτῷ· Ῥαββί,\FnParseFormGloss{c}{masculine singular}{ῥαββί}{Rabbi} οἴδαμεν\FnParse{d}{perfect active indicative 1st plural} ὅτι ἀπὸ θεοῦ ἐλήλυθας\FnParse{e}{perfect active indicative 2nd singular} διδάσκαλος· οὐδεὶς γὰρ ⸂δύναται\FnParse{f}{present middle indicative 3rd singular} ταῦτα τὰ σημεῖα⸃ ποιεῖν\FnParse{g}{present active infinitive} ἃ σὺ ποιεῖς,\FnParse{h}{present active indicative 2nd singular} ἐὰν μὴ ᾖ\FnParse{i}{present active subjunctive 3rd singular} ὁ θεὸς μετ’ αὐτοῦ. 
\MainTextVerseMark{3}{3}⸀ἀπεκρίθη\FnParse{a}{aorist passive indicative 3rd singular} Ἰησοῦς καὶ εἶπεν\FnParse{b}{aorist active indicative 3rd singular} αὐτῷ· Ἀμὴν ἀμὴν λέγω\FnParse{c}{present active indicative 1st singular} σοι, ἐὰν μή τις γεννηθῇ\FnParse{d}{aorist passive subjunctive 3rd singular} ἄνωθεν,\FnFormGloss{e}{ἄνωθεν}{from above} οὐ δύναται\FnParse{f}{present middle indicative 3rd singular} ἰδεῖν\FnParse{g}{aorist active infinitive} τὴν βασιλείαν τοῦ θεοῦ. 
\MainTextVerseMark{3}{4}λέγει\FnParse{a}{present active indicative 3rd singular} πρὸς αὐτὸν ⸀ὁ Νικόδημος·\FnParseFormGloss{b}{nominative masculine singular}{Νικόδημος}{Nicodemus} Πῶς δύναται\FnParse{c}{present middle indicative 3rd singular} ἄνθρωπος γεννηθῆναι\FnParse{d}{aorist passive infinitive} γέρων\FnParseFormGloss{e}{nominative masculine singular}{γέρων}{old person} ὤν;\FnParse{f}{present active participle nominative masculine singular} μὴ δύναται\FnParse{g}{present middle indicative 3rd singular} εἰς τὴν κοιλίαν\FnParseFormGloss{h}{accusative feminine singular}{κοιλία}{any and all internal organs} τῆς μητρὸς αὐτοῦ δεύτερον εἰσελθεῖν\FnParse{i}{aorist active infinitive} καὶ γεννηθῆναι;\FnParse{j}{aorist passive infinitive} 
\MainTextVerseMark{3}{5}⸀ἀπεκρίθη\FnParse{a}{aorist passive indicative 3rd singular} Ἰησοῦς· Ἀμὴν ἀμὴν λέγω\FnParse{b}{present active indicative 1st singular} σοι, ἐὰν μή τις γεννηθῇ\FnParse{c}{aorist passive subjunctive 3rd singular} ἐξ ὕδατος καὶ πνεύματος, οὐ δύναται\FnParse{d}{present middle indicative 3rd singular} εἰσελθεῖν\FnParse{e}{aorist active infinitive} εἰς τὴν βασιλείαν τοῦ θεοῦ. 
\MainTextVerseMark{3}{6}τὸ γεγεννημένον\FnParse{a}{perfect passive participle nominative neuter singular} ἐκ τῆς σαρκὸς σάρξ ἐστιν,\FnParse{b}{present active indicative 3rd singular} καὶ τὸ γεγεννημένον\FnParse{c}{perfect passive participle nominative neuter singular} ἐκ τοῦ πνεύματος πνεῦμά ἐστιν.\FnParse{d}{present active indicative 3rd singular} 
\MainTextVerseMark{3}{7}μὴ θαυμάσῃς\FnParse{a}{aorist active subjunctive 2nd singular} ὅτι εἶπόν\FnParse{b}{aorist active indicative 1st singular} σοι Δεῖ\FnParse{c}{present active indicative 3rd singular} ὑμᾶς γεννηθῆναι\FnParse{d}{aorist passive infinitive} ἄνωθεν.\FnFormGloss{e}{ἄνωθεν}{from above} 
\MainTextVerseMark{3}{8}τὸ πνεῦμα ὅπου θέλει\FnParse{a}{present active indicative 3rd singular} πνεῖ,\FnParseFormGloss{b}{present active indicative 3rd singular}{πνέω}{to blow} καὶ τὴν φωνὴν αὐτοῦ ἀκούεις,\FnParse{c}{present active indicative 2nd singular} ἀλλ’ οὐκ οἶδας\FnParse{d}{perfect active indicative 2nd singular} πόθεν\FnFormGloss{e}{πόθεν}{from where} ἔρχεται\FnParse{f}{present middle indicative 3rd singular} καὶ ποῦ ὑπάγει·\FnParse{g}{present active indicative 3rd singular} οὕτως ἐστὶν\FnParse{h}{present active indicative 3rd singular} πᾶς ὁ γεγεννημένος\FnParse{i}{perfect passive participle nominative masculine singular} ἐκ τοῦ πνεύματος. 
\MainTextVerseMark{3}{9}ἀπεκρίθη\FnParse{a}{aorist passive indicative 3rd singular} Νικόδημος\FnParseFormGloss{b}{nominative masculine singular}{Νικόδημος}{Nicodemus} καὶ εἶπεν\FnParse{c}{aorist active indicative 3rd singular} αὐτῷ· Πῶς δύναται\FnParse{d}{present middle indicative 3rd singular} ταῦτα γενέσθαι;\FnParse{e}{aorist middle infinitive} 
\MainTextVerseMark{3}{10}ἀπεκρίθη\FnParse{a}{aorist passive indicative 3rd singular} Ἰησοῦς καὶ εἶπεν\FnParse{b}{aorist active indicative 3rd singular} αὐτῷ· Σὺ εἶ\FnParse{c}{present active indicative 2nd singular} ὁ διδάσκαλος τοῦ Ἰσραὴλ καὶ ταῦτα οὐ γινώσκεις;\FnParse{d}{present active indicative 2nd singular} 
\MainTextVerseMark{3}{11}ἀμὴν ἀμὴν λέγω\FnParse{a}{present active indicative 1st singular} σοι ὅτι ὃ οἴδαμεν\FnParse{b}{perfect active indicative 1st plural} λαλοῦμεν\FnParse{c}{present active indicative 1st plural} καὶ ὃ ἑωράκαμεν\FnParse{d}{perfect active indicative 1st plural} μαρτυροῦμεν,\FnParse{e}{present active indicative 1st plural} καὶ τὴν μαρτυρίαν ἡμῶν οὐ λαμβάνετε.\FnParse{f}{present active indicative 2nd plural} 
\MainTextVerseMark{3}{12}εἰ τὰ ἐπίγεια\FnParseFormGloss{a}{accusative neuter plural}{ἐπίγειος}{being on the earth} εἶπον\FnParse{b}{aorist active indicative 1st singular} ὑμῖν καὶ οὐ πιστεύετε,\FnParse{c}{present active indicative 2nd plural} πῶς ἐὰν εἴπω\FnParse{d}{aorist active subjunctive 1st singular} ὑμῖν τὰ ἐπουράνια\FnParseFormGloss{e}{accusative neuter plural}{ἐπουράνιος}{heavenly} πιστεύσετε;\FnParse{f}{future active indicative 2nd plural} 
\MainTextVerseMark{3}{13}καὶ οὐδεὶς ἀναβέβηκεν\FnParse{a}{perfect active indicative 3rd singular} εἰς τὸν οὐρανὸν εἰ μὴ ὁ ἐκ τοῦ οὐρανοῦ καταβάς,\FnParse{b}{aorist active participle nominative masculine singular} ὁ υἱὸς τοῦ ⸀ἀνθρώπου. 
\MainTextVerseMark{3}{14}καὶ καθὼς Μωϋσῆς ὕψωσεν\FnParseFormGloss{a}{aorist active indicative 3rd singular}{ὑψόω}{to lift up} τὸν ὄφιν\FnParseFormGloss{b}{accusative masculine singular}{ὄφις}{snake} ἐν τῇ ἐρήμῳ, οὕτως ὑψωθῆναι\FnParseFormGloss{c}{aorist passive infinitive}{ὑψόω}{to lift up} δεῖ\FnParse{d}{present active indicative 3rd singular} τὸν υἱὸν τοῦ ἀνθρώπου, 
\MainTextVerseMark{3}{15}ἵνα πᾶς ὁ πιστεύων\FnParse{a}{present active participle nominative masculine singular} ⸂ἐν αὐτῷ⸃ ἔχῃ\FnParse{b}{present active subjunctive 3rd singular} ζωὴν αἰώνιον. 
\MainTextVerseMark{3}{16}Οὕτως γὰρ ἠγάπησεν\FnParse{a}{aorist active indicative 3rd singular} ὁ θεὸς τὸν κόσμον ὥστε τὸν ⸀υἱὸν τὸν μονογενῆ\FnParseFormGloss{b}{accusative masculine singular}{μονογενής}{one and only} ἔδωκεν,\FnParse{c}{aorist active indicative 3rd singular} ἵνα πᾶς ὁ πιστεύων\FnParse{d}{present active participle nominative masculine singular} εἰς αὐτὸν μὴ ἀπόληται\FnParse{e}{aorist middle subjunctive 3rd singular} ἀλλὰ ἔχῃ\FnParse{f}{present active subjunctive 3rd singular} ζωὴν αἰώνιον. 
\MainTextVerseMark{3}{17}οὐ γὰρ ἀπέστειλεν\FnParse{a}{aorist active indicative 3rd singular} ὁ θεὸς τὸν ⸀υἱὸν εἰς τὸν κόσμον ἵνα κρίνῃ\FnParse{b}{aorist active subjunctive 3rd singular} τὸν κόσμον, ἀλλ’ ἵνα σωθῇ\FnParse{c}{aorist passive subjunctive 3rd singular} ὁ κόσμος δι’ αὐτοῦ. 
\MainTextVerseMark{3}{18}ὁ πιστεύων\FnParse{a}{present active participle nominative masculine singular} εἰς αὐτὸν οὐ κρίνεται·\FnParse{b}{present passive indicative 3rd singular} ὁ ⸀δὲ μὴ πιστεύων\FnParse{c}{present active participle nominative masculine singular} ἤδη κέκριται,\FnParse{d}{perfect passive indicative 3rd singular} ὅτι μὴ πεπίστευκεν\FnParse{e}{perfect active indicative 3rd singular} εἰς τὸ ὄνομα τοῦ μονογενοῦς\FnParseFormGloss{f}{genitive masculine singular}{μονογενής}{one and only} υἱοῦ τοῦ θεοῦ. 
\MainTextVerseMark{3}{19}αὕτη δέ ἐστιν\FnParse{a}{present active indicative 3rd singular} ἡ κρίσις ὅτι τὸ φῶς ἐλήλυθεν\FnParse{b}{perfect active indicative 3rd singular} εἰς τὸν κόσμον καὶ ἠγάπησαν\FnParse{c}{aorist active indicative 3rd plural} οἱ ἄνθρωποι μᾶλλον τὸ σκότος ἢ τὸ φῶς, ἦν\FnParse{d}{imperfect active indicative 3rd singular} γὰρ ⸂αὐτῶν πονηρὰ⸃ τὰ ἔργα. 
\MainTextVerseMark{3}{20}πᾶς γὰρ ὁ φαῦλα\FnParseFormGloss{a}{accusative neuter plural}{φαῦλος}{evil} πράσσων\FnParse{b}{present active participle nominative masculine singular} μισεῖ\FnParse{c}{present active indicative 3rd singular} τὸ φῶς καὶ οὐκ ἔρχεται\FnParse{d}{present middle indicative 3rd singular} πρὸς τὸ φῶς, ἵνα μὴ ἐλεγχθῇ\FnParseFormGloss{e}{aorist passive subjunctive 3rd singular}{ἐλέγχω}{to expose} τὰ ἔργα αὐτοῦ· 
\MainTextVerseMark{3}{21}ὁ δὲ ποιῶν\FnParse{a}{present active participle nominative masculine singular} τὴν ἀλήθειαν ἔρχεται\FnParse{b}{present middle indicative 3rd singular} πρὸς τὸ φῶς, ἵνα φανερωθῇ\FnParse{c}{aorist passive subjunctive 3rd singular} αὐτοῦ τὰ ἔργα ὅτι ἐν θεῷ ἐστιν\FnParse{d}{present active indicative 3rd singular} εἰργασμένα.\FnParse{e}{perfect passive participle nominative neuter plural} 
\MainTextVerseMark{3}{22}Μετὰ ταῦτα ἦλθεν\FnParse{a}{aorist active indicative 3rd singular} ὁ Ἰησοῦς καὶ οἱ μαθηταὶ αὐτοῦ εἰς τὴν Ἰουδαίαν γῆν, καὶ ἐκεῖ διέτριβεν\FnParseFormGloss{b}{imperfect active indicative 3rd singular}{διατρίβω}{to stay} μετ’ αὐτῶν καὶ ἐβάπτιζεν.\FnParse{c}{imperfect active indicative 3rd singular} 
\MainTextVerseMark{3}{23}ἦν\FnParse{a}{imperfect active indicative 3rd singular} δὲ καὶ ⸀ὁ Ἰωάννης βαπτίζων\FnParse{b}{present active participle nominative masculine singular} ἐν Αἰνὼν\FnParseFormGloss{c}{dative feminine singular}{Αἰνών}{Aenon} ἐγγὺς τοῦ Σαλείμ,\FnParseFormGloss{d}{genitive neuter singular}{Σαλείμ}{Salamis} ὅτι ὕδατα πολλὰ ἦν\FnParse{e}{imperfect active indicative 3rd singular} ἐκεῖ, καὶ παρεγίνοντο\FnParse{f}{imperfect middle indicative 3rd plural} καὶ ἐβαπτίζοντο·\FnParse{g}{imperfect passive indicative 3rd plural} 
\MainTextVerseMark{3}{24}οὔπω\FnFormGloss{a}{οὔπω}{not yet} γὰρ ἦν\FnParse{b}{imperfect active indicative 3rd singular} βεβλημένος\FnParse{c}{perfect passive participle nominative masculine singular} εἰς τὴν φυλακὴν ⸀ὁ Ἰωάννης. 
\MainTextVerseMark{3}{25}Ἐγένετο\FnParse{a}{aorist middle indicative 3rd singular} οὖν ζήτησις\FnParseFormGloss{b}{nominative feminine singular}{ζήτησις}{argument} ἐκ τῶν μαθητῶν Ἰωάννου μετὰ Ἰουδαίου περὶ καθαρισμοῦ.\FnParseFormGloss{c}{genitive masculine singular}{καθαρισμός}{cleansing} 
\MainTextVerseMark{3}{26}καὶ ἦλθον\FnParse{a}{aorist active indicative 3rd plural} πρὸς τὸν Ἰωάννην καὶ εἶπαν\FnParse{b}{aorist active indicative 3rd plural} αὐτῷ· Ῥαββί,\FnParseFormGloss{c}{masculine singular}{ῥαββί}{Rabbi} ὃς ἦν\FnParse{d}{imperfect active indicative 3rd singular} μετὰ σοῦ πέραν\FnFormGloss{e}{πέραν}{on the other side} τοῦ Ἰορδάνου,\FnParseFormGloss{f}{genitive masculine singular}{Ἰορδάνης}{Jordan} ᾧ σὺ μεμαρτύρηκας,\FnParse{g}{perfect active indicative 2nd singular} ἴδε οὗτος βαπτίζει\FnParse{h}{present active indicative 3rd singular} καὶ πάντες ἔρχονται\FnParse{i}{present middle indicative 3rd plural} πρὸς αὐτόν. 
\MainTextVerseMark{3}{27}ἀπεκρίθη\FnParse{a}{aorist passive indicative 3rd singular} Ἰωάννης καὶ εἶπεν·\FnParse{b}{aorist active indicative 3rd singular} Οὐ δύναται\FnParse{c}{present middle indicative 3rd singular} ἄνθρωπος λαμβάνειν\FnParse{d}{present active infinitive} ⸂οὐδὲ ἓν⸃ ἐὰν μὴ ᾖ\FnParse{e}{present active subjunctive 3rd singular} δεδομένον\FnParse{f}{perfect passive participle nominative neuter singular} αὐτῷ ἐκ τοῦ οὐρανοῦ. 
\MainTextVerseMark{3}{28}αὐτοὶ ὑμεῖς ⸀μοι μαρτυρεῖτε\FnParse{a}{present active indicative 2nd plural} ὅτι ⸀εἶπον·\FnParse{b}{aorist active indicative 1st singular} Οὐκ εἰμὶ\FnParse{c}{present active indicative 1st singular} ἐγὼ ὁ χριστός, ἀλλ’ ὅτι Ἀπεσταλμένος\FnParse{d}{perfect passive participle nominative masculine singular} εἰμὶ\FnParse{e}{present active indicative 1st singular} ἔμπροσθεν ἐκείνου. 
\MainTextVerseMark{3}{29}ὁ ἔχων\FnParse{a}{present active participle nominative masculine singular} τὴν νύμφην\FnParseFormGloss{b}{accusative feminine singular}{νύμφη}{bride} νυμφίος\FnParseFormGloss{c}{nominative masculine singular}{νυμφίος}{bridegroom} ἐστίν·\FnParse{d}{present active indicative 3rd singular} ὁ δὲ φίλος\FnParseFormGloss{e}{nominative masculine singular}{φίλος}{friendly} τοῦ νυμφίου\FnParseFormGloss{f}{genitive masculine singular}{νυμφίος}{bridegroom} ὁ ἑστηκὼς\FnParse{g}{perfect active participle nominative masculine singular} καὶ ἀκούων\FnParse{h}{present active participle nominative masculine singular} αὐτοῦ, χαρᾷ χαίρει\FnParse{i}{present active indicative 3rd singular} διὰ τὴν φωνὴν τοῦ νυμφίου.\FnParseFormGloss{j}{genitive masculine singular}{νυμφίος}{bridegroom} αὕτη οὖν ἡ χαρὰ ἡ ἐμὴ πεπλήρωται.\FnParse{k}{perfect passive indicative 3rd singular} 
\MainTextVerseMark{3}{30}ἐκεῖνον δεῖ\FnParse{a}{present active indicative 3rd singular} αὐξάνειν,\FnParseFormGloss{b}{present active infinitive}{αὐξάνω}{to cause to grow} ἐμὲ δὲ ἐλαττοῦσθαι.\FnParseFormGloss{c}{present middle infinitive}{ἐλαττόω}{to make lower than} 
\MainTextVerseMark{3}{31}Ὁ ἄνωθεν\FnFormGloss{a}{ἄνωθεν}{from above} ἐρχόμενος\FnParse{b}{present middle participle nominative masculine singular} ἐπάνω\FnFormGloss{c}{ἐπάνω}{above} πάντων ἐστίν.\FnParse{d}{present active indicative 3rd singular} ὁ ὢν\FnParse{e}{present active participle nominative masculine singular} ἐκ τῆς γῆς ἐκ τῆς γῆς ἐστιν\FnParse{f}{present active indicative 3rd singular} καὶ ἐκ τῆς γῆς λαλεῖ·\FnParse{g}{present active indicative 3rd singular} ὁ ἐκ τοῦ οὐρανοῦ ἐρχόμενος\FnParse{h}{present middle participle nominative masculine singular} ἐπάνω\FnFormGloss{i}{ἐπάνω}{above} πάντων ἐστίν·\FnParse{j}{present active indicative 3rd singular} 
\MainTextVerseMark{3}{32}⸀ὃ ἑώρακεν\FnParse{a}{perfect active indicative 3rd singular} καὶ ἤκουσεν\FnParse{b}{aorist active indicative 3rd singular} τοῦτο μαρτυρεῖ,\FnParse{c}{present active indicative 3rd singular} καὶ τὴν μαρτυρίαν αὐτοῦ οὐδεὶς λαμβάνει.\FnParse{d}{present active indicative 3rd singular} 
\MainTextVerseMark{3}{33}ὁ λαβὼν\FnParse{a}{aorist active participle nominative masculine singular} αὐτοῦ τὴν μαρτυρίαν ἐσφράγισεν\FnParseFormGloss{b}{aorist active indicative 3rd singular}{σφραγίζω}{to seal} ὅτι ὁ θεὸς ἀληθής\FnParseFormGloss{c}{nominative masculine singular}{ἀληθής}{true} ἐστιν.\FnParse{d}{present active indicative 3rd singular} 
\MainTextVerseMark{3}{34}ὃν γὰρ ἀπέστειλεν\FnParse{a}{aorist active indicative 3rd singular} ὁ θεὸς τὰ ῥήματα τοῦ θεοῦ λαλεῖ,\FnParse{b}{present active indicative 3rd singular} οὐ γὰρ ἐκ μέτρου\FnParseFormGloss{c}{genitive neuter singular}{μέτρον}{measure} ⸀δίδωσιν\FnParse{d}{present active indicative 3rd singular} τὸ πνεῦμα. 
\MainTextVerseMark{3}{35}ὁ πατὴρ ἀγαπᾷ\FnParse{a}{present active indicative 3rd singular} τὸν υἱόν, καὶ πάντα δέδωκεν\FnParse{b}{perfect active indicative 3rd singular} ἐν τῇ χειρὶ αὐτοῦ. 
\MainTextVerseMark{3}{36}ὁ πιστεύων\FnParse{a}{present active participle nominative masculine singular} εἰς τὸν υἱὸν ἔχει\FnParse{b}{present active indicative 3rd singular} ζωὴν αἰώνιον· ὁ δὲ ἀπειθῶν\FnParseFormGloss{c}{present active participle nominative masculine singular}{ἀπειθέω}{to disobey} τῷ υἱῷ οὐκ ὄψεται\FnParse{d}{future middle indicative 3rd singular} ζωήν, ἀλλ’ ἡ ὀργὴ τοῦ θεοῦ μένει\FnParse{e}{present active indicative 3rd singular} ἐπ’ αὐτόν. 
\ChapterHeader{4}{Chapter 4}

\MainTextVerseMark{4}{1}Ὡς οὖν ἔγνω\FnParse{a}{aorist active indicative 3rd singular} ὁ ⸀Ἰησοῦς ὅτι ἤκουσαν\FnParse{b}{aorist active indicative 3rd plural} οἱ Φαρισαῖοι ὅτι Ἰησοῦς πλείονας μαθητὰς ποιεῖ\FnParse{c}{present active indicative 3rd singular} καὶ βαπτίζει\FnParse{d}{present active indicative 3rd singular} ἢ Ἰωάννης— 
\MainTextVerseMark{4}{2}καίτοιγε\FnFormGloss{a}{καίτοιγε}{although} Ἰησοῦς αὐτὸς οὐκ ἐβάπτιζεν\FnParse{b}{imperfect active indicative 3rd singular} ἀλλ’ οἱ μαθηταὶ αὐτοῦ— 
\MainTextVerseMark{4}{3}ἀφῆκεν\FnParse{a}{aorist active indicative 3rd singular} τὴν Ἰουδαίαν καὶ ἀπῆλθεν\FnParse{b}{aorist active indicative 3rd singular} ⸀πάλιν εἰς τὴν Γαλιλαίαν. 
\MainTextVerseMark{4}{4}ἔδει\FnParse{a}{imperfect active indicative 3rd singular} δὲ αὐτὸν διέρχεσθαι\FnParse{b}{present middle infinitive} διὰ τῆς Σαμαρείας.\FnParseFormGloss{c}{genitive feminine singular}{Σαμάρεια}{Samaria} 
\MainTextVerseMark{4}{5}ἔρχεται\FnParse{a}{present middle indicative 3rd singular} οὖν εἰς πόλιν τῆς Σαμαρείας\FnParseFormGloss{b}{genitive feminine singular}{Σαμάρεια}{Samaria} λεγομένην\FnParse{c}{present passive participle accusative feminine singular} Συχὰρ\FnParseFormGloss{d}{accusative feminine singular}{Συχάρ}{Sychar} πλησίον\FnFormGloss{e}{πλησίον}{near} τοῦ χωρίου\FnParseFormGloss{f}{genitive neuter singular}{χωρίον}{place} ὃ ἔδωκεν\FnParse{g}{aorist active indicative 3rd singular} Ἰακὼβ\FnParseFormGloss{h}{nominative masculine singular}{Ἰακώβ}{Jacob} ⸀τῷ Ἰωσὴφ τῷ υἱῷ αὐτοῦ· 
\MainTextVerseMark{4}{6}ἦν\FnParse{a}{imperfect active indicative 3rd singular} δὲ ἐκεῖ πηγὴ\FnParseFormGloss{b}{nominative feminine singular}{πηγή}{spring} τοῦ Ἰακώβ.\FnParseFormGloss{c}{genitive masculine singular}{Ἰακώβ}{Jacob} ὁ οὖν Ἰησοῦς κεκοπιακὼς\FnParseFormGloss{d}{perfect active participle nominative masculine singular}{κοπιάω}{to work} ἐκ τῆς ὁδοιπορίας\FnParseFormGloss{e}{genitive feminine singular}{ὁδοιπορία}{journey} ἐκαθέζετο\FnParseFormGloss{f}{imperfect middle indicative 3rd singular}{καθέζομαι}{to sit down} οὕτως ἐπὶ τῇ πηγῇ·\FnParseFormGloss{g}{dative feminine singular}{πηγή}{spring} ὥρα ἦν\FnParse{h}{imperfect active indicative 3rd singular} ⸀ὡς ἕκτη.\FnParseFormGloss{i}{nominative feminine singular}{ἕκτος}{sixth} 
\MainTextVerseMark{4}{7}Ἔρχεται\FnParse{a}{present middle indicative 3rd singular} γυνὴ ἐκ τῆς Σαμαρείας\FnParseFormGloss{b}{genitive feminine singular}{Σαμάρεια}{Samaria} ἀντλῆσαι\FnParseFormGloss{c}{aorist active infinitive}{ἀντλέω}{to draw} ὕδωρ. λέγει\FnParse{d}{present active indicative 3rd singular} αὐτῇ ὁ Ἰησοῦς· Δός\FnParse{e}{aorist active imperative 2nd singular} μοι πεῖν·\FnParse{f}{aorist active infinitive} 
\MainTextVerseMark{4}{8}οἱ γὰρ μαθηταὶ αὐτοῦ ἀπεληλύθεισαν\FnParse{a}{pluperfect active indicative 3rd plural} εἰς τὴν πόλιν, ἵνα τροφὰς\FnParseFormGloss{b}{accusative feminine plural}{τροφή}{food} ἀγοράσωσιν.\FnParse{c}{aorist active subjunctive 3rd plural} 
\MainTextVerseMark{4}{9}λέγει\FnParse{a}{present active indicative 3rd singular} οὖν αὐτῷ ἡ γυνὴ ἡ Σαμαρῖτις·\FnParseFormGloss{b}{nominative feminine singular}{Σαμαρῖτις}{Samaritan} Πῶς σὺ Ἰουδαῖος ὢν\FnParse{c}{present active participle nominative masculine singular} παρ’ ἐμοῦ πεῖν\FnParse{d}{aorist active infinitive} αἰτεῖς\FnParse{e}{present active indicative 2nd singular} ⸂γυναικὸς Σαμαρίτιδος\FnParseFormGloss{f}{genitive feminine singular}{Σαμαρῖτις}{Samaritan} οὔσης⸃;\FnParse{g}{present active participle genitive feminine singular} οὐ γὰρ συγχρῶνται\FnParseFormGloss{h}{present middle indicative 3rd plural}{συγχράομαι}{to associate with} Ἰουδαῖοι Σαμαρίταις.\FnParseFormGloss{i}{dative masculine plural}{Σαμαρίτης}{Samaritan} 
\MainTextVerseMark{4}{10}ἀπεκρίθη\FnParse{a}{aorist passive indicative 3rd singular} Ἰησοῦς καὶ εἶπεν\FnParse{b}{aorist active indicative 3rd singular} αὐτῇ· Εἰ ᾔδεις\FnParse{c}{pluperfect active indicative 2nd singular} τὴν δωρεὰν\FnParseFormGloss{d}{accusative feminine singular}{δωρεά}{gift} τοῦ θεοῦ καὶ τίς ἐστιν\FnParse{e}{present active indicative 3rd singular} ὁ λέγων\FnParse{f}{present active participle nominative masculine singular} σοι· Δός\FnParse{g}{aorist active imperative 2nd singular} μοι πεῖν,\FnParse{h}{aorist active infinitive} σὺ ἂν ᾔτησας\FnParse{i}{aorist active indicative 2nd singular} αὐτὸν καὶ ἔδωκεν\FnParse{j}{aorist active indicative 3rd singular} ἄν σοι ὕδωρ ζῶν.\FnParse{k}{present active participle accusative neuter singular} 
\MainTextVerseMark{4}{11}λέγει\FnParse{a}{present active indicative 3rd singular} αὐτῷ ⸂ἡ γυνή⸃· Κύριε, οὔτε ἄντλημα\FnParseFormGloss{b}{accusative neuter singular}{ἄντλημα}{container to draw water with} ἔχεις\FnParse{c}{present active indicative 2nd singular} καὶ τὸ φρέαρ\FnParseFormGloss{d}{nominative neuter singular}{φρέαρ}{well} ἐστὶν\FnParse{e}{present active indicative 3rd singular} βαθύ·\FnParseFormGloss{f}{nominative neuter singular}{βαθύς}{deep} πόθεν\FnFormGloss{g}{πόθεν}{from where} οὖν ἔχεις\FnParse{h}{present active indicative 2nd singular} τὸ ὕδωρ τὸ ζῶν;\FnParse{i}{present active participle accusative neuter singular} 
\MainTextVerseMark{4}{12}μὴ σὺ μείζων εἶ\FnParse{a}{present active indicative 2nd singular} τοῦ πατρὸς ἡμῶν Ἰακώβ,\FnParseFormGloss{b}{genitive masculine singular}{Ἰακώβ}{Jacob} ὃς ἔδωκεν\FnParse{c}{aorist active indicative 3rd singular} ἡμῖν τὸ φρέαρ\FnParseFormGloss{d}{accusative neuter singular}{φρέαρ}{well} καὶ αὐτὸς ἐξ αὐτοῦ ἔπιεν\FnParse{e}{aorist active indicative 3rd singular} καὶ οἱ υἱοὶ αὐτοῦ καὶ τὰ θρέμματα\FnParseFormGloss{f}{nominative neuter plural}{θρέμμα}{livestock} αὐτοῦ; 
\MainTextVerseMark{4}{13}ἀπεκρίθη\FnParse{a}{aorist passive indicative 3rd singular} Ἰησοῦς καὶ εἶπεν\FnParse{b}{aorist active indicative 3rd singular} αὐτῇ· Πᾶς ὁ πίνων\FnParse{c}{present active participle nominative masculine singular} ἐκ τοῦ ὕδατος τούτου διψήσει\FnParseFormGloss{d}{future active indicative 3rd singular}{διψάω}{to be thirsty} πάλιν· 
\MainTextVerseMark{4}{14}ὃς δ’ ἂν πίῃ\FnParse{a}{aorist active subjunctive 3rd singular} ἐκ τοῦ ὕδατος οὗ ἐγὼ δώσω\FnParse{b}{future active indicative 1st singular} αὐτῷ, οὐ μὴ ⸀διψήσει\FnParseFormGloss{c}{future active indicative 3rd singular}{διψάω}{to be thirsty} εἰς τὸν αἰῶνα, ἀλλὰ τὸ ὕδωρ ὃ δώσω\FnParse{d}{future active indicative 1st singular} αὐτῷ γενήσεται\FnParse{e}{future middle indicative 3rd singular} ἐν αὐτῷ πηγὴ\FnParseFormGloss{f}{nominative feminine singular}{πηγή}{spring} ὕδατος ἁλλομένου\FnParseFormGloss{g}{present middle participle genitive neuter singular}{ἅλλομαι}{to jump up} εἰς ζωὴν αἰώνιον. 
\MainTextVerseMark{4}{15}λέγει\FnParse{a}{present active indicative 3rd singular} πρὸς αὐτὸν ἡ γυνή· Κύριε, δός\FnParse{b}{aorist active imperative 2nd singular} μοι τοῦτο τὸ ὕδωρ, ἵνα μὴ διψῶ\FnParseFormGloss{c}{present active subjunctive 1st singular}{διψάω}{to be thirsty} μηδὲ ⸀διέρχωμαι\FnParse{d}{present middle subjunctive 1st singular} ἐνθάδε\FnFormGloss{e}{ἐνθάδε}{here} ἀντλεῖν.\FnParseFormGloss{f}{present active infinitive}{ἀντλέω}{to draw} 
\MainTextVerseMark{4}{16}Λέγει\FnParse{a}{present active indicative 3rd singular} ⸀αὐτῇ· Ὕπαγε\FnParse{b}{present active imperative 2nd singular} φώνησον\FnParse{c}{aorist active imperative 2nd singular} ⸂τὸν ἄνδρα σου⸃ καὶ ἐλθὲ\FnParse{d}{aorist active imperative 2nd singular} ἐνθάδε.\FnFormGloss{e}{ἐνθάδε}{here} 
\MainTextVerseMark{4}{17}ἀπεκρίθη\FnParse{a}{aorist passive indicative 3rd singular} ἡ γυνὴ καὶ εἶπεν\FnParse{b}{aorist active indicative 3rd singular} ⸀αὐτῷ· Οὐκ ἔχω\FnParse{c}{present active indicative 1st singular} ἄνδρα. λέγει\FnParse{d}{present active indicative 3rd singular} αὐτῇ ὁ Ἰησοῦς· Καλῶς εἶπας\FnParse{e}{aorist active indicative 2nd singular} ὅτι Ἄνδρα οὐκ ἔχω·\FnParse{f}{present active indicative 1st singular} 
\MainTextVerseMark{4}{18}πέντε γὰρ ἄνδρας ἔσχες,\FnParse{a}{aorist active indicative 2nd singular} καὶ νῦν ὃν ἔχεις\FnParse{b}{present active indicative 2nd singular} οὐκ ἔστιν\FnParse{c}{present active indicative 3rd singular} σου ἀνήρ· τοῦτο ἀληθὲς\FnParseFormGloss{d}{accusative neuter singular}{ἀληθής}{true} εἴρηκας.\FnParse{e}{perfect active indicative 2nd singular} 
\MainTextVerseMark{4}{19}λέγει\FnParse{a}{present active indicative 3rd singular} αὐτῷ ἡ γυνή· Κύριε, θεωρῶ\FnParse{b}{present active indicative 1st singular} ὅτι προφήτης εἶ\FnParse{c}{present active indicative 2nd singular} σύ. 
\MainTextVerseMark{4}{20}οἱ πατέρες ἡμῶν ἐν τῷ ὄρει τούτῳ προσεκύνησαν·\FnParse{a}{aorist active indicative 3rd plural} καὶ ὑμεῖς λέγετε\FnParse{b}{present active indicative 2nd plural} ὅτι ἐν Ἱεροσολύμοις ἐστὶν\FnParse{c}{present active indicative 3rd singular} ὁ τόπος ὅπου ⸂προσκυνεῖν\FnParse{d}{present active infinitive} δεῖ⸃.\FnParse{e}{present active indicative 3rd singular} 
\MainTextVerseMark{4}{21}λέγει\FnParse{a}{present active indicative 3rd singular} αὐτῇ ὁ Ἰησοῦς· ⸂Πίστευέ\FnParse{b}{present active imperative 2nd singular} μοι, γύναι⸃, ὅτι ἔρχεται\FnParse{c}{present middle indicative 3rd singular} ὥρα ὅτε οὔτε ἐν τῷ ὄρει τούτῳ οὔτε ἐν Ἱεροσολύμοις προσκυνήσετε\FnParse{d}{future active indicative 2nd plural} τῷ πατρί. 
\MainTextVerseMark{4}{22}ὑμεῖς προσκυνεῖτε\FnParse{a}{present active indicative 2nd plural} ὃ οὐκ οἴδατε,\FnParse{b}{perfect active indicative 2nd plural} ἡμεῖς προσκυνοῦμεν\FnParse{c}{present active indicative 1st plural} ὃ οἴδαμεν,\FnParse{d}{perfect active indicative 1st plural} ὅτι ἡ σωτηρία ἐκ τῶν Ἰουδαίων ἐστίν·\FnParse{e}{present active indicative 3rd singular} 
\MainTextVerseMark{4}{23}ἀλλὰ ἔρχεται\FnParse{a}{present middle indicative 3rd singular} ὥρα καὶ νῦν ἐστιν,\FnParse{b}{present active indicative 3rd singular} ὅτε οἱ ἀληθινοὶ\FnParseFormGloss{c}{nominative masculine plural}{ἀληθινός}{true} προσκυνηταὶ\FnParseFormGloss{d}{nominative masculine plural}{προσκυνητής}{worshiper} προσκυνήσουσιν\FnParse{e}{future active indicative 3rd plural} τῷ πατρὶ ἐν πνεύματι καὶ ἀληθείᾳ, καὶ γὰρ ὁ πατὴρ τοιούτους ζητεῖ\FnParse{f}{present active indicative 3rd singular} τοὺς προσκυνοῦντας\FnParse{g}{present active participle accusative masculine plural} αὐτόν· 
\MainTextVerseMark{4}{24}πνεῦμα ὁ θεός, καὶ τοὺς προσκυνοῦντας\FnParse{a}{present active participle accusative masculine plural} αὐτὸν ἐν πνεύματι καὶ ἀληθείᾳ δεῖ\FnParse{b}{present active indicative 3rd singular} προσκυνεῖν.\FnParse{c}{present active infinitive} 
\MainTextVerseMark{4}{25}λέγει\FnParse{a}{present active indicative 3rd singular} αὐτῷ ἡ γυνή· Οἶδα\FnParse{b}{perfect active indicative 1st singular} ὅτι Μεσσίας\FnParseFormGloss{c}{nominative masculine singular}{Μεσσίας}{Messiah} ἔρχεται,\FnParse{d}{present middle indicative 3rd singular} ὁ λεγόμενος\FnParse{e}{present passive participle nominative masculine singular} χριστός· ὅταν ἔλθῃ\FnParse{f}{aorist active subjunctive 3rd singular} ἐκεῖνος, ἀναγγελεῖ\FnParseFormGloss{g}{future active indicative 3rd singular}{ἀναγγέλλω}{to tell} ἡμῖν ⸀ἅπαντα. 
\MainTextVerseMark{4}{26}λέγει\FnParse{a}{present active indicative 3rd singular} αὐτῇ ὁ Ἰησοῦς· Ἐγώ εἰμι,\FnParse{b}{present active indicative 1st singular} ὁ λαλῶν\FnParse{c}{present active participle nominative masculine singular} σοι. 
\MainTextVerseMark{4}{27}Καὶ ἐπὶ τούτῳ ἦλθαν\FnParse{a}{aorist active indicative 3rd plural} οἱ μαθηταὶ αὐτοῦ, καὶ ⸀ἐθαύμαζον\FnParse{b}{imperfect active indicative 3rd plural} ὅτι μετὰ γυναικὸς ἐλάλει·\FnParse{c}{imperfect active indicative 3rd singular} οὐδεὶς μέντοι\FnFormGloss{d}{μέντοι}{but} εἶπεν·\FnParse{e}{aorist active indicative 3rd singular} Τί ζητεῖς;\FnParse{f}{present active indicative 2nd singular} ἢ τί λαλεῖς\FnParse{g}{present active indicative 2nd singular} μετ’ αὐτῆς; 
\MainTextVerseMark{4}{28}ἀφῆκεν\FnParse{a}{aorist active indicative 3rd singular} οὖν τὴν ὑδρίαν\FnParseFormGloss{b}{accusative feminine singular}{ὑδρία}{water jar} αὐτῆς ἡ γυνὴ καὶ ἀπῆλθεν\FnParse{c}{aorist active indicative 3rd singular} εἰς τὴν πόλιν καὶ λέγει\FnParse{d}{present active indicative 3rd singular} τοῖς ἀνθρώποις· 
\MainTextVerseMark{4}{29}Δεῦτε\FnFormGloss{a}{δεῦτε}{come} ἴδετε\FnParse{b}{aorist active imperative 2nd plural} ἄνθρωπον ὃς εἶπέ\FnParse{c}{aorist active indicative 3rd singular} μοι πάντα ⸀ὅσα ἐποίησα·\FnParse{d}{aorist active indicative 1st singular} μήτι\FnFormGloss{e}{μήτι}{[interrogative particle]} οὗτός ἐστιν\FnParse{f}{present active indicative 3rd singular} ὁ χριστός; 
\MainTextVerseMark{4}{30}ἐξῆλθον\FnParse{a}{aorist active indicative 3rd plural} ἐκ τῆς πόλεως καὶ ἤρχοντο\FnParse{b}{imperfect middle indicative 3rd plural} πρὸς αὐτόν. 
\MainTextVerseMark{4}{31}⸀Ἐν τῷ μεταξὺ\FnFormGloss{a}{μεταξύ}{between} ἠρώτων\FnParse{b}{imperfect active indicative 3rd plural} αὐτὸν οἱ μαθηταὶ λέγοντες·\FnParse{c}{present active participle nominative masculine plural} Ῥαββί,\FnParseFormGloss{d}{masculine singular}{ῥαββί}{Rabbi} φάγε.\FnParse{e}{aorist active imperative 2nd singular} 
\MainTextVerseMark{4}{32}ὁ δὲ εἶπεν\FnParse{a}{aorist active indicative 3rd singular} αὐτοῖς· Ἐγὼ βρῶσιν\FnParseFormGloss{b}{accusative feminine singular}{βρῶσις}{consumable} ἔχω\FnParse{c}{present active indicative 1st singular} φαγεῖν\FnParse{d}{aorist active infinitive} ἣν ὑμεῖς οὐκ οἴδατε.\FnParse{e}{perfect active indicative 2nd plural} 
\MainTextVerseMark{4}{33}ἔλεγον\FnParse{a}{imperfect active indicative 3rd plural} οὖν οἱ μαθηταὶ πρὸς ἀλλήλους· Μή τις ἤνεγκεν\FnParse{b}{aorist active indicative 3rd singular} αὐτῷ φαγεῖν;\FnParse{c}{aorist active infinitive} 
\MainTextVerseMark{4}{34}λέγει\FnParse{a}{present active indicative 3rd singular} αὐτοῖς ὁ Ἰησοῦς· Ἐμὸν βρῶμά\FnParseFormGloss{b}{nominative neuter singular}{βρῶμα}{food} ἐστιν\FnParse{c}{present active indicative 3rd singular} ἵνα ⸀ποιήσω\FnParse{d}{aorist active subjunctive 1st singular} τὸ θέλημα τοῦ πέμψαντός\FnParse{e}{aorist active participle genitive masculine singular} με καὶ τελειώσω\FnParseFormGloss{f}{aorist active subjunctive 1st singular}{τελειόω}{to perfect} αὐτοῦ τὸ ἔργον. 
\MainTextVerseMark{4}{35}οὐχ ὑμεῖς λέγετε\FnParse{a}{present active indicative 2nd plural} ὅτι Ἔτι τετράμηνός\FnParseFormGloss{b}{nominative masculine singular}{τετράμηνος}{four months} ἐστιν\FnParse{c}{present active indicative 3rd singular} καὶ ὁ θερισμὸς\FnParseFormGloss{d}{nominative masculine singular}{θερισμός}{harvest} ἔρχεται;\FnParse{e}{present middle indicative 3rd singular} ἰδοὺ λέγω\FnParse{f}{present active indicative 1st singular} ὑμῖν, ἐπάρατε\FnParseFormGloss{g}{aorist active imperative 2nd plural}{ἐπαίρω}{to lift up} τοὺς ὀφθαλμοὺς ὑμῶν καὶ θεάσασθε\FnParseFormGloss{h}{aorist middle imperative 2nd plural}{θεάομαι}{to see} τὰς χώρας\FnParseFormGloss{i}{accusative feminine plural}{χώρα}{country} ὅτι λευκαί\FnParseFormGloss{j}{nominative feminine plural}{λευκός}{white} εἰσιν\FnParse{k}{present active indicative 3rd plural} πρὸς θερισμόν·\FnParseFormGloss{l}{accusative masculine singular}{θερισμός}{harvest} ἤδη 
\MainTextVerseMark{4}{36}⸀ὁ θερίζων\FnParseFormGloss{a}{present active participle nominative masculine singular}{θερίζω}{to reap} μισθὸν\FnParseFormGloss{b}{accusative masculine singular}{μισθός}{wage} λαμβάνει\FnParse{c}{present active indicative 3rd singular} καὶ συνάγει\FnParse{d}{present active indicative 3rd singular} καρπὸν εἰς ζωὴν αἰώνιον, ⸀ἵνα ὁ σπείρων\FnParse{e}{present active participle nominative masculine singular} ὁμοῦ\FnFormGloss{f}{ὁμοῦ}{together} χαίρῃ\FnParse{g}{present active subjunctive 3rd singular} καὶ ὁ θερίζων.\FnParseFormGloss{h}{present active participle nominative masculine singular}{θερίζω}{to reap} 
\MainTextVerseMark{4}{37}ἐν γὰρ τούτῳ ὁ λόγος ⸀ἐστὶν\FnParse{a}{present active indicative 3rd singular} ἀληθινὸς\FnParseFormGloss{b}{nominative masculine singular}{ἀληθινός}{true} ὅτι Ἄλλος ἐστὶν\FnParse{c}{present active indicative 3rd singular} ὁ σπείρων\FnParse{d}{present active participle nominative masculine singular} καὶ ἄλλος ὁ θερίζων·\FnParseFormGloss{e}{present active participle nominative masculine singular}{θερίζω}{to reap} 
\MainTextVerseMark{4}{38}ἐγὼ ἀπέστειλα\FnParse{a}{aorist active indicative 1st singular} ὑμᾶς θερίζειν\FnParseFormGloss{b}{present active infinitive}{θερίζω}{to reap} ὃ οὐχ ὑμεῖς κεκοπιάκατε·\FnParseFormGloss{c}{perfect active indicative 2nd plural}{κοπιάω}{to work} ἄλλοι κεκοπιάκασιν,\FnParseFormGloss{d}{perfect active indicative 3rd plural}{κοπιάω}{to work} καὶ ὑμεῖς εἰς τὸν κόπον\FnParseFormGloss{e}{accusative masculine singular}{κόπος}{labor} αὐτῶν εἰσεληλύθατε.\FnParse{f}{perfect active indicative 2nd plural} 
\MainTextVerseMark{4}{39}Ἐκ δὲ τῆς πόλεως ἐκείνης πολλοὶ ἐπίστευσαν\FnParse{a}{aorist active indicative 3rd plural} εἰς αὐτὸν τῶν Σαμαριτῶν\FnParseFormGloss{b}{genitive masculine plural}{Σαμαρίτης}{Samaritan} διὰ τὸν λόγον τῆς γυναικὸς μαρτυρούσης\FnParse{c}{present active participle genitive feminine singular} ὅτι Εἶπέν\FnParse{d}{aorist active indicative 3rd singular} μοι πάντα ⸀ἃ ἐποίησα.\FnParse{e}{aorist active indicative 1st singular} 
\MainTextVerseMark{4}{40}ὡς οὖν ἦλθον\FnParse{a}{aorist active indicative 3rd plural} πρὸς αὐτὸν οἱ Σαμαρῖται,\FnParseFormGloss{b}{nominative masculine plural}{Σαμαρίτης}{Samaritan} ἠρώτων\FnParse{c}{imperfect active indicative 3rd plural} αὐτὸν μεῖναι\FnParse{d}{aorist active infinitive} παρ’ αὐτοῖς· καὶ ἔμεινεν\FnParse{e}{aorist active indicative 3rd singular} ἐκεῖ δύο ἡμέρας. 
\MainTextVerseMark{4}{41}καὶ πολλῷ πλείους ἐπίστευσαν\FnParse{a}{aorist active indicative 3rd plural} διὰ τὸν λόγον αὐτοῦ, 
\MainTextVerseMark{4}{42}τῇ τε γυναικὶ ἔλεγον\FnParse{a}{imperfect active indicative 3rd plural} ὅτι Οὐκέτι διὰ τὴν σὴν\FnParseFormGloss{b}{accusative feminine singular}{σός}{your} λαλιὰν\FnParseFormGloss{c}{accusative feminine singular}{λαλιά}{speech} πιστεύομεν·\FnParse{d}{present active indicative 1st plural} αὐτοὶ γὰρ ἀκηκόαμεν,\FnParse{e}{perfect active indicative 1st plural} καὶ οἴδαμεν\FnParse{f}{perfect active indicative 1st plural} ὅτι οὗτός ἐστιν\FnParse{g}{present active indicative 3rd singular} ἀληθῶς\FnFormGloss{h}{ἀληθῶς}{truly} ὁ σωτὴρ\FnParseFormGloss{i}{nominative masculine singular}{σωτήρ}{Savior} τοῦ ⸀κόσμου. 
\MainTextVerseMark{4}{43}Μετὰ δὲ τὰς δύο ἡμέρας ἐξῆλθεν\FnParse{a}{aorist active indicative 3rd singular} ⸀ἐκεῖθεν\FnFormGloss{b}{ἐκεῖθεν}{from there} εἰς τὴν Γαλιλαίαν· 
\MainTextVerseMark{4}{44}αὐτὸς ⸀γὰρ Ἰησοῦς ἐμαρτύρησεν\FnParse{a}{aorist active indicative 3rd singular} ὅτι προφήτης ἐν τῇ ἰδίᾳ πατρίδι\FnParseFormGloss{b}{dative feminine singular}{πατρίς}{hometown} τιμὴν οὐκ ἔχει.\FnParse{c}{present active indicative 3rd singular} 
\MainTextVerseMark{4}{45}ὅτε οὖν ἦλθεν\FnParse{a}{aorist active indicative 3rd singular} εἰς τὴν Γαλιλαίαν, ἐδέξαντο\FnParse{b}{aorist middle indicative 3rd plural} αὐτὸν οἱ Γαλιλαῖοι,\FnParseFormGloss{c}{nominative masculine plural}{Γαλιλαῖος}{Galilean} πάντα ἑωρακότες\FnParse{d}{perfect active participle nominative masculine plural} ⸀ὅσα ἐποίησεν\FnParse{e}{aorist active indicative 3rd singular} ἐν Ἱεροσολύμοις ἐν τῇ ἑορτῇ,\FnParseFormGloss{f}{dative feminine singular}{ἑορτή}{feast} καὶ αὐτοὶ γὰρ ἦλθον\FnParse{g}{aorist active indicative 3rd plural} εἰς τὴν ἑορτήν.\FnParseFormGloss{h}{accusative feminine singular}{ἑορτή}{feast} 
\MainTextVerseMark{4}{46}Ἦλθεν\FnParse{a}{aorist active indicative 3rd singular} οὖν ⸀πάλιν εἰς τὴν Κανὰ\FnParseFormGloss{b}{accusative feminine singular}{Κανά}{Cana} τῆς Γαλιλαίας, ὅπου ἐποίησεν\FnParse{c}{aorist active indicative 3rd singular} τὸ ὕδωρ οἶνον. καὶ ἦν\FnParse{d}{imperfect active indicative 3rd singular} τις βασιλικὸς\FnParseFormGloss{e}{nominative masculine singular}{βασιλικός}{royal} οὗ ὁ υἱὸς ἠσθένει\FnParse{f}{imperfect active indicative 3rd singular} ἐν Καφαρναούμ.\FnParseFormGloss{g}{dative feminine singular}{Καφαρναούμ}{Capernaum} 
\MainTextVerseMark{4}{47}οὗτος ἀκούσας\FnParse{a}{aorist active participle nominative masculine singular} ὅτι Ἰησοῦς ἥκει\FnParseFormGloss{b}{present active indicative 3rd singular}{ἥκω}{to come} ἐκ τῆς Ἰουδαίας εἰς τὴν Γαλιλαίαν ἀπῆλθεν\FnParse{c}{aorist active indicative 3rd singular} πρὸς αὐτὸν καὶ ⸀ἠρώτα\FnParse{d}{imperfect active indicative 3rd singular} ἵνα καταβῇ\FnParse{e}{aorist active subjunctive 3rd singular} καὶ ἰάσηται\FnParseFormGloss{f}{aorist middle subjunctive 3rd singular}{ἰάομαι}{to heal} αὐτοῦ τὸν υἱόν, ἤμελλεν\FnParse{g}{imperfect active indicative 3rd singular} γὰρ ἀποθνῄσκειν.\FnParse{h}{present active infinitive} 
\MainTextVerseMark{4}{48}εἶπεν\FnParse{a}{aorist active indicative 3rd singular} οὖν ὁ Ἰησοῦς πρὸς αὐτόν· Ἐὰν μὴ σημεῖα καὶ τέρατα\FnParseFormGloss{b}{accusative neuter plural}{τέρας}{wonder} ἴδητε,\FnParse{c}{aorist active subjunctive 2nd plural} οὐ μὴ πιστεύσητε.\FnParse{d}{aorist active subjunctive 2nd plural} 
\MainTextVerseMark{4}{49}λέγει\FnParse{a}{present active indicative 3rd singular} πρὸς αὐτὸν ὁ βασιλικός·\FnParseFormGloss{b}{nominative masculine singular}{βασιλικός}{royal} Κύριε, κατάβηθι\FnParse{c}{aorist active imperative 2nd singular} πρὶν\FnFormGloss{d}{πρίν}{before} ἀποθανεῖν\FnParse{e}{aorist active infinitive} τὸ παιδίον μου. 
\MainTextVerseMark{4}{50}λέγει\FnParse{a}{present active indicative 3rd singular} αὐτῷ ὁ Ἰησοῦς· Πορεύου·\FnParse{b}{present middle imperative 2nd singular} ὁ υἱός σου ζῇ.\FnParse{c}{present active indicative 3rd singular} ⸀ἐπίστευσεν\FnParse{d}{aorist active indicative 3rd singular} ὁ ἄνθρωπος τῷ λόγῳ ⸀ὃν εἶπεν\FnParse{e}{aorist active indicative 3rd singular} αὐτῷ ὁ Ἰησοῦς καὶ ἐπορεύετο.\FnParse{f}{imperfect middle indicative 3rd singular} 
\MainTextVerseMark{4}{51}ἤδη δὲ αὐτοῦ καταβαίνοντος\FnParse{a}{present active participle genitive masculine singular} οἱ δοῦλοι αὐτοῦ ⸀ὑπήντησαν\FnParseFormGloss{b}{aorist active indicative 3rd plural}{ὑπαντάω}{to go out to meet} ⸀αὐτῷ λέγοντες\FnParse{c}{present active participle nominative masculine plural} ὅτι ὁ παῖς\FnParseFormGloss{d}{nominative masculine singular}{παῖς}{boy} ⸀αὐτοῦ ζῇ.\FnParse{e}{present active indicative 3rd singular} 
\MainTextVerseMark{4}{52}ἐπύθετο\FnParseFormGloss{a}{aorist middle indicative 3rd singular}{πυνθάνομαι}{to ask} οὖν ⸂τὴν ὥραν παρ’ αὐτῶν⸃ ἐν ᾗ κομψότερον\FnParseFormGloss{b}{accusative neuter singular}{κομψότερον}{better} ἔσχεν·\FnParse{c}{aorist active indicative 3rd singular} ⸂εἶπαν\FnParse{d}{aorist active indicative 3rd plural} οὖν⸃ αὐτῷ ὅτι ⸀Ἐχθὲς\FnFormGloss{e}{ἐχθές}{yesterday} ὥραν ἑβδόμην\FnParseFormGloss{f}{accusative feminine singular}{ἕβδομος}{seventh} ἀφῆκεν\FnParse{g}{aorist active indicative 3rd singular} αὐτὸν ὁ πυρετός.\FnParseFormGloss{h}{nominative masculine singular}{πυρετός}{fever} 
\MainTextVerseMark{4}{53}ἔγνω\FnParse{a}{aorist active indicative 3rd singular} οὖν ὁ πατὴρ ⸀ὅτι ἐκείνῃ τῇ ὥρᾳ ἐν ᾗ εἶπεν\FnParse{b}{aorist active indicative 3rd singular} αὐτῷ ὁ ⸀Ἰησοῦς· Ὁ υἱός σου ζῇ,\FnParse{c}{present active indicative 3rd singular} καὶ ἐπίστευσεν\FnParse{d}{aorist active indicative 3rd singular} αὐτὸς καὶ ἡ οἰκία αὐτοῦ ὅλη. 
\MainTextVerseMark{4}{54}τοῦτο ⸀δὲ πάλιν δεύτερον σημεῖον ἐποίησεν\FnParse{a}{aorist active indicative 3rd singular} ὁ Ἰησοῦς ἐλθὼν\FnParse{b}{aorist active participle nominative masculine singular} ἐκ τῆς Ἰουδαίας εἰς τὴν Γαλιλαίαν. 
\ChapterHeader{5}{Chapter 5}

\MainTextVerseMark{5}{1}Μετὰ ταῦτα ⸀ἦν\FnParse{a}{imperfect active indicative 3rd singular} ἑορτὴ\FnParseFormGloss{b}{nominative feminine singular}{ἑορτή}{feast} τῶν Ἰουδαίων, καὶ ⸀ἀνέβη\FnParse{c}{aorist active indicative 3rd singular} Ἰησοῦς εἰς Ἱεροσόλυμα. 
\MainTextVerseMark{5}{2}ἔστιν\FnParse{a}{present active indicative 3rd singular} δὲ ἐν τοῖς Ἱεροσολύμοις ἐπὶ τῇ προβατικῇ\FnParseFormGloss{b}{dative feminine singular}{προβατικός}{of sheep} κολυμβήθρα\FnParseFormGloss{c}{nominative feminine singular}{κολυμβήθρα}{pool} ἡ ἐπιλεγομένη\FnParseFormGloss{d}{present passive participle nominative feminine singular}{ἐπιλέγω}{to be called} Ἑβραϊστὶ\FnFormGloss{e}{Ἑβραϊστί}{in Aramaic} ⸀Βηθεσδά,\FnParseFormGloss{f}{nominative feminine singular}{Βηθεσδά}{Bethesda} πέντε στοὰς\FnParseFormGloss{g}{accusative feminine plural}{στοά}{covered colonnade} ἔχουσα·\FnParse{h}{present active participle nominative feminine singular} 
\MainTextVerseMark{5}{3}ἐν ταύταις κατέκειτο\FnParseFormGloss{a}{imperfect middle indicative 3rd singular}{κατάκειμαι}{to lie down} ⸀πλῆθος τῶν ἀσθενούντων,\FnParse{b}{present active participle genitive masculine plural} τυφλῶν, χωλῶν,\FnParseFormGloss{c}{genitive masculine plural}{χωλός}{lame} ⸀ξηρῶν.\FnParseFormGloss{d}{genitive masculine plural}{ξηρός}{dried up} 
\MainTextVerseMark{5}{5}ἦν\FnParse{a}{imperfect active indicative 3rd singular} δέ τις ἄνθρωπος ἐκεῖ ⸀τριάκοντα\FnParseFormGloss{b}{accusative neuter plural}{τριάκοντα}{thirty} ὀκτὼ\FnParseFormGloss{c}{accusative neuter plural}{ὀκτώ}{eight} ἔτη ἔχων\FnParse{d}{present active participle nominative masculine singular} ἐν τῇ ἀσθενείᾳ\FnParseFormGloss{e}{dative feminine singular}{ἀσθένεια}{weakness} ⸀αὐτοῦ· 
\MainTextVerseMark{5}{6}τοῦτον ἰδὼν\FnParse{a}{aorist active participle nominative masculine singular} ὁ Ἰησοῦς κατακείμενον,\FnParseFormGloss{b}{present middle participle accusative masculine singular}{κατάκειμαι}{to lie down} καὶ γνοὺς\FnParse{c}{aorist active participle nominative masculine singular} ὅτι πολὺν ἤδη χρόνον ἔχει,\FnParse{d}{present active indicative 3rd singular} λέγει\FnParse{e}{present active indicative 3rd singular} αὐτῷ· Θέλεις\FnParse{f}{present active indicative 2nd singular} ὑγιὴς\FnParseFormGloss{g}{nominative masculine singular}{ὑγιής}{healthy} γενέσθαι;\FnParse{h}{aorist middle infinitive} 
\MainTextVerseMark{5}{7}ἀπεκρίθη\FnParse{a}{aorist passive indicative 3rd singular} αὐτῷ ὁ ἀσθενῶν·\FnParse{b}{present active participle nominative masculine singular} Κύριε, ἄνθρωπον οὐκ ἔχω\FnParse{c}{present active indicative 1st singular} ἵνα ὅταν ταραχθῇ\FnParseFormGloss{d}{aorist passive subjunctive 3rd singular}{ταράσσω}{to trouble} τὸ ὕδωρ βάλῃ\FnParse{e}{aorist active subjunctive 3rd singular} με εἰς τὴν κολυμβήθραν·\FnParseFormGloss{f}{accusative feminine singular}{κολυμβήθρα}{pool} ἐν ᾧ δὲ ἔρχομαι\FnParse{g}{present middle indicative 1st singular} ἐγὼ ἄλλος πρὸ ἐμοῦ καταβαίνει.\FnParse{h}{present active indicative 3rd singular} 
\MainTextVerseMark{5}{8}λέγει\FnParse{a}{present active indicative 3rd singular} αὐτῷ ὁ Ἰησοῦς· ⸀Ἔγειρε\FnParse{b}{present active imperative 2nd singular} ἆρον\FnParse{c}{aorist active imperative 2nd singular} τὸν κράβαττόν\FnParseFormGloss{d}{accusative masculine singular}{κράβαττος}{bed} σου καὶ περιπάτει.\FnParse{e}{present active imperative 2nd singular} 
\MainTextVerseMark{5}{9}καὶ εὐθέως ἐγένετο\FnParse{a}{aorist middle indicative 3rd singular} ὑγιὴς\FnParseFormGloss{b}{nominative masculine singular}{ὑγιής}{healthy} ὁ ἄνθρωπος καὶ ἦρε\FnParse{c}{aorist active indicative 3rd singular} τὸν κράβαττον\FnParseFormGloss{d}{accusative masculine singular}{κράβαττος}{bed} αὐτοῦ καὶ περιεπάτει.\FnParse{e}{imperfect active indicative 3rd singular} Ἦν\FnParse{f}{imperfect active indicative 3rd singular} δὲ σάββατον ἐν ἐκείνῃ τῇ ἡμέρᾳ. 
\MainTextVerseMark{5}{10}ἔλεγον\FnParse{a}{imperfect active indicative 3rd plural} οὖν οἱ Ἰουδαῖοι τῷ τεθεραπευμένῳ·\FnParse{b}{perfect passive participle dative masculine singular} Σάββατόν ἐστιν,\FnParse{c}{present active indicative 3rd singular} ⸀καὶ οὐκ ἔξεστίν\FnParse{d}{present active indicative 3rd singular} σοι ἆραι\FnParse{e}{aorist active infinitive} τὸν ⸀κράβαττον.\FnParseFormGloss{f}{accusative masculine singular}{κράβαττος}{bed} 
\MainTextVerseMark{5}{11}⸂ὃς δὲ⸃ ἀπεκρίθη\FnParse{a}{aorist passive indicative 3rd singular} αὐτοῖς· Ὁ ποιήσας\FnParse{b}{aorist active participle nominative masculine singular} με ὑγιῆ\FnParseFormGloss{c}{accusative masculine singular}{ὑγιής}{healthy} ἐκεῖνός μοι εἶπεν\FnParse{d}{aorist active indicative 3rd singular} Ἆρον\FnParse{e}{aorist active imperative 2nd singular} τὸν κράβαττόν\FnParseFormGloss{f}{accusative masculine singular}{κράβαττος}{bed} σου καὶ περιπάτει.\FnParse{g}{present active imperative 2nd singular} 
\MainTextVerseMark{5}{12}ἠρώτησαν\FnParse{a}{aorist active indicative 3rd plural} ⸀οὖν αὐτόν· Τίς ἐστιν\FnParse{b}{present active indicative 3rd singular} ὁ ἄνθρωπος ὁ εἰπών\FnParse{c}{aorist active participle nominative masculine singular} σοι· ⸀Ἆρον\FnParse{d}{aorist active imperative 2nd singular} καὶ περιπάτει;\FnParse{e}{present active imperative 2nd singular} 
\MainTextVerseMark{5}{13}ὁ δὲ ἰαθεὶς\FnParseFormGloss{a}{aorist passive participle nominative masculine singular}{ἰάομαι}{to heal} οὐκ ᾔδει\FnParse{b}{pluperfect active indicative 3rd singular} τίς ἐστιν,\FnParse{c}{present active indicative 3rd singular} ὁ γὰρ Ἰησοῦς ἐξένευσεν\FnParseFormGloss{d}{aorist active indicative 3rd singular}{ἐκνεύω}{to slip away} ὄχλου ὄντος\FnParse{e}{present active participle genitive masculine singular} ἐν τῷ τόπῳ. 
\MainTextVerseMark{5}{14}μετὰ ταῦτα εὑρίσκει\FnParse{a}{present active indicative 3rd singular} αὐτὸν ὁ Ἰησοῦς ἐν τῷ ἱερῷ καὶ εἶπεν\FnParse{b}{aorist active indicative 3rd singular} αὐτῷ· Ἴδε ὑγιὴς\FnParseFormGloss{c}{nominative masculine singular}{ὑγιής}{healthy} γέγονας·\FnParse{d}{perfect active indicative 2nd singular} μηκέτι\FnFormGloss{e}{μηκέτι}{no longer} ἁμάρτανε,\FnParse{f}{present active imperative 2nd singular} ἵνα μὴ χεῖρόν\FnParseFormGloss{g}{nominative neuter singular}{χείρων}{worse} ⸂σοί τι⸃ γένηται.\FnParse{h}{aorist middle subjunctive 3rd singular} 
\MainTextVerseMark{5}{15}ἀπῆλθεν\FnParse{a}{aorist active indicative 3rd singular} ὁ ἄνθρωπος καὶ ⸀ἀνήγγειλεν\FnParseFormGloss{b}{aorist active indicative 3rd singular}{ἀναγγέλλω}{to tell} τοῖς Ἰουδαίοις ὅτι Ἰησοῦς ἐστιν\FnParse{c}{present active indicative 3rd singular} ὁ ποιήσας\FnParse{d}{aorist active participle nominative masculine singular} αὐτὸν ὑγιῆ.\FnParseFormGloss{e}{accusative masculine singular}{ὑγιής}{healthy} 
\MainTextVerseMark{5}{16}καὶ διὰ τοῦτο ἐδίωκον\FnParse{a}{imperfect active indicative 3rd plural} ⸂οἱ Ἰουδαῖοι τὸν Ἰησοῦν⸃ ὅτι ταῦτα ἐποίει\FnParse{b}{imperfect active indicative 3rd singular} ἐν σαββάτῳ. 
\MainTextVerseMark{5}{17}ὁ ⸀δὲ ἀπεκρίνατο\FnParse{a}{aorist middle indicative 3rd singular} αὐτοῖς· Ὁ πατήρ μου ἕως ἄρτι ἐργάζεται\FnParse{b}{present middle indicative 3rd singular} κἀγὼ ἐργάζομαι.\FnParse{c}{present middle indicative 1st singular} 
\MainTextVerseMark{5}{18}διὰ τοῦτο οὖν μᾶλλον ἐζήτουν\FnParse{a}{imperfect active indicative 3rd plural} αὐτὸν οἱ Ἰουδαῖοι ἀποκτεῖναι\FnParse{b}{aorist active infinitive} ὅτι οὐ μόνον ἔλυε\FnParse{c}{imperfect active indicative 3rd singular} τὸ σάββατον, ἀλλὰ καὶ πατέρα ἴδιον ἔλεγε\FnParse{d}{imperfect active indicative 3rd singular} τὸν θεόν, ἴσον\FnParseFormGloss{e}{accusative masculine singular}{ἴσος}{equal} ἑαυτὸν ποιῶν\FnParse{f}{present active participle nominative masculine singular} τῷ θεῷ. 
\MainTextVerseMark{5}{19}Ἀπεκρίνατο\FnParse{a}{aorist middle indicative 3rd singular} οὖν ὁ Ἰησοῦς καὶ ⸀ἔλεγεν\FnParse{b}{imperfect active indicative 3rd singular} αὐτοῖς· Ἀμὴν ἀμὴν λέγω\FnParse{c}{present active indicative 1st singular} ὑμῖν, οὐ δύναται\FnParse{d}{present middle indicative 3rd singular} ὁ υἱὸς ποιεῖν\FnParse{e}{present active infinitive} ἀφ’ ἑαυτοῦ οὐδὲν ⸀ἐὰν μή τι βλέπῃ\FnParse{f}{present active subjunctive 3rd singular} τὸν πατέρα ποιοῦντα·\FnParse{g}{present active participle accusative masculine singular} ἃ γὰρ ἂν ἐκεῖνος ποιῇ,\FnParse{h}{present active subjunctive 3rd singular} ταῦτα καὶ ὁ υἱὸς ὁμοίως ποιεῖ.\FnParse{i}{present active indicative 3rd singular} 
\MainTextVerseMark{5}{20}ὁ γὰρ πατὴρ φιλεῖ\FnParseFormGloss{a}{present active indicative 3rd singular}{φιλέω}{to love} τὸν υἱὸν καὶ πάντα δείκνυσιν\FnParse{b}{present active indicative 3rd singular} αὐτῷ ἃ αὐτὸς ποιεῖ,\FnParse{c}{present active indicative 3rd singular} καὶ μείζονα τούτων δείξει\FnParse{d}{future active indicative 3rd singular} αὐτῷ ἔργα, ἵνα ὑμεῖς θαυμάζητε.\FnParse{e}{present active subjunctive 2nd plural} 
\MainTextVerseMark{5}{21}ὥσπερ γὰρ ὁ πατὴρ ἐγείρει\FnParse{a}{present active indicative 3rd singular} τοὺς νεκροὺς καὶ ζῳοποιεῖ,\FnParseFormGloss{b}{present active indicative 3rd singular}{ζῳοποιέω}{to make alive} οὕτως καὶ ὁ υἱὸς οὓς θέλει\FnParse{c}{present active indicative 3rd singular} ζῳοποιεῖ.\FnParseFormGloss{d}{present active indicative 3rd singular}{ζῳοποιέω}{to make alive} 
\MainTextVerseMark{5}{22}οὐδὲ γὰρ ὁ πατὴρ κρίνει\FnParse{a}{present active indicative 3rd singular} οὐδένα, ἀλλὰ τὴν κρίσιν πᾶσαν δέδωκεν\FnParse{b}{perfect active indicative 3rd singular} τῷ υἱῷ, 
\MainTextVerseMark{5}{23}ἵνα πάντες τιμῶσι\FnParseFormGloss{a}{present active subjunctive 3rd plural}{τιμάω}{to honor} τὸν υἱὸν καθὼς τιμῶσι\FnParseFormGloss{b}{present active indicative 3rd plural}{τιμάω}{to honor} τὸν πατέρα. ὁ μὴ τιμῶν\FnParseFormGloss{c}{present active participle nominative masculine singular}{τιμάω}{to honor} τὸν υἱὸν οὐ τιμᾷ\FnParseFormGloss{d}{present active indicative 3rd singular}{τιμάω}{to honor} τὸν πατέρα τὸν πέμψαντα\FnParse{e}{aorist active participle accusative masculine singular} αὐτόν. 
\MainTextVerseMark{5}{24}Ἀμὴν ἀμὴν λέγω\FnParse{a}{present active indicative 1st singular} ὑμῖν ὅτι ὁ τὸν λόγον μου ἀκούων\FnParse{b}{present active participle nominative masculine singular} καὶ πιστεύων\FnParse{c}{present active participle nominative masculine singular} τῷ πέμψαντί\FnParse{d}{aorist active participle dative masculine singular} με ἔχει\FnParse{e}{present active indicative 3rd singular} ζωὴν αἰώνιον, καὶ εἰς κρίσιν οὐκ ἔρχεται\FnParse{f}{present middle indicative 3rd singular} ἀλλὰ μεταβέβηκεν\FnParseFormGloss{g}{perfect active indicative 3rd singular}{μεταβαίνω}{to go on} ἐκ τοῦ θανάτου εἰς τὴν ζωήν. 
\MainTextVerseMark{5}{25}Ἀμὴν ἀμὴν λέγω\FnParse{a}{present active indicative 1st singular} ὑμῖν ὅτι ἔρχεται\FnParse{b}{present middle indicative 3rd singular} ὥρα καὶ νῦν ἐστιν\FnParse{c}{present active indicative 3rd singular} ὅτε οἱ νεκροὶ ⸀ἀκούσουσιν\FnParse{d}{future active indicative 3rd plural} τῆς φωνῆς τοῦ υἱοῦ τοῦ θεοῦ καὶ οἱ ἀκούσαντες\FnParse{e}{aorist active participle nominative masculine plural} ⸀ζήσουσιν.\FnParse{f}{future active indicative 3rd plural} 
\MainTextVerseMark{5}{26}ὥσπερ γὰρ ὁ πατὴρ ἔχει\FnParse{a}{present active indicative 3rd singular} ζωὴν ἐν ἑαυτῷ, οὕτως ⸂καὶ τῷ υἱῷ ἔδωκεν⸃\FnParse{b}{aorist active indicative 3rd singular} ζωὴν ἔχειν\FnParse{c}{present active infinitive} ἐν ἑαυτῷ· 
\MainTextVerseMark{5}{27}καὶ ἐξουσίαν ἔδωκεν\FnParse{a}{aorist active indicative 3rd singular} ⸀αὐτῷ κρίσιν ποιεῖν,\FnParse{b}{present active infinitive} ὅτι υἱὸς ἀνθρώπου ἐστίν.\FnParse{c}{present active indicative 3rd singular} 
\MainTextVerseMark{5}{28}μὴ θαυμάζετε\FnParse{a}{present active imperative 2nd plural} τοῦτο, ὅτι ἔρχεται\FnParse{b}{present middle indicative 3rd singular} ὥρα ἐν ᾗ πάντες οἱ ἐν τοῖς μνημείοις ⸀ἀκούσουσιν\FnParse{c}{future active indicative 3rd plural} τῆς φωνῆς αὐτοῦ 
\MainTextVerseMark{5}{29}καὶ ἐκπορεύσονται\FnParse{a}{future middle indicative 3rd plural} οἱ τὰ ἀγαθὰ ποιήσαντες\FnParse{b}{aorist active participle nominative masculine plural} εἰς ἀνάστασιν ζωῆς, οἱ ⸀δὲ τὰ φαῦλα\FnParseFormGloss{c}{accusative neuter plural}{φαῦλος}{evil} πράξαντες\FnParse{d}{aorist active participle nominative masculine plural} εἰς ἀνάστασιν κρίσεως. 
\MainTextVerseMark{5}{30}Οὐ δύναμαι\FnParse{a}{present middle indicative 1st singular} ἐγὼ ποιεῖν\FnParse{b}{present active infinitive} ἀπ’ ἐμαυτοῦ οὐδέν· καθὼς ἀκούω\FnParse{c}{present active indicative 1st singular} κρίνω,\FnParse{d}{present active indicative 1st singular} καὶ ἡ κρίσις ἡ ἐμὴ δικαία ἐστίν,\FnParse{e}{present active indicative 3rd singular} ὅτι οὐ ζητῶ\FnParse{f}{present active indicative 1st singular} τὸ θέλημα τὸ ἐμὸν ἀλλὰ τὸ θέλημα τοῦ πέμψαντός\FnParse{g}{aorist active participle genitive masculine singular} ⸀με. 
\MainTextVerseMark{5}{31}Ἐὰν ἐγὼ μαρτυρῶ\FnParse{a}{present active subjunctive 1st singular} περὶ ἐμαυτοῦ, ἡ μαρτυρία μου οὐκ ἔστιν\FnParse{b}{present active indicative 3rd singular} ἀληθής·\FnParseFormGloss{c}{nominative feminine singular}{ἀληθής}{true} 
\MainTextVerseMark{5}{32}ἄλλος ἐστὶν\FnParse{a}{present active indicative 3rd singular} ὁ μαρτυρῶν\FnParse{b}{present active participle nominative masculine singular} περὶ ἐμοῦ, καὶ οἶδα\FnParse{c}{perfect active indicative 1st singular} ὅτι ἀληθής\FnParseFormGloss{d}{nominative feminine singular}{ἀληθής}{true} ἐστιν\FnParse{e}{present active indicative 3rd singular} ἡ μαρτυρία ἣν μαρτυρεῖ\FnParse{f}{present active indicative 3rd singular} περὶ ἐμοῦ. 
\MainTextVerseMark{5}{33}ὑμεῖς ἀπεστάλκατε\FnParse{a}{perfect active indicative 2nd plural} πρὸς Ἰωάννην, καὶ μεμαρτύρηκε\FnParse{b}{perfect active indicative 3rd singular} τῇ ἀληθείᾳ· 
\MainTextVerseMark{5}{34}ἐγὼ δὲ οὐ παρὰ ἀνθρώπου τὴν μαρτυρίαν λαμβάνω,\FnParse{a}{present active indicative 1st singular} ἀλλὰ ταῦτα λέγω\FnParse{b}{present active indicative 1st singular} ἵνα ὑμεῖς σωθῆτε.\FnParse{c}{aorist passive subjunctive 2nd plural} 
\MainTextVerseMark{5}{35}ἐκεῖνος ἦν\FnParse{a}{imperfect active indicative 3rd singular} ὁ λύχνος\FnParseFormGloss{b}{nominative masculine singular}{λύχνος}{lamp} ὁ καιόμενος\FnParseFormGloss{c}{present passive participle nominative masculine singular}{καίω}{to light} καὶ φαίνων,\FnParse{d}{present active participle nominative masculine singular} ὑμεῖς δὲ ἠθελήσατε\FnParse{e}{aorist active indicative 2nd plural} ἀγαλλιαθῆναι\FnParseFormGloss{f}{aorist passive infinitive}{ἀγαλλιάω}{to be filled with delight} πρὸς ὥραν ἐν τῷ φωτὶ αὐτοῦ· 
\MainTextVerseMark{5}{36}ἐγὼ δὲ ἔχω\FnParse{a}{present active indicative 1st singular} τὴν μαρτυρίαν ⸀μείζω τοῦ Ἰωάννου, τὰ γὰρ ἔργα ἃ ⸀δέδωκέν\FnParse{b}{perfect active indicative 3rd singular} μοι ὁ πατὴρ ἵνα τελειώσω\FnParseFormGloss{c}{aorist active subjunctive 1st singular}{τελειόω}{to perfect} αὐτά, αὐτὰ τὰ ἔργα ⸀ἃ ποιῶ,\FnParse{d}{present active indicative 1st singular} μαρτυρεῖ\FnParse{e}{present active indicative 3rd singular} περὶ ἐμοῦ ὅτι ὁ πατήρ με ἀπέσταλκεν,\FnParse{f}{perfect active indicative 3rd singular} 
\MainTextVerseMark{5}{37}καὶ ὁ πέμψας\FnParse{a}{aorist active participle nominative masculine singular} με πατὴρ ⸀ἐκεῖνος μεμαρτύρηκεν\FnParse{b}{perfect active indicative 3rd singular} περὶ ἐμοῦ. οὔτε φωνὴν αὐτοῦ ⸂πώποτε\FnFormGloss{c}{πώποτε}{ever} ἀκηκόατε⸃\FnParse{d}{perfect active indicative 2nd plural} οὔτε εἶδος\FnParseFormGloss{e}{accusative neuter singular}{εἶδος}{form} αὐτοῦ ἑωράκατε,\FnParse{f}{perfect active indicative 2nd plural} 
\MainTextVerseMark{5}{38}καὶ τὸν λόγον αὐτοῦ οὐκ ἔχετε\FnParse{a}{present active indicative 2nd plural} ⸂ἐν ὑμῖν μένοντα⸃,\FnParse{b}{present active participle accusative masculine singular} ὅτι ὃν ἀπέστειλεν\FnParse{c}{aorist active indicative 3rd singular} ἐκεῖνος τούτῳ ὑμεῖς οὐ πιστεύετε.\FnParse{d}{present active indicative 2nd plural} 
\MainTextVerseMark{5}{39}Ἐραυνᾶτε\FnParseFormGloss{a}{present active indicative 2nd plural}{ἐραυνάω}{to search} τὰς γραφάς, ὅτι ὑμεῖς δοκεῖτε\FnParse{b}{present active indicative 2nd plural} ἐν αὐταῖς ζωὴν αἰώνιον ἔχειν·\FnParse{c}{present active infinitive} καὶ ἐκεῖναί εἰσιν\FnParse{d}{present active indicative 3rd plural} αἱ μαρτυροῦσαι\FnParse{e}{present active participle nominative feminine plural} περὶ ἐμοῦ· 
\MainTextVerseMark{5}{40}καὶ οὐ θέλετε\FnParse{a}{present active indicative 2nd plural} ἐλθεῖν\FnParse{b}{aorist active infinitive} πρός με ἵνα ζωὴν ἔχητε.\FnParse{c}{present active subjunctive 2nd plural} 
\MainTextVerseMark{5}{41}δόξαν παρὰ ἀνθρώπων οὐ λαμβάνω,\FnParse{a}{present active indicative 1st singular} 
\MainTextVerseMark{5}{42}ἀλλὰ ἔγνωκα\FnParse{a}{perfect active indicative 1st singular} ὑμᾶς ὅτι τὴν ἀγάπην τοῦ θεοῦ οὐκ ἔχετε\FnParse{b}{present active indicative 2nd plural} ἐν ἑαυτοῖς. 
\MainTextVerseMark{5}{43}ἐγὼ ἐλήλυθα\FnParse{a}{perfect active indicative 1st singular} ἐν τῷ ὀνόματι τοῦ πατρός μου καὶ οὐ λαμβάνετέ\FnParse{b}{present active indicative 2nd plural} με· ἐὰν ἄλλος ἔλθῃ\FnParse{c}{aorist active subjunctive 3rd singular} ἐν τῷ ὀνόματι τῷ ἰδίῳ, ἐκεῖνον λήμψεσθε.\FnParse{d}{future middle indicative 2nd plural} 
\MainTextVerseMark{5}{44}πῶς δύνασθε\FnParse{a}{present middle indicative 2nd plural} ὑμεῖς πιστεῦσαι,\FnParse{b}{aorist active infinitive} δόξαν παρ’ ἀλλήλων λαμβάνοντες,\FnParse{c}{present active participle nominative masculine plural} καὶ τὴν δόξαν τὴν παρὰ τοῦ μόνου θεοῦ οὐ ζητεῖτε;\FnParse{d}{present active indicative 2nd plural} 
\MainTextVerseMark{5}{45}μὴ δοκεῖτε\FnParse{a}{present active imperative 2nd plural} ὅτι ἐγὼ κατηγορήσω\FnParseFormGloss{b}{future active indicative 1st singular}{κατηγορέω}{to accuse} ὑμῶν πρὸς τὸν πατέρα· ἔστιν\FnParse{c}{present active indicative 3rd singular} ὁ κατηγορῶν\FnParseFormGloss{d}{present active participle nominative masculine singular}{κατηγορέω}{to accuse} ὑμῶν Μωϋσῆς, εἰς ὃν ὑμεῖς ἠλπίκατε.\FnParse{e}{perfect active indicative 2nd plural} 
\MainTextVerseMark{5}{46}εἰ γὰρ ἐπιστεύετε\FnParse{a}{imperfect active indicative 2nd plural} Μωϋσεῖ, ἐπιστεύετε\FnParse{b}{imperfect active indicative 2nd plural} ἂν ἐμοί, περὶ γὰρ ἐμοῦ ἐκεῖνος ἔγραψεν.\FnParse{c}{aorist active indicative 3rd singular} 
\MainTextVerseMark{5}{47}εἰ δὲ τοῖς ἐκείνου γράμμασιν\FnParseFormGloss{a}{dative neuter plural}{γράμμα}{letter} οὐ πιστεύετε,\FnParse{b}{present active indicative 2nd plural} πῶς τοῖς ἐμοῖς ῥήμασιν πιστεύσετε;\FnParse{c}{future active indicative 2nd plural} 
\ChapterHeader{6}{Chapter 6}

\MainTextVerseMark{6}{1}Μετὰ ταῦτα ἀπῆλθεν\FnParse{a}{aorist active indicative 3rd singular} ὁ Ἰησοῦς πέραν\FnFormGloss{b}{πέραν}{on the other side} τῆς θαλάσσης τῆς Γαλιλαίας τῆς Τιβεριάδος.\FnParseFormGloss{c}{genitive feminine singular}{Τιβεριάς}{Tiberias} 
\MainTextVerseMark{6}{2}⸂ἠκολούθει\FnParse{a}{imperfect active indicative 3rd singular} δὲ⸃ αὐτῷ ὄχλος πολύς, ὅτι ⸀ἐθεώρουν\FnParse{b}{imperfect active indicative 3rd plural} τὰ σημεῖα ἃ ἐποίει\FnParse{c}{imperfect active indicative 3rd singular} ἐπὶ τῶν ἀσθενούντων.\FnParse{d}{present active participle genitive masculine plural} 
\MainTextVerseMark{6}{3}ἀνῆλθεν\FnParseFormGloss{a}{aorist active indicative 3rd singular}{ἀνέρχομαι}{to go up} δὲ εἰς τὸ ⸀ὄρος Ἰησοῦς, καὶ ἐκεῖ ἐκάθητο\FnParse{b}{imperfect middle indicative 3rd singular} μετὰ τῶν μαθητῶν αὐτοῦ. 
\MainTextVerseMark{6}{4}ἦν\FnParse{a}{imperfect active indicative 3rd singular} δὲ ἐγγὺς τὸ πάσχα,\FnParseFormGloss{b}{nominative neuter singular}{πάσχα}{Passover} ἡ ἑορτὴ\FnParseFormGloss{c}{nominative feminine singular}{ἑορτή}{feast} τῶν Ἰουδαίων. 
\MainTextVerseMark{6}{5}ἐπάρας\FnParseFormGloss{a}{aorist active participle nominative masculine singular}{ἐπαίρω}{to lift up} οὖν ⸂τοὺς ὀφθαλμοὺς ὁ Ἰησοῦς⸃ καὶ θεασάμενος\FnParseFormGloss{b}{aorist middle participle nominative masculine singular}{θεάομαι}{to see} ὅτι πολὺς ὄχλος ἔρχεται\FnParse{c}{present middle indicative 3rd singular} πρὸς αὐτὸν λέγει\FnParse{d}{present active indicative 3rd singular} ⸀πρὸς Φίλιππον· Πόθεν\FnFormGloss{e}{πόθεν}{from where} ⸀ἀγοράσωμεν\FnParse{f}{aorist active subjunctive 1st plural} ἄρτους ἵνα φάγωσιν\FnParse{g}{aorist active subjunctive 3rd plural} οὗτοι; 
\MainTextVerseMark{6}{6}τοῦτο δὲ ἔλεγεν\FnParse{a}{imperfect active indicative 3rd singular} πειράζων\FnParse{b}{present active participle nominative masculine singular} αὐτόν, αὐτὸς γὰρ ᾔδει\FnParse{c}{pluperfect active indicative 3rd singular} τί ἔμελλεν\FnParse{d}{imperfect active indicative 3rd singular} ποιεῖν.\FnParse{e}{present active infinitive} 
\MainTextVerseMark{6}{7}ἀπεκρίθη\FnParse{a}{aorist passive indicative 3rd singular} ⸀αὐτῷ Φίλιππος· Διακοσίων\FnParseFormGloss{b}{genitive neuter plural}{διακόσιοι}{two hundred} δηναρίων\FnParseFormGloss{c}{genitive neuter plural}{δηνάριον}{denarius} ἄρτοι οὐκ ἀρκοῦσιν\FnParseFormGloss{d}{present active indicative 3rd plural}{ἀρκέω}{to be content} αὐτοῖς ἵνα ⸀ἕκαστος βραχύ\FnParseFormGloss{e}{accusative neuter singular}{βραχύς}{little} ⸀τι λάβῃ.\FnParse{f}{aorist active subjunctive 3rd singular} 
\MainTextVerseMark{6}{8}λέγει\FnParse{a}{present active indicative 3rd singular} αὐτῷ εἷς ἐκ τῶν μαθητῶν αὐτοῦ, Ἀνδρέας\FnParseFormGloss{b}{nominative masculine singular}{Ἀνδρέας}{Andrew} ὁ ἀδελφὸς Σίμωνος Πέτρου· 
\MainTextVerseMark{6}{9}Ἔστιν\FnParse{a}{present active indicative 3rd singular} ⸀παιδάριον\FnParseFormGloss{b}{nominative neuter singular}{παιδάριον}{little boy} ὧδε ⸀ὃς ἔχει\FnParse{c}{present active indicative 3rd singular} πέντε ἄρτους κριθίνους\FnParseFormGloss{d}{accusative masculine plural}{κρίθινος}{made of barley} καὶ δύο ὀψάρια·\FnParseFormGloss{e}{accusative neuter plural}{ὀψάριον}{fish} ἀλλὰ ταῦτα τί ἐστιν\FnParse{f}{present active indicative 3rd singular} εἰς τοσούτους;\FnParseFormGloss{g}{accusative masculine plural}{τοσοῦτος}{so great} 
\MainTextVerseMark{6}{10}⸀εἶπεν\FnParse{a}{aorist active indicative 3rd singular} ὁ Ἰησοῦς· Ποιήσατε\FnParse{b}{aorist active imperative 2nd plural} τοὺς ἀνθρώπους ἀναπεσεῖν.\FnParseFormGloss{c}{aorist active infinitive}{ἀναπίπτω}{recline} ἦν\FnParse{d}{imperfect active indicative 3rd singular} δὲ χόρτος\FnParseFormGloss{e}{nominative masculine singular}{χόρτος}{grass} πολὺς ἐν τῷ τόπῳ. ἀνέπεσαν\FnParseFormGloss{f}{aorist active indicative 3rd plural}{ἀναπίπτω}{recline} οὖν οἱ ἄνδρες τὸν ἀριθμὸν\FnParseFormGloss{g}{accusative masculine singular}{ἀριθμός}{number} ⸀ὡς πεντακισχίλιοι.\FnParseFormGloss{h}{nominative masculine plural}{πεντακισχίλιοι}{five thousand} 
\MainTextVerseMark{6}{11}ἔλαβεν\FnParse{a}{aorist active indicative 3rd singular} ⸀οὖν τοὺς ἄρτους ὁ Ἰησοῦς καὶ εὐχαριστήσας\FnParse{b}{aorist active participle nominative masculine singular} ⸀διέδωκεν\FnParseFormGloss{c}{aorist active indicative 3rd singular}{διαδίδωμι}{to distribute} τοῖς ἀνακειμένοις,\FnParseFormGloss{d}{present middle participle dative masculine plural}{ἀνάκειμαι}{to recline for a meal} ὁμοίως καὶ ἐκ τῶν ὀψαρίων\FnParseFormGloss{e}{genitive neuter plural}{ὀψάριον}{fish} ὅσον ἤθελον.\FnParse{f}{imperfect active indicative 3rd plural} 
\MainTextVerseMark{6}{12}ὡς δὲ ἐνεπλήσθησαν\FnParseFormGloss{a}{aorist passive indicative 3rd plural}{ἐμπίμπλημι}{to fill} λέγει\FnParse{b}{present active indicative 3rd singular} τοῖς μαθηταῖς αὐτοῦ· Συναγάγετε\FnParse{c}{aorist active imperative 2nd plural} τὰ περισσεύσαντα\FnParse{d}{aorist active participle accusative neuter plural} κλάσματα,\FnParseFormGloss{e}{accusative neuter plural}{κλάσμα}{broken piece} ἵνα μή τι ἀπόληται.\FnParse{f}{aorist middle subjunctive 3rd singular} 
\MainTextVerseMark{6}{13}συνήγαγον\FnParse{a}{aorist active indicative 3rd plural} οὖν, καὶ ἐγέμισαν\FnParseFormGloss{b}{aorist active indicative 3rd plural}{γεμίζω}{to fill} δώδεκα κοφίνους\FnParseFormGloss{c}{accusative masculine plural}{κόφινος}{a large basket} κλασμάτων\FnParseFormGloss{d}{genitive neuter plural}{κλάσμα}{broken piece} ἐκ τῶν πέντε ἄρτων τῶν κριθίνων\FnParseFormGloss{e}{genitive masculine plural}{κρίθινος}{made of barley} ἃ ⸀ἐπερίσσευσαν\FnParse{f}{aorist active indicative 3rd plural} τοῖς βεβρωκόσιν.\FnParseFormGloss{g}{perfect active participle dative masculine plural}{βιβρώσκω}{to eat} 
\MainTextVerseMark{6}{14}οἱ οὖν ἄνθρωποι ἰδόντες\FnParse{a}{aorist active participle nominative masculine plural} ⸂ὃ ἐποίησεν\FnParse{b}{aorist active indicative 3rd singular} σημεῖον⸃ ⸀ἔλεγον\FnParse{c}{imperfect active indicative 3rd plural} ὅτι Οὗτός ἐστιν\FnParse{d}{present active indicative 3rd singular} ἀληθῶς\FnFormGloss{e}{ἀληθῶς}{truly} ὁ προφήτης ὁ ἐρχόμενος\FnParse{f}{present middle participle nominative masculine singular} εἰς τὸν κόσμον. 
\MainTextVerseMark{6}{15}Ἰησοῦς οὖν γνοὺς\FnParse{a}{aorist active participle nominative masculine singular} ὅτι μέλλουσιν\FnParse{b}{present active indicative 3rd plural} ἔρχεσθαι\FnParse{c}{present middle infinitive} καὶ ἁρπάζειν\FnParseFormGloss{d}{present active infinitive}{ἁρπάζω}{to catch} αὐτὸν ἵνα ⸀ποιήσωσιν\FnParse{e}{aorist active subjunctive 3rd plural} βασιλέα ἀνεχώρησεν\FnParseFormGloss{f}{aorist active indicative 3rd singular}{ἀναχωρέω}{to withdraw} ⸀πάλιν εἰς τὸ ὄρος αὐτὸς μόνος. 
\MainTextVerseMark{6}{16}Ὡς δὲ ὀψία\FnParseFormGloss{a}{nominative feminine singular}{ὀψία}{evening} ἐγένετο\FnParse{b}{aorist middle indicative 3rd singular} κατέβησαν\FnParse{c}{aorist active indicative 3rd plural} οἱ μαθηταὶ αὐτοῦ ἐπὶ τὴν θάλασσαν, 
\MainTextVerseMark{6}{17}καὶ ἐμβάντες\FnParseFormGloss{a}{aorist active participle nominative masculine plural}{ἐμβαίνω}{to get into} ⸀εἰς πλοῖον ἤρχοντο\FnParse{b}{imperfect middle indicative 3rd plural} πέραν\FnFormGloss{c}{πέραν}{on the other side} τῆς θαλάσσης εἰς Καφαρναούμ.\FnParseFormGloss{d}{accusative feminine singular}{Καφαρναούμ}{Capernaum} καὶ σκοτία\FnParseFormGloss{e}{nominative feminine singular}{σκοτία}{darkness} ἤδη ἐγεγόνει\FnParse{f}{pluperfect active indicative 3rd singular} καὶ ⸀οὔπω\FnFormGloss{g}{οὔπω}{not yet} ἐληλύθει\FnParse{h}{pluperfect active indicative 3rd singular} πρὸς αὐτοὺς ὁ Ἰησοῦς, 
\MainTextVerseMark{6}{18}ἥ τε θάλασσα ἀνέμου μεγάλου πνέοντος\FnParseFormGloss{a}{present active participle genitive masculine singular}{πνέω}{to blow} ⸀διεγείρετο.\FnParseFormGloss{b}{imperfect passive indicative 3rd singular}{διεγείρω}{to get up} 
\MainTextVerseMark{6}{19}ἐληλακότες\FnParseFormGloss{a}{perfect active participle nominative masculine plural}{ἐλαύνω}{to row} οὖν ὡς σταδίους\FnParseFormGloss{b}{accusative masculine plural}{στάδιος}{arena} εἴκοσι\FnParseFormGloss{c}{accusative masculine plural}{εἴκοσι(ν)}{twenty} πέντε ἢ τριάκοντα\FnParseFormGloss{d}{accusative masculine plural}{τριάκοντα}{thirty} θεωροῦσιν\FnParse{e}{present active indicative 3rd plural} τὸν Ἰησοῦν περιπατοῦντα\FnParse{f}{present active participle accusative masculine singular} ἐπὶ τῆς θαλάσσης καὶ ἐγγὺς τοῦ πλοίου γινόμενον,\FnParse{g}{present middle participle accusative masculine singular} καὶ ἐφοβήθησαν.\FnParse{h}{aorist passive indicative 3rd plural} 
\MainTextVerseMark{6}{20}ὁ δὲ λέγει\FnParse{a}{present active indicative 3rd singular} αὐτοῖς· Ἐγώ εἰμι,\FnParse{b}{present active indicative 1st singular} μὴ φοβεῖσθε.\FnParse{c}{present middle imperative 2nd plural} 
\MainTextVerseMark{6}{21}ἤθελον\FnParse{a}{imperfect active indicative 3rd plural} οὖν λαβεῖν\FnParse{b}{aorist active infinitive} αὐτὸν εἰς τὸ πλοῖον, καὶ εὐθέως ⸂ἐγένετο\FnParse{c}{aorist middle indicative 3rd singular} τὸ πλοῖον⸃ ἐπὶ τῆς γῆς εἰς ἣν ὑπῆγον.\FnParse{d}{imperfect active indicative 3rd plural} 
\MainTextVerseMark{6}{22}Τῇ ἐπαύριον\FnFormGloss{a}{ἐπαύριον}{the next day} ὁ ὄχλος ὁ ἑστηκὼς\FnParse{b}{perfect active participle nominative masculine singular} πέραν\FnFormGloss{c}{πέραν}{on the other side} τῆς θαλάσσης ⸀εἶδον\FnParse{d}{aorist active indicative 3rd plural} ὅτι πλοιάριον\FnParseFormGloss{e}{nominative neuter singular}{πλοιάριον}{boat} ἄλλο οὐκ ἦν\FnParse{f}{imperfect active indicative 3rd singular} ἐκεῖ εἰ μὴ ⸀ἕν, καὶ ὅτι οὐ συνεισῆλθεν\FnParseFormGloss{g}{aorist active indicative 3rd singular}{συνεισέρχομαι}{to enter together with} τοῖς μαθηταῖς αὐτοῦ ὁ Ἰησοῦς εἰς τὸ ⸀πλοῖον ἀλλὰ μόνοι οἱ μαθηταὶ αὐτοῦ ἀπῆλθον·\FnParse{h}{aorist active indicative 3rd plural} 
\MainTextVerseMark{6}{23}⸀ἀλλὰ ἦλθεν\FnParse{a}{aorist active indicative 3rd singular} ⸀πλοιάρια\FnParseFormGloss{b}{nominative neuter plural}{πλοιάριον}{boat} ἐκ Τιβεριάδος\FnParseFormGloss{c}{genitive feminine singular}{Τιβεριάς}{Tiberias} ἐγγὺς τοῦ τόπου ὅπου ἔφαγον\FnParse{d}{aorist active indicative 3rd plural} τὸν ἄρτον εὐχαριστήσαντος\FnParse{e}{aorist active participle genitive masculine singular} τοῦ κυρίου. 
\MainTextVerseMark{6}{24}ὅτε οὖν εἶδεν\FnParse{a}{aorist active indicative 3rd singular} ὁ ὄχλος ὅτι Ἰησοῦς οὐκ ἔστιν\FnParse{b}{present active indicative 3rd singular} ἐκεῖ οὐδὲ οἱ μαθηταὶ αὐτοῦ, ἐνέβησαν\FnParseFormGloss{c}{aorist active indicative 3rd plural}{ἐμβαίνω}{to get into} αὐτοὶ εἰς τὰ ⸀πλοιάρια\FnParseFormGloss{d}{accusative neuter plural}{πλοιάριον}{boat} καὶ ἦλθον\FnParse{e}{aorist active indicative 3rd plural} εἰς Καφαρναοὺμ\FnParseFormGloss{f}{accusative feminine singular}{Καφαρναούμ}{Capernaum} ζητοῦντες\FnParse{g}{present active participle nominative masculine plural} τὸν Ἰησοῦν. 
\MainTextVerseMark{6}{25}Καὶ εὑρόντες\FnParse{a}{aorist active participle nominative masculine plural} αὐτὸν πέραν\FnFormGloss{b}{πέραν}{on the other side} τῆς θαλάσσης εἶπον\FnParse{c}{aorist active indicative 3rd plural} αὐτῷ· Ῥαββί,\FnParseFormGloss{d}{masculine singular}{ῥαββί}{Rabbi} πότε\FnFormGloss{e}{πότε}{when? how long?} ὧδε γέγονας;\FnParse{f}{perfect active indicative 2nd singular} 
\MainTextVerseMark{6}{26}ἀπεκρίθη\FnParse{a}{aorist passive indicative 3rd singular} αὐτοῖς ὁ Ἰησοῦς καὶ εἶπεν·\FnParse{b}{aorist active indicative 3rd singular} Ἀμὴν ἀμὴν λέγω\FnParse{c}{present active indicative 1st singular} ὑμῖν, ζητεῖτέ\FnParse{d}{present active indicative 2nd plural} με οὐχ ὅτι εἴδετε\FnParse{e}{aorist active indicative 2nd plural} σημεῖα ἀλλ’ ὅτι ἐφάγετε\FnParse{f}{aorist active indicative 2nd plural} ἐκ τῶν ἄρτων καὶ ἐχορτάσθητε·\FnParseFormGloss{g}{aorist passive indicative 2nd plural}{χορτάζω}{to feed} 
\MainTextVerseMark{6}{27}ἐργάζεσθε\FnParse{a}{present middle imperative 2nd plural} μὴ τὴν βρῶσιν\FnParseFormGloss{b}{accusative feminine singular}{βρῶσις}{consumable} τὴν ἀπολλυμένην\FnParse{c}{present middle participle accusative feminine singular} ἀλλὰ τὴν βρῶσιν\FnParseFormGloss{d}{accusative feminine singular}{βρῶσις}{consumable} τὴν μένουσαν\FnParse{e}{present active participle accusative feminine singular} εἰς ζωὴν αἰώνιον, ἣν ὁ υἱὸς τοῦ ἀνθρώπου ὑμῖν δώσει,\FnParse{f}{future active indicative 3rd singular} τοῦτον γὰρ ὁ πατὴρ ἐσφράγισεν\FnParseFormGloss{g}{aorist active indicative 3rd singular}{σφραγίζω}{to seal} ὁ θεός. 
\MainTextVerseMark{6}{28}εἶπον\FnParse{a}{aorist active indicative 3rd plural} οὖν πρὸς αὐτόν· Τί ποιῶμεν\FnParse{b}{present active subjunctive 1st plural} ἵνα ἐργαζώμεθα\FnParse{c}{present middle subjunctive 1st plural} τὰ ἔργα τοῦ θεοῦ; 
\MainTextVerseMark{6}{29}ἀπεκρίθη\FnParse{a}{aorist passive indicative 3rd singular} ⸀ὁ Ἰησοῦς καὶ εἶπεν\FnParse{b}{aorist active indicative 3rd singular} αὐτοῖς· Τοῦτό ἐστιν\FnParse{c}{present active indicative 3rd singular} τὸ ἔργον τοῦ θεοῦ ἵνα ⸀πιστεύητε\FnParse{d}{present active subjunctive 2nd plural} εἰς ὃν ἀπέστειλεν\FnParse{e}{aorist active indicative 3rd singular} ἐκεῖνος. 
\MainTextVerseMark{6}{30}εἶπον\FnParse{a}{aorist active indicative 3rd plural} οὖν αὐτῷ· Τί οὖν ποιεῖς\FnParse{b}{present active indicative 2nd singular} σὺ σημεῖον, ἵνα ἴδωμεν\FnParse{c}{aorist active subjunctive 1st plural} καὶ πιστεύσωμέν\FnParse{d}{aorist active subjunctive 1st plural} σοι; τί ἐργάζῃ;\FnParse{e}{present middle indicative 2nd singular} 
\MainTextVerseMark{6}{31}οἱ πατέρες ἡμῶν τὸ μάννα\FnParseFormGloss{a}{accusative neuter singular}{μάννα}{manna} ἔφαγον\FnParse{b}{aorist active indicative 3rd plural} ἐν τῇ ἐρήμῳ, καθώς ἐστιν\FnParse{c}{present active indicative 3rd singular} γεγραμμένον·\FnParse{d}{perfect passive participle nominative neuter singular} Ἄρτον ἐκ τοῦ οὐρανοῦ ἔδωκεν\FnParse{e}{aorist active indicative 3rd singular} αὐτοῖς φαγεῖν.\FnParse{f}{aorist active infinitive} 
\MainTextVerseMark{6}{32}εἶπεν\FnParse{a}{aorist active indicative 3rd singular} οὖν αὐτοῖς ὁ Ἰησοῦς· Ἀμὴν ἀμὴν λέγω\FnParse{b}{present active indicative 1st singular} ὑμῖν, οὐ Μωϋσῆς ⸀δέδωκεν\FnParse{c}{perfect active indicative 3rd singular} ὑμῖν τὸν ἄρτον ἐκ τοῦ οὐρανοῦ, ἀλλ’ ὁ πατήρ μου δίδωσιν\FnParse{d}{present active indicative 3rd singular} ὑμῖν τὸν ἄρτον ἐκ τοῦ οὐρανοῦ τὸν ἀληθινόν·\FnParseFormGloss{e}{accusative masculine singular}{ἀληθινός}{true} 
\MainTextVerseMark{6}{33}ὁ γὰρ ἄρτος τοῦ θεοῦ ἐστιν\FnParse{a}{present active indicative 3rd singular} ὁ καταβαίνων\FnParse{b}{present active participle nominative masculine singular} ἐκ τοῦ οὐρανοῦ καὶ ζωὴν διδοὺς\FnParse{c}{present active participle nominative masculine singular} τῷ κόσμῳ. 
\MainTextVerseMark{6}{34}εἶπον\FnParse{a}{aorist active indicative 3rd plural} οὖν πρὸς αὐτόν· Κύριε, πάντοτε δὸς\FnParse{b}{aorist active imperative 2nd singular} ἡμῖν τὸν ἄρτον τοῦτον. 
\MainTextVerseMark{6}{35}⸀Εἶπεν\FnParse{a}{aorist active indicative 3rd singular} αὐτοῖς ὁ Ἰησοῦς· Ἐγώ εἰμι\FnParse{b}{present active indicative 1st singular} ὁ ἄρτος τῆς ζωῆς· ὁ ἐρχόμενος\FnParse{c}{present middle participle nominative masculine singular} πρὸς ἐμὲ οὐ μὴ πεινάσῃ,\FnParseFormGloss{d}{aorist active subjunctive 3rd singular}{πεινάω}{to be hungry} καὶ ὁ πιστεύων\FnParse{e}{present active participle nominative masculine singular} εἰς ⸀ἐμὲ οὐ μὴ ⸀διψήσει\FnParseFormGloss{f}{future active indicative 3rd singular}{διψάω}{to be thirsty} πώποτε.\FnFormGloss{g}{πώποτε}{ever} 
\MainTextVerseMark{6}{36}ἀλλ’ εἶπον\FnParse{a}{aorist active indicative 1st singular} ὑμῖν ὅτι καὶ ἑωράκατέ\FnParse{b}{perfect active indicative 2nd plural} με καὶ οὐ πιστεύετε.\FnParse{c}{present active indicative 2nd plural} 
\MainTextVerseMark{6}{37}πᾶν ὃ δίδωσίν\FnParse{a}{present active indicative 3rd singular} μοι ὁ πατὴρ πρὸς ἐμὲ ἥξει,\FnParseFormGloss{b}{future active indicative 3rd singular}{ἥκω}{to come} καὶ τὸν ἐρχόμενον\FnParse{c}{present middle participle accusative masculine singular} πρός ⸀με οὐ μὴ ἐκβάλω\FnParse{d}{aorist active subjunctive 1st singular} ἔξω, 
\MainTextVerseMark{6}{38}ὅτι καταβέβηκα\FnParse{a}{perfect active indicative 1st singular} ⸀ἀπὸ τοῦ οὐρανοῦ οὐχ ἵνα ποιῶ\FnParse{b}{present active subjunctive 1st singular} τὸ θέλημα τὸ ἐμὸν ἀλλὰ τὸ θέλημα τοῦ πέμψαντός\FnParse{c}{aorist active participle genitive masculine singular} με· 
\MainTextVerseMark{6}{39}τοῦτο δέ ἐστιν\FnParse{a}{present active indicative 3rd singular} τὸ θέλημα τοῦ πέμψαντός\FnParse{b}{aorist active participle genitive masculine singular} ⸀με ἵνα πᾶν ὃ δέδωκέν\FnParse{c}{perfect active indicative 3rd singular} μοι μὴ ἀπολέσω\FnParse{d}{aorist active subjunctive 1st singular} ἐξ αὐτοῦ ἀλλὰ ἀναστήσω\FnParse{e}{aorist active subjunctive 1st singular} ⸀αὐτὸ τῇ ἐσχάτῃ ἡμέρᾳ. 
\MainTextVerseMark{6}{40}τοῦτο ⸀γάρ ἐστιν\FnParse{a}{present active indicative 3rd singular} τὸ θέλημα τοῦ ⸂πατρός μου⸃ ἵνα πᾶς ὁ θεωρῶν\FnParse{b}{present active participle nominative masculine singular} τὸν υἱὸν καὶ πιστεύων\FnParse{c}{present active participle nominative masculine singular} εἰς αὐτὸν ἔχῃ\FnParse{d}{present active subjunctive 3rd singular} ζωὴν αἰώνιον, καὶ ἀναστήσω\FnParse{e}{future active indicative 1st singular} αὐτὸν ⸀ἐγὼ τῇ ἐσχάτῃ ἡμέρᾳ. 
\MainTextVerseMark{6}{41}Ἐγόγγυζον\FnParseFormGloss{a}{imperfect active indicative 3rd plural}{γογγύζω}{to grumble} οὖν οἱ Ἰουδαῖοι περὶ αὐτοῦ ὅτι εἶπεν·\FnParse{b}{aorist active indicative 3rd singular} Ἐγώ εἰμι\FnParse{c}{present active indicative 1st singular} ὁ ἄρτος ὁ καταβὰς\FnParse{d}{aorist active participle nominative masculine singular} ἐκ τοῦ οὐρανοῦ, 
\MainTextVerseMark{6}{42}καὶ ἔλεγον·\FnParse{a}{imperfect active indicative 3rd plural} ⸀Οὐχ οὗτός ἐστιν\FnParse{b}{present active indicative 3rd singular} Ἰησοῦς ὁ υἱὸς Ἰωσήφ, οὗ ἡμεῖς οἴδαμεν\FnParse{c}{perfect active indicative 1st plural} τὸν πατέρα καὶ τὴν μητέρα; πῶς ⸀νῦν ⸀λέγει\FnParse{d}{present active indicative 3rd singular} ὅτι Ἐκ τοῦ οὐρανοῦ καταβέβηκα;\FnParse{e}{perfect active indicative 1st singular} 
\MainTextVerseMark{6}{43}⸀ἀπεκρίθη\FnParse{a}{aorist passive indicative 3rd singular} Ἰησοῦς καὶ εἶπεν\FnParse{b}{aorist active indicative 3rd singular} αὐτοῖς· Μὴ γογγύζετε\FnParseFormGloss{c}{present active imperative 2nd plural}{γογγύζω}{to grumble} μετ’ ἀλλήλων. 
\MainTextVerseMark{6}{44}οὐδεὶς δύναται\FnParse{a}{present middle indicative 3rd singular} ἐλθεῖν\FnParse{b}{aorist active infinitive} πρός ⸀με ἐὰν μὴ ὁ πατὴρ ὁ πέμψας\FnParse{c}{aorist active participle nominative masculine singular} με ἑλκύσῃ\FnParseFormGloss{d}{aorist active subjunctive 3rd singular}{ἕλκω}{to drag} αὐτόν, κἀγὼ ἀναστήσω\FnParse{e}{future active indicative 1st singular} αὐτὸν ἐν τῇ ἐσχάτῃ ἡμέρᾳ. 
\MainTextVerseMark{6}{45}ἔστιν\FnParse{a}{present active indicative 3rd singular} γεγραμμένον\FnParse{b}{perfect passive participle nominative neuter singular} ἐν τοῖς προφήταις· Καὶ ἔσονται\FnParse{c}{future middle indicative 3rd plural} πάντες διδακτοὶ\FnParseFormGloss{d}{nominative masculine plural}{διδακτός}{taught} θεοῦ· ⸀πᾶς ὁ ⸀ἀκούσας\FnParse{e}{aorist active participle nominative masculine singular} παρὰ τοῦ πατρὸς καὶ μαθὼν\FnParseFormGloss{f}{aorist active participle nominative masculine singular}{μανθάνω}{to learn} ἔρχεται\FnParse{g}{present middle indicative 3rd singular} πρὸς ⸀ἐμέ. 
\MainTextVerseMark{6}{46}οὐχ ὅτι τὸν πατέρα ⸂ἑώρακέν\FnParse{a}{perfect active indicative 3rd singular} τις⸃ εἰ μὴ ὁ ὢν\FnParse{b}{present active participle nominative masculine singular} παρὰ τοῦ θεοῦ, οὗτος ἑώρακεν\FnParse{c}{perfect active indicative 3rd singular} τὸν πατέρα. 
\MainTextVerseMark{6}{47}ἀμὴν ἀμὴν λέγω\FnParse{a}{present active indicative 1st singular} ὑμῖν, ὁ ⸀πιστεύων\FnParse{b}{present active participle nominative masculine singular} ἔχει\FnParse{c}{present active indicative 3rd singular} ζωὴν αἰώνιον. 
\MainTextVerseMark{6}{48}ἐγώ εἰμι\FnParse{a}{present active indicative 1st singular} ὁ ἄρτος τῆς ζωῆς· 
\MainTextVerseMark{6}{49}οἱ πατέρες ὑμῶν ἔφαγον\FnParse{a}{aorist active indicative 3rd plural} ⸂ἐν τῇ ἐρήμῳ τὸ μάννα⸃\FnParseFormGloss{b}{accusative neuter singular}{μάννα}{manna} καὶ ἀπέθανον·\FnParse{c}{aorist active indicative 3rd plural} 
\MainTextVerseMark{6}{50}οὗτός ἐστιν\FnParse{a}{present active indicative 3rd singular} ὁ ἄρτος ὁ ἐκ τοῦ οὐρανοῦ καταβαίνων\FnParse{b}{present active participle nominative masculine singular} ἵνα τις ἐξ αὐτοῦ φάγῃ\FnParse{c}{aorist active subjunctive 3rd singular} καὶ μὴ ἀποθάνῃ·\FnParse{d}{aorist active subjunctive 3rd singular} 
\MainTextVerseMark{6}{51}ἐγώ εἰμι\FnParse{a}{present active indicative 1st singular} ὁ ἄρτος ὁ ζῶν\FnParse{b}{present active participle nominative masculine singular} ὁ ἐκ τοῦ οὐρανοῦ καταβάς·\FnParse{c}{aorist active participle nominative masculine singular} ἐάν τις φάγῃ\FnParse{d}{aorist active subjunctive 3rd singular} ἐκ τούτου τοῦ ἄρτου ⸀ζήσει\FnParse{e}{future active indicative 3rd singular} εἰς τὸν αἰῶνα, καὶ ὁ ἄρτος δὲ ὃν ἐγὼ δώσω\FnParse{f}{future active indicative 1st singular} ἡ σάρξ μού ⸀ἐστιν\FnParse{g}{present active indicative 3rd singular} ὑπὲρ τῆς τοῦ κόσμου ζωῆς. 
\MainTextVerseMark{6}{52}Ἐμάχοντο\FnParseFormGloss{a}{imperfect middle indicative 3rd plural}{μάχομαι}{to fight} οὖν πρὸς ἀλλήλους οἱ Ἰουδαῖοι λέγοντες·\FnParse{b}{present active participle nominative masculine plural} Πῶς δύναται\FnParse{c}{present middle indicative 3rd singular} οὗτος ἡμῖν δοῦναι\FnParse{d}{aorist active infinitive} τὴν σάρκα ⸀αὐτοῦ φαγεῖν;\FnParse{e}{aorist active infinitive} 
\MainTextVerseMark{6}{53}εἶπεν\FnParse{a}{aorist active indicative 3rd singular} οὖν αὐτοῖς ὁ Ἰησοῦς· Ἀμὴν ἀμὴν λέγω\FnParse{b}{present active indicative 1st singular} ὑμῖν, ἐὰν μὴ φάγητε\FnParse{c}{aorist active subjunctive 2nd plural} τὴν σάρκα τοῦ υἱοῦ τοῦ ἀνθρώπου καὶ πίητε\FnParse{d}{aorist active subjunctive 2nd plural} αὐτοῦ τὸ αἷμα, οὐκ ἔχετε\FnParse{e}{present active indicative 2nd plural} ζωὴν ἐν ἑαυτοῖς. 
\MainTextVerseMark{6}{54}ὁ τρώγων\FnParseFormGloss{a}{present active participle nominative masculine singular}{τρώγω}{to eat} μου τὴν σάρκα καὶ πίνων\FnParse{b}{present active participle nominative masculine singular} μου τὸ αἷμα ἔχει\FnParse{c}{present active indicative 3rd singular} ζωὴν αἰώνιον, κἀγὼ ἀναστήσω\FnParse{d}{future active indicative 1st singular} αὐτὸν τῇ ἐσχάτῃ ἡμέρᾳ. 
\MainTextVerseMark{6}{55}ἡ γὰρ σάρξ μου ⸀ἀληθής\FnParseFormGloss{a}{nominative feminine singular}{ἀληθής}{true} ἐστι\FnParse{b}{present active indicative 3rd singular} βρῶσις,\FnParseFormGloss{c}{nominative feminine singular}{βρῶσις}{consumable} καὶ τὸ αἷμά μου ⸁ἀληθής\FnParseFormGloss{d}{nominative feminine singular}{ἀληθής}{true} ἐστι\FnParse{e}{present active indicative 3rd singular} πόσις.\FnParseFormGloss{f}{nominative feminine singular}{πόσις}{drinking} 
\MainTextVerseMark{6}{56}ὁ τρώγων\FnParseFormGloss{a}{present active participle nominative masculine singular}{τρώγω}{to eat} μου τὴν σάρκα καὶ πίνων\FnParse{b}{present active participle nominative masculine singular} μου τὸ αἷμα ἐν ἐμοὶ μένει\FnParse{c}{present active indicative 3rd singular} κἀγὼ ἐν αὐτῷ. 
\MainTextVerseMark{6}{57}καθὼς ἀπέστειλέν\FnParse{a}{aorist active indicative 3rd singular} με ὁ ζῶν\FnParse{b}{present active participle nominative masculine singular} πατὴρ κἀγὼ ζῶ\FnParse{c}{present active indicative 1st singular} διὰ τὸν πατέρα, καὶ ὁ τρώγων\FnParseFormGloss{d}{present active participle nominative masculine singular}{τρώγω}{to eat} με κἀκεῖνος\FnParseFormGloss{e}{nominative masculine singular}{κἀκεῖνος}{and that one} ⸀ζήσει\FnParse{f}{future active indicative 3rd singular} δι’ ἐμέ. 
\MainTextVerseMark{6}{58}οὗτός ἐστιν\FnParse{a}{present active indicative 3rd singular} ὁ ἄρτος ὁ ⸀ἐξ οὐρανοῦ καταβάς,\FnParse{b}{aorist active participle nominative masculine singular} οὐ καθὼς ἔφαγον\FnParse{c}{aorist active indicative 3rd plural} οἱ ⸀πατέρες καὶ ἀπέθανον·\FnParse{d}{aorist active indicative 3rd plural} ὁ τρώγων\FnParseFormGloss{e}{present active participle nominative masculine singular}{τρώγω}{to eat} τοῦτον τὸν ἄρτον ⸀ζήσει\FnParse{f}{future active indicative 3rd singular} εἰς τὸν αἰῶνα. 
\MainTextVerseMark{6}{59}ταῦτα εἶπεν\FnParse{a}{aorist active indicative 3rd singular} ἐν συναγωγῇ διδάσκων\FnParse{b}{present active participle nominative masculine singular} ἐν Καφαρναούμ.\FnParseFormGloss{c}{dative feminine singular}{Καφαρναούμ}{Capernaum} 
\MainTextVerseMark{6}{60}Πολλοὶ οὖν ἀκούσαντες\FnParse{a}{aorist active participle nominative masculine plural} ἐκ τῶν μαθητῶν αὐτοῦ εἶπαν·\FnParse{b}{aorist active indicative 3rd plural} Σκληρός\FnParseFormGloss{c}{nominative masculine singular}{σκληρός}{hard} ἐστιν\FnParse{d}{present active indicative 3rd singular} ⸂ὁ λόγος οὗτος⸃· τίς δύναται\FnParse{e}{present middle indicative 3rd singular} αὐτοῦ ἀκούειν;\FnParse{f}{present active infinitive} 
\MainTextVerseMark{6}{61}εἰδὼς\FnParse{a}{perfect active participle nominative masculine singular} δὲ ὁ Ἰησοῦς ἐν ἑαυτῷ ὅτι γογγύζουσιν\FnParseFormGloss{b}{present active indicative 3rd plural}{γογγύζω}{to grumble} περὶ τούτου οἱ μαθηταὶ αὐτοῦ εἶπεν\FnParse{c}{aorist active indicative 3rd singular} αὐτοῖς· Τοῦτο ὑμᾶς σκανδαλίζει;\FnParse{d}{present active indicative 3rd singular} 
\MainTextVerseMark{6}{62}ἐὰν οὖν θεωρῆτε\FnParse{a}{present active subjunctive 2nd plural} τὸν υἱὸν τοῦ ἀνθρώπου ἀναβαίνοντα\FnParse{b}{present active participle accusative masculine singular} ὅπου ἦν\FnParse{c}{imperfect active indicative 3rd singular} τὸ πρότερον;\FnParseFormGloss{d}{accusative neuter singular}{πρότερος}{before} 
\MainTextVerseMark{6}{63}τὸ πνεῦμά ἐστιν\FnParse{a}{present active indicative 3rd singular} τὸ ζῳοποιοῦν,\FnParseFormGloss{b}{present active participle nominative neuter singular}{ζῳοποιέω}{to make alive} ἡ σὰρξ οὐκ ὠφελεῖ\FnParseFormGloss{c}{present active indicative 3rd singular}{ὠφελέω}{to be of good use} οὐδέν· τὰ ῥήματα ἃ ἐγὼ ⸀λελάληκα\FnParse{d}{perfect active indicative 1st singular} ὑμῖν πνεῦμά ἐστιν\FnParse{e}{present active indicative 3rd singular} καὶ ζωή ἐστιν.\FnParse{f}{present active indicative 3rd singular} 
\MainTextVerseMark{6}{64}ἀλλὰ εἰσὶν\FnParse{a}{present active indicative 3rd plural} ἐξ ὑμῶν τινες οἳ οὐ πιστεύουσιν.\FnParse{b}{present active indicative 3rd plural} ᾔδει\FnParse{c}{pluperfect active indicative 3rd singular} γὰρ ἐξ ἀρχῆς ὁ Ἰησοῦς τίνες εἰσὶν\FnParse{d}{present active indicative 3rd plural} οἱ μὴ πιστεύοντες\FnParse{e}{present active participle nominative masculine plural} καὶ τίς ἐστιν\FnParse{f}{present active indicative 3rd singular} ὁ παραδώσων\FnParse{g}{future active participle nominative masculine singular} αὐτόν. 
\MainTextVerseMark{6}{65}καὶ ἔλεγεν·\FnParse{a}{imperfect active indicative 3rd singular} Διὰ τοῦτο εἴρηκα\FnParse{b}{perfect active indicative 1st singular} ὑμῖν ὅτι οὐδεὶς δύναται\FnParse{c}{present middle indicative 3rd singular} ἐλθεῖν\FnParse{d}{aorist active infinitive} πρός με ἐὰν μὴ ᾖ\FnParse{e}{present active subjunctive 3rd singular} δεδομένον\FnParse{f}{perfect passive participle nominative neuter singular} αὐτῷ ἐκ τοῦ ⸀πατρός. 
\MainTextVerseMark{6}{66}Ἐκ τούτου πολλοὶ ⸀ἐκ ⸂τῶν μαθητῶν αὐτοῦ ἀπῆλθον⸃\FnParse{a}{aorist active indicative 3rd plural} εἰς τὰ ὀπίσω καὶ οὐκέτι μετ’ αὐτοῦ περιεπάτουν.\FnParse{b}{imperfect active indicative 3rd plural} 
\MainTextVerseMark{6}{67}εἶπεν\FnParse{a}{aorist active indicative 3rd singular} οὖν ὁ Ἰησοῦς τοῖς δώδεκα· Μὴ καὶ ὑμεῖς θέλετε\FnParse{b}{present active indicative 2nd plural} ὑπάγειν;\FnParse{c}{present active infinitive} 
\MainTextVerseMark{6}{68}⸀ἀπεκρίθη\FnParse{a}{aorist passive indicative 3rd singular} αὐτῷ Σίμων Πέτρος· Κύριε, πρὸς τίνα ἀπελευσόμεθα;\FnParse{b}{future middle indicative 1st plural} ῥήματα ζωῆς αἰωνίου ἔχεις,\FnParse{c}{present active indicative 2nd singular} 
\MainTextVerseMark{6}{69}καὶ ἡμεῖς πεπιστεύκαμεν\FnParse{a}{perfect active indicative 1st plural} καὶ ἐγνώκαμεν\FnParse{b}{perfect active indicative 1st plural} ὅτι σὺ εἶ\FnParse{c}{present active indicative 2nd singular} ὁ ⸀ἅγιος τοῦ ⸀θεοῦ. 
\MainTextVerseMark{6}{70}ἀπεκρίθη\FnParse{a}{aorist passive indicative 3rd singular} αὐτοῖς ὁ Ἰησοῦς· Οὐκ ἐγὼ ὑμᾶς τοὺς δώδεκα ἐξελεξάμην;\FnParseFormGloss{b}{aorist middle indicative 1st singular}{ἐκλέγομαι}{to chose} καὶ ἐξ ὑμῶν εἷς διάβολός ἐστιν.\FnParse{c}{present active indicative 3rd singular} 
\MainTextVerseMark{6}{71}ἔλεγεν\FnParse{a}{imperfect active indicative 3rd singular} δὲ τὸν Ἰούδαν Σίμωνος ⸀Ἰσκαριώτου·\FnParseFormGloss{b}{genitive masculine singular}{Ἰσκαριώτης}{Iscariot} οὗτος γὰρ ἔμελλεν\FnParse{c}{imperfect active indicative 3rd singular} ⸂παραδιδόναι\FnParse{d}{present active infinitive} αὐτόν⸃, ⸀εἷς ἐκ τῶν δώδεκα. 
\ChapterHeader{7}{Chapter 7}

\MainTextVerseMark{7}{1}Καὶ ⸂μετὰ ταῦτα περιεπάτει\FnParse{a}{imperfect active indicative 3rd singular} ὁ Ἰησοῦς⸃ ἐν τῇ Γαλιλαίᾳ, οὐ γὰρ ἤθελεν\FnParse{b}{imperfect active indicative 3rd singular} ἐν τῇ Ἰουδαίᾳ περιπατεῖν,\FnParse{c}{present active infinitive} ὅτι ἐζήτουν\FnParse{d}{imperfect active indicative 3rd plural} αὐτὸν οἱ Ἰουδαῖοι ἀποκτεῖναι.\FnParse{e}{aorist active infinitive} 
\MainTextVerseMark{7}{2}ἦν\FnParse{a}{imperfect active indicative 3rd singular} δὲ ἐγγὺς ἡ ἑορτὴ\FnParseFormGloss{b}{nominative feminine singular}{ἑορτή}{feast} τῶν Ἰουδαίων ἡ σκηνοπηγία.\FnParseFormGloss{c}{nominative feminine singular}{σκηνοπηγία}{Tabernacles} 
\MainTextVerseMark{7}{3}εἶπον\FnParse{a}{aorist active indicative 3rd plural} οὖν πρὸς αὐτὸν οἱ ἀδελφοὶ αὐτοῦ· Μετάβηθι\FnParseFormGloss{b}{aorist active imperative 2nd singular}{μεταβαίνω}{to go on} ἐντεῦθεν\FnFormGloss{c}{ἐντεῦθεν}{from here} καὶ ὕπαγε\FnParse{d}{present active imperative 2nd singular} εἰς τὴν Ἰουδαίαν, ἵνα καὶ οἱ μαθηταί σου ⸀θεωρήσουσιν\FnParse{e}{future active indicative 3rd plural} ⸂σοῦ τὰ ἔργα⸃ ἃ ποιεῖς·\FnParse{f}{present active indicative 2nd singular} 
\MainTextVerseMark{7}{4}οὐδεὶς γάρ ⸂τι ἐν κρυπτῷ⸃\FnParseFormGloss{a}{dative neuter singular}{κρυπτός}{hidden} ποιεῖ\FnParse{b}{present active indicative 3rd singular} καὶ ζητεῖ\FnParse{c}{present active indicative 3rd singular} αὐτὸς ἐν παρρησίᾳ εἶναι·\FnParse{d}{present active infinitive} εἰ ταῦτα ποιεῖς,\FnParse{e}{present active indicative 2nd singular} φανέρωσον\FnParse{f}{aorist active imperative 2nd singular} σεαυτὸν τῷ κόσμῳ. 
\MainTextVerseMark{7}{5}οὐδὲ γὰρ οἱ ἀδελφοὶ αὐτοῦ ἐπίστευον\FnParse{a}{imperfect active indicative 3rd plural} εἰς αὐτόν. 
\MainTextVerseMark{7}{6}λέγει\FnParse{a}{present active indicative 3rd singular} οὖν αὐτοῖς ὁ Ἰησοῦς· Ὁ καιρὸς ὁ ἐμὸς οὔπω\FnFormGloss{b}{οὔπω}{not yet} πάρεστιν,\FnParseFormGloss{c}{present active indicative 3rd singular}{πάρειμι}{to be present} ὁ δὲ καιρὸς ὁ ὑμέτερος\FnParseFormGloss{d}{nominative masculine singular}{ὑμέτερος}{your} πάντοτέ ἐστιν\FnParse{e}{present active indicative 3rd singular} ἕτοιμος.\FnParseFormGloss{f}{nominative masculine singular}{ἕτοιμος}{ready} 
\MainTextVerseMark{7}{7}οὐ δύναται\FnParse{a}{present middle indicative 3rd singular} ὁ κόσμος μισεῖν\FnParse{b}{present active infinitive} ὑμᾶς, ἐμὲ δὲ μισεῖ,\FnParse{c}{present active indicative 3rd singular} ὅτι ἐγὼ μαρτυρῶ\FnParse{d}{present active indicative 1st singular} περὶ αὐτοῦ ὅτι τὰ ἔργα αὐτοῦ πονηρά ἐστιν.\FnParse{e}{present active indicative 3rd singular} 
\MainTextVerseMark{7}{8}ὑμεῖς ἀνάβητε\FnParse{a}{aorist active imperative 2nd plural} εἰς τὴν ⸀ἑορτήν·\FnParseFormGloss{b}{accusative feminine singular}{ἑορτή}{feast} ἐγὼ ⸀οὐκ ἀναβαίνω\FnParse{c}{present active indicative 1st singular} εἰς τὴν ἑορτὴν\FnParseFormGloss{d}{accusative feminine singular}{ἑορτή}{feast} ταύτην, ὅτι ὁ ⸂ἐμὸς καιρὸς⸃ οὔπω\FnFormGloss{e}{οὔπω}{not yet} πεπλήρωται.\FnParse{f}{perfect passive indicative 3rd singular} 
\MainTextVerseMark{7}{9}ταῦτα ⸀δὲ εἰπὼν\FnParse{a}{aorist active participle nominative masculine singular} ⸀αὐτὸς ἔμεινεν\FnParse{b}{aorist active indicative 3rd singular} ἐν τῇ Γαλιλαίᾳ. 
\MainTextVerseMark{7}{10}Ὡς δὲ ἀνέβησαν\FnParse{a}{aorist active indicative 3rd plural} οἱ ἀδελφοὶ αὐτοῦ ⸂εἰς τὴν ἑορτήν,\FnParseFormGloss{b}{accusative feminine singular}{ἑορτή}{feast} τότε καὶ αὐτὸς ἀνέβη⸃,\FnParse{c}{aorist active indicative 3rd singular} οὐ φανερῶς\FnFormGloss{d}{φανερῶς}{manifestly} ἀλλὰ ὡς ἐν κρυπτῷ.\FnParseFormGloss{e}{dative neuter singular}{κρυπτός}{hidden} 
\MainTextVerseMark{7}{11}οἱ οὖν Ἰουδαῖοι ἐζήτουν\FnParse{a}{imperfect active indicative 3rd plural} αὐτὸν ἐν τῇ ἑορτῇ\FnParseFormGloss{b}{dative feminine singular}{ἑορτή}{feast} καὶ ἔλεγον·\FnParse{c}{imperfect active indicative 3rd plural} Ποῦ ἐστιν\FnParse{d}{present active indicative 3rd singular} ἐκεῖνος; 
\MainTextVerseMark{7}{12}καὶ γογγυσμὸς\FnParseFormGloss{a}{nominative masculine singular}{γογγυσμός}{complaint} ⸂περὶ αὐτοῦ ἦν\FnParse{b}{imperfect active indicative 3rd singular} πολὺς⸃ ἐν τοῖς ὄχλοις· οἱ μὲν ἔλεγον\FnParse{c}{imperfect active indicative 3rd plural} ὅτι Ἀγαθός ἐστιν,\FnParse{d}{present active indicative 3rd singular} ἄλλοι ⸀δὲ ἔλεγον·\FnParse{e}{imperfect active indicative 3rd plural} Οὔ,\FnFormGloss{f}{οὔ}{no} ἀλλὰ πλανᾷ\FnParse{g}{present active indicative 3rd singular} τὸν ὄχλον. 
\MainTextVerseMark{7}{13}οὐδεὶς μέντοι\FnFormGloss{a}{μέντοι}{but} παρρησίᾳ ἐλάλει\FnParse{b}{imperfect active indicative 3rd singular} περὶ αὐτοῦ διὰ τὸν φόβον τῶν Ἰουδαίων. 
\MainTextVerseMark{7}{14}Ἤδη δὲ τῆς ἑορτῆς\FnParseFormGloss{a}{genitive feminine singular}{ἑορτή}{feast} μεσούσης\FnParseFormGloss{b}{present active participle genitive feminine singular}{μεσόω}{to be halfway through} ⸀ἀνέβη\FnParse{c}{aorist active indicative 3rd singular} Ἰησοῦς εἰς τὸ ἱερὸν καὶ ἐδίδασκεν.\FnParse{d}{imperfect active indicative 3rd singular} 
\MainTextVerseMark{7}{15}⸂ἐθαύμαζον\FnParse{a}{imperfect active indicative 3rd plural} οὖν⸃ οἱ Ἰουδαῖοι λέγοντες·\FnParse{b}{present active participle nominative masculine plural} Πῶς οὗτος γράμματα\FnParseFormGloss{c}{accusative neuter plural}{γράμμα}{letter} οἶδεν\FnParse{d}{perfect active indicative 3rd singular} μὴ μεμαθηκώς;\FnParseFormGloss{e}{perfect active participle nominative masculine singular}{μανθάνω}{to learn} 
\MainTextVerseMark{7}{16}ἀπεκρίθη\FnParse{a}{aorist passive indicative 3rd singular} οὖν αὐτοῖς ⸀ὁ Ἰησοῦς καὶ εἶπεν·\FnParse{b}{aorist active indicative 3rd singular} Ἡ ἐμὴ διδαχὴ οὐκ ἔστιν\FnParse{c}{present active indicative 3rd singular} ἐμὴ ἀλλὰ τοῦ πέμψαντός\FnParse{d}{aorist active participle genitive masculine singular} με· 
\MainTextVerseMark{7}{17}ἐάν τις θέλῃ\FnParse{a}{present active subjunctive 3rd singular} τὸ θέλημα αὐτοῦ ποιεῖν,\FnParse{b}{present active infinitive} γνώσεται\FnParse{c}{future middle indicative 3rd singular} περὶ τῆς διδαχῆς πότερον\FnFormGloss{d}{πότερον}{whether} ἐκ τοῦ θεοῦ ἐστιν\FnParse{e}{present active indicative 3rd singular} ἢ ἐγὼ ἀπ’ ἐμαυτοῦ λαλῶ.\FnParse{f}{present active indicative 1st singular} 
\MainTextVerseMark{7}{18}ὁ ἀφ’ ἑαυτοῦ λαλῶν\FnParse{a}{present active participle nominative masculine singular} τὴν δόξαν τὴν ἰδίαν ζητεῖ·\FnParse{b}{present active indicative 3rd singular} ὁ δὲ ζητῶν\FnParse{c}{present active participle nominative masculine singular} τὴν δόξαν τοῦ πέμψαντος\FnParse{d}{aorist active participle genitive masculine singular} αὐτὸν οὗτος ἀληθής\FnParseFormGloss{e}{nominative masculine singular}{ἀληθής}{true} ἐστιν\FnParse{f}{present active indicative 3rd singular} καὶ ἀδικία\FnParseFormGloss{g}{nominative feminine singular}{ἀδικία}{wickedness} ἐν αὐτῷ οὐκ ἔστιν.\FnParse{h}{present active indicative 3rd singular} 
\MainTextVerseMark{7}{19}Οὐ Μωϋσῆς ⸀δέδωκεν\FnParse{a}{perfect active indicative 3rd singular} ὑμῖν τὸν νόμον; καὶ οὐδεὶς ἐξ ὑμῶν ποιεῖ\FnParse{b}{present active indicative 3rd singular} τὸν νόμον. τί με ζητεῖτε\FnParse{c}{present active indicative 2nd plural} ἀποκτεῖναι;\FnParse{d}{aorist active infinitive} 
\MainTextVerseMark{7}{20}ἀπεκρίθη\FnParse{a}{aorist passive indicative 3rd singular} ὁ ⸀ὄχλος· Δαιμόνιον ἔχεις·\FnParse{b}{present active indicative 2nd singular} τίς σε ζητεῖ\FnParse{c}{present active indicative 3rd singular} ἀποκτεῖναι;\FnParse{d}{aorist active infinitive} 
\MainTextVerseMark{7}{21}ἀπεκρίθη\FnParse{a}{aorist passive indicative 3rd singular} Ἰησοῦς καὶ εἶπεν\FnParse{b}{aorist active indicative 3rd singular} αὐτοῖς· Ἓν ἔργον ἐποίησα\FnParse{c}{aorist active indicative 1st singular} καὶ πάντες θαυμάζετε.\FnParse{d}{present active indicative 2nd plural} 
\MainTextVerseMark{7}{22}διὰ τοῦτο Μωϋσῆς δέδωκεν\FnParse{a}{perfect active indicative 3rd singular} ὑμῖν τὴν περιτομήν— οὐχ ὅτι ἐκ τοῦ Μωϋσέως ἐστὶν\FnParse{b}{present active indicative 3rd singular} ἀλλ’ ἐκ τῶν πατέρων— καὶ ἐν σαββάτῳ περιτέμνετε\FnParseFormGloss{c}{present active indicative 2nd plural}{περιτέμνω}{to circumcise} ἄνθρωπον. 
\MainTextVerseMark{7}{23}εἰ περιτομὴν ⸀λαμβάνει\FnParse{a}{present active indicative 3rd singular} ἄνθρωπος ἐν σαββάτῳ ἵνα μὴ λυθῇ\FnParse{b}{aorist passive subjunctive 3rd singular} ὁ νόμος Μωϋσέως, ἐμοὶ χολᾶτε\FnParseFormGloss{c}{present active indicative 2nd plural}{χολάω}{to be angry} ὅτι ὅλον ἄνθρωπον ὑγιῆ\FnParseFormGloss{d}{accusative masculine singular}{ὑγιής}{healthy} ἐποίησα\FnParse{e}{aorist active indicative 1st singular} ἐν σαββάτῳ; 
\MainTextVerseMark{7}{24}μὴ κρίνετε\FnParse{a}{present active imperative 2nd plural} κατ’ ὄψιν,\FnParseFormGloss{b}{accusative feminine singular}{ὄψις}{face} ἀλλὰ τὴν δικαίαν κρίσιν ⸀κρίνετε.\FnParse{c}{present active imperative 2nd plural} 
\MainTextVerseMark{7}{25}Ἔλεγον\FnParse{a}{imperfect active indicative 3rd plural} οὖν τινες ἐκ τῶν Ἱεροσολυμιτῶν·\FnParseFormGloss{b}{genitive masculine plural}{Ἱεροσολυμίτης}{inhabitant of Jerusalem} Οὐχ οὗτός ἐστιν\FnParse{c}{present active indicative 3rd singular} ὃν ζητοῦσιν\FnParse{d}{present active indicative 3rd plural} ἀποκτεῖναι;\FnParse{e}{aorist active infinitive} 
\MainTextVerseMark{7}{26}καὶ ἴδε παρρησίᾳ λαλεῖ\FnParse{a}{present active indicative 3rd singular} καὶ οὐδὲν αὐτῷ λέγουσιν·\FnParse{b}{present active indicative 3rd plural} μήποτε\FnFormGloss{c}{μήποτε}{never} ἀληθῶς\FnFormGloss{d}{ἀληθῶς}{truly} ἔγνωσαν\FnParse{e}{aorist active indicative 3rd plural} οἱ ἄρχοντες ὅτι οὗτός ⸀ἐστιν\FnParse{f}{present active indicative 3rd singular} ὁ χριστός; 
\MainTextVerseMark{7}{27}ἀλλὰ τοῦτον οἴδαμεν\FnParse{a}{perfect active indicative 1st plural} πόθεν\FnFormGloss{b}{πόθεν}{from where} ἐστίν·\FnParse{c}{present active indicative 3rd singular} ὁ δὲ χριστὸς ὅταν ἔρχηται\FnParse{d}{present middle subjunctive 3rd singular} οὐδεὶς γινώσκει\FnParse{e}{present active indicative 3rd singular} πόθεν\FnFormGloss{f}{πόθεν}{from where} ἐστίν.\FnParse{g}{present active indicative 3rd singular} 
\MainTextVerseMark{7}{28}ἔκραξεν\FnParse{a}{aorist active indicative 3rd singular} οὖν ἐν τῷ ἱερῷ διδάσκων\FnParse{b}{present active participle nominative masculine singular} ὁ Ἰησοῦς καὶ λέγων·\FnParse{c}{present active participle nominative masculine singular} Κἀμὲ οἴδατε\FnParse{d}{perfect active indicative 2nd plural} καὶ οἴδατε\FnParse{e}{perfect active indicative 2nd plural} πόθεν\FnFormGloss{f}{πόθεν}{from where} εἰμί·\FnParse{g}{present active indicative 1st singular} καὶ ἀπ’ ἐμαυτοῦ οὐκ ἐλήλυθα,\FnParse{h}{perfect active indicative 1st singular} ἀλλ’ ἔστιν\FnParse{i}{present active indicative 3rd singular} ἀληθινὸς\FnParseFormGloss{j}{nominative masculine singular}{ἀληθινός}{true} ὁ πέμψας\FnParse{k}{aorist active participle nominative masculine singular} με, ὃν ὑμεῖς οὐκ οἴδατε·\FnParse{l}{perfect active indicative 2nd plural} 
\MainTextVerseMark{7}{29}ἐγὼ οἶδα\FnParse{a}{perfect active indicative 1st singular} αὐτόν, ὅτι παρ’ αὐτοῦ εἰμι\FnParse{b}{present active indicative 1st singular} κἀκεῖνός\FnParseFormGloss{c}{nominative masculine singular}{κἀκεῖνος}{and that one} με ἀπέστειλεν.\FnParse{d}{aorist active indicative 3rd singular} 
\MainTextVerseMark{7}{30}ἐζήτουν\FnParse{a}{imperfect active indicative 3rd plural} οὖν αὐτὸν πιάσαι,\FnParseFormGloss{b}{aorist active infinitive}{πιάζω}{to seize} καὶ οὐδεὶς ἐπέβαλεν\FnParseFormGloss{c}{aorist active indicative 3rd singular}{ἐπιβάλλω}{to throw over} ἐπ’ αὐτὸν τὴν χεῖρα, ὅτι οὔπω\FnFormGloss{d}{οὔπω}{not yet} ἐληλύθει\FnParse{e}{pluperfect active indicative 3rd singular} ἡ ὥρα αὐτοῦ. 
\MainTextVerseMark{7}{31}⸂ἐκ τοῦ ὄχλου δὲ πολλοὶ⸃ ἐπίστευσαν\FnParse{a}{aorist active indicative 3rd plural} εἰς αὐτόν, καὶ ⸀ἔλεγον·\FnParse{b}{imperfect active indicative 3rd plural} Ὁ χριστὸς ὅταν ἔλθῃ\FnParse{c}{aorist active subjunctive 3rd singular} ⸀μὴ πλείονα ⸀σημεῖα ποιήσει\FnParse{d}{future active indicative 3rd singular} ὧν οὗτος ἐποίησεν;\FnParse{e}{aorist active indicative 3rd singular} 
\MainTextVerseMark{7}{32}Ἤκουσαν\FnParse{a}{aorist active indicative 3rd plural} οἱ Φαρισαῖοι τοῦ ὄχλου γογγύζοντος\FnParseFormGloss{b}{present active participle genitive masculine singular}{γογγύζω}{to grumble} περὶ αὐτοῦ ταῦτα, καὶ ἀπέστειλαν\FnParse{c}{aorist active indicative 3rd plural} ⸂οἱ ἀρχιερεῖς καὶ οἱ Φαρισαῖοι ὑπηρέτας⸃\FnParseFormGloss{d}{accusative masculine plural}{ὑπηρέτης}{servant} ἵνα πιάσωσιν\FnParseFormGloss{e}{aorist active subjunctive 3rd plural}{πιάζω}{to seize} αὐτόν. 
\MainTextVerseMark{7}{33}εἶπεν\FnParse{a}{aorist active indicative 3rd singular} οὖν ὁ Ἰησοῦς· Ἔτι ⸂χρόνον μικρὸν⸃ μεθ’ ὑμῶν εἰμι\FnParse{b}{present active indicative 1st singular} καὶ ὑπάγω\FnParse{c}{present active indicative 1st singular} πρὸς τὸν πέμψαντά\FnParse{d}{aorist active participle accusative masculine singular} με. 
\MainTextVerseMark{7}{34}ζητήσετέ\FnParse{a}{future active indicative 2nd plural} με καὶ οὐχ ⸀εὑρήσετέ,\FnParse{b}{future active indicative 2nd plural} καὶ ὅπου εἰμὶ\FnParse{c}{present active indicative 1st singular} ἐγὼ ὑμεῖς οὐ δύνασθε\FnParse{d}{present middle indicative 2nd plural} ἐλθεῖν.\FnParse{e}{aorist active infinitive} 
\MainTextVerseMark{7}{35}εἶπον\FnParse{a}{aorist active indicative 3rd plural} οὖν οἱ Ἰουδαῖοι πρὸς ἑαυτούς· Ποῦ οὗτος μέλλει\FnParse{b}{present active indicative 3rd singular} πορεύεσθαι\FnParse{c}{present middle infinitive} ὅτι ἡμεῖς οὐχ εὑρήσομεν\FnParse{d}{future active indicative 1st plural} αὐτόν; μὴ εἰς τὴν διασπορὰν\FnParseFormGloss{e}{accusative feminine singular}{διασπορά}{scattering} τῶν Ἑλλήνων\FnParseFormGloss{f}{genitive masculine plural}{Ἕλλην}{Greek} μέλλει\FnParse{g}{present active indicative 3rd singular} πορεύεσθαι\FnParse{h}{present middle infinitive} καὶ διδάσκειν\FnParse{i}{present active infinitive} τοὺς Ἕλληνας;\FnParseFormGloss{j}{accusative masculine plural}{Ἕλλην}{Greek} 
\MainTextVerseMark{7}{36}τίς ἐστιν\FnParse{a}{present active indicative 3rd singular} ⸂ὁ λόγος οὗτος⸃ ὃν εἶπε·\FnParse{b}{aorist active indicative 3rd singular} Ζητήσετέ\FnParse{c}{future active indicative 2nd plural} με καὶ οὐχ ⸀εὑρήσετέ,\FnParse{d}{future active indicative 2nd plural} καὶ ὅπου εἰμὶ\FnParse{e}{present active indicative 1st singular} ἐγὼ ὑμεῖς οὐ δύνασθε\FnParse{f}{present middle indicative 2nd plural} ἐλθεῖν;\FnParse{g}{aorist active infinitive} 
\MainTextVerseMark{7}{37}Ἐν δὲ τῇ ἐσχάτῃ ἡμέρᾳ τῇ μεγάλῃ τῆς ἑορτῆς\FnParseFormGloss{a}{genitive feminine singular}{ἑορτή}{feast} εἱστήκει\FnParse{b}{pluperfect active indicative 3rd singular} ὁ Ἰησοῦς, καὶ ἔκραξεν\FnParse{c}{aorist active indicative 3rd singular} λέγων·\FnParse{d}{present active participle nominative masculine singular} Ἐάν τις διψᾷ\FnParseFormGloss{e}{present active subjunctive 3rd singular}{διψάω}{to be thirsty} ἐρχέσθω\FnParse{f}{present middle imperative 3rd singular} πρός με καὶ πινέτω.\FnParse{g}{present active imperative 3rd singular} 
\MainTextVerseMark{7}{38}ὁ πιστεύων\FnParse{a}{present active participle nominative masculine singular} εἰς ἐμέ, καθὼς εἶπεν\FnParse{b}{aorist active indicative 3rd singular} ἡ γραφή, ποταμοὶ\FnParseFormGloss{c}{nominative masculine plural}{ποταμός}{river} ἐκ τῆς κοιλίας\FnParseFormGloss{d}{genitive feminine singular}{κοιλία}{any and all internal organs} αὐτοῦ ῥεύσουσιν\FnParseFormGloss{e}{future active indicative 3rd plural}{ῥέω}{to flow} ὕδατος ζῶντος.\FnParse{f}{present active participle genitive neuter singular} 
\MainTextVerseMark{7}{39}τοῦτο δὲ εἶπεν\FnParse{a}{aorist active indicative 3rd singular} περὶ τοῦ πνεύματος ⸀οὗ ἔμελλον\FnParse{b}{imperfect active indicative 3rd plural} λαμβάνειν\FnParse{c}{present active infinitive} οἱ ⸀πιστεύσαντες\FnParse{d}{aorist active participle nominative masculine plural} εἰς αὐτόν· οὔπω\FnFormGloss{e}{οὔπω}{not yet} γὰρ ἦν\FnParse{f}{imperfect active indicative 3rd singular} ⸀πνεῦμα, ὅτι Ἰησοῦς ⸀οὐδέπω\FnFormGloss{g}{οὐδέπω}{not yet} ἐδοξάσθη.\FnParse{h}{aorist passive indicative 3rd singular} 
\MainTextVerseMark{7}{40}⸂Ἐκ τοῦ ὄχλου οὖν⸃ ἀκούσαντες\FnParse{a}{aorist active participle nominative masculine plural} ⸂τῶν λόγων τούτων⸃ ⸀ἔλεγον·\FnParse{b}{imperfect active indicative 3rd plural} Οὗτός ἐστιν\FnParse{c}{present active indicative 3rd singular} ἀληθῶς\FnFormGloss{d}{ἀληθῶς}{truly} ὁ προφήτης· 
\MainTextVerseMark{7}{41}ἄλλοι ἔλεγον·\FnParse{a}{imperfect active indicative 3rd plural} Οὗτός ἐστιν\FnParse{b}{present active indicative 3rd singular} ὁ χριστός· ⸂οἱ δὲ⸃ ἔλεγον·\FnParse{c}{imperfect active indicative 3rd plural} Μὴ γὰρ ἐκ τῆς Γαλιλαίας ὁ χριστὸς ἔρχεται;\FnParse{d}{present middle indicative 3rd singular} 
\MainTextVerseMark{7}{42}⸀οὐχ ἡ γραφὴ εἶπεν\FnParse{a}{aorist active indicative 3rd singular} ὅτι ἐκ τοῦ σπέρματος Δαυὶδ, καὶ ἀπὸ Βηθλέεμ\FnParseFormGloss{b}{genitive feminine singular}{Βηθλέεμ}{Bethlehem} τῆς κώμης\FnParseFormGloss{c}{genitive feminine singular}{κώμη}{a village} ὅπου ἦν\FnParse{d}{imperfect active indicative 3rd singular} Δαυὶδ, ⸂ἔρχεται\FnParse{e}{present middle indicative 3rd singular} ὁ χριστός⸃; 
\MainTextVerseMark{7}{43}σχίσμα\FnParseFormGloss{a}{nominative neuter singular}{σχίσμα}{tear} οὖν ⸂ἐγένετο\FnParse{b}{aorist middle indicative 3rd singular} ἐν τῷ ὄχλῳ⸃ δι’ αὐτόν. 
\MainTextVerseMark{7}{44}τινὲς δὲ ἤθελον\FnParse{a}{imperfect active indicative 3rd plural} ἐξ αὐτῶν πιάσαι\FnParseFormGloss{b}{aorist active infinitive}{πιάζω}{to seize} αὐτόν, ἀλλ’ οὐδεὶς ⸀ἐπέβαλεν\FnParseFormGloss{c}{aorist active indicative 3rd singular}{ἐπιβάλλω}{to throw over} ἐπ’ αὐτὸν τὰς χεῖρας. 
\MainTextVerseMark{7}{45}Ἦλθον\FnParse{a}{aorist active indicative 3rd plural} οὖν οἱ ὑπηρέται\FnParseFormGloss{b}{nominative masculine plural}{ὑπηρέτης}{servant} πρὸς τοὺς ἀρχιερεῖς καὶ Φαρισαίους, καὶ εἶπον\FnParse{c}{aorist active indicative 3rd plural} αὐτοῖς ἐκεῖνοι· Διὰ τί οὐκ ἠγάγετε\FnParse{d}{aorist active indicative 2nd plural} αὐτόν; 
\MainTextVerseMark{7}{46}ἀπεκρίθησαν\FnParse{a}{aorist passive indicative 3rd plural} οἱ ὑπηρέται·\FnParseFormGloss{b}{nominative masculine plural}{ὑπηρέτης}{servant} Οὐδέποτε\FnFormGloss{c}{οὐδέποτε}{never} ⸂ἐλάλησεν\FnParse{d}{aorist active indicative 3rd singular} οὕτως⸃ ⸀ἄνθρωπος. 
\MainTextVerseMark{7}{47}ἀπεκρίθησαν\FnParse{a}{aorist passive indicative 3rd plural} οὖν αὐτοῖς οἱ Φαρισαῖοι· Μὴ καὶ ὑμεῖς πεπλάνησθε;\FnParse{b}{perfect passive indicative 2nd plural} 
\MainTextVerseMark{7}{48}μή τις ἐκ τῶν ἀρχόντων ἐπίστευσεν\FnParse{a}{aorist active indicative 3rd singular} εἰς αὐτὸν ἢ ἐκ τῶν Φαρισαίων; 
\MainTextVerseMark{7}{49}ἀλλὰ ὁ ὄχλος οὗτος ὁ μὴ γινώσκων\FnParse{a}{present active participle nominative masculine singular} τὸν νόμον ⸀ἐπάρατοί\FnParseFormGloss{b}{nominative masculine plural}{ἐπάρατος}{accursed} εἰσιν.\FnParse{c}{present active indicative 3rd plural} 
\MainTextVerseMark{7}{50}λέγει\FnParse{a}{present active indicative 3rd singular} Νικόδημος\FnParseFormGloss{b}{nominative masculine singular}{Νικόδημος}{Nicodemus} πρὸς αὐτούς, ὁ ⸀ἐλθὼν\FnParse{c}{aorist active participle nominative masculine singular} πρὸς αὐτὸν ⸀πρότερον,\FnParseFormGloss{d}{accusative neuter singular}{πρότερος}{before} εἷς ὢν\FnParse{e}{present active participle nominative masculine singular} ἐξ αὐτῶν· 
\MainTextVerseMark{7}{51}Μὴ ὁ νόμος ἡμῶν κρίνει\FnParse{a}{present active indicative 3rd singular} τὸν ἄνθρωπον ἐὰν μὴ ἀκούσῃ\FnParse{b}{aorist active subjunctive 3rd singular} ⸂πρῶτον παρ’ αὐτοῦ⸃ καὶ γνῷ\FnParse{c}{aorist active subjunctive 3rd singular} τί ποιεῖ;\FnParse{d}{present active indicative 3rd singular} 
\MainTextVerseMark{7}{52}ἀπεκρίθησαν\FnParse{a}{aorist passive indicative 3rd plural} καὶ εἶπαν\FnParse{b}{aorist active indicative 3rd plural} αὐτῷ· Μὴ καὶ σὺ ἐκ τῆς Γαλιλαίας εἶ;\FnParse{c}{present active indicative 2nd singular} ἐραύνησον\FnParseFormGloss{d}{aorist active imperative 2nd singular}{ἐραυνάω}{to search} καὶ ἴδε\FnParse{e}{aorist active imperative 2nd singular} ὅτι ⸂ἐκ τῆς Γαλιλαίας προφήτης⸃ οὐκ ⸀ἐγείρεται.\FnParse{f}{present passive indicative 3rd singular} 
\ChapterHeader{8}{Chapter 8}

\MainTextVerseMark{8}{12}Πάλιν 
\MainTextVerseMark{8}{12}οὖν αὐτοῖς ⸂ἐλάλησεν\FnParse{a}{aorist active indicative 3rd singular} ὁ Ἰησοῦς⸃ λέγων·\FnParse{b}{present active participle nominative masculine singular} Ἐγώ εἰμι\FnParse{c}{present active indicative 1st singular} τὸ φῶς τοῦ κόσμου· ὁ ἀκολουθῶν\FnParse{d}{present active participle nominative masculine singular} ⸀ἐμοὶ οὐ μὴ περιπατήσῃ\FnParse{e}{aorist active subjunctive 3rd singular} ἐν τῇ σκοτίᾳ,\FnParseFormGloss{f}{dative feminine singular}{σκοτία}{darkness} ἀλλ’ ἕξει\FnParse{g}{future active indicative 3rd singular} τὸ φῶς τῆς ζωῆς. 
\MainTextVerseMark{8}{13}εἶπον\FnParse{a}{aorist active indicative 3rd plural} οὖν αὐτῷ οἱ Φαρισαῖοι· Σὺ περὶ σεαυτοῦ μαρτυρεῖς·\FnParse{b}{present active indicative 2nd singular} ἡ μαρτυρία σου οὐκ ἔστιν\FnParse{c}{present active indicative 3rd singular} ἀληθής.\FnParseFormGloss{d}{nominative feminine singular}{ἀληθής}{true} 
\MainTextVerseMark{8}{14}ἀπεκρίθη\FnParse{a}{aorist passive indicative 3rd singular} Ἰησοῦς καὶ εἶπεν\FnParse{b}{aorist active indicative 3rd singular} αὐτοῖς· Κἂν\FnFormGloss{c}{κἄν}{and if} ἐγὼ μαρτυρῶ\FnParse{d}{present active subjunctive 1st singular} περὶ ἐμαυτοῦ, ἀληθής\FnParseFormGloss{e}{nominative feminine singular}{ἀληθής}{true} ἐστιν\FnParse{f}{present active indicative 3rd singular} ἡ μαρτυρία μου, ὅτι οἶδα\FnParse{g}{perfect active indicative 1st singular} πόθεν\FnFormGloss{h}{πόθεν}{from where} ἦλθον\FnParse{i}{aorist active indicative 1st singular} καὶ ποῦ ὑπάγω·\FnParse{j}{present active indicative 1st singular} ὑμεῖς δὲ οὐκ οἴδατε\FnParse{k}{perfect active indicative 2nd plural} πόθεν\FnFormGloss{l}{πόθεν}{from where} ἔρχομαι\FnParse{m}{present middle indicative 1st singular} ⸀ἢ ποῦ ὑπάγω.\FnParse{n}{present active indicative 1st singular} 
\MainTextVerseMark{8}{15}ὑμεῖς κατὰ τὴν σάρκα κρίνετε,\FnParse{a}{present active indicative 2nd plural} ἐγὼ οὐ κρίνω\FnParse{b}{present active indicative 1st singular} οὐδένα. 
\MainTextVerseMark{8}{16}καὶ ἐὰν κρίνω\FnParse{a}{present active subjunctive 1st singular} δὲ ἐγώ, ἡ κρίσις ἡ ἐμὴ ⸀ἀληθινή\FnParseFormGloss{b}{nominative feminine singular}{ἀληθινός}{true} ἐστιν,\FnParse{c}{present active indicative 3rd singular} ὅτι μόνος οὐκ εἰμί,\FnParse{d}{present active indicative 1st singular} ἀλλ’ ἐγὼ καὶ ὁ πέμψας\FnParse{e}{aorist active participle nominative masculine singular} με πατήρ. 
\MainTextVerseMark{8}{17}καὶ ἐν τῷ νόμῳ δὲ τῷ ὑμετέρῳ\FnParseFormGloss{a}{dative masculine singular}{ὑμέτερος}{your} γέγραπται\FnParse{b}{perfect passive indicative 3rd singular} ὅτι δύο ἀνθρώπων ἡ μαρτυρία ἀληθής\FnParseFormGloss{c}{nominative feminine singular}{ἀληθής}{true} ἐστιν.\FnParse{d}{present active indicative 3rd singular} 
\MainTextVerseMark{8}{18}ἐγώ εἰμι\FnParse{a}{present active indicative 1st singular} ὁ μαρτυρῶν\FnParse{b}{present active participle nominative masculine singular} περὶ ἐμαυτοῦ καὶ μαρτυρεῖ\FnParse{c}{present active indicative 3rd singular} περὶ ἐμοῦ ὁ πέμψας\FnParse{d}{aorist active participle nominative masculine singular} με πατήρ. 
\MainTextVerseMark{8}{19}ἔλεγον\FnParse{a}{imperfect active indicative 3rd plural} οὖν αὐτῷ· Ποῦ ἐστιν\FnParse{b}{present active indicative 3rd singular} ὁ πατήρ σου; ἀπεκρίθη\FnParse{c}{aorist passive indicative 3rd singular} Ἰησοῦς· Οὔτε ἐμὲ οἴδατε\FnParse{d}{perfect active indicative 2nd plural} οὔτε τὸν πατέρα μου· εἰ ἐμὲ ᾔδειτε,\FnParse{e}{pluperfect active indicative 2nd plural} καὶ τὸν πατέρα μου ⸂ἂν ᾔδειτε⸃.\FnParse{f}{pluperfect active indicative 2nd plural} 
\MainTextVerseMark{8}{20}ταῦτα τὰ ῥήματα ⸀ἐλάλησεν\FnParse{a}{aorist active indicative 3rd singular} ἐν τῷ γαζοφυλακίῳ\FnParseFormGloss{b}{dative neuter singular}{γαζοφυλάκιον}{treasury} διδάσκων\FnParse{c}{present active participle nominative masculine singular} ἐν τῷ ἱερῷ· καὶ οὐδεὶς ἐπίασεν\FnParseFormGloss{d}{aorist active indicative 3rd singular}{πιάζω}{to seize} αὐτόν, ὅτι οὔπω\FnFormGloss{e}{οὔπω}{not yet} ἐληλύθει\FnParse{f}{pluperfect active indicative 3rd singular} ἡ ὥρα αὐτοῦ. 
\MainTextVerseMark{8}{21}Εἶπεν\FnParse{a}{aorist active indicative 3rd singular} οὖν πάλιν ⸀αὐτοῖς· Ἐγὼ ὑπάγω\FnParse{b}{present active indicative 1st singular} καὶ ζητήσετέ\FnParse{c}{future active indicative 2nd plural} με, καὶ ἐν τῇ ἁμαρτίᾳ ὑμῶν ἀποθανεῖσθε·\FnParse{d}{future middle indicative 2nd plural} ὅπου ἐγὼ ὑπάγω\FnParse{e}{present active indicative 1st singular} ὑμεῖς οὐ δύνασθε\FnParse{f}{present middle indicative 2nd plural} ἐλθεῖν.\FnParse{g}{aorist active infinitive} 
\MainTextVerseMark{8}{22}ἔλεγον\FnParse{a}{imperfect active indicative 3rd plural} οὖν οἱ Ἰουδαῖοι· Μήτι\FnFormGloss{b}{μήτι}{[interrogative particle]} ἀποκτενεῖ\FnParse{c}{future active indicative 3rd singular} ἑαυτὸν ὅτι λέγει·\FnParse{d}{present active indicative 3rd singular} Ὅπου ἐγὼ ὑπάγω\FnParse{e}{present active indicative 1st singular} ὑμεῖς οὐ δύνασθε\FnParse{f}{present middle indicative 2nd plural} ἐλθεῖν;\FnParse{g}{aorist active infinitive} 
\MainTextVerseMark{8}{23}καὶ ⸀ἔλεγεν\FnParse{a}{imperfect active indicative 3rd singular} αὐτοῖς· Ὑμεῖς ἐκ τῶν κάτω\FnFormGloss{b}{κάτω}{below} ἐστέ,\FnParse{c}{present active indicative 2nd plural} ἐγὼ ἐκ τῶν ἄνω\FnFormGloss{d}{ἄνω}{above} εἰμί·\FnParse{e}{present active indicative 1st singular} ὑμεῖς ἐκ ⸂τούτου τοῦ κόσμου⸃ ἐστέ,\FnParse{f}{present active indicative 2nd plural} ἐγὼ οὐκ εἰμὶ\FnParse{g}{present active indicative 1st singular} ἐκ τοῦ κόσμου τούτου. 
\MainTextVerseMark{8}{24}εἶπον\FnParse{a}{aorist active indicative 1st singular} οὖν ὑμῖν ὅτι ἀποθανεῖσθε\FnParse{b}{future middle indicative 2nd plural} ἐν ταῖς ἁμαρτίαις ὑμῶν· ἐὰν γὰρ μὴ πιστεύσητε\FnParse{c}{aorist active subjunctive 2nd plural} ὅτι ἐγώ εἰμι,\FnParse{d}{present active indicative 1st singular} ἀποθανεῖσθε\FnParse{e}{future middle indicative 2nd plural} ἐν ταῖς ἁμαρτίαις ὑμῶν. 
\MainTextVerseMark{8}{25}ἔλεγον\FnParse{a}{imperfect active indicative 3rd plural} οὖν αὐτῷ· Σὺ τίς εἶ;\FnParse{b}{present active indicative 2nd singular} ⸀εἶπεν\FnParse{c}{aorist active indicative 3rd singular} αὐτοῖς ὁ Ἰησοῦς· Τὴν ἀρχὴν ⸂ὅ τι⸃ καὶ λαλῶ\FnParse{d}{present active indicative 1st singular} ὑμῖν; 
\MainTextVerseMark{8}{26}πολλὰ ἔχω\FnParse{a}{present active indicative 1st singular} περὶ ὑμῶν λαλεῖν\FnParse{b}{present active infinitive} καὶ κρίνειν·\FnParse{c}{present active infinitive} ἀλλ’ ὁ πέμψας\FnParse{d}{aorist active participle nominative masculine singular} με ἀληθής\FnParseFormGloss{e}{nominative masculine singular}{ἀληθής}{true} ἐστιν,\FnParse{f}{present active indicative 3rd singular} κἀγὼ ἃ ἤκουσα\FnParse{g}{aorist active indicative 1st singular} παρ’ αὐτοῦ ταῦτα ⸀λαλῶ\FnParse{h}{present active indicative 1st singular} εἰς τὸν κόσμον. 
\MainTextVerseMark{8}{27}οὐκ ἔγνωσαν\FnParse{a}{aorist active indicative 3rd plural} ὅτι τὸν πατέρα αὐτοῖς ἔλεγεν.\FnParse{b}{imperfect active indicative 3rd singular} 
\MainTextVerseMark{8}{28}εἶπεν\FnParse{a}{aorist active indicative 3rd singular} ⸀οὖν ὁ Ἰησοῦς· Ὅταν ὑψώσητε\FnParseFormGloss{b}{aorist active subjunctive 2nd plural}{ὑψόω}{to lift up} τὸν υἱὸν τοῦ ἀνθρώπου, τότε γνώσεσθε\FnParse{c}{future middle indicative 2nd plural} ὅτι ἐγώ εἰμι,\FnParse{d}{present active indicative 1st singular} καὶ ἀπ’ ἐμαυτοῦ ποιῶ\FnParse{e}{present active indicative 1st singular} οὐδέν, ἀλλὰ καθὼς ἐδίδαξέν\FnParse{f}{aorist active indicative 3rd singular} με ὁ ⸀πατὴρ ταῦτα λαλῶ.\FnParse{g}{present active indicative 1st singular} 
\MainTextVerseMark{8}{29}καὶ ὁ πέμψας\FnParse{a}{aorist active participle nominative masculine singular} με μετ’ ἐμοῦ ἐστιν·\FnParse{b}{present active indicative 3rd singular} οὐκ ἀφῆκέν\FnParse{c}{aorist active indicative 3rd singular} με ⸀μόνον, ὅτι ἐγὼ τὰ ἀρεστὰ\FnParseFormGloss{d}{accusative neuter plural}{ἀρεστός}{pleasing} αὐτῷ ποιῶ\FnParse{e}{present active indicative 1st singular} πάντοτε. 
\MainTextVerseMark{8}{30}ταῦτα αὐτοῦ λαλοῦντος\FnParse{a}{present active participle genitive masculine singular} πολλοὶ ἐπίστευσαν\FnParse{b}{aorist active indicative 3rd plural} εἰς αὐτόν. 
\MainTextVerseMark{8}{31}Ἔλεγεν\FnParse{a}{imperfect active indicative 3rd singular} οὖν ὁ Ἰησοῦς πρὸς τοὺς πεπιστευκότας\FnParse{b}{perfect active participle accusative masculine plural} αὐτῷ Ἰουδαίους· Ἐὰν ὑμεῖς μείνητε\FnParse{c}{aorist active subjunctive 2nd plural} ἐν τῷ λόγῳ τῷ ἐμῷ, ἀληθῶς\FnFormGloss{d}{ἀληθῶς}{truly} μαθηταί μού ἐστε,\FnParse{e}{present active indicative 2nd plural} 
\MainTextVerseMark{8}{32}καὶ γνώσεσθε\FnParse{a}{future middle indicative 2nd plural} τὴν ἀλήθειαν, καὶ ἡ ἀλήθεια ἐλευθερώσει\FnParseFormGloss{b}{future active indicative 3rd singular}{ἐλευθερόω}{to set free} ὑμᾶς. 
\MainTextVerseMark{8}{33}ἀπεκρίθησαν\FnParse{a}{aorist passive indicative 3rd plural} ⸂πρὸς αὐτόν⸃· Σπέρμα Ἀβραάμ ἐσμεν\FnParse{b}{present active indicative 1st plural} καὶ οὐδενὶ δεδουλεύκαμεν\FnParseFormGloss{c}{perfect active indicative 1st plural}{δουλεύω}{to serve} πώποτε·\FnFormGloss{d}{πώποτε}{ever} πῶς σὺ λέγεις\FnParse{e}{present active indicative 2nd singular} ὅτι Ἐλεύθεροι\FnParseFormGloss{f}{nominative masculine plural}{ἐλεύθερος}{free} γενήσεσθε;\FnParse{g}{future middle indicative 2nd plural} 
\MainTextVerseMark{8}{34}Ἀπεκρίθη\FnParse{a}{aorist passive indicative 3rd singular} αὐτοῖς ὁ Ἰησοῦς· Ἀμὴν ἀμὴν λέγω\FnParse{b}{present active indicative 1st singular} ὑμῖν ὅτι πᾶς ὁ ποιῶν\FnParse{c}{present active participle nominative masculine singular} τὴν ἁμαρτίαν δοῦλός ἐστιν\FnParse{d}{present active indicative 3rd singular} τῆς ἁμαρτίας· 
\MainTextVerseMark{8}{35}ὁ δὲ δοῦλος οὐ μένει\FnParse{a}{present active indicative 3rd singular} ἐν τῇ οἰκίᾳ εἰς τὸν αἰῶνα· ⸀ὁ υἱὸς μένει\FnParse{b}{present active indicative 3rd singular} εἰς τὸν αἰῶνα. 
\MainTextVerseMark{8}{36}ἐὰν οὖν ὁ υἱὸς ὑμᾶς ἐλευθερώσῃ,\FnParseFormGloss{a}{aorist active subjunctive 3rd singular}{ἐλευθερόω}{to set free} ὄντως\FnFormGloss{b}{ὄντως}{really} ἐλεύθεροι\FnParseFormGloss{c}{nominative masculine plural}{ἐλεύθερος}{free} ἔσεσθε.\FnParse{d}{future middle indicative 2nd plural} 
\MainTextVerseMark{8}{37}οἶδα\FnParse{a}{perfect active indicative 1st singular} ὅτι σπέρμα Ἀβραάμ ἐστε·\FnParse{b}{present active indicative 2nd plural} ἀλλὰ ζητεῖτέ\FnParse{c}{present active indicative 2nd plural} με ἀποκτεῖναι,\FnParse{d}{aorist active infinitive} ὅτι ὁ λόγος ὁ ἐμὸς οὐ χωρεῖ\FnParseFormGloss{e}{present active indicative 3rd singular}{χωρέω}{to go} ἐν ὑμῖν. 
\MainTextVerseMark{8}{38}⸂ἃ ἐγὼ⸃ ἑώρακα\FnParse{a}{perfect active indicative 1st singular} παρὰ τῷ ⸀πατρὶ λαλῶ·\FnParse{b}{present active indicative 1st singular} καὶ ὑμεῖς οὖν ⸂ἃ ἠκούσατε⸃\FnParse{c}{aorist active indicative 2nd plural} παρὰ ⸂τοῦ πατρὸς⸃ ποιεῖτε.\FnParse{d}{present active indicative 2nd plural} 
\MainTextVerseMark{8}{39}Ἀπεκρίθησαν\FnParse{a}{aorist passive indicative 3rd plural} καὶ εἶπαν\FnParse{b}{aorist active indicative 3rd plural} αὐτῷ· Ὁ πατὴρ ἡμῶν Ἀβραάμ ἐστιν.\FnParse{c}{present active indicative 3rd singular} λέγει\FnParse{d}{present active indicative 3rd singular} αὐτοῖς ὁ Ἰησοῦς· Εἰ τέκνα τοῦ Ἀβραάμ ⸀ἐστε,\FnParse{e}{present active indicative 2nd plural} τὰ ἔργα τοῦ Ἀβραὰμ ⸀ἐποιεῖτε·\FnParse{f}{imperfect active indicative 2nd plural} 
\MainTextVerseMark{8}{40}νῦν δὲ ζητεῖτέ\FnParse{a}{present active indicative 2nd plural} με ἀποκτεῖναι,\FnParse{b}{aorist active infinitive} ἄνθρωπον ὃς τὴν ἀλήθειαν ὑμῖν λελάληκα\FnParse{c}{perfect active indicative 1st singular} ἣν ἤκουσα\FnParse{d}{aorist active indicative 1st singular} παρὰ τοῦ θεοῦ· τοῦτο Ἀβραὰμ οὐκ ἐποίησεν.\FnParse{e}{aorist active indicative 3rd singular} 
\MainTextVerseMark{8}{41}ὑμεῖς ποιεῖτε\FnParse{a}{present active indicative 2nd plural} τὰ ἔργα τοῦ πατρὸς ὑμῶν. ⸀εἶπαν\FnParse{b}{aorist active indicative 3rd plural} αὐτῷ· Ἡμεῖς ἐκ πορνείας\FnParseFormGloss{c}{genitive feminine singular}{πορνεία}{sexual immorality} ⸂οὐ γεγεννήμεθα⸃·\FnParse{d}{perfect passive indicative 1st plural} ἕνα πατέρα ἔχομεν\FnParse{e}{present active indicative 1st plural} τὸν θεόν. 
\MainTextVerseMark{8}{42}⸀εἶπεν\FnParse{a}{aorist active indicative 3rd singular} αὐτοῖς ὁ Ἰησοῦς· Εἰ ὁ θεὸς πατὴρ ὑμῶν ἦν\FnParse{b}{imperfect active indicative 3rd singular} ἠγαπᾶτε\FnParse{c}{imperfect active indicative 2nd plural} ἂν ἐμέ, ἐγὼ γὰρ ἐκ τοῦ θεοῦ ἐξῆλθον\FnParse{d}{aorist active indicative 1st singular} καὶ ἥκω·\FnParseFormGloss{e}{present active indicative 1st singular}{ἥκω}{to come} οὐδὲ γὰρ ἀπ’ ἐμαυτοῦ ἐλήλυθα,\FnParse{f}{perfect active indicative 1st singular} ἀλλ’ ἐκεῖνός με ἀπέστειλεν.\FnParse{g}{aorist active indicative 3rd singular} 
\MainTextVerseMark{8}{43}διὰ τί τὴν λαλιὰν\FnParseFormGloss{a}{accusative feminine singular}{λαλιά}{speech} τὴν ἐμὴν οὐ γινώσκετε;\FnParse{b}{present active indicative 2nd plural} ὅτι οὐ δύνασθε\FnParse{c}{present middle indicative 2nd plural} ἀκούειν\FnParse{d}{present active infinitive} τὸν λόγον τὸν ἐμόν. 
\MainTextVerseMark{8}{44}ὑμεῖς ἐκ τοῦ πατρὸς τοῦ διαβόλου ἐστὲ\FnParse{a}{present active indicative 2nd plural} καὶ τὰς ἐπιθυμίας τοῦ πατρὸς ὑμῶν θέλετε\FnParse{b}{present active indicative 2nd plural} ποιεῖν.\FnParse{c}{present active infinitive} ἐκεῖνος ἀνθρωποκτόνος\FnParseFormGloss{d}{nominative masculine singular}{ἀνθρωποκτόνος}{murderer} ἦν\FnParse{e}{imperfect active indicative 3rd singular} ἀπ’ ἀρχῆς, καὶ ἐν τῇ ἀληθείᾳ ⸂οὐκ ἔστηκεν⸃,\FnParse{f}{imperfect active indicative 3rd singular} ὅτι οὐκ ἔστιν\FnParse{g}{present active indicative 3rd singular} ἀλήθεια ἐν αὐτῷ. ὅταν λαλῇ\FnParse{h}{present active subjunctive 3rd singular} τὸ ψεῦδος,\FnParseFormGloss{i}{accusative neuter singular}{ψεῦδος}{lie} ἐκ τῶν ἰδίων λαλεῖ,\FnParse{j}{present active indicative 3rd singular} ὅτι ψεύστης\FnParseFormGloss{k}{nominative masculine singular}{ψεύστης}{liar} ἐστὶν\FnParse{l}{present active indicative 3rd singular} καὶ ὁ πατὴρ αὐτοῦ. 
\MainTextVerseMark{8}{45}ἐγὼ δὲ ὅτι τὴν ἀλήθειαν λέγω,\FnParse{a}{present active indicative 1st singular} οὐ πιστεύετέ\FnParse{b}{present active indicative 2nd plural} μοι. 
\MainTextVerseMark{8}{46}τίς ἐξ ὑμῶν ἐλέγχει\FnParseFormGloss{a}{present active indicative 3rd singular}{ἐλέγχω}{to expose} με περὶ ἁμαρτίας; ⸀εἰ ἀλήθειαν λέγω,\FnParse{b}{present active indicative 1st singular} διὰ τί ὑμεῖς οὐ πιστεύετέ\FnParse{c}{present active indicative 2nd plural} μοι; 
\MainTextVerseMark{8}{47}ὁ ὢν\FnParse{a}{present active participle nominative masculine singular} ἐκ τοῦ θεοῦ τὰ ῥήματα τοῦ θεοῦ ἀκούει·\FnParse{b}{present active indicative 3rd singular} διὰ τοῦτο ὑμεῖς οὐκ ἀκούετε\FnParse{c}{present active indicative 2nd plural} ὅτι ἐκ τοῦ θεοῦ οὐκ ἐστέ.\FnParse{d}{present active indicative 2nd plural} 
\MainTextVerseMark{8}{48}⸀Ἀπεκρίθησαν\FnParse{a}{aorist passive indicative 3rd plural} οἱ Ἰουδαῖοι καὶ εἶπαν\FnParse{b}{aorist active indicative 3rd plural} αὐτῷ· Οὐ καλῶς λέγομεν\FnParse{c}{present active indicative 1st plural} ἡμεῖς ὅτι Σαμαρίτης\FnParseFormGloss{d}{nominative masculine singular}{Σαμαρίτης}{Samaritan} εἶ\FnParse{e}{present active indicative 2nd singular} σὺ καὶ δαιμόνιον ἔχεις;\FnParse{f}{present active indicative 2nd singular} 
\MainTextVerseMark{8}{49}ἀπεκρίθη\FnParse{a}{aorist passive indicative 3rd singular} Ἰησοῦς· Ἐγὼ δαιμόνιον οὐκ ἔχω,\FnParse{b}{present active indicative 1st singular} ἀλλὰ τιμῶ\FnParseFormGloss{c}{present active indicative 1st singular}{τιμάω}{to honor} τὸν πατέρα μου, καὶ ὑμεῖς ἀτιμάζετέ\FnParseFormGloss{d}{present active indicative 2nd plural}{ἀτιμάζω}{to dishonor} με. 
\MainTextVerseMark{8}{50}ἐγὼ δὲ οὐ ζητῶ\FnParse{a}{present active indicative 1st singular} τὴν δόξαν μου· ἔστιν\FnParse{b}{present active indicative 3rd singular} ὁ ζητῶν\FnParse{c}{present active participle nominative masculine singular} καὶ κρίνων.\FnParse{d}{present active participle nominative masculine singular} 
\MainTextVerseMark{8}{51}ἀμὴν ἀμὴν λέγω\FnParse{a}{present active indicative 1st singular} ὑμῖν, ἐάν τις τὸν ⸂ἐμὸν λόγον⸃ τηρήσῃ,\FnParse{b}{aorist active subjunctive 3rd singular} θάνατον οὐ μὴ θεωρήσῃ\FnParse{c}{aorist active subjunctive 3rd singular} εἰς τὸν αἰῶνα. 
\MainTextVerseMark{8}{52}⸀εἶπον\FnParse{a}{aorist active indicative 3rd plural} αὐτῷ οἱ Ἰουδαῖοι· Νῦν ἐγνώκαμεν\FnParse{b}{perfect active indicative 1st plural} ὅτι δαιμόνιον ἔχεις.\FnParse{c}{present active indicative 2nd singular} Ἀβραὰμ ἀπέθανεν\FnParse{d}{aorist active indicative 3rd singular} καὶ οἱ προφῆται, καὶ σὺ λέγεις·\FnParse{e}{present active indicative 2nd singular} Ἐάν τις τὸν λόγον μου τηρήσῃ,\FnParse{f}{aorist active subjunctive 3rd singular} οὐ μὴ γεύσηται\FnParseFormGloss{g}{aorist middle subjunctive 3rd singular}{γεύομαι}{to taste} θανάτου εἰς τὸν αἰῶνα· 
\MainTextVerseMark{8}{53}μὴ σὺ μείζων εἶ\FnParse{a}{present active indicative 2nd singular} τοῦ πατρὸς ἡμῶν Ἀβραάμ, ὅστις ἀπέθανεν;\FnParse{b}{aorist active indicative 3rd singular} καὶ οἱ προφῆται ἀπέθανον·\FnParse{c}{aorist active indicative 3rd plural} τίνα ⸀σεαυτὸν ποιεῖς;\FnParse{d}{present active indicative 2nd singular} 
\MainTextVerseMark{8}{54}ἀπεκρίθη\FnParse{a}{aorist passive indicative 3rd singular} Ἰησοῦς· Ἐὰν ἐγὼ ⸀δοξάσω\FnParse{b}{aorist active subjunctive 1st singular} ἐμαυτόν, ἡ δόξα μου οὐδέν ἐστιν·\FnParse{c}{present active indicative 3rd singular} ἔστιν\FnParse{d}{present active indicative 3rd singular} ὁ πατήρ μου ὁ δοξάζων\FnParse{e}{present active participle nominative masculine singular} με, ὃν ὑμεῖς λέγετε\FnParse{f}{present active indicative 2nd plural} ὅτι θεὸς ⸀ἡμῶν ἐστιν,\FnParse{g}{present active indicative 3rd singular} 
\MainTextVerseMark{8}{55}καὶ οὐκ ἐγνώκατε\FnParse{a}{perfect active indicative 2nd plural} αὐτόν, ἐγὼ δὲ οἶδα\FnParse{b}{perfect active indicative 1st singular} αὐτόν· κἂν\FnFormGloss{c}{κἄν}{and if} εἴπω\FnParse{d}{aorist active subjunctive 1st singular} ὅτι οὐκ οἶδα\FnParse{e}{perfect active indicative 1st singular} αὐτόν, ἔσομαι\FnParse{f}{future middle indicative 1st singular} ὅμοιος ⸀ὑμῖν ψεύστης·\FnParseFormGloss{g}{nominative masculine singular}{ψεύστης}{liar} ἀλλὰ οἶδα\FnParse{h}{perfect active indicative 1st singular} αὐτὸν καὶ τὸν λόγον αὐτοῦ τηρῶ.\FnParse{i}{present active indicative 1st singular} 
\MainTextVerseMark{8}{56}Ἀβραὰμ ὁ πατὴρ ὑμῶν ἠγαλλιάσατο\FnParseFormGloss{a}{aorist middle indicative 3rd singular}{ἀγαλλιάω}{to be filled with delight} ἵνα ἴδῃ\FnParse{b}{aorist active subjunctive 3rd singular} τὴν ἡμέραν τὴν ἐμήν, καὶ εἶδεν\FnParse{c}{aorist active indicative 3rd singular} καὶ ἐχάρη.\FnParse{d}{aorist passive indicative 3rd singular} 
\MainTextVerseMark{8}{57}εἶπον\FnParse{a}{aorist active indicative 3rd plural} οὖν οἱ Ἰουδαῖοι πρὸς αὐτόν· Πεντήκοντα\FnParseFormGloss{b}{accusative neuter plural}{πεντήκοντα}{fifty} ἔτη οὔπω\FnFormGloss{c}{οὔπω}{not yet} ἔχεις\FnParse{d}{present active indicative 2nd singular} καὶ Ἀβραὰμ ἑώρακας;\FnParse{e}{perfect active indicative 2nd singular} 
\MainTextVerseMark{8}{58}εἶπεν\FnParse{a}{aorist active indicative 3rd singular} ⸀αὐτοῖς Ἰησοῦς· Ἀμὴν ἀμὴν λέγω\FnParse{b}{present active indicative 1st singular} ὑμῖν, πρὶν\FnFormGloss{c}{πρίν}{before} Ἀβραὰμ γενέσθαι\FnParse{d}{aorist middle infinitive} ἐγὼ εἰμί.\FnParse{e}{present active indicative 1st singular} 
\MainTextVerseMark{8}{59}ἦραν\FnParse{a}{aorist active indicative 3rd plural} οὖν λίθους ἵνα βάλωσιν\FnParse{b}{aorist active subjunctive 3rd plural} ἐπ’ αὐτόν· Ἰησοῦς δὲ ἐκρύβη\FnParseFormGloss{c}{aorist passive indicative 3rd singular}{κρύπτω}{to hide} καὶ ἐξῆλθεν\FnParse{d}{aorist active indicative 3rd singular} ἐκ τοῦ ⸀ἱεροῦ. 
\ChapterHeader{9}{Chapter 9}

\MainTextVerseMark{9}{1}Καὶ παράγων\FnParseFormGloss{a}{present active participle nominative masculine singular}{παράγω}{to pass by} εἶδεν\FnParse{b}{aorist active indicative 3rd singular} ἄνθρωπον τυφλὸν ἐκ γενετῆς.\FnParseFormGloss{c}{genitive feminine singular}{γενετή}{birth} 
\MainTextVerseMark{9}{2}καὶ ἠρώτησαν\FnParse{a}{aorist active indicative 3rd plural} αὐτὸν οἱ μαθηταὶ αὐτοῦ λέγοντες·\FnParse{b}{present active participle nominative masculine plural} Ῥαββί,\FnParseFormGloss{c}{masculine singular}{ῥαββί}{Rabbi} τίς ἥμαρτεν,\FnParse{d}{aorist active indicative 3rd singular} οὗτος ἢ οἱ γονεῖς\FnParseFormGloss{e}{nominative masculine plural}{γονεύς}{parents} αὐτοῦ, ἵνα τυφλὸς γεννηθῇ;\FnParse{f}{aorist passive subjunctive 3rd singular} 
\MainTextVerseMark{9}{3}ἀπεκρίθη\FnParse{a}{aorist passive indicative 3rd singular} Ἰησοῦς· Οὔτε οὗτος ἥμαρτεν\FnParse{b}{aorist active indicative 3rd singular} οὔτε οἱ γονεῖς\FnParseFormGloss{c}{nominative masculine plural}{γονεύς}{parents} αὐτοῦ, ἀλλ’ ἵνα φανερωθῇ\FnParse{d}{aorist passive subjunctive 3rd singular} τὰ ἔργα τοῦ θεοῦ ἐν αὐτῷ. 
\MainTextVerseMark{9}{4}⸀ἡμᾶς δεῖ\FnParse{a}{present active indicative 3rd singular} ἐργάζεσθαι\FnParse{b}{present middle infinitive} τὰ ἔργα τοῦ πέμψαντός\FnParse{c}{aorist active participle genitive masculine singular} με ἕως ἡμέρα ἐστίν·\FnParse{d}{present active indicative 3rd singular} ἔρχεται\FnParse{e}{present middle indicative 3rd singular} νὺξ ὅτε οὐδεὶς δύναται\FnParse{f}{present middle indicative 3rd singular} ἐργάζεσθαι.\FnParse{g}{present middle infinitive} 
\MainTextVerseMark{9}{5}ὅταν ἐν τῷ κόσμῳ ὦ,\FnParse{a}{present active subjunctive 1st singular} φῶς εἰμι\FnParse{b}{present active indicative 1st singular} τοῦ κόσμου. 
\MainTextVerseMark{9}{6}ταῦτα εἰπὼν\FnParse{a}{aorist active participle nominative masculine singular} ἔπτυσεν\FnParseFormGloss{b}{aorist active indicative 3rd singular}{πτύω}{to spit} χαμαὶ\FnFormGloss{c}{χαμαί}{to the ground} καὶ ἐποίησεν\FnParse{d}{aorist active indicative 3rd singular} πηλὸν\FnParseFormGloss{e}{accusative masculine singular}{πηλός}{mud} ἐκ τοῦ πτύσματος,\FnParseFormGloss{f}{genitive neuter singular}{πτύσμα}{saliva} καὶ ⸀ἐπέχρισεν\FnParseFormGloss{g}{aorist active indicative 3rd singular}{ἐπιχρίω}{to put on} ⸀αὐτοῦ τὸν πηλὸν\FnParseFormGloss{h}{accusative masculine singular}{πηλός}{mud} ἐπὶ τοὺς ⸀ὀφθαλμούς, 
\MainTextVerseMark{9}{7}καὶ εἶπεν\FnParse{a}{aorist active indicative 3rd singular} αὐτῷ· Ὕπαγε\FnParse{b}{present active imperative 2nd singular} νίψαι\FnParseFormGloss{c}{aorist middle imperative 2nd singular}{νίπτω}{to wash} εἰς τὴν κολυμβήθραν\FnParseFormGloss{d}{accusative feminine singular}{κολυμβήθρα}{pool} τοῦ Σιλωάμ\FnParseFormGloss{e}{genitive masculine singular}{Σιλωάμ}{Siloam} (ὃ ἑρμηνεύεται\FnParseFormGloss{f}{present passive indicative 3rd singular}{ἑρμηνεύω}{to translate} Ἀπεσταλμένος).\FnParse{g}{perfect passive participle nominative masculine singular} ἀπῆλθεν\FnParse{h}{aorist active indicative 3rd singular} οὖν καὶ ἐνίψατο,\FnParseFormGloss{i}{aorist middle indicative 3rd singular}{νίπτω}{to wash} καὶ ἦλθεν\FnParse{j}{aorist active indicative 3rd singular} βλέπων.\FnParse{k}{present active participle nominative masculine singular} 
\MainTextVerseMark{9}{8}οἱ οὖν γείτονες\FnParseFormGloss{a}{nominative masculine plural}{γείτων}{neighbor} καὶ οἱ θεωροῦντες\FnParse{b}{present active participle nominative masculine plural} αὐτὸν τὸ πρότερον\FnParseFormGloss{c}{accusative neuter singular}{πρότερος}{before} ὅτι ⸀προσαίτης\FnParseFormGloss{d}{nominative masculine singular}{προσαίτης}{beggar} ἦν\FnParse{e}{imperfect active indicative 3rd singular} ἔλεγον·\FnParse{f}{imperfect active indicative 3rd plural} Οὐχ οὗτός ἐστιν\FnParse{g}{present active indicative 3rd singular} ὁ καθήμενος\FnParse{h}{present middle participle nominative masculine singular} καὶ προσαιτῶν;\FnParseFormGloss{i}{present active participle nominative masculine singular}{προσαιτέω}{to beg} 
\MainTextVerseMark{9}{9}ἄλλοι ἔλεγον\FnParse{a}{imperfect active indicative 3rd plural} ὅτι Οὗτός ἐστιν·\FnParse{b}{present active indicative 3rd singular} ἄλλοι ⸂ἔλεγον·\FnParse{c}{imperfect active indicative 3rd plural} Οὐχί, ἀλλὰ⸃ ὅμοιος αὐτῷ ἐστιν.\FnParse{d}{present active indicative 3rd singular} ἐκεῖνος ἔλεγεν\FnParse{e}{imperfect active indicative 3rd singular} ὅτι Ἐγώ εἰμι.\FnParse{f}{present active indicative 1st singular} 
\MainTextVerseMark{9}{10}ἔλεγον\FnParse{a}{imperfect active indicative 3rd plural} οὖν αὐτῷ· ⸀Πῶς ἠνεῴχθησάν\FnParse{b}{aorist passive indicative 3rd plural} σου οἱ ὀφθαλμοί; 
\MainTextVerseMark{9}{11}ἀπεκρίθη\FnParse{a}{aorist passive indicative 3rd singular} ⸀ἐκεῖνος· ⸂Ὁ ἄνθρωπος ὁ⸃ λεγόμενος\FnParse{b}{present passive participle nominative masculine singular} Ἰησοῦς πηλὸν\FnParseFormGloss{c}{accusative masculine singular}{πηλός}{mud} ἐποίησεν\FnParse{d}{aorist active indicative 3rd singular} καὶ ἐπέχρισέν\FnParseFormGloss{e}{aorist active indicative 3rd singular}{ἐπιχρίω}{to put on} μου τοὺς ὀφθαλμοὺς καὶ εἶπέν\FnParse{f}{aorist active indicative 3rd singular} μοι ⸀ὅτι Ὕπαγε\FnParse{g}{present active imperative 2nd singular} εἰς ⸀τὸν Σιλωὰμ\FnParseFormGloss{h}{accusative masculine singular}{Σιλωάμ}{Siloam} καὶ νίψαι·\FnParseFormGloss{i}{aorist middle imperative 2nd singular}{νίπτω}{to wash} ἀπελθὼν\FnParse{j}{aorist active participle nominative masculine singular} ⸀οὖν καὶ νιψάμενος\FnParseFormGloss{k}{aorist middle participle nominative masculine singular}{νίπτω}{to wash} ἀνέβλεψα.\FnParseFormGloss{l}{aorist active indicative 1st singular}{ἀναβλέπω}{to look up} 
\MainTextVerseMark{9}{12}⸂καὶ εἶπαν⸃\FnParse{a}{aorist active indicative 3rd plural} αὐτῷ· Ποῦ ἐστιν\FnParse{b}{present active indicative 3rd singular} ἐκεῖνος; λέγει·\FnParse{c}{present active indicative 3rd singular} Οὐκ οἶδα.\FnParse{d}{perfect active indicative 1st singular} 
\MainTextVerseMark{9}{13}Ἄγουσιν\FnParse{a}{present active indicative 3rd plural} αὐτὸν πρὸς τοὺς Φαρισαίους τόν ποτε\FnFormGloss{b}{ποτέ}{once} τυφλόν. 
\MainTextVerseMark{9}{14}ἦν\FnParse{a}{imperfect active indicative 3rd singular} δὲ σάββατον ⸂ἐν ᾗ ἡμέρᾳ⸃ τὸν πηλὸν\FnParseFormGloss{b}{accusative masculine singular}{πηλός}{mud} ἐποίησεν\FnParse{c}{aorist active indicative 3rd singular} ὁ Ἰησοῦς καὶ ἀνέῳξεν\FnParse{d}{aorist active indicative 3rd singular} αὐτοῦ τοὺς ὀφθαλμούς. 
\MainTextVerseMark{9}{15}πάλιν οὖν ἠρώτων\FnParse{a}{imperfect active indicative 3rd plural} αὐτὸν καὶ οἱ Φαρισαῖοι πῶς ἀνέβλεψεν.\FnParseFormGloss{b}{aorist active indicative 3rd singular}{ἀναβλέπω}{to look up} ὁ δὲ εἶπεν\FnParse{c}{aorist active indicative 3rd singular} αὐτοῖς· Πηλὸν\FnParseFormGloss{d}{accusative masculine singular}{πηλός}{mud} ἐπέθηκέν\FnParse{e}{aorist active indicative 3rd singular} μου ἐπὶ τοὺς ὀφθαλμούς, καὶ ἐνιψάμην,\FnParseFormGloss{f}{aorist middle indicative 1st singular}{νίπτω}{to wash} καὶ βλέπω.\FnParse{g}{present active indicative 1st singular} 
\MainTextVerseMark{9}{16}ἔλεγον\FnParse{a}{imperfect active indicative 3rd plural} οὖν ἐκ τῶν Φαρισαίων τινές· ⸂Οὐκ ἔστιν\FnParse{b}{present active indicative 3rd singular} οὗτος παρὰ θεοῦ ὁ ἄνθρωπος⸃, ὅτι τὸ σάββατον οὐ τηρεῖ.\FnParse{c}{present active indicative 3rd singular} ⸀ἄλλοι ἔλεγον·\FnParse{d}{imperfect active indicative 3rd plural} Πῶς δύναται\FnParse{e}{present middle indicative 3rd singular} ἄνθρωπος ἁμαρτωλὸς τοιαῦτα σημεῖα ποιεῖν;\FnParse{f}{present active infinitive} καὶ σχίσμα\FnParseFormGloss{g}{nominative neuter singular}{σχίσμα}{tear} ἦν\FnParse{h}{imperfect active indicative 3rd singular} ἐν αὐτοῖς. 
\MainTextVerseMark{9}{17}λέγουσιν\FnParse{a}{present active indicative 3rd plural} ⸀οὖν τῷ τυφλῷ πάλιν· ⸂Τί σὺ⸃ λέγεις\FnParse{b}{present active indicative 2nd singular} περὶ αὐτοῦ, ὅτι ἠνέῳξέν\FnParse{c}{aorist active indicative 3rd singular} σου τοὺς ὀφθαλμούς; ὁ δὲ εἶπεν\FnParse{d}{aorist active indicative 3rd singular} ὅτι Προφήτης ἐστίν.\FnParse{e}{present active indicative 3rd singular} 
\MainTextVerseMark{9}{18}Οὐκ ἐπίστευσαν\FnParse{a}{aorist active indicative 3rd plural} οὖν οἱ Ἰουδαῖοι περὶ αὐτοῦ ὅτι ⸂ἦν\FnParse{b}{imperfect active indicative 3rd singular} τυφλὸς⸃ καὶ ἀνέβλεψεν,\FnParseFormGloss{c}{aorist active indicative 3rd singular}{ἀναβλέπω}{to look up} ἕως ὅτου ἐφώνησαν\FnParse{d}{aorist active indicative 3rd plural} τοὺς γονεῖς\FnParseFormGloss{e}{accusative masculine plural}{γονεύς}{parents} αὐτοῦ ⸂τοῦ ἀναβλέψαντος⸃\FnParseFormGloss{f}{aorist active participle genitive masculine singular}{ἀναβλέπω}{to look up} 
\MainTextVerseMark{9}{19}καὶ ἠρώτησαν\FnParse{a}{aorist active indicative 3rd plural} αὐτοὺς λέγοντες·\FnParse{b}{present active participle nominative masculine plural} Οὗτός ἐστιν\FnParse{c}{present active indicative 3rd singular} ὁ υἱὸς ὑμῶν, ὃν ὑμεῖς λέγετε\FnParse{d}{present active indicative 2nd plural} ὅτι τυφλὸς ἐγεννήθη;\FnParse{e}{aorist passive indicative 3rd singular} πῶς οὖν ⸂βλέπει\FnParse{f}{present active indicative 3rd singular} ἄρτι⸃; 
\MainTextVerseMark{9}{20}ἀπεκρίθησαν\FnParse{a}{aorist passive indicative 3rd plural} ⸀οὖν οἱ γονεῖς\FnParseFormGloss{b}{nominative masculine plural}{γονεύς}{parents} αὐτοῦ καὶ εἶπαν·\FnParse{c}{aorist active indicative 3rd plural} Οἴδαμεν\FnParse{d}{perfect active indicative 1st plural} ὅτι οὗτός ἐστιν\FnParse{e}{present active indicative 3rd singular} ὁ υἱὸς ἡμῶν καὶ ὅτι τυφλὸς ἐγεννήθη·\FnParse{f}{aorist passive indicative 3rd singular} 
\MainTextVerseMark{9}{21}πῶς δὲ νῦν βλέπει\FnParse{a}{present active indicative 3rd singular} οὐκ οἴδαμεν,\FnParse{b}{perfect active indicative 1st plural} ἢ τίς ἤνοιξεν\FnParse{c}{aorist active indicative 3rd singular} αὐτοῦ τοὺς ὀφθαλμοὺς ἡμεῖς οὐκ οἴδαμεν·\FnParse{d}{perfect active indicative 1st plural} ⸂αὐτὸν ἐρωτήσατε,\FnParse{e}{aorist active imperative 2nd plural} ἡλικίαν\FnParseFormGloss{f}{accusative feminine singular}{ἡλικία}{life} ἔχει⸃,\FnParse{g}{present active indicative 3rd singular} αὐτὸς περὶ ἑαυτοῦ λαλήσει.\FnParse{h}{future active indicative 3rd singular} 
\MainTextVerseMark{9}{22}ταῦτα εἶπαν\FnParse{a}{aorist active indicative 3rd plural} οἱ γονεῖς\FnParseFormGloss{b}{nominative masculine plural}{γονεύς}{parents} αὐτοῦ ὅτι ἐφοβοῦντο\FnParse{c}{imperfect middle indicative 3rd plural} τοὺς Ἰουδαίους, ἤδη γὰρ συνετέθειντο\FnParseFormGloss{d}{pluperfect middle indicative 3rd plural}{συντίθεμαι}{to agree} οἱ Ἰουδαῖοι ἵνα ἐάν τις αὐτὸν ὁμολογήσῃ\FnParseFormGloss{e}{aorist active subjunctive 3rd singular}{ὁμολογέω}{to confess} χριστόν, ἀποσυνάγωγος\FnParseFormGloss{f}{nominative masculine singular}{ἀποσυνάγωγος}{put out of the synagogue} γένηται.\FnParse{g}{aorist middle subjunctive 3rd singular} 
\MainTextVerseMark{9}{23}διὰ τοῦτο οἱ γονεῖς\FnParseFormGloss{a}{nominative masculine plural}{γονεύς}{parents} αὐτοῦ εἶπαν\FnParse{b}{aorist active indicative 3rd plural} ὅτι Ἡλικίαν\FnParseFormGloss{c}{accusative feminine singular}{ἡλικία}{life} ἔχει,\FnParse{d}{present active indicative 3rd singular} αὐτὸν ⸀ἐπερωτήσατε.\FnParse{e}{aorist active imperative 2nd plural} 
\MainTextVerseMark{9}{24}Ἐφώνησαν\FnParse{a}{aorist active indicative 3rd plural} οὖν ⸂τὸν ἄνθρωπον ἐκ δευτέρου⸃ ὃς ἦν\FnParse{b}{imperfect active indicative 3rd singular} τυφλὸς καὶ εἶπαν\FnParse{c}{aorist active indicative 3rd plural} αὐτῷ· Δὸς\FnParse{d}{aorist active imperative 2nd singular} δόξαν τῷ θεῷ· ἡμεῖς οἴδαμεν\FnParse{e}{perfect active indicative 1st plural} ὅτι ⸂οὗτος ὁ ἄνθρωπος⸃ ἁμαρτωλός ἐστιν.\FnParse{f}{present active indicative 3rd singular} 
\MainTextVerseMark{9}{25}ἀπεκρίθη\FnParse{a}{aorist passive indicative 3rd singular} οὖν ⸀ἐκεῖνος· Εἰ ἁμαρτωλός ἐστιν\FnParse{b}{present active indicative 3rd singular} οὐκ οἶδα·\FnParse{c}{perfect active indicative 1st singular} ἓν οἶδα\FnParse{d}{perfect active indicative 1st singular} ὅτι τυφλὸς ὢν\FnParse{e}{present active participle nominative masculine singular} ἄρτι βλέπω.\FnParse{f}{present active indicative 1st singular} 
\MainTextVerseMark{9}{26}εἶπον\FnParse{a}{aorist active indicative 3rd plural} ⸂οὖν αὐτῷ⸃· Τί ἐποίησέν\FnParse{b}{aorist active indicative 3rd singular} σοι; πῶς ἤνοιξέν\FnParse{c}{aorist active indicative 3rd singular} σου τοὺς ὀφθαλμούς; 
\MainTextVerseMark{9}{27}ἀπεκρίθη\FnParse{a}{aorist passive indicative 3rd singular} αὐτοῖς· Εἶπον\FnParse{b}{aorist active indicative 1st singular} ὑμῖν ἤδη καὶ οὐκ ἠκούσατε·\FnParse{c}{aorist active indicative 2nd plural} τί πάλιν θέλετε\FnParse{d}{present active indicative 2nd plural} ἀκούειν;\FnParse{e}{present active infinitive} μὴ καὶ ὑμεῖς θέλετε\FnParse{f}{present active indicative 2nd plural} αὐτοῦ μαθηταὶ γενέσθαι;\FnParse{g}{aorist middle infinitive} 
\MainTextVerseMark{9}{28}⸀ἐλοιδόρησαν\FnParseFormGloss{a}{aorist active indicative 3rd plural}{λοιδορέω}{to insult} αὐτὸν καὶ εἶπον·\FnParse{b}{aorist active indicative 3rd plural} Σὺ ⸂μαθητὴς εἶ⸃\FnParse{c}{present active indicative 2nd singular} ἐκείνου, ἡμεῖς δὲ τοῦ Μωϋσέως ἐσμὲν\FnParse{d}{present active indicative 1st plural} μαθηταί· 
\MainTextVerseMark{9}{29}ἡμεῖς οἴδαμεν\FnParse{a}{perfect active indicative 1st plural} ὅτι Μωϋσεῖ λελάληκεν\FnParse{b}{perfect active indicative 3rd singular} ὁ θεός, τοῦτον δὲ οὐκ οἴδαμεν\FnParse{c}{perfect active indicative 1st plural} πόθεν\FnFormGloss{d}{πόθεν}{from where} ἐστίν.\FnParse{e}{present active indicative 3rd singular} 
\MainTextVerseMark{9}{30}ἀπεκρίθη\FnParse{a}{aorist passive indicative 3rd singular} ὁ ἄνθρωπος καὶ εἶπεν\FnParse{b}{aorist active indicative 3rd singular} αὐτοῖς· Ἐν ⸂τούτῳ γὰρ τὸ⸃ θαυμαστόν\FnParseFormGloss{c}{nominative neuter singular}{θαυμαστός}{wonderful} ἐστιν\FnParse{d}{present active indicative 3rd singular} ὅτι ὑμεῖς οὐκ οἴδατε\FnParse{e}{perfect active indicative 2nd plural} πόθεν\FnFormGloss{f}{πόθεν}{from where} ἐστίν,\FnParse{g}{present active indicative 3rd singular} καὶ ἤνοιξέν\FnParse{h}{aorist active indicative 3rd singular} μου τοὺς ὀφθαλμούς. 
\MainTextVerseMark{9}{31}⸀οἴδαμεν\FnParse{a}{perfect active indicative 1st plural} ὅτι ⸂ἁμαρτωλῶν ὁ θεὸς⸃ οὐκ ἀκούει,\FnParse{b}{present active indicative 3rd singular} ἀλλ’ ἐάν τις θεοσεβὴς\FnParseFormGloss{c}{nominative masculine singular}{θεοσεβής}{godly} ᾖ\FnParse{d}{present active subjunctive 3rd singular} καὶ τὸ θέλημα αὐτοῦ ποιῇ\FnParse{e}{present active subjunctive 3rd singular} τούτου ἀκούει.\FnParse{f}{present active indicative 3rd singular} 
\MainTextVerseMark{9}{32}ἐκ τοῦ αἰῶνος οὐκ ἠκούσθη\FnParse{a}{aorist passive indicative 3rd singular} ὅτι ἠνέῳξέν\FnParse{b}{aorist active indicative 3rd singular} τις ὀφθαλμοὺς τυφλοῦ γεγεννημένου·\FnParse{c}{perfect passive participle genitive masculine singular} 
\MainTextVerseMark{9}{33}εἰ μὴ ἦν\FnParse{a}{imperfect active indicative 3rd singular} οὗτος παρὰ θεοῦ, οὐκ ἠδύνατο\FnParse{b}{imperfect middle indicative 3rd singular} ποιεῖν\FnParse{c}{present active infinitive} οὐδέν. 
\MainTextVerseMark{9}{34}ἀπεκρίθησαν\FnParse{a}{aorist passive indicative 3rd plural} καὶ εἶπαν\FnParse{b}{aorist active indicative 3rd plural} αὐτῷ· Ἐν ἁμαρτίαις σὺ ἐγεννήθης\FnParse{c}{aorist passive indicative 2nd singular} ὅλος, καὶ σὺ διδάσκεις\FnParse{d}{present active indicative 2nd singular} ἡμᾶς; καὶ ἐξέβαλον\FnParse{e}{aorist active indicative 3rd plural} αὐτὸν ἔξω. 
\MainTextVerseMark{9}{35}⸀Ἤκουσεν\FnParse{a}{aorist active indicative 3rd singular} Ἰησοῦς ὅτι ἐξέβαλον\FnParse{b}{aorist active indicative 3rd plural} αὐτὸν ἔξω, καὶ εὑρὼν\FnParse{c}{aorist active participle nominative masculine singular} αὐτὸν ⸀εἶπεν·\FnParse{d}{aorist active indicative 3rd singular} Σὺ πιστεύεις\FnParse{e}{present active indicative 2nd singular} εἰς τὸν υἱὸν τοῦ ⸀ἀνθρώπου; 
\MainTextVerseMark{9}{36}ἀπεκρίθη\FnParse{a}{aorist passive indicative 3rd singular} ἐκεῖνος καὶ εἶπεν·\FnParse{b}{aorist active indicative 3rd singular} Καὶ τίς ἐστιν,\FnParse{c}{present active indicative 3rd singular} κύριε, ἵνα πιστεύσω\FnParse{d}{aorist active subjunctive 1st singular} εἰς αὐτόν; 
\MainTextVerseMark{9}{37}⸀εἶπεν\FnParse{a}{aorist active indicative 3rd singular} αὐτῷ ὁ Ἰησοῦς· Καὶ ἑώρακας\FnParse{b}{perfect active indicative 2nd singular} αὐτὸν καὶ ὁ λαλῶν\FnParse{c}{present active participle nominative masculine singular} μετὰ σοῦ ἐκεῖνός ἐστιν.\FnParse{d}{present active indicative 3rd singular} 
\MainTextVerseMark{9}{38}ὁ δὲ ἔφη·\FnParse{a}{imperfect active indicative 3rd singular} Πιστεύω,\FnParse{b}{present active indicative 1st singular} κύριε· καὶ προσεκύνησεν\FnParse{c}{aorist active indicative 3rd singular} αὐτῷ. 
\MainTextVerseMark{9}{39}καὶ εἶπεν\FnParse{a}{aorist active indicative 3rd singular} ὁ Ἰησοῦς· Εἰς κρίμα\FnParseFormGloss{b}{accusative neuter singular}{κρίμα}{judgment} ἐγὼ εἰς τὸν κόσμον τοῦτον ἦλθον,\FnParse{c}{aorist active indicative 1st singular} ἵνα οἱ μὴ βλέποντες\FnParse{d}{present active participle nominative masculine plural} βλέπωσιν\FnParse{e}{present active subjunctive 3rd plural} καὶ οἱ βλέποντες\FnParse{f}{present active participle nominative masculine plural} τυφλοὶ γένωνται.\FnParse{g}{aorist middle subjunctive 3rd plural} 
\MainTextVerseMark{9}{40}⸀ἤκουσαν\FnParse{a}{aorist active indicative 3rd plural} ἐκ τῶν Φαρισαίων ταῦτα οἱ ⸂μετ’ αὐτοῦ ὄντες⸃,\FnParse{b}{present active participle nominative masculine plural} καὶ εἶπον\FnParse{c}{aorist active indicative 3rd plural} αὐτῷ· Μὴ καὶ ἡμεῖς τυφλοί ἐσμεν;\FnParse{d}{present active indicative 1st plural} 
\MainTextVerseMark{9}{41}εἶπεν\FnParse{a}{aorist active indicative 3rd singular} αὐτοῖς ὁ Ἰησοῦς· Εἰ τυφλοὶ ἦτε,\FnParse{b}{imperfect active indicative 2nd plural} οὐκ ἂν εἴχετε\FnParse{c}{imperfect active indicative 2nd plural} ἁμαρτίαν· νῦν δὲ λέγετε\FnParse{d}{present active indicative 2nd plural} ὅτι Βλέπομεν·\FnParse{e}{present active indicative 1st plural} ⸀ἡ ἁμαρτία ὑμῶν μένει.\FnParse{f}{present active indicative 3rd singular} 
\ChapterHeader{10}{Chapter 10}

\MainTextVerseMark{10}{1}Ἀμὴν ἀμὴν λέγω\FnParse{a}{present active indicative 1st singular} ὑμῖν, ὁ μὴ εἰσερχόμενος\FnParse{b}{present middle participle nominative masculine singular} διὰ τῆς θύρας εἰς τὴν αὐλὴν\FnParseFormGloss{c}{accusative feminine singular}{αὐλή}{palace} τῶν προβάτων ἀλλὰ ἀναβαίνων\FnParse{d}{present active participle nominative masculine singular} ἀλλαχόθεν\FnFormGloss{e}{ἀλλαχόθεν}{from another way} ἐκεῖνος κλέπτης\FnParseFormGloss{f}{nominative masculine singular}{κλέπτης}{thief} ἐστὶν\FnParse{g}{present active indicative 3rd singular} καὶ λῃστής·\FnParseFormGloss{h}{nominative masculine singular}{λῃστής}{robber} 
\MainTextVerseMark{10}{2}ὁ δὲ εἰσερχόμενος\FnParse{a}{present middle participle nominative masculine singular} διὰ τῆς θύρας ποιμήν\FnParseFormGloss{b}{nominative masculine singular}{ποιμήν}{shepherd} ἐστιν\FnParse{c}{present active indicative 3rd singular} τῶν προβάτων. 
\MainTextVerseMark{10}{3}τούτῳ ὁ θυρωρὸς\FnParseFormGloss{a}{nominative masculine singular}{θυρωρός}{doorkeeper} ἀνοίγει,\FnParse{b}{present active indicative 3rd singular} καὶ τὰ πρόβατα τῆς φωνῆς αὐτοῦ ἀκούει\FnParse{c}{present active indicative 3rd singular} καὶ τὰ ἴδια πρόβατα ⸀φωνεῖ\FnParse{d}{present active indicative 3rd singular} κατ’ ὄνομα καὶ ἐξάγει\FnParseFormGloss{e}{present active indicative 3rd singular}{ἐξάγω}{to lead out} αὐτά. 
\MainTextVerseMark{10}{4}⸀ὅταν τὰ ἴδια ⸀πάντα ἐκβάλῃ,\FnParse{a}{aorist active subjunctive 3rd singular} ἔμπροσθεν αὐτῶν πορεύεται,\FnParse{b}{present middle indicative 3rd singular} καὶ τὰ πρόβατα αὐτῷ ἀκολουθεῖ,\FnParse{c}{present active indicative 3rd singular} ὅτι οἴδασιν\FnParse{d}{perfect active indicative 3rd plural} τὴν φωνὴν αὐτοῦ· 
\MainTextVerseMark{10}{5}ἀλλοτρίῳ\FnParseFormGloss{a}{dative masculine singular}{ἀλλότριος}{belonging to another} δὲ οὐ μὴ ⸀ἀκολουθήσουσιν\FnParse{b}{future active indicative 3rd plural} ἀλλὰ φεύξονται\FnParseFormGloss{c}{future middle indicative 3rd plural}{φεύγω}{flee} ἀπ’ αὐτοῦ, ὅτι οὐκ οἴδασι\FnParse{d}{perfect active indicative 3rd plural} τῶν ἀλλοτρίων\FnParseFormGloss{e}{genitive masculine plural}{ἀλλότριος}{belonging to another} τὴν φωνήν. 
\MainTextVerseMark{10}{6}ταύτην τὴν παροιμίαν\FnParseFormGloss{a}{accusative feminine singular}{παροιμία}{figure of speech} εἶπεν\FnParse{b}{aorist active indicative 3rd singular} αὐτοῖς ὁ Ἰησοῦς· ἐκεῖνοι δὲ οὐκ ἔγνωσαν\FnParse{c}{aorist active indicative 3rd plural} τίνα ⸀ἦν\FnParse{d}{imperfect active indicative 3rd singular} ἃ ἐλάλει\FnParse{e}{imperfect active indicative 3rd singular} αὐτοῖς. 
\MainTextVerseMark{10}{7}Εἶπεν\FnParse{a}{aorist active indicative 3rd singular} οὖν πάλιν ⸀αὐτοῖς ὁ Ἰησοῦς· Ἀμὴν ἀμὴν λέγω\FnParse{b}{present active indicative 1st singular} ὑμῖν ⸀ὅτι ἐγώ εἰμι\FnParse{c}{present active indicative 1st singular} ἡ θύρα τῶν προβάτων. 
\MainTextVerseMark{10}{8}πάντες ὅσοι ἦλθον\FnParse{a}{aorist active indicative 3rd plural} ⸂πρὸ ἐμοῦ⸃ κλέπται\FnParseFormGloss{b}{nominative masculine plural}{κλέπτης}{thief} εἰσὶν\FnParse{c}{present active indicative 3rd plural} καὶ λῃσταί·\FnParseFormGloss{d}{nominative masculine plural}{λῃστής}{robber} ἀλλ’ οὐκ ἤκουσαν\FnParse{e}{aorist active indicative 3rd plural} αὐτῶν τὰ πρόβατα. 
\MainTextVerseMark{10}{9}ἐγώ εἰμι\FnParse{a}{present active indicative 1st singular} ἡ θύρα· δι’ ἐμοῦ ἐάν τις εἰσέλθῃ\FnParse{b}{aorist active subjunctive 3rd singular} σωθήσεται\FnParse{c}{future passive indicative 3rd singular} καὶ εἰσελεύσεται\FnParse{d}{future middle indicative 3rd singular} καὶ ἐξελεύσεται\FnParse{e}{future middle indicative 3rd singular} καὶ νομὴν\FnParseFormGloss{f}{accusative feminine singular}{νομή}{pasture} εὑρήσει.\FnParse{g}{future active indicative 3rd singular} 
\MainTextVerseMark{10}{10}ὁ κλέπτης\FnParseFormGloss{a}{nominative masculine singular}{κλέπτης}{thief} οὐκ ἔρχεται\FnParse{b}{present middle indicative 3rd singular} εἰ μὴ ἵνα κλέψῃ\FnParseFormGloss{c}{aorist active subjunctive 3rd singular}{κλέπτω}{steal} καὶ θύσῃ\FnParseFormGloss{d}{aorist active subjunctive 3rd singular}{θύω}{to kill} καὶ ἀπολέσῃ·\FnParse{e}{aorist active subjunctive 3rd singular} ἐγὼ ἦλθον\FnParse{f}{aorist active indicative 1st singular} ἵνα ζωὴν ἔχωσιν\FnParse{g}{present active subjunctive 3rd plural} καὶ περισσὸν\FnParseFormGloss{h}{accusative neuter singular}{περισσός}{exceeding} ἔχωσιν.\FnParse{i}{present active subjunctive 3rd plural} 
\MainTextVerseMark{10}{11}Ἐγώ εἰμι\FnParse{a}{present active indicative 1st singular} ὁ ποιμὴν\FnParseFormGloss{b}{nominative masculine singular}{ποιμήν}{shepherd} ὁ καλός· ὁ ποιμὴν\FnParseFormGloss{c}{nominative masculine singular}{ποιμήν}{shepherd} ὁ καλὸς τὴν ψυχὴν αὐτοῦ τίθησιν\FnParse{d}{present active indicative 3rd singular} ὑπὲρ τῶν προβάτων· 
\MainTextVerseMark{10}{12}ὁ ⸀μισθωτὸς\FnParseFormGloss{a}{nominative masculine singular}{μισθωτός}{hired worker} καὶ οὐκ ὢν\FnParse{b}{present active participle nominative masculine singular} ποιμήν,\FnParseFormGloss{c}{nominative masculine singular}{ποιμήν}{shepherd} οὗ οὐκ ⸀ἔστιν\FnParse{d}{present active indicative 3rd singular} τὰ πρόβατα ἴδια, θεωρεῖ\FnParse{e}{present active indicative 3rd singular} τὸν λύκον\FnParseFormGloss{f}{accusative masculine singular}{λύκος}{wolf} ἐρχόμενον\FnParse{g}{present middle participle accusative masculine singular} καὶ ἀφίησιν\FnParse{h}{present active indicative 3rd singular} τὰ πρόβατα καὶ φεύγει—\FnParseFormGloss{i}{present active indicative 3rd singular}{φεύγω}{flee} καὶ ὁ λύκος\FnParseFormGloss{j}{nominative masculine singular}{λύκος}{wolf} ἁρπάζει\FnParseFormGloss{k}{present active indicative 3rd singular}{ἁρπάζω}{to catch} αὐτὰ καὶ ⸀σκορπίζει—\FnParseFormGloss{l}{present active indicative 3rd singular}{σκορπίζω}{to scatter} 
\MainTextVerseMark{10}{13}⸀ὅτι μισθωτός\FnParseFormGloss{a}{nominative masculine singular}{μισθωτός}{hired worker} ἐστιν\FnParse{b}{present active indicative 3rd singular} καὶ οὐ μέλει\FnParseFormGloss{c}{present active indicative 3rd singular}{μέλει}{it is a care} αὐτῷ περὶ τῶν προβάτων. 
\MainTextVerseMark{10}{14}ἐγώ εἰμι\FnParse{a}{present active indicative 1st singular} ὁ ποιμὴν\FnParseFormGloss{b}{nominative masculine singular}{ποιμήν}{shepherd} ὁ καλός, καὶ γινώσκω\FnParse{c}{present active indicative 1st singular} τὰ ἐμὰ καὶ ⸂γινώσκουσί\FnParse{d}{present active indicative 3rd plural} με τὰ ἐμά⸃, 
\MainTextVerseMark{10}{15}καθὼς γινώσκει\FnParse{a}{present active indicative 3rd singular} με ὁ πατὴρ κἀγὼ γινώσκω\FnParse{b}{present active indicative 1st singular} τὸν πατέρα, καὶ τὴν ψυχήν μου τίθημι\FnParse{c}{present active indicative 1st singular} ὑπὲρ τῶν προβάτων. 
\MainTextVerseMark{10}{16}καὶ ἄλλα πρόβατα ἔχω\FnParse{a}{present active indicative 1st singular} ἃ οὐκ ἔστιν\FnParse{b}{present active indicative 3rd singular} ἐκ τῆς αὐλῆς\FnParseFormGloss{c}{genitive feminine singular}{αὐλή}{palace} ταύτης· κἀκεῖνα\FnParseFormGloss{d}{accusative neuter plural}{κἀκεῖνος}{and that one} ⸂δεῖ\FnParse{e}{present active indicative 3rd singular} με⸃ ἀγαγεῖν,\FnParse{f}{aorist active infinitive} καὶ τῆς φωνῆς μου ἀκούσουσιν,\FnParse{g}{future active indicative 3rd plural} καὶ ⸀γενήσονται\FnParse{h}{future middle indicative 3rd plural} μία ποίμνη,\FnParseFormGloss{i}{nominative feminine singular}{ποίμνη}{flock} εἷς ποιμήν.\FnParseFormGloss{j}{nominative masculine singular}{ποιμήν}{shepherd} 
\MainTextVerseMark{10}{17}διὰ τοῦτό ⸂με ὁ πατὴρ⸃ ἀγαπᾷ\FnParse{a}{present active indicative 3rd singular} ὅτι ἐγὼ τίθημι\FnParse{b}{present active indicative 1st singular} τὴν ψυχήν μου, ἵνα πάλιν λάβω\FnParse{c}{aorist active subjunctive 1st singular} αὐτήν. 
\MainTextVerseMark{10}{18}οὐδεὶς ⸀αἴρει\FnParse{a}{present active indicative 3rd singular} αὐτὴν ἀπ’ ἐμοῦ, ἀλλ’ ἐγὼ τίθημι\FnParse{b}{present active indicative 1st singular} αὐτὴν ἀπ’ ἐμαυτοῦ. ἐξουσίαν ἔχω\FnParse{c}{present active indicative 1st singular} θεῖναι\FnParse{d}{aorist active infinitive} αὐτήν, καὶ ἐξουσίαν ἔχω\FnParse{e}{present active indicative 1st singular} πάλιν λαβεῖν\FnParse{f}{aorist active infinitive} αὐτήν· ταύτην τὴν ἐντολὴν ἔλαβον\FnParse{g}{aorist active indicative 1st singular} παρὰ τοῦ πατρός μου. 
\MainTextVerseMark{10}{19}⸀Σχίσμα\FnParseFormGloss{a}{nominative neuter singular}{σχίσμα}{tear} πάλιν ἐγένετο\FnParse{b}{aorist middle indicative 3rd singular} ἐν τοῖς Ἰουδαίοις διὰ τοὺς λόγους τούτους. 
\MainTextVerseMark{10}{20}ἔλεγον\FnParse{a}{imperfect active indicative 3rd plural} δὲ πολλοὶ ἐξ αὐτῶν· Δαιμόνιον ἔχει\FnParse{b}{present active indicative 3rd singular} καὶ μαίνεται·\FnParseFormGloss{c}{present middle indicative 3rd singular}{μαίνομαι}{to rave} τί αὐτοῦ ἀκούετε;\FnParse{d}{present active indicative 2nd plural} 
\MainTextVerseMark{10}{21}⸀ἄλλοι ἔλεγον·\FnParse{a}{imperfect active indicative 3rd plural} Ταῦτα τὰ ῥήματα οὐκ ἔστιν\FnParse{b}{present active indicative 3rd singular} δαιμονιζομένου·\FnParseFormGloss{c}{present middle participle genitive masculine singular}{δαιμονίζομαι}{to be demon-possessed} μὴ δαιμόνιον δύναται\FnParse{d}{present middle indicative 3rd singular} τυφλῶν ὀφθαλμοὺς ⸀ἀνοῖξαι;\FnParse{e}{aorist active infinitive} 
\MainTextVerseMark{10}{22}Ἐγένετο\FnParse{a}{aorist middle indicative 3rd singular} ⸀τότε τὰ ἐγκαίνια\FnParseFormGloss{b}{nominative neuter plural}{ἐγκαίνια}{Feast of Dedication} ἐν ⸀τοῖς Ἱεροσολύμοις· ⸀χειμὼν\FnParseFormGloss{c}{nominative masculine singular}{χειμών}{winter} ἦν,\FnParse{d}{imperfect active indicative 3rd singular} 
\MainTextVerseMark{10}{23}καὶ περιεπάτει\FnParse{a}{imperfect active indicative 3rd singular} ὁ Ἰησοῦς ἐν τῷ ἱερῷ ἐν τῇ στοᾷ\FnParseFormGloss{b}{dative feminine singular}{στοά}{covered colonnade} ⸀τοῦ Σολομῶνος.\FnParseFormGloss{c}{genitive masculine singular}{Σολομών}{Solomon} 
\MainTextVerseMark{10}{24}ἐκύκλωσαν\FnParseFormGloss{a}{aorist active indicative 3rd plural}{κυκλόω}{gather around} οὖν αὐτὸν οἱ Ἰουδαῖοι καὶ ἔλεγον\FnParse{b}{imperfect active indicative 3rd plural} αὐτῷ· Ἕως πότε\FnFormGloss{c}{πότε}{when? how long?} τὴν ψυχὴν ἡμῶν αἴρεις;\FnParse{d}{present active indicative 2nd singular} εἰ σὺ εἶ\FnParse{e}{present active indicative 2nd singular} ὁ χριστός, εἰπὲ\FnParse{f}{aorist active imperative 2nd singular} ἡμῖν παρρησίᾳ. 
\MainTextVerseMark{10}{25}ἀπεκρίθη\FnParse{a}{aorist passive indicative 3rd singular} αὐτοῖς ὁ Ἰησοῦς· Εἶπον\FnParse{b}{aorist active indicative 1st singular} ὑμῖν καὶ οὐ πιστεύετε·\FnParse{c}{present active indicative 2nd plural} τὰ ἔργα ἃ ἐγὼ ποιῶ\FnParse{d}{present active indicative 1st singular} ἐν τῷ ὀνόματι τοῦ πατρός μου ταῦτα μαρτυρεῖ\FnParse{e}{present active indicative 3rd singular} περὶ ἐμοῦ· 
\MainTextVerseMark{10}{26}ἀλλὰ ὑμεῖς οὐ πιστεύετε,\FnParse{a}{present active indicative 2nd plural} ⸂ὅτι οὐκ⸃ ἐστὲ\FnParse{b}{present active indicative 2nd plural} ἐκ τῶν προβάτων τῶν ⸀ἐμῶν. 
\MainTextVerseMark{10}{27}τὰ πρόβατα τὰ ἐμὰ τῆς φωνῆς μου ⸀ἀκούουσιν,\FnParse{a}{present active indicative 3rd plural} κἀγὼ γινώσκω\FnParse{b}{present active indicative 1st singular} αὐτά, καὶ ἀκολουθοῦσίν\FnParse{c}{present active indicative 3rd plural} μοι, 
\MainTextVerseMark{10}{28}κἀγὼ ⸂δίδωμι\FnParse{a}{present active indicative 1st singular} αὐτοῖς ζωὴν αἰώνιον⸃, καὶ οὐ μὴ ἀπόλωνται\FnParse{b}{aorist middle subjunctive 3rd plural} εἰς τὸν αἰῶνα, καὶ οὐχ ἁρπάσει\FnParseFormGloss{c}{future active indicative 3rd singular}{ἁρπάζω}{to catch} τις αὐτὰ ἐκ τῆς χειρός μου. 
\MainTextVerseMark{10}{29}ὁ πατήρ μου ⸀ὃ δέδωκέν\FnParse{a}{perfect active indicative 3rd singular} ⸀μοι ⸂πάντων μεῖζων⸃ ἐστιν,\FnParse{b}{present active indicative 3rd singular} καὶ οὐδεὶς δύναται\FnParse{c}{present middle indicative 3rd singular} ἁρπάζειν\FnParseFormGloss{d}{present active infinitive}{ἁρπάζω}{to catch} ἐκ τῆς χειρὸς τοῦ ⸀πατρός. 
\MainTextVerseMark{10}{30}ἐγὼ καὶ ὁ πατὴρ ἕν ἐσμεν.\FnParse{a}{present active indicative 1st plural} 
\MainTextVerseMark{10}{31}Ἐβάστασαν\FnParseFormGloss{a}{aorist active indicative 3rd plural}{βαστάζω}{to carry} ⸀οὖν πάλιν λίθους οἱ Ἰουδαῖοι ἵνα λιθάσωσιν\FnParseFormGloss{b}{aorist active subjunctive 3rd plural}{λιθάζω}{to stone} αὐτόν. 
\MainTextVerseMark{10}{32}ἀπεκρίθη\FnParse{a}{aorist passive indicative 3rd singular} αὐτοῖς ὁ Ἰησοῦς· Πολλὰ ⸂ἔργα καλὰ ἔδειξα\FnParse{b}{aorist active indicative 1st singular} ὑμῖν⸃ ἐκ τοῦ ⸀πατρός· διὰ ποῖον αὐτῶν ἔργον ⸂ἐμὲ λιθάζετε⸃;\FnParseFormGloss{c}{present active indicative 2nd plural}{λιθάζω}{to stone} 
\MainTextVerseMark{10}{33}ἀπεκρίθησαν\FnParse{a}{aorist passive indicative 3rd plural} αὐτῷ οἱ ⸀Ἰουδαῖοι· Περὶ καλοῦ ἔργου οὐ λιθάζομέν\FnParseFormGloss{b}{present active indicative 1st plural}{λιθάζω}{to stone} σε ἀλλὰ περὶ βλασφημίας,\FnParseFormGloss{c}{genitive feminine singular}{βλασφημία}{blasphemy} καὶ ὅτι σὺ ἄνθρωπος ὢν\FnParse{d}{present active participle nominative masculine singular} ποιεῖς\FnParse{e}{present active indicative 2nd singular} σεαυτὸν θεόν. 
\MainTextVerseMark{10}{34}ἀπεκρίθη\FnParse{a}{aorist passive indicative 3rd singular} αὐτοῖς ὁ Ἰησοῦς· Οὐκ ἔστιν\FnParse{b}{present active indicative 3rd singular} γεγραμμένον\FnParse{c}{perfect passive participle nominative neuter singular} ἐν τῷ νόμῳ ὑμῶν ⸀ὅτι Ἐγὼ εἶπα·\FnParse{d}{aorist active indicative 1st singular} Θεοί ἐστε;\FnParse{e}{present active indicative 2nd plural} 
\MainTextVerseMark{10}{35}εἰ ἐκείνους εἶπεν\FnParse{a}{aorist active indicative 3rd singular} θεοὺς πρὸς οὓς ὁ λόγος τοῦ θεοῦ ἐγένετο,\FnParse{b}{aorist middle indicative 3rd singular} καὶ οὐ δύναται\FnParse{c}{present middle indicative 3rd singular} λυθῆναι\FnParse{d}{aorist passive infinitive} ἡ γραφή, 
\MainTextVerseMark{10}{36}ὃν ὁ πατὴρ ἡγίασεν\FnParseFormGloss{a}{aorist active indicative 3rd singular}{ἁγιάζω}{to sanctify} καὶ ἀπέστειλεν\FnParse{b}{aorist active indicative 3rd singular} εἰς τὸν κόσμον ὑμεῖς λέγετε\FnParse{c}{present active indicative 2nd plural} ὅτι Βλασφημεῖς,\FnParse{d}{present active indicative 2nd singular} ὅτι εἶπον·\FnParse{e}{aorist active indicative 1st singular} Υἱὸς τοῦ θεοῦ εἰμι;\FnParse{f}{present active indicative 1st singular} 
\MainTextVerseMark{10}{37}εἰ οὐ ποιῶ\FnParse{a}{present active indicative 1st singular} τὰ ἔργα τοῦ πατρός μου, μὴ πιστεύετέ\FnParse{b}{present active imperative 2nd plural} μοι· 
\MainTextVerseMark{10}{38}εἰ δὲ ποιῶ,\FnParse{a}{present active indicative 1st singular} κἂν\FnFormGloss{b}{κἄν}{and if} ἐμοὶ μὴ πιστεύητε\FnParse{c}{present active subjunctive 2nd plural} τοῖς ἔργοις ⸀πιστεύετε,\FnParse{d}{present active imperative 2nd plural} ἵνα γνῶτε\FnParse{e}{aorist active subjunctive 2nd plural} καὶ ⸀γινώσκητε\FnParse{f}{present active subjunctive 2nd plural} ὅτι ἐν ἐμοὶ ὁ πατὴρ κἀγὼ ἐν ⸂τῷ πατρί⸃. 
\MainTextVerseMark{10}{39}ἐζήτουν\FnParse{a}{imperfect active indicative 3rd plural} οὖν ⸂πάλιν αὐτὸν⸃ πιάσαι·\FnParseFormGloss{b}{aorist active infinitive}{πιάζω}{to seize} καὶ ἐξῆλθεν\FnParse{c}{aorist active indicative 3rd singular} ἐκ τῆς χειρὸς αὐτῶν. 
\MainTextVerseMark{10}{40}Καὶ ἀπῆλθεν\FnParse{a}{aorist active indicative 3rd singular} πάλιν πέραν\FnFormGloss{b}{πέραν}{on the other side} τοῦ Ἰορδάνου\FnParseFormGloss{c}{genitive masculine singular}{Ἰορδάνης}{Jordan} εἰς τὸν τόπον ὅπου ἦν\FnParse{d}{imperfect active indicative 3rd singular} Ἰωάννης τὸ πρῶτον βαπτίζων,\FnParse{e}{present active participle nominative masculine singular} καὶ ἔμεινεν\FnParse{f}{aorist active indicative 3rd singular} ἐκεῖ. 
\MainTextVerseMark{10}{41}καὶ πολλοὶ ἦλθον\FnParse{a}{aorist active indicative 3rd plural} πρὸς αὐτὸν καὶ ἔλεγον\FnParse{b}{imperfect active indicative 3rd plural} ὅτι Ἰωάννης μὲν σημεῖον ἐποίησεν\FnParse{c}{aorist active indicative 3rd singular} οὐδέν, πάντα δὲ ὅσα εἶπεν\FnParse{d}{aorist active indicative 3rd singular} Ἰωάννης περὶ τούτου ἀληθῆ\FnParseFormGloss{e}{nominative neuter plural}{ἀληθής}{true} ἦν.\FnParse{f}{imperfect active indicative 3rd singular} 
\MainTextVerseMark{10}{42}καὶ ⸂πολλοὶ ἐπίστευσαν\FnParse{a}{aorist active indicative 3rd plural} εἰς αὐτὸν ἐκεῖ⸃. 
\ChapterHeader{11}{Chapter 11}

\MainTextVerseMark{11}{1}Ἦν\FnParse{a}{imperfect active indicative 3rd singular} δέ τις ἀσθενῶν,\FnParse{b}{present active participle nominative masculine singular} Λάζαρος\FnParseFormGloss{c}{nominative masculine singular}{Λάζαρος}{Lazarus} ἀπὸ Βηθανίας\FnParseFormGloss{d}{genitive feminine singular}{Βηθανία}{Bethany} ἐκ τῆς κώμης\FnParseFormGloss{e}{genitive feminine singular}{κώμη}{a village} Μαρίας\FnParseFormGloss{f}{genitive feminine singular}{Μαρία}{Mary} καὶ Μάρθας\FnParseFormGloss{g}{genitive feminine singular}{Μάρθα}{Martha} τῆς ἀδελφῆς\FnParseFormGloss{h}{genitive feminine singular}{ἀδελφή}{sister} αὐτῆς. 
\MainTextVerseMark{11}{2}ἦν\FnParse{a}{imperfect active indicative 3rd singular} δὲ ⸀Μαριὰμ\FnParseFormGloss{b}{nominative feminine singular}{Μαριάμ}{Mary} ἡ ἀλείψασα\FnParseFormGloss{c}{aorist active participle nominative feminine singular}{ἀλείφω}{to pour on} τὸν κύριον μύρῳ\FnParseFormGloss{d}{dative neuter singular}{μύρον}{perfume} καὶ ἐκμάξασα\FnParseFormGloss{e}{aorist active participle nominative feminine singular}{ἐκμάσσω}{to wipe off} τοὺς πόδας αὐτοῦ ταῖς θριξὶν\FnParseFormGloss{f}{dative feminine plural}{θρίξ}{hair} αὐτῆς, ἧς ὁ ἀδελφὸς Λάζαρος\FnParseFormGloss{g}{nominative masculine singular}{Λάζαρος}{Lazarus} ἠσθένει.\FnParse{h}{imperfect active indicative 3rd singular} 
\MainTextVerseMark{11}{3}ἀπέστειλαν\FnParse{a}{aorist active indicative 3rd plural} οὖν αἱ ἀδελφαὶ\FnParseFormGloss{b}{nominative feminine plural}{ἀδελφή}{sister} πρὸς αὐτὸν λέγουσαι·\FnParse{c}{present active participle nominative feminine plural} Κύριε, ἴδε ὃν φιλεῖς\FnParseFormGloss{d}{present active indicative 2nd singular}{φιλέω}{to love} ἀσθενεῖ.\FnParse{e}{present active indicative 3rd singular} 
\MainTextVerseMark{11}{4}ἀκούσας\FnParse{a}{aorist active participle nominative masculine singular} δὲ ὁ Ἰησοῦς εἶπεν·\FnParse{b}{aorist active indicative 3rd singular} Αὕτη ἡ ἀσθένεια\FnParseFormGloss{c}{nominative feminine singular}{ἀσθένεια}{weakness} οὐκ ἔστιν\FnParse{d}{present active indicative 3rd singular} πρὸς θάνατον ἀλλ’ ὑπὲρ τῆς δόξης τοῦ θεοῦ ἵνα δοξασθῇ\FnParse{e}{aorist passive subjunctive 3rd singular} ὁ υἱὸς τοῦ θεοῦ δι’ αὐτῆς. 
\MainTextVerseMark{11}{5}ἠγάπα\FnParse{a}{imperfect active indicative 3rd singular} δὲ ὁ Ἰησοῦς τὴν Μάρθαν\FnParseFormGloss{b}{accusative feminine singular}{Μάρθα}{Martha} καὶ τὴν ἀδελφὴν\FnParseFormGloss{c}{accusative feminine singular}{ἀδελφή}{sister} αὐτῆς καὶ τὸν Λάζαρον.\FnParseFormGloss{d}{accusative masculine singular}{Λάζαρος}{Lazarus} 
\MainTextVerseMark{11}{6}ὡς οὖν ἤκουσεν\FnParse{a}{aorist active indicative 3rd singular} ὅτι ἀσθενεῖ,\FnParse{b}{present active indicative 3rd singular} τότε μὲν ἔμεινεν\FnParse{c}{aorist active indicative 3rd singular} ἐν ᾧ ἦν\FnParse{d}{imperfect active indicative 3rd singular} τόπῳ δύο ἡμέρας· 
\MainTextVerseMark{11}{7}ἔπειτα\FnFormGloss{a}{ἔπειτα}{then} μετὰ τοῦτο λέγει\FnParse{b}{present active indicative 3rd singular} τοῖς μαθηταῖς· Ἄγωμεν\FnParse{c}{present active subjunctive 1st plural} εἰς τὴν Ἰουδαίαν πάλιν. 
\MainTextVerseMark{11}{8}λέγουσιν\FnParse{a}{present active indicative 3rd plural} αὐτῷ οἱ μαθηταί· Ῥαββί,\FnParseFormGloss{b}{masculine singular}{ῥαββί}{Rabbi} νῦν ἐζήτουν\FnParse{c}{imperfect active indicative 3rd plural} σε λιθάσαι\FnParseFormGloss{d}{aorist active infinitive}{λιθάζω}{to stone} οἱ Ἰουδαῖοι, καὶ πάλιν ὑπάγεις\FnParse{e}{present active indicative 2nd singular} ἐκεῖ; 
\MainTextVerseMark{11}{9}ἀπεκρίθη\FnParse{a}{aorist passive indicative 3rd singular} Ἰησοῦς· Οὐχὶ δώδεκα ⸂ὧραί εἰσιν⸃\FnParse{b}{present active indicative 3rd plural} τῆς ἡμέρας; ἐάν τις περιπατῇ\FnParse{c}{present active subjunctive 3rd singular} ἐν τῇ ἡμέρᾳ, οὐ προσκόπτει,\FnParseFormGloss{d}{present active indicative 3rd singular}{προσκόπτω}{to strike} ὅτι τὸ φῶς τοῦ κόσμου τούτου βλέπει·\FnParse{e}{present active indicative 3rd singular} 
\MainTextVerseMark{11}{10}ἐὰν δέ τις περιπατῇ\FnParse{a}{present active subjunctive 3rd singular} ἐν τῇ νυκτί, προσκόπτει,\FnParseFormGloss{b}{present active indicative 3rd singular}{προσκόπτω}{to strike} ὅτι τὸ φῶς οὐκ ἔστιν\FnParse{c}{present active indicative 3rd singular} ἐν αὐτῷ. 
\MainTextVerseMark{11}{11}ταῦτα εἶπεν,\FnParse{a}{aorist active indicative 3rd singular} καὶ μετὰ τοῦτο λέγει\FnParse{b}{present active indicative 3rd singular} αὐτοῖς· Λάζαρος\FnParseFormGloss{c}{nominative masculine singular}{Λάζαρος}{Lazarus} ὁ φίλος\FnParseFormGloss{d}{nominative masculine singular}{φίλος}{friendly} ἡμῶν κεκοίμηται,\FnParseFormGloss{e}{perfect middle indicative 3rd singular}{κοιμάομαι}{to fall asleep} ἀλλὰ πορεύομαι\FnParse{f}{present middle indicative 1st singular} ἵνα ἐξυπνίσω\FnParseFormGloss{g}{aorist active subjunctive 1st singular}{ἐξυπνίζω}{to wake up} αὐτόν. 
\MainTextVerseMark{11}{12}εἶπαν\FnParse{a}{aorist active indicative 3rd plural} οὖν οἱ μαθηταὶ ⸀αὐτῷ· Κύριε, εἰ κεκοίμηται\FnParseFormGloss{b}{perfect middle indicative 3rd singular}{κοιμάομαι}{to fall asleep} σωθήσεται.\FnParse{c}{future passive indicative 3rd singular} 
\MainTextVerseMark{11}{13}εἰρήκει\FnParse{a}{pluperfect active indicative 3rd singular} δὲ ὁ Ἰησοῦς περὶ τοῦ θανάτου αὐτοῦ. ἐκεῖνοι δὲ ἔδοξαν\FnParse{b}{aorist active indicative 3rd plural} ὅτι περὶ τῆς κοιμήσεως\FnParseFormGloss{c}{genitive feminine singular}{κοίμησις}{sleep} τοῦ ὕπνου\FnParseFormGloss{d}{genitive masculine singular}{ὕπνος}{sleep} λέγει.\FnParse{e}{present active indicative 3rd singular} 
\MainTextVerseMark{11}{14}τότε οὖν εἶπεν\FnParse{a}{aorist active indicative 3rd singular} αὐτοῖς ὁ Ἰησοῦς παρρησίᾳ· Λάζαρος\FnParseFormGloss{b}{nominative masculine singular}{Λάζαρος}{Lazarus} ἀπέθανεν,\FnParse{c}{aorist active indicative 3rd singular} 
\MainTextVerseMark{11}{15}καὶ χαίρω\FnParse{a}{present active indicative 1st singular} δι’ ὑμᾶς, ἵνα πιστεύσητε,\FnParse{b}{aorist active subjunctive 2nd plural} ὅτι οὐκ ἤμην\FnParse{c}{imperfect middle indicative 1st singular} ἐκεῖ· ἀλλὰ ἄγωμεν\FnParse{d}{present active subjunctive 1st plural} πρὸς αὐτόν. 
\MainTextVerseMark{11}{16}εἶπεν\FnParse{a}{aorist active indicative 3rd singular} οὖν Θωμᾶς\FnParseFormGloss{b}{nominative masculine singular}{Θωμᾶς}{Thomas} ὁ λεγόμενος\FnParse{c}{present passive participle nominative masculine singular} Δίδυμος\FnParseFormGloss{d}{nominative masculine singular}{Δίδυμος}{Didymus} τοῖς συμμαθηταῖς·\FnParseFormGloss{e}{dative masculine plural}{συμμαθητής}{fellow disciple} Ἄγωμεν\FnParse{f}{present active subjunctive 1st plural} καὶ ἡμεῖς ἵνα ἀποθάνωμεν\FnParse{g}{aorist active subjunctive 1st plural} μετ’ αὐτοῦ. 
\MainTextVerseMark{11}{17}Ἐλθὼν\FnParse{a}{aorist active participle nominative masculine singular} οὖν ὁ Ἰησοῦς εὗρεν\FnParse{b}{aorist active indicative 3rd singular} αὐτὸν τέσσαρας ⸂ἤδη ἡμέρας⸃ ἔχοντα\FnParse{c}{present active participle accusative masculine singular} ἐν τῷ μνημείῳ. 
\MainTextVerseMark{11}{18}ἦν\FnParse{a}{imperfect active indicative 3rd singular} δὲ ⸀ἡ Βηθανία\FnParseFormGloss{b}{nominative feminine singular}{Βηθανία}{Bethany} ἐγγὺς τῶν Ἱεροσολύμων ὡς ἀπὸ σταδίων\FnParseFormGloss{c}{genitive masculine plural}{στάδιος}{arena} δεκαπέντε.\FnParseFormGloss{d}{genitive masculine plural}{δεκαπέντε}{fifteen} 
\MainTextVerseMark{11}{19}⸂πολλοὶ δὲ⸃ ἐκ τῶν Ἰουδαίων ἐληλύθεισαν\FnParse{a}{pluperfect active indicative 3rd plural} πρὸς ⸀τὴν Μάρθαν\FnParseFormGloss{b}{accusative feminine singular}{Μάρθα}{Martha} καὶ ⸀Μαριὰμ\FnParseFormGloss{c}{accusative feminine singular}{Μαριάμ}{Mary} ἵνα παραμυθήσωνται\FnParseFormGloss{d}{aorist middle subjunctive 3rd plural}{παραμυθέομαι}{to comfort} αὐτὰς περὶ τοῦ ⸀ἀδελφοῦ. 
\MainTextVerseMark{11}{20}ἡ οὖν Μάρθα\FnParseFormGloss{a}{nominative feminine singular}{Μάρθα}{Martha} ὡς ἤκουσεν\FnParse{b}{aorist active indicative 3rd singular} ὅτι Ἰησοῦς ἔρχεται\FnParse{c}{present middle indicative 3rd singular} ὑπήντησεν\FnParseFormGloss{d}{aorist active indicative 3rd singular}{ὑπαντάω}{to go out to meet} αὐτῷ· ⸀Μαρία\FnParseFormGloss{e}{nominative feminine singular}{Μαρία}{Mary} δὲ ἐν τῷ οἴκῳ ἐκαθέζετο.\FnParseFormGloss{f}{imperfect middle indicative 3rd singular}{καθέζομαι}{to sit down} 
\MainTextVerseMark{11}{21}εἶπεν\FnParse{a}{aorist active indicative 3rd singular} οὖν ⸀ἡ Μάρθα\FnParseFormGloss{b}{nominative feminine singular}{Μάρθα}{Martha} πρὸς ⸀τὸν Ἰησοῦν· Κύριε, εἰ ἦς\FnParse{c}{imperfect active indicative 2nd singular} ὧδε ⸂οὐκ ἂν ἀπέθανεν\FnParse{d}{aorist active indicative 3rd singular} ὁ ἀδελφός μου⸃· 
\MainTextVerseMark{11}{22}⸀καὶ νῦν οἶδα\FnParse{a}{perfect active indicative 1st singular} ὅτι ὅσα ἂν αἰτήσῃ\FnParse{b}{aorist middle subjunctive 2nd singular} τὸν θεὸν δώσει\FnParse{c}{future active indicative 3rd singular} σοι ὁ θεός. 
\MainTextVerseMark{11}{23}λέγει\FnParse{a}{present active indicative 3rd singular} αὐτῇ ὁ Ἰησοῦς· Ἀναστήσεται\FnParse{b}{future middle indicative 3rd singular} ὁ ἀδελφός σου. 
\MainTextVerseMark{11}{24}λέγει\FnParse{a}{present active indicative 3rd singular} αὐτῷ ⸀ἡ Μάρθα·\FnParseFormGloss{b}{nominative feminine singular}{Μάρθα}{Martha} Οἶδα\FnParse{c}{perfect active indicative 1st singular} ὅτι ἀναστήσεται\FnParse{d}{future middle indicative 3rd singular} ἐν τῇ ἀναστάσει ἐν τῇ ἐσχάτῃ ἡμέρᾳ. 
\MainTextVerseMark{11}{25}εἶπεν\FnParse{a}{aorist active indicative 3rd singular} αὐτῇ ὁ Ἰησοῦς· Ἐγώ εἰμι\FnParse{b}{present active indicative 1st singular} ἡ ἀνάστασις καὶ ἡ ζωή· ὁ πιστεύων\FnParse{c}{present active participle nominative masculine singular} εἰς ἐμὲ κἂν\FnFormGloss{d}{κἄν}{and if} ἀποθάνῃ\FnParse{e}{aorist active subjunctive 3rd singular} ζήσεται,\FnParse{f}{future middle indicative 3rd singular} 
\MainTextVerseMark{11}{26}καὶ πᾶς ὁ ζῶν\FnParse{a}{present active participle nominative masculine singular} καὶ πιστεύων\FnParse{b}{present active participle nominative masculine singular} εἰς ἐμὲ οὐ μὴ ἀποθάνῃ\FnParse{c}{aorist active subjunctive 3rd singular} εἰς τὸν αἰῶνα· πιστεύεις\FnParse{d}{present active indicative 2nd singular} τοῦτο; 
\MainTextVerseMark{11}{27}λέγει\FnParse{a}{present active indicative 3rd singular} αὐτῷ· Ναί, κύριε· ἐγὼ πεπίστευκα\FnParse{b}{perfect active indicative 1st singular} ὅτι σὺ εἶ\FnParse{c}{present active indicative 2nd singular} ὁ χριστὸς ὁ υἱὸς τοῦ θεοῦ ὁ εἰς τὸν κόσμον ἐρχόμενος.\FnParse{d}{present middle participle nominative masculine singular} 
\MainTextVerseMark{11}{28}Καὶ ⸀τοῦτο εἰποῦσα\FnParse{a}{aorist active participle nominative feminine singular} ἀπῆλθεν\FnParse{b}{aorist active indicative 3rd singular} καὶ ἐφώνησεν\FnParse{c}{aorist active indicative 3rd singular} ⸀Μαριὰμ\FnParseFormGloss{d}{accusative feminine singular}{Μαριάμ}{Mary} τὴν ἀδελφὴν\FnParseFormGloss{e}{accusative feminine singular}{ἀδελφή}{sister} αὐτῆς λάθρᾳ\FnFormGloss{f}{λάθρᾳ}{secretly} εἰποῦσα·\FnParse{g}{aorist active participle nominative feminine singular} Ὁ διδάσκαλος πάρεστιν\FnParseFormGloss{h}{present active indicative 3rd singular}{πάρειμι}{to be present} καὶ φωνεῖ\FnParse{i}{present active indicative 3rd singular} σε. 
\MainTextVerseMark{11}{29}ἐκείνη ⸀δὲ ὡς ἤκουσεν\FnParse{a}{aorist active indicative 3rd singular} ⸀ἠγέρθη\FnParse{b}{aorist passive indicative 3rd singular} ταχὺ\FnParseFormGloss{c}{accusative neuter singular}{ταχύ}{quick} καὶ ⸀ἤρχετο\FnParse{d}{imperfect middle indicative 3rd singular} πρὸς αὐτόν· 
\MainTextVerseMark{11}{30}οὔπω\FnFormGloss{a}{οὔπω}{not yet} δὲ ἐληλύθει\FnParse{b}{pluperfect active indicative 3rd singular} ὁ Ἰησοῦς εἰς τὴν κώμην,\FnParseFormGloss{c}{accusative feminine singular}{κώμη}{a village} ἀλλ’ ἦν\FnParse{d}{imperfect active indicative 3rd singular} ⸀ἔτι ἐν τῷ τόπῳ ὅπου ὑπήντησεν\FnParseFormGloss{e}{aorist active indicative 3rd singular}{ὑπαντάω}{to go out to meet} αὐτῷ ἡ Μάρθα.\FnParseFormGloss{f}{nominative feminine singular}{Μάρθα}{Martha} 
\MainTextVerseMark{11}{31}οἱ οὖν Ἰουδαῖοι οἱ ὄντες\FnParse{a}{present active participle nominative masculine plural} μετ’ αὐτῆς ἐν τῇ οἰκίᾳ καὶ παραμυθούμενοι\FnParseFormGloss{b}{present middle participle nominative masculine plural}{παραμυθέομαι}{to comfort} αὐτήν, ἰδόντες\FnParse{c}{aorist active participle nominative masculine plural} τὴν ⸀Μαριὰμ\FnParseFormGloss{d}{accusative feminine singular}{Μαριάμ}{Mary} ὅτι ταχέως\FnFormGloss{e}{ταχέως}{quickly} ἀνέστη\FnParse{f}{aorist active indicative 3rd singular} καὶ ἐξῆλθεν,\FnParse{g}{aorist active indicative 3rd singular} ἠκολούθησαν\FnParse{h}{aorist active indicative 3rd plural} αὐτῇ ⸀δόξαντες\FnParse{i}{aorist active participle nominative masculine plural} ὅτι ὑπάγει\FnParse{j}{present active indicative 3rd singular} εἰς τὸ μνημεῖον ἵνα κλαύσῃ\FnParse{k}{aorist active subjunctive 3rd singular} ἐκεῖ. 
\MainTextVerseMark{11}{32}ἡ οὖν ⸀Μαριὰμ\FnParseFormGloss{a}{nominative feminine singular}{Μαριάμ}{Mary} ὡς ἦλθεν\FnParse{b}{aorist active indicative 3rd singular} ὅπου ⸀ἦν\FnParse{c}{imperfect active indicative 3rd singular} Ἰησοῦς ἰδοῦσα\FnParse{d}{aorist active participle nominative feminine singular} αὐτὸν ἔπεσεν\FnParse{e}{aorist active indicative 3rd singular} αὐτοῦ ⸀πρὸς τοὺς πόδας, λέγουσα\FnParse{f}{present active participle nominative feminine singular} αὐτῷ· Κύριε, εἰ ἦς\FnParse{g}{imperfect active indicative 2nd singular} ὧδε οὐκ ἄν ⸂μου ἀπέθανεν⸃\FnParse{h}{aorist active indicative 3rd singular} ὁ ἀδελφός. 
\MainTextVerseMark{11}{33}Ἰησοῦς οὖν ὡς εἶδεν\FnParse{a}{aorist active indicative 3rd singular} αὐτὴν κλαίουσαν\FnParse{b}{present active participle accusative feminine singular} καὶ τοὺς συνελθόντας\FnParse{c}{aorist active participle accusative masculine plural} αὐτῇ Ἰουδαίους κλαίοντας\FnParse{d}{present active participle accusative masculine plural} ἐνεβριμήσατο\FnParseFormGloss{e}{aorist middle indicative 3rd singular}{ἐμβριμάομαι}{to warn sternly} τῷ πνεύματι καὶ ἐτάραξεν\FnParseFormGloss{f}{aorist active indicative 3rd singular}{ταράσσω}{to trouble} ἑαυτόν, 
\MainTextVerseMark{11}{34}καὶ εἶπεν·\FnParse{a}{aorist active indicative 3rd singular} Ποῦ τεθείκατε\FnParse{b}{perfect active indicative 2nd plural} αὐτόν; λέγουσιν\FnParse{c}{present active indicative 3rd plural} αὐτῷ· Κύριε, ἔρχου\FnParse{d}{present middle imperative 2nd singular} καὶ ἴδε.\FnParse{e}{aorist active imperative 2nd singular} 
\MainTextVerseMark{11}{35}ἐδάκρυσεν\FnParseFormGloss{a}{aorist active indicative 3rd singular}{δακρύω}{to weep} ὁ Ἰησοῦς. 
\MainTextVerseMark{11}{36}ἔλεγον\FnParse{a}{imperfect active indicative 3rd plural} οὖν οἱ Ἰουδαῖοι· Ἴδε πῶς ἐφίλει\FnParseFormGloss{b}{imperfect active indicative 3rd singular}{φιλέω}{to love} αὐτόν. 
\MainTextVerseMark{11}{37}τινὲς δὲ ἐξ αὐτῶν εἶπαν·\FnParse{a}{aorist active indicative 3rd plural} Οὐκ ἐδύνατο\FnParse{b}{imperfect middle indicative 3rd singular} οὗτος ὁ ἀνοίξας\FnParse{c}{aorist active participle nominative masculine singular} τοὺς ὀφθαλμοὺς τοῦ τυφλοῦ ποιῆσαι\FnParse{d}{aorist active infinitive} ἵνα καὶ οὗτος μὴ ἀποθάνῃ;\FnParse{e}{aorist active subjunctive 3rd singular} 
\MainTextVerseMark{11}{38}Ἰησοῦς οὖν πάλιν ἐμβριμώμενος\FnParseFormGloss{a}{present middle participle nominative masculine singular}{ἐμβριμάομαι}{to warn sternly} ἐν ἑαυτῷ ἔρχεται\FnParse{b}{present middle indicative 3rd singular} εἰς τὸ μνημεῖον· ἦν\FnParse{c}{imperfect active indicative 3rd singular} δὲ σπήλαιον,\FnParseFormGloss{d}{nominative neuter singular}{σπήλαιον}{den} καὶ λίθος ἐπέκειτο\FnParseFormGloss{e}{imperfect middle indicative 3rd singular}{ἐπίκειμαι}{to lay upon} ἐπ’ αὐτῷ. 
\MainTextVerseMark{11}{39}λέγει\FnParse{a}{present active indicative 3rd singular} ὁ Ἰησοῦς· Ἄρατε\FnParse{b}{aorist active imperative 2nd plural} τὸν λίθον. λέγει\FnParse{c}{present active indicative 3rd singular} αὐτῷ ἡ ἀδελφὴ\FnParseFormGloss{d}{nominative feminine singular}{ἀδελφή}{sister} τοῦ ⸀τετελευτηκότος\FnParseFormGloss{e}{perfect active participle genitive masculine singular}{τελευτάω}{to die} Μάρθα·\FnParseFormGloss{f}{nominative feminine singular}{Μάρθα}{Martha} Κύριε, ἤδη ὄζει,\FnParseFormGloss{g}{present active indicative 3rd singular}{ὄζω}{to give off a bad odor} τεταρταῖος\FnParseFormGloss{h}{nominative masculine singular}{τεταρταῖος}{fourth} γάρ ἐστιν.\FnParse{i}{present active indicative 3rd singular} 
\MainTextVerseMark{11}{40}λέγει\FnParse{a}{present active indicative 3rd singular} αὐτῇ ὁ Ἰησοῦς· Οὐκ εἶπόν\FnParse{b}{aorist active indicative 1st singular} σοι ὅτι ἐὰν πιστεύσῃς\FnParse{c}{aorist active subjunctive 2nd singular} ⸀ὄψῃ\FnParse{d}{future middle indicative 2nd singular} τὴν δόξαν τοῦ θεοῦ; 
\MainTextVerseMark{11}{41}ἦραν\FnParse{a}{aorist active indicative 3rd plural} οὖν τὸν ⸀λίθον. ὁ δὲ Ἰησοῦς ἦρεν\FnParse{b}{aorist active indicative 3rd singular} τοὺς ὀφθαλμοὺς ἄνω\FnFormGloss{c}{ἄνω}{above} καὶ εἶπεν·\FnParse{d}{aorist active indicative 3rd singular} Πάτερ, εὐχαριστῶ\FnParse{e}{present active indicative 1st singular} σοι ὅτι ἤκουσάς\FnParse{f}{aorist active indicative 2nd singular} μου, 
\MainTextVerseMark{11}{42}ἐγὼ δὲ ᾔδειν\FnParse{a}{pluperfect active indicative 1st singular} ὅτι πάντοτέ μου ἀκούεις·\FnParse{b}{present active indicative 2nd singular} ἀλλὰ διὰ τὸν ὄχλον τὸν περιεστῶτα\FnParseFormGloss{c}{perfect active participle accusative masculine singular}{περιΐστημι}{to stand around} εἶπον,\FnParse{d}{aorist active indicative 1st singular} ἵνα πιστεύσωσιν\FnParse{e}{aorist active subjunctive 3rd plural} ὅτι σύ με ἀπέστειλας.\FnParse{f}{aorist active indicative 2nd singular} 
\MainTextVerseMark{11}{43}καὶ ταῦτα εἰπὼν\FnParse{a}{aorist active participle nominative masculine singular} φωνῇ μεγάλῃ ἐκραύγασεν·\FnParseFormGloss{b}{aorist active indicative 3rd singular}{κραυγάζω}{to shout} Λάζαρε,\FnParseFormGloss{c}{masculine singular}{Λάζαρος}{Lazarus} δεῦρο\FnFormGloss{d}{δεῦρο}{come} ἔξω. 
\MainTextVerseMark{11}{44}⸀ἐξῆλθεν\FnParse{a}{aorist active indicative 3rd singular} ὁ τεθνηκὼς\FnParseFormGloss{b}{perfect active participle nominative masculine singular}{θνῄσκω}{to have died} δεδεμένος\FnParse{c}{perfect passive participle nominative masculine singular} τοὺς πόδας καὶ τὰς χεῖρας κειρίαις,\FnParseFormGloss{d}{dative feminine plural}{κειρία}{strip of linen} καὶ ἡ ὄψις\FnParseFormGloss{e}{nominative feminine singular}{ὄψις}{face} αὐτοῦ σουδαρίῳ\FnParseFormGloss{f}{dative neuter singular}{σουδάριον}{piece of cloth} περιεδέδετο.\FnParseFormGloss{g}{pluperfect passive indicative 3rd singular}{περιδέω}{to wrap around} λέγει\FnParse{h}{present active indicative 3rd singular} ⸂αὐτοῖς ὁ Ἰησοῦς⸃· Λύσατε\FnParse{i}{aorist active imperative 2nd plural} αὐτὸν καὶ ἄφετε\FnParse{j}{aorist active imperative 2nd plural} ⸀αὐτὸν ὑπάγειν.\FnParse{k}{present active infinitive} 
\MainTextVerseMark{11}{45}Πολλοὶ οὖν ἐκ τῶν Ἰουδαίων, οἱ ἐλθόντες\FnParse{a}{aorist active participle nominative masculine plural} πρὸς τὴν ⸀Μαριὰμ\FnParseFormGloss{b}{accusative feminine singular}{Μαριάμ}{Mary} καὶ θεασάμενοι\FnParseFormGloss{c}{aorist middle participle nominative masculine plural}{θεάομαι}{to see} ⸀ἃ ⸀ἐποίησεν,\FnParse{d}{aorist active indicative 3rd singular} ἐπίστευσαν\FnParse{e}{aorist active indicative 3rd plural} εἰς αὐτόν· 
\MainTextVerseMark{11}{46}τινὲς δὲ ἐξ αὐτῶν ἀπῆλθον\FnParse{a}{aorist active indicative 3rd plural} πρὸς τοὺς Φαρισαίους καὶ εἶπαν\FnParse{b}{aorist active indicative 3rd plural} αὐτοῖς ἃ ⸀ἐποίησεν\FnParse{c}{aorist active indicative 3rd singular} Ἰησοῦς. 
\MainTextVerseMark{11}{47}συνήγαγον\FnParse{a}{aorist active indicative 3rd plural} οὖν οἱ ἀρχιερεῖς καὶ οἱ Φαρισαῖοι συνέδριον,\FnParseFormGloss{b}{accusative neuter singular}{συνέδριον}{Sanhedrin} καὶ ἔλεγον·\FnParse{c}{imperfect active indicative 3rd plural} Τί ποιοῦμεν\FnParse{d}{present active indicative 1st plural} ὅτι οὗτος ὁ ἄνθρωπος πολλὰ ⸂ποιεῖ\FnParse{e}{present active indicative 3rd singular} σημεῖα⸃; 
\MainTextVerseMark{11}{48}ἐὰν ἀφῶμεν\FnParse{a}{aorist active subjunctive 1st plural} αὐτὸν οὕτως, πάντες πιστεύσουσιν\FnParse{b}{future active indicative 3rd plural} εἰς αὐτόν, καὶ ἐλεύσονται\FnParse{c}{future middle indicative 3rd plural} οἱ Ῥωμαῖοι\FnParseFormGloss{d}{nominative masculine plural}{Ῥωμαῖος}{Roman} καὶ ἀροῦσιν\FnParse{e}{future active indicative 3rd plural} ἡμῶν καὶ τὸν τόπον καὶ τὸ ἔθνος. 
\MainTextVerseMark{11}{49}εἷς δέ τις ἐξ αὐτῶν Καϊάφας,\FnParseFormGloss{a}{nominative masculine singular}{Καϊάφας}{Caiaphas} ἀρχιερεὺς ὢν\FnParse{b}{present active participle nominative masculine singular} τοῦ ἐνιαυτοῦ\FnParseFormGloss{c}{genitive masculine singular}{ἐνιαυτός}{year} ἐκείνου, εἶπεν\FnParse{d}{aorist active indicative 3rd singular} αὐτοῖς· Ὑμεῖς οὐκ οἴδατε\FnParse{e}{perfect active indicative 2nd plural} οὐδέν, 
\MainTextVerseMark{11}{50}οὐδὲ ⸀λογίζεσθε\FnParse{a}{present middle indicative 2nd plural} ὅτι συμφέρει\FnParseFormGloss{b}{present active indicative 3rd singular}{συμφέρω}{to bring together} ⸀ὑμῖν ἵνα εἷς ἄνθρωπος ἀποθάνῃ\FnParse{c}{aorist active subjunctive 3rd singular} ὑπὲρ τοῦ λαοῦ καὶ μὴ ὅλον τὸ ἔθνος ἀπόληται.\FnParse{d}{aorist middle subjunctive 3rd singular} 
\MainTextVerseMark{11}{51}τοῦτο δὲ ἀφ’ ἑαυτοῦ οὐκ εἶπεν,\FnParse{a}{aorist active indicative 3rd singular} ἀλλὰ ἀρχιερεὺς ὢν\FnParse{b}{present active participle nominative masculine singular} τοῦ ἐνιαυτοῦ\FnParseFormGloss{c}{genitive masculine singular}{ἐνιαυτός}{year} ἐκείνου ἐπροφήτευσεν\FnParseFormGloss{d}{aorist active indicative 3rd singular}{προφητεύω}{to prophesy} ὅτι ἔμελλεν\FnParse{e}{imperfect active indicative 3rd singular} Ἰησοῦς ἀποθνῄσκειν\FnParse{f}{present active infinitive} ὑπὲρ τοῦ ἔθνους, 
\MainTextVerseMark{11}{52}καὶ οὐχ ὑπὲρ τοῦ ἔθνους μόνον, ἀλλ’ ἵνα καὶ τὰ τέκνα τοῦ θεοῦ τὰ διεσκορπισμένα\FnParseFormGloss{a}{perfect passive participle accusative neuter plural}{διασκορπίζω}{to scatter} συναγάγῃ\FnParse{b}{aorist active subjunctive 3rd singular} εἰς ἕν. 
\MainTextVerseMark{11}{53}ἀπ’ ἐκείνης οὖν τῆς ἡμέρας ⸀ἐβουλεύσαντο\FnParseFormGloss{a}{aorist middle indicative 3rd plural}{βουλεύομαι}{to make plans} ἵνα ἀποκτείνωσιν\FnParse{b}{aorist active subjunctive 3rd plural} αὐτόν. 
\MainTextVerseMark{11}{54}⸂Ὁ οὖν Ἰησοῦς⸃ οὐκέτι παρρησίᾳ περιεπάτει\FnParse{a}{imperfect active indicative 3rd singular} ἐν τοῖς Ἰουδαίοις, ἀλλὰ ἀπῆλθεν\FnParse{b}{aorist active indicative 3rd singular} ἐκεῖθεν\FnFormGloss{c}{ἐκεῖθεν}{from there} εἰς τὴν χώραν\FnParseFormGloss{d}{accusative feminine singular}{χώρα}{country} ἐγγὺς τῆς ἐρήμου, εἰς Ἐφραὶμ\FnParseFormGloss{e}{accusative masculine singular}{Ἐφραίμ}{Ephraim} λεγομένην\FnParse{f}{present passive participle accusative feminine singular} πόλιν, κἀκεῖ\FnFormGloss{g}{κἀκεῖ}{and there} ⸀ἔμεινεν\FnParse{h}{aorist active indicative 3rd singular} μετὰ τῶν ⸀μαθητῶν. 
\MainTextVerseMark{11}{55}Ἦν\FnParse{a}{imperfect active indicative 3rd singular} δὲ ἐγγὺς τὸ πάσχα\FnParseFormGloss{b}{nominative neuter singular}{πάσχα}{Passover} τῶν Ἰουδαίων, καὶ ἀνέβησαν\FnParse{c}{aorist active indicative 3rd plural} πολλοὶ εἰς Ἱεροσόλυμα ἐκ τῆς χώρας\FnParseFormGloss{d}{genitive feminine singular}{χώρα}{country} πρὸ τοῦ πάσχα\FnParseFormGloss{e}{genitive neuter singular}{πάσχα}{Passover} ἵνα ἁγνίσωσιν\FnParseFormGloss{f}{aorist active subjunctive 3rd plural}{ἁγνίζω}{to purify} ἑαυτούς. 
\MainTextVerseMark{11}{56}ἐζήτουν\FnParse{a}{imperfect active indicative 3rd plural} οὖν τὸν Ἰησοῦν καὶ ἔλεγον\FnParse{b}{imperfect active indicative 3rd plural} μετ’ ἀλλήλων ἐν τῷ ἱερῷ ἑστηκότες·\FnParse{c}{perfect active participle nominative masculine plural} Τί δοκεῖ\FnParse{d}{present active indicative 3rd singular} ὑμῖν; ὅτι οὐ μὴ ἔλθῃ\FnParse{e}{aorist active subjunctive 3rd singular} εἰς τὴν ἑορτήν;\FnParseFormGloss{f}{accusative feminine singular}{ἑορτή}{feast} 
\MainTextVerseMark{11}{57}δεδώκεισαν\FnParse{a}{pluperfect active indicative 3rd plural} ⸀δὲ οἱ ἀρχιερεῖς καὶ οἱ Φαρισαῖοι ⸀ἐντολὰς ἵνα ἐάν τις γνῷ\FnParse{b}{aorist active subjunctive 3rd singular} ποῦ ἐστιν\FnParse{c}{present active indicative 3rd singular} μηνύσῃ,\FnParseFormGloss{d}{aorist active subjunctive 3rd singular}{μηνύω}{to inform} ὅπως πιάσωσιν\FnParseFormGloss{e}{aorist active subjunctive 3rd plural}{πιάζω}{to seize} αὐτόν. 
\ChapterHeader{12}{Chapter 12}

\MainTextVerseMark{12}{1}Ὁ οὖν Ἰησοῦς πρὸ ἓξ\FnParseFormGloss{a}{genitive feminine plural}{ἕξ}{six} ἡμερῶν τοῦ πάσχα\FnParseFormGloss{b}{genitive neuter singular}{πάσχα}{Passover} ἦλθεν\FnParse{c}{aorist active indicative 3rd singular} εἰς Βηθανίαν,\FnParseFormGloss{d}{accusative feminine singular}{Βηθανία}{Bethany} ὅπου ἦν\FnParse{e}{imperfect active indicative 3rd singular} ⸀Λάζαρος,\FnParseFormGloss{f}{nominative masculine singular}{Λάζαρος}{Lazarus} ὃν ἤγειρεν\FnParse{g}{aorist active indicative 3rd singular} ἐκ νεκρῶν ⸀Ἰησοῦς. 
\MainTextVerseMark{12}{2}ἐποίησαν\FnParse{a}{aorist active indicative 3rd plural} οὖν αὐτῷ δεῖπνον\FnParseFormGloss{b}{accusative neuter singular}{δεῖπνον}{banquet} ἐκεῖ, καὶ ἡ Μάρθα\FnParseFormGloss{c}{nominative feminine singular}{Μάρθα}{Martha} διηκόνει,\FnParse{d}{imperfect active indicative 3rd singular} ὁ δὲ Λάζαρος\FnParseFormGloss{e}{nominative masculine singular}{Λάζαρος}{Lazarus} εἷς ἦν\FnParse{f}{imperfect active indicative 3rd singular} ⸀ἐκ τῶν ἀνακειμένων\FnParseFormGloss{g}{present middle participle genitive masculine plural}{ἀνάκειμαι}{to recline for a meal} σὺν αὐτῷ· 
\MainTextVerseMark{12}{3}ἡ οὖν ⸀Μαριὰμ\FnParseFormGloss{a}{nominative feminine singular}{Μαριάμ}{Mary} λαβοῦσα\FnParse{b}{aorist active participle nominative feminine singular} λίτραν\FnParseFormGloss{c}{accusative feminine singular}{λίτρα}{pound} μύρου\FnParseFormGloss{d}{genitive neuter singular}{μύρον}{perfume} νάρδου\FnParseFormGloss{e}{genitive feminine singular}{νάρδος}{nard} πιστικῆς\FnParseFormGloss{f}{genitive feminine singular}{πιστικός}{pure} πολυτίμου\FnParseFormGloss{g}{genitive feminine singular}{πολύτιμος}{expensive} ἤλειψεν\FnParseFormGloss{h}{aorist active indicative 3rd singular}{ἀλείφω}{to pour on} τοὺς πόδας τοῦ Ἰησοῦ καὶ ἐξέμαξεν\FnParseFormGloss{i}{aorist active indicative 3rd singular}{ἐκμάσσω}{to wipe off} ταῖς θριξὶν\FnParseFormGloss{j}{dative feminine plural}{θρίξ}{hair} αὐτῆς τοὺς πόδας αὐτοῦ· ἡ δὲ οἰκία ἐπληρώθη\FnParse{k}{aorist passive indicative 3rd singular} ἐκ τῆς ὀσμῆς\FnParseFormGloss{l}{genitive feminine singular}{ὀσμή}{fragrance} τοῦ μύρου.\FnParseFormGloss{m}{genitive neuter singular}{μύρον}{perfume} 
\MainTextVerseMark{12}{4}λέγει\FnParse{a}{present active indicative 3rd singular} ⸀δὲ ⸂Ἰούδας ὁ Ἰσκαριώτης\FnParseFormGloss{b}{nominative masculine singular}{Ἰσκαριώτης}{Iscariot} εἷς τῶν μαθητῶν αὐτοῦ⸃, ὁ μέλλων\FnParse{c}{present active participle nominative masculine singular} αὐτὸν παραδιδόναι·\FnParse{d}{present active infinitive} 
\MainTextVerseMark{12}{5}Διὰ τί τοῦτο τὸ μύρον\FnParseFormGloss{a}{nominative neuter singular}{μύρον}{perfume} οὐκ ἐπράθη\FnParseFormGloss{b}{aorist passive indicative 3rd singular}{πιπράσκω}{to sell} τριακοσίων\FnParseFormGloss{c}{genitive neuter plural}{τριακόσιοι}{three hundred} δηναρίων\FnParseFormGloss{d}{genitive neuter plural}{δηνάριον}{denarius} καὶ ἐδόθη\FnParse{e}{aorist passive indicative 3rd singular} πτωχοῖς; 
\MainTextVerseMark{12}{6}εἶπεν\FnParse{a}{aorist active indicative 3rd singular} δὲ τοῦτο οὐχ ὅτι περὶ τῶν πτωχῶν ἔμελεν\FnParseFormGloss{b}{imperfect active indicative 3rd singular}{μέλει}{it is a care} αὐτῷ, ἀλλ’ ὅτι κλέπτης\FnParseFormGloss{c}{nominative masculine singular}{κλέπτης}{thief} ἦν\FnParse{d}{imperfect active indicative 3rd singular} καὶ τὸ γλωσσόκομον\FnParseFormGloss{e}{accusative neuter singular}{γλωσσόκομον}{container for money} ⸀ἔχων\FnParse{f}{present active participle nominative masculine singular} τὰ βαλλόμενα\FnParse{g}{present passive participle accusative neuter plural} ἐβάσταζεν.\FnParseFormGloss{h}{imperfect active indicative 3rd singular}{βαστάζω}{to carry} 
\MainTextVerseMark{12}{7}εἶπεν\FnParse{a}{aorist active indicative 3rd singular} οὖν ὁ Ἰησοῦς· Ἄφες\FnParse{b}{aorist active imperative 2nd singular} αὐτήν, ⸀ἵνα εἰς τὴν ἡμέραν τοῦ ἐνταφιασμοῦ\FnParseFormGloss{c}{genitive masculine singular}{ἐνταφιασμός}{preparation for burial} μου ⸀τηρήσῃ\FnParse{d}{aorist active subjunctive 3rd singular} αὐτό· 
\MainTextVerseMark{12}{8}τοὺς πτωχοὺς γὰρ πάντοτε ἔχετε\FnParse{a}{present active indicative 2nd plural} μεθ’ ἑαυτῶν, ἐμὲ δὲ οὐ πάντοτε ἔχετε.\FnParse{b}{present active indicative 2nd plural} 
\MainTextVerseMark{12}{9}Ἔγνω\FnParse{a}{aorist active indicative 3rd singular} ⸀οὖν ὄχλος πολὺς ἐκ τῶν Ἰουδαίων ὅτι ἐκεῖ ἐστιν,\FnParse{b}{present active indicative 3rd singular} καὶ ἦλθον\FnParse{c}{aorist active indicative 3rd plural} οὐ διὰ τὸν Ἰησοῦν μόνον, ἀλλ’ ἵνα καὶ τὸν Λάζαρον\FnParseFormGloss{d}{accusative masculine singular}{Λάζαρος}{Lazarus} ἴδωσιν\FnParse{e}{aorist active subjunctive 3rd plural} ὃν ἤγειρεν\FnParse{f}{aorist active indicative 3rd singular} ἐκ νεκρῶν. 
\MainTextVerseMark{12}{10}ἐβουλεύσαντο\FnParseFormGloss{a}{aorist middle indicative 3rd plural}{βουλεύομαι}{to make plans} δὲ οἱ ἀρχιερεῖς ἵνα καὶ τὸν Λάζαρον\FnParseFormGloss{b}{accusative masculine singular}{Λάζαρος}{Lazarus} ἀποκτείνωσιν,\FnParse{c}{aorist active subjunctive 3rd plural} 
\MainTextVerseMark{12}{11}ὅτι πολλοὶ δι’ αὐτὸν ὑπῆγον\FnParse{a}{imperfect active indicative 3rd plural} τῶν Ἰουδαίων καὶ ἐπίστευον\FnParse{b}{imperfect active indicative 3rd plural} εἰς τὸν Ἰησοῦν. 
\MainTextVerseMark{12}{12}Τῇ ἐπαύριον\FnFormGloss{a}{ἐπαύριον}{the next day} ⸀ὁ ὄχλος πολὺς ὁ ἐλθὼν\FnParse{b}{aorist active participle nominative masculine singular} εἰς τὴν ἑορτήν,\FnParseFormGloss{c}{accusative feminine singular}{ἑορτή}{feast} ἀκούσαντες\FnParse{d}{aorist active participle nominative masculine plural} ὅτι ἔρχεται\FnParse{e}{present middle indicative 3rd singular} ⸁ὁ Ἰησοῦς εἰς Ἱεροσόλυμα, 
\MainTextVerseMark{12}{13}ἔλαβον\FnParse{a}{aorist active indicative 3rd plural} τὰ βαΐα\FnParseFormGloss{b}{accusative neuter plural}{βάϊον}{branch} τῶν φοινίκων\FnParseFormGloss{c}{genitive masculine plural}{φοῖνιξ}{palm tree} καὶ ἐξῆλθον\FnParse{d}{aorist active indicative 3rd plural} εἰς ὑπάντησιν\FnParseFormGloss{e}{accusative feminine singular}{ὑπάντησις}{meeting} αὐτῷ, καὶ ⸀ἐκραύγαζον·\FnParseFormGloss{f}{imperfect active indicative 3rd plural}{κραυγάζω}{to shout} Ὡσαννά,\FnFormGloss{g}{ὡσαννά}{Hosanna!} εὐλογημένος\FnParse{h}{perfect passive participle nominative masculine singular} ὁ ἐρχόμενος\FnParse{i}{present middle participle nominative masculine singular} ἐν ὀνόματι κυρίου, ⸂καὶ ὁ⸃ βασιλεὺς τοῦ Ἰσραήλ. 
\MainTextVerseMark{12}{14}εὑρὼν\FnParse{a}{aorist active participle nominative masculine singular} δὲ ὁ Ἰησοῦς ὀνάριον\FnParseFormGloss{b}{accusative neuter singular}{ὀνάριον}{young donkey} ἐκάθισεν\FnParse{c}{aorist active indicative 3rd singular} ἐπ’ αὐτό, καθώς ἐστιν\FnParse{d}{present active indicative 3rd singular} γεγραμμένον·\FnParse{e}{perfect passive participle nominative neuter singular} 
\MainTextVerseMark{12}{15}Μὴ φοβοῦ,\FnParse{a}{present middle imperative 2nd singular} θυγάτηρ\FnParseFormGloss{b}{feminine singular}{θυγάτηρ}{daughter} Σιών·\FnParseFormGloss{c}{genitive feminine singular}{Σιών}{Zion} ἰδοὺ ὁ βασιλεύς σου ἔρχεται,\FnParse{d}{present middle indicative 3rd singular} καθήμενος\FnParse{e}{present middle participle nominative masculine singular} ἐπὶ πῶλον\FnParseFormGloss{f}{accusative masculine singular}{πῶλος}{colt} ὄνου.\FnParseFormGloss{g}{genitive feminine singular}{ὄνος}{donkey} 
\MainTextVerseMark{12}{16}⸀ταῦτα οὐκ ἔγνωσαν\FnParse{a}{aorist active indicative 3rd plural} ⸂αὐτοῦ οἱ μαθηταὶ⸃ τὸ πρῶτον, ἀλλ’ ὅτε ἐδοξάσθη\FnParse{b}{aorist passive indicative 3rd singular} Ἰησοῦς τότε ἐμνήσθησαν\FnParseFormGloss{c}{aorist passive indicative 3rd plural}{μιμνῄσκομαι}{to remember} ὅτι ταῦτα ἦν\FnParse{d}{imperfect active indicative 3rd singular} ἐπ’ αὐτῷ γεγραμμένα\FnParse{e}{perfect passive participle nominative neuter plural} καὶ ταῦτα ἐποίησαν\FnParse{f}{aorist active indicative 3rd plural} αὐτῷ. 
\MainTextVerseMark{12}{17}ἐμαρτύρει\FnParse{a}{imperfect active indicative 3rd singular} οὖν ὁ ὄχλος ὁ ὢν\FnParse{b}{present active participle nominative masculine singular} μετ’ αὐτοῦ ὅτε τὸν Λάζαρον\FnParseFormGloss{c}{accusative masculine singular}{Λάζαρος}{Lazarus} ἐφώνησεν\FnParse{d}{aorist active indicative 3rd singular} ἐκ τοῦ μνημείου καὶ ἤγειρεν\FnParse{e}{aorist active indicative 3rd singular} αὐτὸν ἐκ νεκρῶν. 
\MainTextVerseMark{12}{18}διὰ τοῦτο ⸀καὶ ὑπήντησεν\FnParseFormGloss{a}{aorist active indicative 3rd singular}{ὑπαντάω}{to go out to meet} αὐτῷ ὁ ὄχλος ὅτι ⸀ἤκουσαν\FnParse{b}{aorist active indicative 3rd plural} τοῦτο αὐτὸν πεποιηκέναι\FnParse{c}{perfect active infinitive} τὸ σημεῖον. 
\MainTextVerseMark{12}{19}οἱ οὖν Φαρισαῖοι εἶπαν\FnParse{a}{aorist active indicative 3rd plural} πρὸς ἑαυτούς· Θεωρεῖτε\FnParse{b}{present active indicative 2nd plural} ὅτι οὐκ ὠφελεῖτε\FnParseFormGloss{c}{present active indicative 2nd plural}{ὠφελέω}{to be of good use} οὐδέν· ἴδε ὁ ⸀κόσμος ὀπίσω αὐτοῦ ἀπῆλθεν.\FnParse{d}{aorist active indicative 3rd singular} 
\MainTextVerseMark{12}{20}Ἦσαν\FnParse{a}{imperfect active indicative 3rd plural} δὲ ⸂Ἕλληνές\FnParseFormGloss{b}{nominative masculine plural}{Ἕλλην}{Greek} τινες⸃ ἐκ τῶν ἀναβαινόντων\FnParse{c}{present active participle genitive masculine plural} ἵνα προσκυνήσωσιν\FnParse{d}{aorist active subjunctive 3rd plural} ἐν τῇ ἑορτῇ·\FnParseFormGloss{e}{dative feminine singular}{ἑορτή}{feast} 
\MainTextVerseMark{12}{21}οὗτοι οὖν προσῆλθον\FnParse{a}{aorist active indicative 3rd plural} Φιλίππῳ τῷ ἀπὸ Βηθσαϊδὰ\FnParseFormGloss{b}{genitive feminine singular}{Βηθσαϊδά}{Bethsaida} τῆς Γαλιλαίας, καὶ ἠρώτων\FnParse{c}{imperfect active indicative 3rd plural} αὐτὸν λέγοντες·\FnParse{d}{present active participle nominative masculine plural} Κύριε, θέλομεν\FnParse{e}{present active indicative 1st plural} τὸν Ἰησοῦν ἰδεῖν.\FnParse{f}{aorist active infinitive} 
\MainTextVerseMark{12}{22}ἔρχεται\FnParse{a}{present middle indicative 3rd singular} ⸀ὁ Φίλιππος καὶ λέγει\FnParse{b}{present active indicative 3rd singular} τῷ Ἀνδρέᾳ·\FnParseFormGloss{c}{dative masculine singular}{Ἀνδρέας}{Andrew} ⸀ἔρχεται\FnParse{d}{present middle indicative 3rd singular} Ἀνδρέας\FnParseFormGloss{e}{nominative masculine singular}{Ἀνδρέας}{Andrew} καὶ Φίλιππος ⸀καὶ λέγουσιν\FnParse{f}{present active indicative 3rd plural} τῷ Ἰησοῦ. 
\MainTextVerseMark{12}{23}ὁ δὲ Ἰησοῦς ⸀ἀποκρίνεται\FnParse{a}{present middle indicative 3rd singular} αὐτοῖς λέγων·\FnParse{b}{present active participle nominative masculine singular} Ἐλήλυθεν\FnParse{c}{perfect active indicative 3rd singular} ἡ ὥρα ἵνα δοξασθῇ\FnParse{d}{aorist passive subjunctive 3rd singular} ὁ υἱὸς τοῦ ἀνθρώπου. 
\MainTextVerseMark{12}{24}ἀμὴν ἀμὴν λέγω\FnParse{a}{present active indicative 1st singular} ὑμῖν, ἐὰν μὴ ὁ κόκκος\FnParseFormGloss{b}{nominative masculine singular}{κόκκος}{seed} τοῦ σίτου\FnParseFormGloss{c}{genitive masculine singular}{σῖτος}{wheat} πεσὼν\FnParse{d}{aorist active participle nominative masculine singular} εἰς τὴν γῆν ἀποθάνῃ,\FnParse{e}{aorist active subjunctive 3rd singular} αὐτὸς μόνος μένει·\FnParse{f}{present active indicative 3rd singular} ἐὰν δὲ ἀποθάνῃ,\FnParse{g}{aorist active subjunctive 3rd singular} πολὺν καρπὸν φέρει.\FnParse{h}{present active indicative 3rd singular} 
\MainTextVerseMark{12}{25}ὁ φιλῶν\FnParseFormGloss{a}{present active participle nominative masculine singular}{φιλέω}{to love} τὴν ψυχὴν αὐτοῦ ⸀ἀπολλύει\FnParse{b}{present active indicative 3rd singular} αὐτήν, καὶ ὁ μισῶν\FnParse{c}{present active participle nominative masculine singular} τὴν ψυχὴν αὐτοῦ ἐν τῷ κόσμῳ τούτῳ εἰς ζωὴν αἰώνιον φυλάξει\FnParse{d}{future active indicative 3rd singular} αὐτήν. 
\MainTextVerseMark{12}{26}ἐὰν ἐμοί ⸂τις διακονῇ⸃\FnParse{a}{present active subjunctive 3rd singular} ἐμοὶ ἀκολουθείτω,\FnParse{b}{present active imperative 3rd singular} καὶ ὅπου εἰμὶ\FnParse{c}{present active indicative 1st singular} ἐγὼ ἐκεῖ καὶ ὁ διάκονος\FnParseFormGloss{d}{nominative masculine singular}{διάκονος}{servant} ὁ ἐμὸς ἔσται·\FnParse{e}{future middle indicative 3rd singular} ⸀ἐάν τις ἐμοὶ διακονῇ\FnParse{f}{present active subjunctive 3rd singular} τιμήσει\FnParseFormGloss{g}{future active indicative 3rd singular}{τιμάω}{to honor} αὐτὸν ὁ πατήρ. 
\MainTextVerseMark{12}{27}Νῦν ἡ ψυχή μου τετάρακται,\FnParseFormGloss{a}{perfect passive indicative 3rd singular}{ταράσσω}{to trouble} καὶ τί εἴπω;\FnParse{b}{aorist active subjunctive 1st singular} πάτερ, σῶσόν\FnParse{c}{aorist active imperative 2nd singular} με ἐκ τῆς ὥρας ταύτης. ἀλλὰ διὰ τοῦτο ἦλθον\FnParse{d}{aorist active indicative 1st singular} εἰς τὴν ὥραν ταύτην. 
\MainTextVerseMark{12}{28}πάτερ, δόξασόν\FnParse{a}{aorist active imperative 2nd singular} σου τὸ ὄνομα. ἦλθεν\FnParse{b}{aorist active indicative 3rd singular} οὖν φωνὴ ἐκ τοῦ οὐρανοῦ· Καὶ ἐδόξασα\FnParse{c}{aorist active indicative 1st singular} καὶ πάλιν δοξάσω.\FnParse{d}{future active indicative 1st singular} 
\MainTextVerseMark{12}{29}ὁ οὖν ὄχλος ὁ ἑστὼς\FnParse{a}{perfect active participle nominative masculine singular} καὶ ἀκούσας\FnParse{b}{aorist active participle nominative masculine singular} ἔλεγεν\FnParse{c}{imperfect active indicative 3rd singular} βροντὴν\FnParseFormGloss{d}{accusative feminine singular}{βροντή}{thunder} γεγονέναι·\FnParse{e}{perfect active infinitive} ἄλλοι ἔλεγον·\FnParse{f}{imperfect active indicative 3rd plural} Ἄγγελος αὐτῷ λελάληκεν.\FnParse{g}{perfect active indicative 3rd singular} 
\MainTextVerseMark{12}{30}ἀπεκρίθη\FnParse{a}{aorist passive indicative 3rd singular} ⸂Ἰησοῦς καὶ εἶπεν⸃·\FnParse{b}{aorist active indicative 3rd singular} Οὐ δι’ ἐμὲ ⸂ἡ φωνὴ αὕτη⸃ γέγονεν\FnParse{c}{perfect active indicative 3rd singular} ἀλλὰ δι’ ὑμᾶς. 
\MainTextVerseMark{12}{31}νῦν κρίσις ἐστὶν\FnParse{a}{present active indicative 3rd singular} τοῦ κόσμου τούτου, νῦν ὁ ἄρχων τοῦ κόσμου τούτου ἐκβληθήσεται\FnParse{b}{future passive indicative 3rd singular} ἔξω· 
\MainTextVerseMark{12}{32}κἀγὼ ⸀ἐὰν ὑψωθῶ\FnParseFormGloss{a}{aorist passive subjunctive 1st singular}{ὑψόω}{to lift up} ἐκ τῆς γῆς, πάντας ἑλκύσω\FnParseFormGloss{b}{future active indicative 1st singular}{ἕλκω}{to drag} πρὸς ἐμαυτόν. 
\MainTextVerseMark{12}{33}τοῦτο δὲ ἔλεγεν\FnParse{a}{imperfect active indicative 3rd singular} σημαίνων\FnParseFormGloss{b}{present active participle nominative masculine singular}{σημαίνω}{to make known} ποίῳ θανάτῳ ἤμελλεν\FnParse{c}{imperfect active indicative 3rd singular} ἀποθνῄσκειν.\FnParse{d}{present active infinitive} 
\MainTextVerseMark{12}{34}ἀπεκρίθη\FnParse{a}{aorist passive indicative 3rd singular} ⸀οὖν αὐτῷ ὁ ὄχλος· Ἡμεῖς ἠκούσαμεν\FnParse{b}{aorist active indicative 1st plural} ἐκ τοῦ νόμου ὅτι ὁ χριστὸς μένει\FnParse{c}{present active indicative 3rd singular} εἰς τὸν αἰῶνα, καὶ πῶς ⸂λέγεις\FnParse{d}{present active indicative 2nd singular} σὺ⸃ ⸀ὅτι δεῖ\FnParse{e}{present active indicative 3rd singular} ὑψωθῆναι\FnParseFormGloss{f}{aorist passive infinitive}{ὑψόω}{to lift up} τὸν υἱὸν τοῦ ἀνθρώπου; τίς ἐστιν\FnParse{g}{present active indicative 3rd singular} οὗτος ὁ υἱὸς τοῦ ἀνθρώπου; 
\MainTextVerseMark{12}{35}εἶπεν\FnParse{a}{aorist active indicative 3rd singular} οὖν αὐτοῖς ὁ Ἰησοῦς· Ἔτι μικρὸν χρόνον τὸ φῶς ⸂ἐν ὑμῖν⸃ ἐστιν.\FnParse{b}{present active indicative 3rd singular} περιπατεῖτε\FnParse{c}{present active imperative 2nd plural} ⸀ὡς τὸ φῶς ἔχετε,\FnParse{d}{present active indicative 2nd plural} ἵνα μὴ σκοτία\FnParseFormGloss{e}{nominative feminine singular}{σκοτία}{darkness} ὑμᾶς καταλάβῃ,\FnParseFormGloss{f}{aorist active subjunctive 3rd singular}{καταλαμβάνω}{to obtain} καὶ ὁ περιπατῶν\FnParse{g}{present active participle nominative masculine singular} ἐν τῇ σκοτίᾳ\FnParseFormGloss{h}{dative feminine singular}{σκοτία}{darkness} οὐκ οἶδεν\FnParse{i}{perfect active indicative 3rd singular} ποῦ ὑπάγει.\FnParse{j}{present active indicative 3rd singular} 
\MainTextVerseMark{12}{36}⸀ὡς τὸ φῶς ἔχετε,\FnParse{a}{present active indicative 2nd plural} πιστεύετε\FnParse{b}{present active imperative 2nd plural} εἰς τὸ φῶς, ἵνα υἱοὶ φωτὸς γένησθε.\FnParse{c}{aorist middle subjunctive 2nd plural} Ταῦτα ⸀ἐλάλησεν\FnParse{d}{aorist active indicative 3rd singular} Ἰησοῦς, καὶ ἀπελθὼν\FnParse{e}{aorist active participle nominative masculine singular} ἐκρύβη\FnParseFormGloss{f}{aorist passive indicative 3rd singular}{κρύπτω}{to hide} ἀπ’ αὐτῶν. 
\MainTextVerseMark{12}{37}τοσαῦτα\FnParseFormGloss{a}{accusative neuter plural}{τοσοῦτος}{so great} δὲ αὐτοῦ σημεῖα πεποιηκότος\FnParse{b}{perfect active participle genitive masculine singular} ἔμπροσθεν αὐτῶν οὐκ ἐπίστευον\FnParse{c}{imperfect active indicative 3rd plural} εἰς αὐτόν, 
\MainTextVerseMark{12}{38}ἵνα ὁ λόγος Ἠσαΐου\FnParseFormGloss{a}{genitive masculine singular}{Ἠσαΐας}{Isaiah} τοῦ προφήτου πληρωθῇ\FnParse{b}{aorist passive subjunctive 3rd singular} ὃν εἶπεν·\FnParse{c}{aorist active indicative 3rd singular} Κύριε, τίς ἐπίστευσεν\FnParse{d}{aorist active indicative 3rd singular} τῇ ἀκοῇ\FnParseFormGloss{e}{dative feminine singular}{ἀκοή}{hearing} ἡμῶν; καὶ ὁ βραχίων\FnParseFormGloss{f}{nominative masculine singular}{βραχίων}{arm} κυρίου τίνι ἀπεκαλύφθη;\FnParseFormGloss{g}{aorist passive indicative 3rd singular}{ἀποκαλύπτω}{to reveal} 
\MainTextVerseMark{12}{39}διὰ τοῦτο οὐκ ἠδύναντο\FnParse{a}{imperfect middle indicative 3rd plural} πιστεύειν\FnParse{b}{present active infinitive} ὅτι πάλιν εἶπεν\FnParse{c}{aorist active indicative 3rd singular} Ἠσαΐας·\FnParseFormGloss{d}{nominative masculine singular}{Ἠσαΐας}{Isaiah} 
\MainTextVerseMark{12}{40}Τετύφλωκεν\FnParseFormGloss{a}{perfect active indicative 3rd singular}{τυφλόω}{to cause blindness} αὐτῶν τοὺς ὀφθαλμοὺς καὶ ⸀ἐπώρωσεν\FnParseFormGloss{b}{aorist active indicative 3rd singular}{πωρόω}{to harden} αὐτῶν τὴν καρδίαν, ἵνα μὴ ἴδωσιν\FnParse{c}{aorist active subjunctive 3rd plural} τοῖς ὀφθαλμοῖς καὶ νοήσωσιν\FnParseFormGloss{d}{aorist active subjunctive 3rd plural}{νοέω}{to understand} τῇ καρδίᾳ καὶ ⸀στραφῶσιν,\FnParseFormGloss{e}{aorist passive subjunctive 3rd plural}{στρέφω}{to turn} καὶ ⸀ἰάσομαι\FnParseFormGloss{f}{future middle indicative 1st singular}{ἰάομαι}{to heal} αὐτούς. 
\MainTextVerseMark{12}{41}ταῦτα εἶπεν\FnParse{a}{aorist active indicative 3rd singular} Ἠσαΐας\FnParseFormGloss{b}{nominative masculine singular}{Ἠσαΐας}{Isaiah} ⸀ὅτι εἶδεν\FnParse{c}{aorist active indicative 3rd singular} τὴν δόξαν αὐτοῦ, καὶ ἐλάλησεν\FnParse{d}{aorist active indicative 3rd singular} περὶ αὐτοῦ. 
\MainTextVerseMark{12}{42}ὅμως\FnFormGloss{a}{ὅμως}{just as} μέντοι\FnFormGloss{b}{μέντοι}{but} καὶ ἐκ τῶν ἀρχόντων πολλοὶ ἐπίστευσαν\FnParse{c}{aorist active indicative 3rd plural} εἰς αὐτόν, ἀλλὰ διὰ τοὺς Φαρισαίους οὐχ ὡμολόγουν\FnParseFormGloss{d}{imperfect active indicative 3rd plural}{ὁμολογέω}{to confess} ἵνα μὴ ἀποσυνάγωγοι\FnParseFormGloss{e}{nominative masculine plural}{ἀποσυνάγωγος}{put out of the synagogue} γένωνται,\FnParse{f}{aorist middle subjunctive 3rd plural} 
\MainTextVerseMark{12}{43}ἠγάπησαν\FnParse{a}{aorist active indicative 3rd plural} γὰρ τὴν δόξαν τῶν ἀνθρώπων μᾶλλον ἤπερ\FnFormGloss{b}{ἤπερ}{than} τὴν δόξαν τοῦ θεοῦ. 
\MainTextVerseMark{12}{44}Ἰησοῦς δὲ ἔκραξεν\FnParse{a}{aorist active indicative 3rd singular} καὶ εἶπεν·\FnParse{b}{aorist active indicative 3rd singular} Ὁ πιστεύων\FnParse{c}{present active participle nominative masculine singular} εἰς ἐμὲ οὐ πιστεύει\FnParse{d}{present active indicative 3rd singular} εἰς ἐμὲ ἀλλὰ εἰς τὸν πέμψαντά\FnParse{e}{aorist active participle accusative masculine singular} με, 
\MainTextVerseMark{12}{45}καὶ ὁ θεωρῶν\FnParse{a}{present active participle nominative masculine singular} ἐμὲ θεωρεῖ\FnParse{b}{present active indicative 3rd singular} τὸν πέμψαντά\FnParse{c}{aorist active participle accusative masculine singular} με. 
\MainTextVerseMark{12}{46}ἐγὼ φῶς εἰς τὸν κόσμον ἐλήλυθα,\FnParse{a}{perfect active indicative 1st singular} ἵνα πᾶς ὁ πιστεύων\FnParse{b}{present active participle nominative masculine singular} εἰς ἐμὲ ἐν τῇ σκοτίᾳ\FnParseFormGloss{c}{dative feminine singular}{σκοτία}{darkness} μὴ μείνῃ.\FnParse{d}{aorist active subjunctive 3rd singular} 
\MainTextVerseMark{12}{47}καὶ ἐάν τίς μου ἀκούσῃ\FnParse{a}{aorist active subjunctive 3rd singular} τῶν ῥημάτων καὶ μὴ ⸀φυλάξῃ,\FnParse{b}{aorist active subjunctive 3rd singular} ἐγὼ οὐ κρίνω\FnParse{c}{present active indicative 1st singular} αὐτόν, οὐ γὰρ ἦλθον\FnParse{d}{aorist active indicative 1st singular} ἵνα κρίνω\FnParse{e}{aorist active subjunctive 1st singular} τὸν κόσμον ἀλλ’ ἵνα σώσω\FnParse{f}{aorist active subjunctive 1st singular} τὸν κόσμον. 
\MainTextVerseMark{12}{48}ὁ ἀθετῶν\FnParseFormGloss{a}{present active participle nominative masculine singular}{ἀθετέω}{to reject} ἐμὲ καὶ μὴ λαμβάνων\FnParse{b}{present active participle nominative masculine singular} τὰ ῥήματά μου ἔχει\FnParse{c}{present active indicative 3rd singular} τὸν κρίνοντα\FnParse{d}{present active participle accusative masculine singular} αὐτόν· ὁ λόγος ὃν ἐλάλησα\FnParse{e}{aorist active indicative 1st singular} ἐκεῖνος κρινεῖ\FnParse{f}{future active indicative 3rd singular} αὐτὸν ἐν τῇ ἐσχάτῃ ἡμέρᾳ· 
\MainTextVerseMark{12}{49}ὅτι ἐγὼ ἐξ ἐμαυτοῦ οὐκ ἐλάλησα,\FnParse{a}{aorist active indicative 1st singular} ἀλλ’ ὁ πέμψας\FnParse{b}{aorist active participle nominative masculine singular} με πατὴρ αὐτός μοι ἐντολὴν ⸀δέδωκεν\FnParse{c}{perfect active indicative 3rd singular} τί εἴπω\FnParse{d}{aorist active subjunctive 1st singular} καὶ τί λαλήσω.\FnParse{e}{aorist active subjunctive 1st singular} 
\MainTextVerseMark{12}{50}καὶ οἶδα\FnParse{a}{perfect active indicative 1st singular} ὅτι ἡ ἐντολὴ αὐτοῦ ζωὴ αἰώνιός ἐστιν.\FnParse{b}{present active indicative 3rd singular} ἃ οὖν ⸂ἐγὼ λαλῶ⸃,\FnParse{c}{present active indicative 1st singular} καθὼς εἴρηκέν\FnParse{d}{perfect active indicative 3rd singular} μοι ὁ πατήρ, οὕτως λαλῶ.\FnParse{e}{present active indicative 1st singular} 
\ChapterHeader{13}{Chapter 13}

\MainTextVerseMark{13}{1}Πρὸ δὲ τῆς ἑορτῆς\FnParseFormGloss{a}{genitive feminine singular}{ἑορτή}{feast} τοῦ πάσχα\FnParseFormGloss{b}{genitive neuter singular}{πάσχα}{Passover} εἰδὼς\FnParse{c}{perfect active participle nominative masculine singular} ὁ Ἰησοῦς ὅτι ⸀ἦλθεν\FnParse{d}{aorist active indicative 3rd singular} αὐτοῦ ἡ ὥρα ἵνα μεταβῇ\FnParseFormGloss{e}{aorist active subjunctive 3rd singular}{μεταβαίνω}{to go on} ἐκ τοῦ κόσμου τούτου πρὸς τὸν πατέρα ἀγαπήσας\FnParse{f}{aorist active participle nominative masculine singular} τοὺς ἰδίους τοὺς ἐν τῷ κόσμῳ εἰς τέλος ἠγάπησεν\FnParse{g}{aorist active indicative 3rd singular} αὐτούς. 
\MainTextVerseMark{13}{2}καὶ δείπνου\FnParseFormGloss{a}{genitive neuter singular}{δεῖπνον}{banquet} ⸀γινομένου,\FnParse{b}{present middle participle genitive neuter singular} τοῦ διαβόλου ἤδη βεβληκότος\FnParse{c}{perfect active participle genitive masculine singular} εἰς τὴν καρδίαν ⸂ἵνα παραδοῖ\FnParse{d}{aorist active subjunctive 3rd singular} αὐτὸν Ἰούδας Σίμωνος ⸀Ἰσκαριώτου⸃,\FnParseFormGloss{e}{genitive masculine singular}{Ἰσκαριώτης}{Iscariot} 
\MainTextVerseMark{13}{3}⸀εἰδὼς\FnParse{a}{perfect active participle nominative masculine singular} ὅτι πάντα ⸀ἔδωκεν\FnParse{b}{aorist active indicative 3rd singular} αὐτῷ ὁ πατὴρ εἰς τὰς χεῖρας, καὶ ὅτι ἀπὸ θεοῦ ἐξῆλθεν\FnParse{c}{aorist active indicative 3rd singular} καὶ πρὸς τὸν θεὸν ὑπάγει,\FnParse{d}{present active indicative 3rd singular} 
\MainTextVerseMark{13}{4}ἐγείρεται\FnParse{a}{present passive indicative 3rd singular} ἐκ τοῦ δείπνου\FnParseFormGloss{b}{genitive neuter singular}{δεῖπνον}{banquet} καὶ τίθησιν\FnParse{c}{present active indicative 3rd singular} τὰ ἱμάτια καὶ λαβὼν\FnParse{d}{aorist active participle nominative masculine singular} λέντιον\FnParseFormGloss{e}{accusative neuter singular}{λέντιον}{towel} διέζωσεν\FnParseFormGloss{f}{aorist active indicative 3rd singular}{διαζώννυμι}{to wrap around} ἑαυτόν· 
\MainTextVerseMark{13}{5}εἶτα\FnFormGloss{a}{εἶτα}{then} βάλλει\FnParse{b}{present active indicative 3rd singular} ὕδωρ εἰς τὸν νιπτῆρα,\FnParseFormGloss{c}{accusative masculine singular}{νιπτήρ}{basin for washing} καὶ ἤρξατο\FnParse{d}{aorist middle indicative 3rd singular} νίπτειν\FnParseFormGloss{e}{present active infinitive}{νίπτω}{to wash} τοὺς πόδας τῶν μαθητῶν καὶ ἐκμάσσειν\FnParseFormGloss{f}{present active infinitive}{ἐκμάσσω}{to wipe off} τῷ λεντίῳ\FnParseFormGloss{g}{dative neuter singular}{λέντιον}{towel} ᾧ ἦν\FnParse{h}{imperfect active indicative 3rd singular} διεζωσμένος.\FnParseFormGloss{i}{perfect middle participle nominative masculine singular}{διαζώννυμι}{to wrap around} 
\MainTextVerseMark{13}{6}ἔρχεται\FnParse{a}{present middle indicative 3rd singular} οὖν πρὸς Σίμωνα Πέτρον. ⸀λέγει\FnParse{b}{present active indicative 3rd singular} ⸀αὐτῷ· Κύριε, σύ μου νίπτεις\FnParseFormGloss{c}{present active indicative 2nd singular}{νίπτω}{to wash} τοὺς πόδας; 
\MainTextVerseMark{13}{7}ἀπεκρίθη\FnParse{a}{aorist passive indicative 3rd singular} Ἰησοῦς καὶ εἶπεν\FnParse{b}{aorist active indicative 3rd singular} αὐτῷ· Ὃ ἐγὼ ποιῶ\FnParse{c}{present active indicative 1st singular} σὺ οὐκ οἶδας\FnParse{d}{perfect active indicative 2nd singular} ἄρτι, γνώσῃ\FnParse{e}{future middle indicative 2nd singular} δὲ μετὰ ταῦτα. 
\MainTextVerseMark{13}{8}λέγει\FnParse{a}{present active indicative 3rd singular} αὐτῷ Πέτρος· Οὐ μὴ νίψῃς\FnParseFormGloss{b}{aorist active subjunctive 2nd singular}{νίπτω}{to wash} ⸂μου τοὺς πόδας⸃ εἰς τὸν αἰῶνα. ἀπεκρίθη\FnParse{c}{aorist passive indicative 3rd singular} ⸂Ἰησοῦς αὐτῷ⸃· Ἐὰν μὴ νίψω\FnParseFormGloss{d}{aorist active subjunctive 1st singular}{νίπτω}{to wash} σε, οὐκ ἔχεις\FnParse{e}{present active indicative 2nd singular} μέρος μετ’ ἐμοῦ. 
\MainTextVerseMark{13}{9}λέγει\FnParse{a}{present active indicative 3rd singular} αὐτῷ Σίμων Πέτρος· Κύριε, μὴ τοὺς πόδας μου μόνον ἀλλὰ καὶ τὰς χεῖρας καὶ τὴν κεφαλήν. 
\MainTextVerseMark{13}{10}λέγει\FnParse{a}{present active indicative 3rd singular} αὐτῷ ⸀ὁ Ἰησοῦς· Ὁ λελουμένος\FnParseFormGloss{b}{perfect middle participle nominative masculine singular}{λούω}{to wash} ⸂οὐκ ἔχει\FnParse{c}{present active indicative 3rd singular} χρείαν⸃ ⸂εἰ μὴ⸃ τοὺς πόδας νίψασθαι,\FnParseFormGloss{d}{aorist middle infinitive}{νίπτω}{to wash} ἀλλ’ ἔστιν\FnParse{e}{present active indicative 3rd singular} καθαρὸς\FnParseFormGloss{f}{nominative masculine singular}{καθαρός}{clean} ὅλος· καὶ ὑμεῖς καθαροί\FnParseFormGloss{g}{nominative masculine plural}{καθαρός}{clean} ἐστε,\FnParse{h}{present active indicative 2nd plural} ἀλλ’ οὐχὶ πάντες. 
\MainTextVerseMark{13}{11}ᾔδει\FnParse{a}{pluperfect active indicative 3rd singular} γὰρ τὸν παραδιδόντα\FnParse{b}{present active participle accusative masculine singular} αὐτόν· διὰ τοῦτο εἶπεν\FnParse{c}{aorist active indicative 3rd singular} ⸀ὅτι Οὐχὶ πάντες καθαροί\FnParseFormGloss{d}{nominative masculine plural}{καθαρός}{clean} ἐστε.\FnParse{e}{present active indicative 2nd plural} 
\MainTextVerseMark{13}{12}Ὅτε οὖν ἔνιψεν\FnParseFormGloss{a}{aorist active indicative 3rd singular}{νίπτω}{to wash} τοὺς πόδας αὐτῶν ⸀καὶ ἔλαβεν\FnParse{b}{aorist active indicative 3rd singular} τὰ ἱμάτια αὐτοῦ ⸂καὶ ἀνέπεσεν⸃,\FnParseFormGloss{c}{aorist active indicative 3rd singular}{ἀναπίπτω}{recline} πάλιν εἶπεν\FnParse{d}{aorist active indicative 3rd singular} αὐτοῖς· Γινώσκετε\FnParse{e}{present active indicative 2nd plural} τί πεποίηκα\FnParse{f}{perfect active indicative 1st singular} ὑμῖν; 
\MainTextVerseMark{13}{13}ὑμεῖς φωνεῖτέ\FnParse{a}{present active indicative 2nd plural} με Ὁ διδάσκαλος καὶ Ὁ κύριος, καὶ καλῶς λέγετε,\FnParse{b}{present active indicative 2nd plural} εἰμὶ\FnParse{c}{present active indicative 1st singular} γάρ. 
\MainTextVerseMark{13}{14}εἰ οὖν ἐγὼ ἔνιψα\FnParseFormGloss{a}{aorist active indicative 1st singular}{νίπτω}{to wash} ὑμῶν τοὺς πόδας ὁ κύριος καὶ ὁ διδάσκαλος, καὶ ὑμεῖς ὀφείλετε\FnParse{b}{present active indicative 2nd plural} ἀλλήλων νίπτειν\FnParseFormGloss{c}{present active infinitive}{νίπτω}{to wash} τοὺς πόδας· 
\MainTextVerseMark{13}{15}ὑπόδειγμα\FnParseFormGloss{a}{accusative neuter singular}{ὑπόδειγμα}{example} γὰρ ἔδωκα\FnParse{b}{aorist active indicative 1st singular} ὑμῖν ἵνα καθὼς ἐγὼ ἐποίησα\FnParse{c}{aorist active indicative 1st singular} ὑμῖν καὶ ὑμεῖς ποιῆτε.\FnParse{d}{present active subjunctive 2nd plural} 
\MainTextVerseMark{13}{16}ἀμὴν ἀμὴν λέγω\FnParse{a}{present active indicative 1st singular} ὑμῖν, οὐκ ἔστιν\FnParse{b}{present active indicative 3rd singular} δοῦλος μείζων τοῦ κυρίου αὐτοῦ οὐδὲ ἀπόστολος μείζων τοῦ πέμψαντος\FnParse{c}{aorist active participle genitive masculine singular} αὐτόν. 
\MainTextVerseMark{13}{17}εἰ ταῦτα οἴδατε,\FnParse{a}{perfect active indicative 2nd plural} μακάριοί ἐστε\FnParse{b}{present active indicative 2nd plural} ἐὰν ποιῆτε\FnParse{c}{present active subjunctive 2nd plural} αὐτά. 
\MainTextVerseMark{13}{18}οὐ περὶ πάντων ὑμῶν λέγω·\FnParse{a}{present active indicative 1st singular} ἐγὼ οἶδα\FnParse{b}{perfect active indicative 1st singular} ⸀τίνας ἐξελεξάμην·\FnParseFormGloss{c}{aorist middle indicative 1st singular}{ἐκλέγομαι}{to chose} ἀλλ’ ἵνα ἡ γραφὴ πληρωθῇ·\FnParse{d}{aorist passive subjunctive 3rd singular} Ὁ τρώγων\FnParseFormGloss{e}{present active participle nominative masculine singular}{τρώγω}{to eat} ⸀μου τὸν ἄρτον ἐπῆρεν\FnParseFormGloss{f}{aorist active indicative 3rd singular}{ἐπαίρω}{to lift up} ἐπ’ ἐμὲ τὴν πτέρναν\FnParseFormGloss{g}{accusative feminine singular}{πτέρνα}{heel} αὐτοῦ. 
\MainTextVerseMark{13}{19}ἀπ’ ἄρτι λέγω\FnParse{a}{present active indicative 1st singular} ὑμῖν πρὸ τοῦ γενέσθαι,\FnParse{b}{aorist middle infinitive} ἵνα ⸂πιστεύσητε\FnParse{c}{aorist active subjunctive 2nd plural} ὅταν γένηται⸃\FnParse{d}{aorist middle subjunctive 3rd singular} ὅτι ἐγώ εἰμι.\FnParse{e}{present active indicative 1st singular} 
\MainTextVerseMark{13}{20}ἀμὴν ἀμὴν λέγω\FnParse{a}{present active indicative 1st singular} ὑμῖν, ὁ λαμβάνων\FnParse{b}{present active participle nominative masculine singular} ⸀ἄν τινα πέμψω\FnParse{c}{aorist active subjunctive 1st singular} ἐμὲ λαμβάνει,\FnParse{d}{present active indicative 3rd singular} ὁ δὲ ἐμὲ λαμβάνων\FnParse{e}{present active participle nominative masculine singular} λαμβάνει\FnParse{f}{present active indicative 3rd singular} τὸν πέμψαντά\FnParse{g}{aorist active participle accusative masculine singular} με. 
\MainTextVerseMark{13}{21}Ταῦτα εἰπὼν\FnParse{a}{aorist active participle nominative masculine singular} ⸀ὁ Ἰησοῦς ἐταράχθη\FnParseFormGloss{b}{aorist passive indicative 3rd singular}{ταράσσω}{to trouble} τῷ πνεύματι καὶ ἐμαρτύρησεν\FnParse{c}{aorist active indicative 3rd singular} καὶ εἶπεν·\FnParse{d}{aorist active indicative 3rd singular} Ἀμὴν ἀμὴν λέγω\FnParse{e}{present active indicative 1st singular} ὑμῖν ὅτι εἷς ἐξ ὑμῶν παραδώσει\FnParse{f}{future active indicative 3rd singular} με. 
\MainTextVerseMark{13}{22}⸀ἔβλεπον\FnParse{a}{imperfect active indicative 3rd plural} εἰς ἀλλήλους οἱ μαθηταὶ ἀπορούμενοι\FnParseFormGloss{b}{present middle participle nominative masculine plural}{ἀπορέω}{to be puzzled} περὶ τίνος λέγει.\FnParse{c}{present active indicative 3rd singular} 
\MainTextVerseMark{13}{23}⸀ἦν\FnParse{a}{imperfect active indicative 3rd singular} ἀνακείμενος\FnParseFormGloss{b}{present middle participle nominative masculine singular}{ἀνάκειμαι}{to recline for a meal} εἷς ⸀ἐκ τῶν μαθητῶν αὐτοῦ ἐν τῷ κόλπῳ\FnParseFormGloss{c}{dative masculine singular}{κόλπος}{side} τοῦ Ἰησοῦ, ὃν ἠγάπα\FnParse{d}{imperfect active indicative 3rd singular} ὁ Ἰησοῦς· 
\MainTextVerseMark{13}{24}νεύει\FnParseFormGloss{a}{present active indicative 3rd singular}{νεύω}{to motion} οὖν τούτῳ Σίμων Πέτρος ⸂πυθέσθαι\FnParseFormGloss{b}{aorist middle infinitive}{πυνθάνομαι}{to ask} τίς ἂν εἴη⸃\FnParse{c}{present active optative 3rd singular} περὶ οὗ λέγει.\FnParse{d}{present active indicative 3rd singular} 
\MainTextVerseMark{13}{25}⸀ἀναπεσὼν\FnParseFormGloss{a}{aorist active participle nominative masculine singular}{ἀναπίπτω}{recline} ἐκεῖνος οὕτως ἐπὶ τὸ στῆθος\FnParseFormGloss{b}{accusative neuter singular}{στῆθος}{chest} τοῦ Ἰησοῦ λέγει\FnParse{c}{present active indicative 3rd singular} αὐτῷ· Κύριε, τίς ἐστιν;\FnParse{d}{present active indicative 3rd singular} 
\MainTextVerseMark{13}{26}⸀ἀποκρίνεται\FnParse{a}{present middle indicative 3rd singular} ὁ Ἰησοῦς· Ἐκεῖνός ἐστιν\FnParse{b}{present active indicative 3rd singular} ᾧ ἐγὼ ⸀βάψω\FnParseFormGloss{c}{future active indicative 1st singular}{βάπτω}{to dip} τὸ ψωμίον\FnParseFormGloss{d}{accusative neuter singular}{ψωμίον}{piece of bread} ⸂καὶ δώσω\FnParse{e}{future active indicative 1st singular} αὐτῷ⸃· ⸂βάψας\FnParseFormGloss{f}{aorist active participle nominative masculine singular}{βάπτω}{to dip} οὖν⸃ τὸ ⸀ψωμίον\FnParseFormGloss{g}{accusative neuter singular}{ψωμίον}{piece of bread} δίδωσιν\FnParse{h}{present active indicative 3rd singular} Ἰούδᾳ Σίμωνος ⸀Ἰσκαριώτου.\FnParseFormGloss{i}{genitive masculine singular}{Ἰσκαριώτης}{Iscariot} 
\MainTextVerseMark{13}{27}καὶ μετὰ τὸ ψωμίον\FnParseFormGloss{a}{accusative neuter singular}{ψωμίον}{piece of bread} τότε εἰσῆλθεν\FnParse{b}{aorist active indicative 3rd singular} εἰς ἐκεῖνον ὁ Σατανᾶς. λέγει\FnParse{c}{present active indicative 3rd singular} οὖν αὐτῷ ⸀ὁ Ἰησοῦς· Ὃ ποιεῖς\FnParse{d}{present active indicative 2nd singular} ποίησον\FnParse{e}{aorist active imperative 2nd singular} τάχιον.\FnParseFormGloss{f}{accusative neuter singular}{τάχιον}{quickly} 
\MainTextVerseMark{13}{28}τοῦτο δὲ οὐδεὶς ἔγνω\FnParse{a}{aorist active indicative 3rd singular} τῶν ἀνακειμένων\FnParseFormGloss{b}{present middle participle genitive masculine plural}{ἀνάκειμαι}{to recline for a meal} πρὸς τί εἶπεν\FnParse{c}{aorist active indicative 3rd singular} αὐτῷ· 
\MainTextVerseMark{13}{29}τινὲς γὰρ ἐδόκουν,\FnParse{a}{imperfect active indicative 3rd plural} ἐπεὶ\FnFormGloss{b}{ἐπεί}{since} τὸ γλωσσόκομον\FnParseFormGloss{c}{accusative neuter singular}{γλωσσόκομον}{container for money} ⸀εἶχεν\FnParse{d}{imperfect active indicative 3rd singular} Ἰούδας, ὅτι λέγει\FnParse{e}{present active indicative 3rd singular} αὐτῷ ⸀ὁ Ἰησοῦς· Ἀγόρασον\FnParse{f}{aorist active imperative 2nd singular} ὧν χρείαν ἔχομεν\FnParse{g}{present active indicative 1st plural} εἰς τὴν ἑορτήν,\FnParseFormGloss{h}{accusative feminine singular}{ἑορτή}{feast} ἢ τοῖς πτωχοῖς ἵνα τι δῷ.\FnParse{i}{aorist active subjunctive 3rd singular} 
\MainTextVerseMark{13}{30}λαβὼν\FnParse{a}{aorist active participle nominative masculine singular} οὖν τὸ ψωμίον\FnParseFormGloss{b}{accusative neuter singular}{ψωμίον}{piece of bread} ἐκεῖνος ⸂ἐξῆλθεν\FnParse{c}{aorist active indicative 3rd singular} εὐθύς⸃. ἦν\FnParse{d}{imperfect active indicative 3rd singular} δὲ νύξ. 
\MainTextVerseMark{13}{31}Ὅτε ⸀οὖν ἐξῆλθεν\FnParse{a}{aorist active indicative 3rd singular} ⸀λέγει\FnParse{b}{present active indicative 3rd singular} Ἰησοῦς· Νῦν ἐδοξάσθη\FnParse{c}{aorist passive indicative 3rd singular} ὁ υἱὸς τοῦ ἀνθρώπου, καὶ ὁ θεὸς ἐδοξάσθη\FnParse{d}{aorist passive indicative 3rd singular} ἐν αὐτῷ· 
\MainTextVerseMark{13}{32}⸂εἰ ὁ θεὸς ἐδοξάσθη\FnParse{a}{aorist passive indicative 3rd singular} ἐν αὐτῷ,⸃ καὶ ὁ θεὸς δοξάσει\FnParse{b}{future active indicative 3rd singular} αὐτὸν ἐν ⸀αὑτῷ, καὶ εὐθὺς δοξάσει\FnParse{c}{future active indicative 3rd singular} αὐτόν. 
\MainTextVerseMark{13}{33}τεκνία,\FnParseFormGloss{a}{neuter plural}{τεκνίον}{dear children} ἔτι μικρὸν μεθ’ ὑμῶν εἰμι·\FnParse{b}{present active indicative 1st singular} ζητήσετέ\FnParse{c}{future active indicative 2nd plural} με, καὶ καθὼς εἶπον\FnParse{d}{aorist active indicative 1st singular} τοῖς Ἰουδαίοις ὅτι Ὅπου ⸂ἐγὼ ὑπάγω⸃\FnParse{e}{present active indicative 1st singular} ὑμεῖς οὐ δύνασθε\FnParse{f}{present middle indicative 2nd plural} ἐλθεῖν,\FnParse{g}{aorist active infinitive} καὶ ὑμῖν λέγω\FnParse{h}{present active indicative 1st singular} ἄρτι. 
\MainTextVerseMark{13}{34}ἐντολὴν καινὴν δίδωμι\FnParse{a}{present active indicative 1st singular} ὑμῖν ἵνα ἀγαπᾶτε\FnParse{b}{present active subjunctive 2nd plural} ἀλλήλους, καθὼς ἠγάπησα\FnParse{c}{aorist active indicative 1st singular} ὑμᾶς ἵνα καὶ ὑμεῖς ἀγαπᾶτε\FnParse{d}{present active subjunctive 2nd plural} ἀλλήλους. 
\MainTextVerseMark{13}{35}ἐν τούτῳ γνώσονται\FnParse{a}{future middle indicative 3rd plural} πάντες ὅτι ἐμοὶ μαθηταί ἐστε,\FnParse{b}{present active indicative 2nd plural} ἐὰν ἀγάπην ἔχητε\FnParse{c}{present active subjunctive 2nd plural} ἐν ἀλλήλοις. 
\MainTextVerseMark{13}{36}Λέγει\FnParse{a}{present active indicative 3rd singular} αὐτῷ Σίμων Πέτρος· Κύριε, ποῦ ὑπάγεις;\FnParse{b}{present active indicative 2nd singular} ⸀ἀπεκρίθη\FnParse{c}{aorist passive indicative 3rd singular} Ἰησοῦς· Ὅπου ὑπάγω\FnParse{d}{present active indicative 1st singular} οὐ δύνασαί\FnParse{e}{present middle indicative 2nd singular} μοι νῦν ἀκολουθῆσαι,\FnParse{f}{aorist active infinitive} ⸂ἀκολουθήσεις\FnParse{g}{future active indicative 2nd singular} δὲ ὕστερον⸃.\FnParseFormGloss{h}{accusative neuter singular}{ὕστερον}{later} 
\MainTextVerseMark{13}{37}λέγει\FnParse{a}{present active indicative 3rd singular} αὐτῷ ⸀ὁ Πέτρος· Κύριε, διὰ τί οὐ δύναμαί\FnParse{b}{present middle indicative 1st singular} σοι ⸀ἀκολουθῆσαι\FnParse{c}{aorist active infinitive} ἄρτι; τὴν ψυχήν μου ὑπὲρ σοῦ θήσω.\FnParse{d}{future active indicative 1st singular} 
\MainTextVerseMark{13}{38}⸀ἀποκρίνεται\FnParse{a}{present middle indicative 3rd singular} Ἰησοῦς· Τὴν ψυχήν σου ὑπὲρ ἐμοῦ θήσεις;\FnParse{b}{future active indicative 2nd singular} ἀμὴν ἀμὴν λέγω\FnParse{c}{present active indicative 1st singular} σοι, οὐ μὴ ἀλέκτωρ\FnParseFormGloss{d}{nominative masculine singular}{ἀλέκτωρ}{rooster} φωνήσῃ\FnParse{e}{aorist active subjunctive 3rd singular} ἕως οὗ ⸀ἀρνήσῃ\FnParse{f}{future middle indicative 2nd singular} με τρίς.\FnFormGloss{g}{τρίς}{three times} 
\ChapterHeader{14}{Chapter 14}

\MainTextVerseMark{14}{1}Μὴ ταρασσέσθω\FnParseFormGloss{a}{present passive imperative 3rd singular}{ταράσσω}{to trouble} ὑμῶν ἡ καρδία· πιστεύετε\FnParse{b}{present active indicative 2nd plural} εἰς τὸν θεόν, καὶ εἰς ἐμὲ πιστεύετε.\FnParse{c}{present active indicative 2nd plural} 
\MainTextVerseMark{14}{2}ἐν τῇ οἰκίᾳ τοῦ πατρός μου μοναὶ\FnParseFormGloss{a}{nominative feminine plural}{μονή}{room} πολλαί εἰσιν·\FnParse{b}{present active indicative 3rd plural} εἰ δὲ μή, εἶπον\FnParse{c}{aorist active indicative 1st singular} ἂν ὑμῖν ⸀ὅτι πορεύομαι\FnParse{d}{present middle indicative 1st singular} ἑτοιμάσαι\FnParse{e}{aorist active infinitive} τόπον ὑμῖν· 
\MainTextVerseMark{14}{3}καὶ ἐὰν πορευθῶ\FnParse{a}{aorist passive subjunctive 1st singular} ⸀καὶ ἑτοιμάσω\FnParse{b}{aorist active subjunctive 1st singular} ⸂τόπον ὑμῖν⸃, πάλιν ἔρχομαι\FnParse{c}{present middle indicative 1st singular} καὶ παραλήμψομαι\FnParse{d}{future middle indicative 1st singular} ὑμᾶς πρὸς ἐμαυτόν, ἵνα ὅπου εἰμὶ\FnParse{e}{present active indicative 1st singular} ἐγὼ καὶ ὑμεῖς ἦτε.\FnParse{f}{present active subjunctive 2nd plural} 
\MainTextVerseMark{14}{4}καὶ ὅπου ἐγὼ ὑπάγω\FnParse{a}{present active indicative 1st singular} οἴδατε\FnParse{b}{perfect active indicative 2nd plural} ⸂τὴν ὁδόν⸃. 
\MainTextVerseMark{14}{5}λέγει\FnParse{a}{present active indicative 3rd singular} αὐτῷ Θωμᾶς·\FnParseFormGloss{b}{nominative masculine singular}{Θωμᾶς}{Thomas} Κύριε, οὐκ οἴδαμεν\FnParse{c}{perfect active indicative 1st plural} ποῦ ὑπάγεις·\FnParse{d}{present active indicative 2nd singular} ⸀πῶς ⸂δυνάμεθα\FnParse{e}{present middle indicative 1st plural} τὴν ὁδὸν εἰδέναι⸃;\FnParse{f}{perfect active infinitive} 
\MainTextVerseMark{14}{6}λέγει\FnParse{a}{present active indicative 3rd singular} αὐτῷ ⸀ὁ Ἰησοῦς· Ἐγώ εἰμι\FnParse{b}{present active indicative 1st singular} ἡ ὁδὸς καὶ ἡ ἀλήθεια καὶ ἡ ζωή· οὐδεὶς ἔρχεται\FnParse{c}{present middle indicative 3rd singular} πρὸς τὸν πατέρα εἰ μὴ δι’ ἐμοῦ. 
\MainTextVerseMark{14}{7}εἰ ⸀ἐγνώκειτέ\FnParse{a}{pluperfect active indicative 2nd plural} με, καὶ τὸν πατέρα μου ⸂ἂν ᾔδειτε⸃·\FnParse{b}{pluperfect active indicative 2nd plural} ⸀ἀπ’ ἄρτι γινώσκετε\FnParse{c}{present active indicative 2nd plural} αὐτὸν καὶ ἑωράκατε\FnParse{d}{perfect active indicative 2nd plural} ⸀αὐτόν. 
\MainTextVerseMark{14}{8}Λέγει\FnParse{a}{present active indicative 3rd singular} αὐτῷ Φίλιππος· Κύριε, δεῖξον\FnParse{b}{aorist active imperative 2nd singular} ἡμῖν τὸν πατέρα, καὶ ἀρκεῖ\FnParseFormGloss{c}{present active indicative 3rd singular}{ἀρκέω}{to be content} ἡμῖν. 
\MainTextVerseMark{14}{9}λέγει\FnParse{a}{present active indicative 3rd singular} αὐτῷ ὁ Ἰησοῦς· ⸂Τοσούτῳ\FnParseFormGloss{b}{dative masculine singular}{τοσοῦτος}{so great} χρόνῳ⸃ μεθ’ ὑμῶν εἰμι\FnParse{c}{present active indicative 1st singular} καὶ οὐκ ἔγνωκάς\FnParse{d}{perfect active indicative 2nd singular} με, Φίλιππε; ὁ ἑωρακὼς\FnParse{e}{perfect active participle nominative masculine singular} ἐμὲ ἑώρακεν\FnParse{f}{perfect active indicative 3rd singular} τὸν πατέρα· ⸀πῶς σὺ λέγεις·\FnParse{g}{present active indicative 2nd singular} Δεῖξον\FnParse{h}{aorist active imperative 2nd singular} ἡμῖν τὸν πατέρα; 
\MainTextVerseMark{14}{10}οὐ πιστεύεις\FnParse{a}{present active indicative 2nd singular} ὅτι ἐγὼ ἐν τῷ πατρὶ καὶ ὁ πατὴρ ἐν ἐμοί ἐστιν;\FnParse{b}{present active indicative 3rd singular} τὰ ῥήματα ἃ ἐγὼ ⸀λέγω\FnParse{c}{present active indicative 1st singular} ὑμῖν ἀπ’ ἐμαυτοῦ οὐ λαλῶ,\FnParse{d}{present active indicative 1st singular} ὁ δὲ ⸀πατὴρ ἐν ἐμοὶ ⸀μένων\FnParse{e}{present active participle nominative masculine singular} ποιεῖ\FnParse{f}{present active indicative 3rd singular} τὰ ἔργα ⸀αὐτοῦ. 
\MainTextVerseMark{14}{11}πιστεύετέ\FnParse{a}{present active imperative 2nd plural} μοι ὅτι ἐγὼ ἐν τῷ πατρὶ καὶ ὁ πατὴρ ἐν ἐμοί· εἰ δὲ μή, διὰ τὰ ἔργα αὐτὰ ⸀πιστεύετε.\FnParse{b}{present active imperative 2nd plural} 
\MainTextVerseMark{14}{12}ἀμὴν ἀμὴν λέγω\FnParse{a}{present active indicative 1st singular} ὑμῖν, ὁ πιστεύων\FnParse{b}{present active participle nominative masculine singular} εἰς ἐμὲ τὰ ἔργα ἃ ἐγὼ ποιῶ\FnParse{c}{present active indicative 1st singular} κἀκεῖνος\FnParseFormGloss{d}{nominative masculine singular}{κἀκεῖνος}{and that one} ποιήσει,\FnParse{e}{future active indicative 3rd singular} καὶ μείζονα τούτων ποιήσει,\FnParse{f}{future active indicative 3rd singular} ὅτι ἐγὼ πρὸς τὸν ⸀πατέρα πορεύομαι·\FnParse{g}{present middle indicative 1st singular} 
\MainTextVerseMark{14}{13}καὶ ὅ τι ἂν αἰτήσητε\FnParse{a}{aorist active subjunctive 2nd plural} ἐν τῷ ὀνόματί μου τοῦτο ποιήσω,\FnParse{b}{future active indicative 1st singular} ἵνα δοξασθῇ\FnParse{c}{aorist passive subjunctive 3rd singular} ὁ πατὴρ ἐν τῷ υἱῷ· 
\MainTextVerseMark{14}{14}ἐάν τι αἰτήσητέ\FnParse{a}{aorist active subjunctive 2nd plural} ⸀με ἐν τῷ ὀνόματί μου ⸀ἐγὼ ποιήσω.\FnParse{b}{future active indicative 1st singular} 
\MainTextVerseMark{14}{15}Ἐὰν ἀγαπᾶτέ\FnParse{a}{present active subjunctive 2nd plural} με, τὰς ἐντολὰς τὰς ἐμὰς ⸀τηρήσετε·\FnParse{b}{future active indicative 2nd plural} 
\MainTextVerseMark{14}{16}κἀγὼ ἐρωτήσω\FnParse{a}{future active indicative 1st singular} τὸν πατέρα καὶ ἄλλον παράκλητον\FnParseFormGloss{b}{accusative masculine singular}{παράκλητος}{counselor} δώσει\FnParse{c}{future active indicative 3rd singular} ὑμῖν ἵνα ⸀ᾖ\FnParse{d}{present active subjunctive 3rd singular} μεθ’ ὑμῶν εἰς τὸν ⸀αἰῶνα, 
\MainTextVerseMark{14}{17}τὸ πνεῦμα τῆς ἀληθείας, ὃ ὁ κόσμος οὐ δύναται\FnParse{a}{present middle indicative 3rd singular} λαβεῖν,\FnParse{b}{aorist active infinitive} ὅτι οὐ θεωρεῖ\FnParse{c}{present active indicative 3rd singular} αὐτὸ οὐδὲ ⸀γινώσκει·\FnParse{d}{present active indicative 3rd singular} ⸀ὑμεῖς γινώσκετε\FnParse{e}{present active indicative 2nd plural} αὐτό, ὅτι παρ’ ὑμῖν μένει\FnParse{f}{present active indicative 3rd singular} καὶ ἐν ὑμῖν ⸀ἔσται.\FnParse{g}{future middle indicative 3rd singular} 
\MainTextVerseMark{14}{18}Οὐκ ἀφήσω\FnParse{a}{future active indicative 1st singular} ὑμᾶς ὀρφανούς,\FnParseFormGloss{b}{accusative masculine plural}{ὀρφανός}{an orphan} ἔρχομαι\FnParse{c}{present middle indicative 1st singular} πρὸς ὑμᾶς. 
\MainTextVerseMark{14}{19}ἔτι μικρὸν καὶ ὁ κόσμος με οὐκέτι θεωρεῖ,\FnParse{a}{present active indicative 3rd singular} ὑμεῖς δὲ θεωρεῖτέ\FnParse{b}{present active indicative 2nd plural} με, ὅτι ἐγὼ ζῶ\FnParse{c}{present active indicative 1st singular} καὶ ὑμεῖς ⸀ζήσετε.\FnParse{d}{future active indicative 2nd plural} 
\MainTextVerseMark{14}{20}ἐν ἐκείνῃ τῇ ἡμέρᾳ ⸂γνώσεσθε\FnParse{a}{future middle indicative 2nd plural} ὑμεῖς⸃ ὅτι ἐγὼ ἐν τῷ πατρί μου καὶ ὑμεῖς ἐν ἐμοὶ κἀγὼ ἐν ὑμῖν. 
\MainTextVerseMark{14}{21}ὁ ἔχων\FnParse{a}{present active participle nominative masculine singular} τὰς ἐντολάς μου καὶ τηρῶν\FnParse{b}{present active participle nominative masculine singular} αὐτὰς ἐκεῖνός ἐστιν\FnParse{c}{present active indicative 3rd singular} ὁ ἀγαπῶν\FnParse{d}{present active participle nominative masculine singular} με· ὁ δὲ ἀγαπῶν\FnParse{e}{present active participle nominative masculine singular} με ἀγαπηθήσεται\FnParse{f}{future passive indicative 3rd singular} ὑπὸ τοῦ πατρός μου, κἀγὼ ἀγαπήσω\FnParse{g}{future active indicative 1st singular} αὐτὸν καὶ ἐμφανίσω\FnParseFormGloss{h}{future active indicative 1st singular}{ἐμφανίζω}{to show} αὐτῷ ἐμαυτόν. 
\MainTextVerseMark{14}{22}λέγει\FnParse{a}{present active indicative 3rd singular} αὐτῷ Ἰούδας, οὐχ ὁ Ἰσκαριώτης·\FnParseFormGloss{b}{nominative masculine singular}{Ἰσκαριώτης}{Iscariot} Κύριε, ⸀τί γέγονεν\FnParse{c}{perfect active indicative 3rd singular} ὅτι ἡμῖν μέλλεις\FnParse{d}{present active indicative 2nd singular} ἐμφανίζειν\FnParseFormGloss{e}{present active infinitive}{ἐμφανίζω}{to show} σεαυτὸν καὶ οὐχὶ τῷ κόσμῳ; 
\MainTextVerseMark{14}{23}ἀπεκρίθη\FnParse{a}{aorist passive indicative 3rd singular} Ἰησοῦς καὶ εἶπεν\FnParse{b}{aorist active indicative 3rd singular} αὐτῷ· Ἐάν τις ἀγαπᾷ\FnParse{c}{present active subjunctive 3rd singular} με τὸν λόγον μου τηρήσει,\FnParse{d}{future active indicative 3rd singular} καὶ ὁ πατήρ μου ἀγαπήσει\FnParse{e}{future active indicative 3rd singular} αὐτόν, καὶ πρὸς αὐτὸν ἐλευσόμεθα\FnParse{f}{future middle indicative 1st plural} καὶ μονὴν\FnParseFormGloss{g}{accusative feminine singular}{μονή}{room} παρ’ αὐτῷ ⸀ποιησόμεθα.\FnParse{h}{future middle indicative 1st plural} 
\MainTextVerseMark{14}{24}ὁ μὴ ἀγαπῶν\FnParse{a}{present active participle nominative masculine singular} με τοὺς λόγους μου οὐ τηρεῖ·\FnParse{b}{present active indicative 3rd singular} καὶ ὁ λόγος ὃν ἀκούετε\FnParse{c}{present active indicative 2nd plural} οὐκ ἔστιν\FnParse{d}{present active indicative 3rd singular} ἐμὸς ἀλλὰ τοῦ πέμψαντός\FnParse{e}{aorist active participle genitive masculine singular} με πατρός. 
\MainTextVerseMark{14}{25}Ταῦτα λελάληκα\FnParse{a}{perfect active indicative 1st singular} ὑμῖν παρ’ ὑμῖν μένων·\FnParse{b}{present active participle nominative masculine singular} 
\MainTextVerseMark{14}{26}ὁ δὲ παράκλητος,\FnParseFormGloss{a}{nominative masculine singular}{παράκλητος}{counselor} τὸ πνεῦμα τὸ ἅγιον ὃ πέμψει\FnParse{b}{future active indicative 3rd singular} ὁ πατὴρ ἐν τῷ ὀνόματί μου, ἐκεῖνος ὑμᾶς διδάξει\FnParse{c}{future active indicative 3rd singular} πάντα καὶ ὑπομνήσει\FnParseFormGloss{d}{future active indicative 3rd singular}{ὑπομιμνῄσκω}{to remind} ὑμᾶς πάντα ἃ εἶπον\FnParse{e}{aorist active indicative 1st singular} ⸀ὑμῖν. 
\MainTextVerseMark{14}{27}εἰρήνην ἀφίημι\FnParse{a}{present active indicative 1st singular} ὑμῖν, εἰρήνην τὴν ἐμὴν δίδωμι\FnParse{b}{present active indicative 1st singular} ὑμῖν· οὐ καθὼς ὁ κόσμος δίδωσιν\FnParse{c}{present active indicative 3rd singular} ἐγὼ δίδωμι\FnParse{d}{present active indicative 1st singular} ὑμῖν. μὴ ταρασσέσθω\FnParseFormGloss{e}{present passive imperative 3rd singular}{ταράσσω}{to trouble} ὑμῶν ἡ καρδία μηδὲ δειλιάτω.\FnParseFormGloss{f}{present active imperative 3rd singular}{δειλιάω}{to be afraid} 
\MainTextVerseMark{14}{28}ἠκούσατε\FnParse{a}{aorist active indicative 2nd plural} ὅτι ἐγὼ εἶπον\FnParse{b}{aorist active indicative 1st singular} ὑμῖν· Ὑπάγω\FnParse{c}{present active indicative 1st singular} καὶ ἔρχομαι\FnParse{d}{present middle indicative 1st singular} πρὸς ὑμᾶς. εἰ ἠγαπᾶτέ\FnParse{e}{imperfect active indicative 2nd plural} με ἐχάρητε\FnParse{f}{aorist passive indicative 2nd plural} ἄν, ⸀ὅτι πορεύομαι\FnParse{g}{present middle indicative 1st singular} πρὸς τὸν πατέρα, ὅτι ὁ ⸀πατὴρ μείζων μού ἐστιν.\FnParse{h}{present active indicative 3rd singular} 
\MainTextVerseMark{14}{29}καὶ νῦν εἴρηκα\FnParse{a}{perfect active indicative 1st singular} ὑμῖν πρὶν\FnFormGloss{b}{πρίν}{before} γενέσθαι,\FnParse{c}{aorist middle infinitive} ἵνα ὅταν γένηται\FnParse{d}{aorist middle subjunctive 3rd singular} πιστεύσητε.\FnParse{e}{aorist active subjunctive 2nd plural} 
\MainTextVerseMark{14}{30}οὐκέτι πολλὰ λαλήσω\FnParse{a}{future active indicative 1st singular} μεθ’ ὑμῶν, ἔρχεται\FnParse{b}{present middle indicative 3rd singular} γὰρ ὁ τοῦ κόσμου ἄρχων· καὶ ἐν ἐμοὶ οὐκ ἔχει\FnParse{c}{present active indicative 3rd singular} οὐδέν, 
\MainTextVerseMark{14}{31}ἀλλ’ ἵνα γνῷ\FnParse{a}{aorist active subjunctive 3rd singular} ὁ κόσμος ὅτι ἀγαπῶ\FnParse{b}{present active indicative 1st singular} τὸν πατέρα, καὶ καθὼς ⸀ἐνετείλατο\FnParseFormGloss{c}{aorist middle indicative 3rd singular}{ἐντέλλομαι}{to command} μοι ὁ πατὴρ οὕτως ποιῶ.\FnParse{d}{present active indicative 1st singular} Ἐγείρεσθε,\FnParse{e}{present passive imperative 2nd plural} ἄγωμεν\FnParse{f}{present active subjunctive 1st plural} ἐντεῦθεν.\FnFormGloss{g}{ἐντεῦθεν}{from here} 
\ChapterHeader{15}{Chapter 15}

\MainTextVerseMark{15}{1}Ἐγώ εἰμι\FnParse{a}{present active indicative 1st singular} ἡ ἄμπελος\FnParseFormGloss{b}{nominative feminine singular}{ἄμπελος}{vine} ἡ ἀληθινή,\FnParseFormGloss{c}{nominative feminine singular}{ἀληθινός}{true} καὶ ὁ πατήρ μου ὁ γεωργός\FnParseFormGloss{d}{nominative masculine singular}{γεωργός}{farmer} ἐστιν·\FnParse{e}{present active indicative 3rd singular} 
\MainTextVerseMark{15}{2}πᾶν κλῆμα\FnParseFormGloss{a}{accusative neuter singular}{κλῆμα}{branch} ἐν ἐμοὶ μὴ φέρον\FnParse{b}{present active participle accusative neuter singular} καρπὸν αἴρει\FnParse{c}{present active indicative 3rd singular} αὐτό, καὶ πᾶν τὸ καρπὸν φέρον\FnParse{d}{present active participle accusative neuter singular} καθαίρει\FnParseFormGloss{e}{present active indicative 3rd singular}{καθαίρω}{to prune} αὐτὸ ἵνα ⸂καρπὸν πλείονα⸃ φέρῃ.\FnParse{f}{present active subjunctive 3rd singular} 
\MainTextVerseMark{15}{3}ἤδη ὑμεῖς καθαροί\FnParseFormGloss{a}{nominative masculine plural}{καθαρός}{clean} ἐστε\FnParse{b}{present active indicative 2nd plural} διὰ τὸν λόγον ὃν λελάληκα\FnParse{c}{perfect active indicative 1st singular} ὑμῖν· 
\MainTextVerseMark{15}{4}μείνατε\FnParse{a}{aorist active imperative 2nd plural} ἐν ἐμοί, κἀγὼ ἐν ὑμῖν. καθὼς τὸ κλῆμα\FnParseFormGloss{b}{nominative neuter singular}{κλῆμα}{branch} οὐ δύναται\FnParse{c}{present middle indicative 3rd singular} καρπὸν φέρειν\FnParse{d}{present active infinitive} ἀφ’ ἑαυτοῦ ἐὰν μὴ ⸀μένῃ\FnParse{e}{present active subjunctive 3rd singular} ἐν τῇ ἀμπέλῳ,\FnParseFormGloss{f}{dative feminine singular}{ἄμπελος}{vine} οὕτως οὐδὲ ὑμεῖς ἐὰν μὴ ἐν ἐμοὶ ⸀μένητε.\FnParse{g}{present active subjunctive 2nd plural} 
\MainTextVerseMark{15}{5}ἐγώ εἰμι\FnParse{a}{present active indicative 1st singular} ἡ ἄμπελος,\FnParseFormGloss{b}{nominative feminine singular}{ἄμπελος}{vine} ὑμεῖς τὰ κλήματα.\FnParseFormGloss{c}{nominative neuter plural}{κλῆμα}{branch} ὁ μένων\FnParse{d}{present active participle nominative masculine singular} ἐν ἐμοὶ κἀγὼ ἐν αὐτῷ οὗτος φέρει\FnParse{e}{present active indicative 3rd singular} καρπὸν πολύν, ὅτι χωρὶς ἐμοῦ οὐ δύνασθε\FnParse{f}{present middle indicative 2nd plural} ποιεῖν\FnParse{g}{present active infinitive} οὐδέν. 
\MainTextVerseMark{15}{6}ἐὰν μή τις ⸀μένῃ\FnParse{a}{present active subjunctive 3rd singular} ἐν ἐμοί, ἐβλήθη\FnParse{b}{aorist passive indicative 3rd singular} ἔξω ὡς τὸ κλῆμα\FnParseFormGloss{c}{nominative neuter singular}{κλῆμα}{branch} καὶ ἐξηράνθη,\FnParseFormGloss{d}{aorist passive indicative 3rd singular}{ξηραίνω}{to wither} καὶ συνάγουσιν\FnParse{e}{present active indicative 3rd plural} αὐτὰ καὶ εἰς τὸ πῦρ βάλλουσιν\FnParse{f}{present active indicative 3rd plural} καὶ καίεται.\FnParseFormGloss{g}{present passive indicative 3rd singular}{καίω}{to light} 
\MainTextVerseMark{15}{7}ἐὰν μείνητε\FnParse{a}{aorist active subjunctive 2nd plural} ἐν ἐμοὶ καὶ τὰ ῥήματά μου ἐν ὑμῖν μείνῃ,\FnParse{b}{aorist active subjunctive 3rd singular} ὃ ἐὰν θέλητε\FnParse{c}{present active subjunctive 2nd plural} ⸀αἰτήσασθε\FnParse{d}{aorist middle imperative 2nd plural} καὶ γενήσεται\FnParse{e}{future middle indicative 3rd singular} ὑμῖν· 
\MainTextVerseMark{15}{8}ἐν τούτῳ ἐδοξάσθη\FnParse{a}{aorist passive indicative 3rd singular} ὁ πατήρ μου ἵνα καρπὸν πολὺν φέρητε\FnParse{b}{present active subjunctive 2nd plural} καὶ ⸀γένησθε\FnParse{c}{aorist middle subjunctive 2nd plural} ἐμοὶ μαθηταί. 
\MainTextVerseMark{15}{9}καθὼς ἠγάπησέν\FnParse{a}{aorist active indicative 3rd singular} με ὁ πατήρ, κἀγὼ ⸂ὑμᾶς ἠγάπησα⸃,\FnParse{b}{aorist active indicative 1st singular} μείνατε\FnParse{c}{aorist active imperative 2nd plural} ἐν τῇ ἀγάπῃ τῇ ἐμῇ. 
\MainTextVerseMark{15}{10}ἐὰν τὰς ἐντολάς μου τηρήσητε,\FnParse{a}{aorist active subjunctive 2nd plural} μενεῖτε\FnParse{b}{future active indicative 2nd plural} ἐν τῇ ἀγάπῃ μου, καθὼς ἐγὼ ⸂τὰς ἐντολὰς τοῦ πατρός μου⸃ τετήρηκα\FnParse{c}{perfect active indicative 1st singular} καὶ μένω\FnParse{d}{present active indicative 1st singular} αὐτοῦ ἐν τῇ ἀγάπῃ. 
\MainTextVerseMark{15}{11}ταῦτα λελάληκα\FnParse{a}{perfect active indicative 1st singular} ὑμῖν ἵνα ἡ χαρὰ ἡ ἐμὴ ἐν ὑμῖν ⸀ᾖ\FnParse{b}{present active subjunctive 3rd singular} καὶ ἡ χαρὰ ὑμῶν πληρωθῇ.\FnParse{c}{aorist passive subjunctive 3rd singular} 
\MainTextVerseMark{15}{12}Αὕτη ἐστὶν\FnParse{a}{present active indicative 3rd singular} ἡ ἐντολὴ ἡ ἐμὴ ἵνα ἀγαπᾶτε\FnParse{b}{present active subjunctive 2nd plural} ἀλλήλους καθὼς ἠγάπησα\FnParse{c}{aorist active indicative 1st singular} ὑμᾶς· 
\MainTextVerseMark{15}{13}μείζονα ταύτης ἀγάπην οὐδεὶς ἔχει,\FnParse{a}{present active indicative 3rd singular} ἵνα τις τὴν ψυχὴν αὐτοῦ θῇ\FnParse{b}{aorist active subjunctive 3rd singular} ὑπὲρ τῶν φίλων\FnParseFormGloss{c}{genitive masculine plural}{φίλος}{friendly} αὐτοῦ. 
\MainTextVerseMark{15}{14}ὑμεῖς φίλοι\FnParseFormGloss{a}{nominative masculine plural}{φίλος}{friendly} μού ἐστε\FnParse{b}{present active indicative 2nd plural} ἐὰν ποιῆτε\FnParse{c}{present active subjunctive 2nd plural} ⸀ἃ ἐγὼ ἐντέλλομαι\FnParseFormGloss{d}{present middle indicative 1st singular}{ἐντέλλομαι}{to command} ὑμῖν. 
\MainTextVerseMark{15}{15}οὐκέτι ⸂λέγω\FnParse{a}{present active indicative 1st singular} ὑμᾶς⸃ δούλους, ὅτι ὁ δοῦλος οὐκ οἶδεν\FnParse{b}{perfect active indicative 3rd singular} τί ποιεῖ\FnParse{c}{present active indicative 3rd singular} αὐτοῦ ὁ κύριος· ὑμᾶς δὲ εἴρηκα\FnParse{d}{perfect active indicative 1st singular} φίλους,\FnParseFormGloss{e}{accusative masculine plural}{φίλος}{friendly} ὅτι πάντα ἃ ἤκουσα\FnParse{f}{aorist active indicative 1st singular} παρὰ τοῦ πατρός μου ἐγνώρισα\FnParseFormGloss{g}{aorist active indicative 1st singular}{γνωρίζω}{to make known} ὑμῖν. 
\MainTextVerseMark{15}{16}οὐχ ὑμεῖς με ἐξελέξασθε,\FnParseFormGloss{a}{aorist middle indicative 2nd plural}{ἐκλέγομαι}{to chose} ἀλλ’ ἐγὼ ἐξελεξάμην\FnParseFormGloss{b}{aorist middle indicative 1st singular}{ἐκλέγομαι}{to chose} ὑμᾶς, καὶ ἔθηκα\FnParse{c}{aorist active indicative 1st singular} ὑμᾶς ἵνα ὑμεῖς ὑπάγητε\FnParse{d}{present active subjunctive 2nd plural} καὶ καρπὸν φέρητε\FnParse{e}{present active subjunctive 2nd plural} καὶ ὁ καρπὸς ὑμῶν μένῃ,\FnParse{f}{present active subjunctive 3rd singular} ἵνα ὅ τι ἂν αἰτήσητε\FnParse{g}{aorist active subjunctive 2nd plural} τὸν πατέρα ἐν τῷ ὀνόματί μου δῷ\FnParse{h}{aorist active subjunctive 3rd singular} ὑμῖν. 
\MainTextVerseMark{15}{17}ταῦτα ἐντέλλομαι\FnParseFormGloss{a}{present middle indicative 1st singular}{ἐντέλλομαι}{to command} ὑμῖν ἵνα ἀγαπᾶτε\FnParse{b}{present active subjunctive 2nd plural} ἀλλήλους. 
\MainTextVerseMark{15}{18}Εἰ ὁ κόσμος ὑμᾶς μισεῖ,\FnParse{a}{present active indicative 3rd singular} γινώσκετε\FnParse{b}{present active indicative 2nd plural} ὅτι ἐμὲ πρῶτον ⸀ὑμῶν μεμίσηκεν.\FnParse{c}{perfect active indicative 3rd singular} 
\MainTextVerseMark{15}{19}εἰ ἐκ τοῦ κόσμου ἦτε,\FnParse{a}{imperfect active indicative 2nd plural} ὁ κόσμος ἂν τὸ ἴδιον ἐφίλει·\FnParseFormGloss{b}{imperfect active indicative 3rd singular}{φιλέω}{to love} ὅτι δὲ ἐκ τοῦ κόσμου οὐκ ἐστέ,\FnParse{c}{present active indicative 2nd plural} ἀλλ’ ἐγὼ ἐξελεξάμην\FnParseFormGloss{d}{aorist middle indicative 1st singular}{ἐκλέγομαι}{to chose} ὑμᾶς ἐκ τοῦ κόσμου, διὰ τοῦτο μισεῖ\FnParse{e}{present active indicative 3rd singular} ὑμᾶς ὁ κόσμος. 
\MainTextVerseMark{15}{20}μνημονεύετε\FnParseFormGloss{a}{present active imperative 2nd plural}{μνημονεύω}{to remember} τοῦ λόγου οὗ ἐγὼ εἶπον\FnParse{b}{aorist active indicative 1st singular} ὑμῖν· Οὐκ ἔστιν\FnParse{c}{present active indicative 3rd singular} δοῦλος μείζων τοῦ κυρίου αὐτοῦ· εἰ ἐμὲ ἐδίωξαν,\FnParse{d}{aorist active indicative 3rd plural} καὶ ὑμᾶς διώξουσιν·\FnParse{e}{future active indicative 3rd plural} εἰ τὸν λόγον μου ἐτήρησαν,\FnParse{f}{aorist active indicative 3rd plural} καὶ τὸν ὑμέτερον\FnParseFormGloss{g}{accusative masculine singular}{ὑμέτερος}{your} τηρήσουσιν.\FnParse{h}{future active indicative 3rd plural} 
\MainTextVerseMark{15}{21}ἀλλὰ ταῦτα πάντα ποιήσουσιν\FnParse{a}{future active indicative 3rd plural} ⸂εἰς ὑμᾶς⸃ διὰ τὸ ὄνομά μου, ὅτι οὐκ οἴδασιν\FnParse{b}{perfect active indicative 3rd plural} τὸν πέμψαντά\FnParse{c}{aorist active participle accusative masculine singular} με. 
\MainTextVerseMark{15}{22}εἰ μὴ ἦλθον\FnParse{a}{aorist active indicative 1st singular} καὶ ἐλάλησα\FnParse{b}{aorist active indicative 1st singular} αὐτοῖς, ἁμαρτίαν οὐκ ⸀εἴχοσαν·\FnParse{c}{imperfect active indicative 3rd plural} νῦν δὲ πρόφασιν\FnParseFormGloss{d}{accusative feminine singular}{πρόφασις}{excuse} οὐκ ἔχουσιν\FnParse{e}{present active indicative 3rd plural} περὶ τῆς ἁμαρτίας αὐτῶν. 
\MainTextVerseMark{15}{23}ὁ ἐμὲ μισῶν\FnParse{a}{present active participle nominative masculine singular} καὶ τὸν πατέρα μου μισεῖ.\FnParse{b}{present active indicative 3rd singular} 
\MainTextVerseMark{15}{24}εἰ τὰ ἔργα μὴ ἐποίησα\FnParse{a}{aorist active indicative 1st singular} ἐν αὐτοῖς ἃ οὐδεὶς ἄλλος ⸀ἐποίησεν,\FnParse{b}{aorist active indicative 3rd singular} ἁμαρτίαν οὐκ ⸀εἴχοσαν·\FnParse{c}{imperfect active indicative 3rd plural} νῦν δὲ καὶ ἑωράκασιν\FnParse{d}{perfect active indicative 3rd plural} καὶ μεμισήκασιν\FnParse{e}{perfect active indicative 3rd plural} καὶ ἐμὲ καὶ τὸν πατέρα μου. 
\MainTextVerseMark{15}{25}ἀλλ’ ἵνα πληρωθῇ\FnParse{a}{aorist passive subjunctive 3rd singular} ὁ λόγος ὁ ⸂ἐν τῷ νόμῳ αὐτῶν γεγραμμένος⸃\FnParse{b}{perfect passive participle nominative masculine singular} ὅτι Ἐμίσησάν\FnParse{c}{aorist active indicative 3rd plural} με δωρεάν.\FnFormGloss{d}{δωρεάν}{freely} 
\MainTextVerseMark{15}{26}⸀Ὅταν ἔλθῃ\FnParse{a}{aorist active subjunctive 3rd singular} ὁ παράκλητος\FnParseFormGloss{b}{nominative masculine singular}{παράκλητος}{counselor} ὃν ἐγὼ πέμψω\FnParse{c}{future active indicative 1st singular} ὑμῖν παρὰ τοῦ πατρός, τὸ πνεῦμα τῆς ἀληθείας ὃ παρὰ τοῦ πατρὸς ἐκπορεύεται,\FnParse{d}{present middle indicative 3rd singular} ἐκεῖνος μαρτυρήσει\FnParse{e}{future active indicative 3rd singular} περὶ ἐμοῦ· 
\MainTextVerseMark{15}{27}καὶ ὑμεῖς δὲ μαρτυρεῖτε,\FnParse{a}{present active indicative 2nd plural} ὅτι ἀπ’ ἀρχῆς μετ’ ἐμοῦ ἐστε.\FnParse{b}{present active indicative 2nd plural} 
\ChapterHeader{16}{Chapter 16}

\MainTextVerseMark{16}{1}Ταῦτα λελάληκα\FnParse{a}{perfect active indicative 1st singular} ὑμῖν ἵνα μὴ σκανδαλισθῆτε.\FnParse{b}{aorist passive subjunctive 2nd plural} 
\MainTextVerseMark{16}{2}ἀποσυναγώγους\FnParseFormGloss{a}{accusative masculine plural}{ἀποσυνάγωγος}{put out of the synagogue} ποιήσουσιν\FnParse{b}{future active indicative 3rd plural} ὑμᾶς· ἀλλ’ ἔρχεται\FnParse{c}{present middle indicative 3rd singular} ὥρα ἵνα πᾶς ὁ ἀποκτείνας\FnParse{d}{aorist active participle nominative masculine singular} ὑμᾶς δόξῃ\FnParse{e}{aorist active subjunctive 3rd singular} λατρείαν\FnParseFormGloss{f}{accusative feminine singular}{λατρεία}{worship} προσφέρειν\FnParse{g}{present active infinitive} τῷ θεῷ. 
\MainTextVerseMark{16}{3}καὶ ταῦτα ποιήσουσιν\FnParse{a}{future active indicative 3rd plural} ὅτι οὐκ ἔγνωσαν\FnParse{b}{aorist active indicative 3rd plural} τὸν πατέρα οὐδὲ ἐμέ. 
\MainTextVerseMark{16}{4}ἀλλὰ ταῦτα λελάληκα\FnParse{a}{perfect active indicative 1st singular} ὑμῖν ἵνα ὅταν ἔλθῃ\FnParse{b}{aorist active subjunctive 3rd singular} ἡ ὥρα ⸀αὐτῶν μνημονεύητε\FnParseFormGloss{c}{present active subjunctive 2nd plural}{μνημονεύω}{to remember} ⸁αὐτῶν ὅτι ἐγὼ εἶπον\FnParse{d}{aorist active indicative 1st singular} ὑμῖν. Ταῦτα δὲ ὑμῖν ἐξ ἀρχῆς οὐκ εἶπον,\FnParse{e}{aorist active indicative 1st singular} ὅτι μεθ’ ὑμῶν ἤμην.\FnParse{f}{imperfect middle indicative 1st singular} 
\MainTextVerseMark{16}{5}νῦν δὲ ὑπάγω\FnParse{a}{present active indicative 1st singular} πρὸς τὸν πέμψαντά\FnParse{b}{aorist active participle accusative masculine singular} με καὶ οὐδεὶς ἐξ ὑμῶν ἐρωτᾷ\FnParse{c}{present active indicative 3rd singular} με· Ποῦ ὑπάγεις;\FnParse{d}{present active indicative 2nd singular} 
\MainTextVerseMark{16}{6}ἀλλ’ ὅτι ταῦτα λελάληκα\FnParse{a}{perfect active indicative 1st singular} ὑμῖν ἡ λύπη\FnParseFormGloss{b}{nominative feminine singular}{λύπη}{sorrow} πεπλήρωκεν\FnParse{c}{perfect active indicative 3rd singular} ὑμῶν τὴν καρδίαν. 
\MainTextVerseMark{16}{7}ἀλλ’ ἐγὼ τὴν ἀλήθειαν λέγω\FnParse{a}{present active indicative 1st singular} ὑμῖν, συμφέρει\FnParseFormGloss{b}{present active indicative 3rd singular}{συμφέρω}{to bring together} ὑμῖν ἵνα ἐγὼ ἀπέλθω.\FnParse{c}{aorist active subjunctive 1st singular} ἐὰν ⸀γὰρ μὴ ἀπέλθω,\FnParse{d}{aorist active subjunctive 1st singular} ὁ παράκλητος\FnParseFormGloss{e}{nominative masculine singular}{παράκλητος}{counselor} ⸂οὐ μὴ ἔλθῃ⸃\FnParse{f}{aorist active subjunctive 3rd singular} πρὸς ὑμᾶς· ἐὰν δὲ πορευθῶ,\FnParse{g}{aorist passive subjunctive 1st singular} πέμψω\FnParse{h}{future active indicative 1st singular} αὐτὸν πρὸς ὑμᾶς. 
\MainTextVerseMark{16}{8}καὶ ἐλθὼν\FnParse{a}{aorist active participle nominative masculine singular} ἐκεῖνος ἐλέγξει\FnParseFormGloss{b}{future active indicative 3rd singular}{ἐλέγχω}{to expose} τὸν κόσμον περὶ ἁμαρτίας καὶ περὶ δικαιοσύνης καὶ περὶ κρίσεως· 
\MainTextVerseMark{16}{9}περὶ ἁμαρτίας μέν, ὅτι οὐ πιστεύουσιν\FnParse{a}{present active indicative 3rd plural} εἰς ἐμέ· 
\MainTextVerseMark{16}{10}περὶ δικαιοσύνης δέ, ὅτι πρὸς τὸν ⸀πατέρα ὑπάγω\FnParse{a}{present active indicative 1st singular} καὶ οὐκέτι θεωρεῖτέ\FnParse{b}{present active indicative 2nd plural} με· 
\MainTextVerseMark{16}{11}περὶ δὲ κρίσεως, ὅτι ὁ ἄρχων τοῦ κόσμου τούτου κέκριται.\FnParse{a}{perfect passive indicative 3rd singular} 
\MainTextVerseMark{16}{12}Ἔτι πολλὰ ἔχω\FnParse{a}{present active indicative 1st singular} ⸂ὑμῖν λέγειν⸃,\FnParse{b}{present active infinitive} ἀλλ’ οὐ δύνασθε\FnParse{c}{present middle indicative 2nd plural} βαστάζειν\FnParseFormGloss{d}{present active infinitive}{βαστάζω}{to carry} ἄρτι· 
\MainTextVerseMark{16}{13}ὅταν δὲ ἔλθῃ\FnParse{a}{aorist active subjunctive 3rd singular} ἐκεῖνος, τὸ πνεῦμα τῆς ἀληθείας, ὁδηγήσει\FnParseFormGloss{b}{future active indicative 3rd singular}{ὁδηγέω}{to lead} ὑμᾶς ⸂ἐν τῇ ἀληθείᾳ πάσῃ⸃, οὐ γὰρ λαλήσει\FnParse{c}{future active indicative 3rd singular} ἀφ’ ἑαυτοῦ, ἀλλ’ ὅσα ⸀ἀκούσει\FnParse{d}{future active indicative 3rd singular} λαλήσει,\FnParse{e}{future active indicative 3rd singular} καὶ τὰ ἐρχόμενα\FnParse{f}{present middle participle accusative neuter plural} ἀναγγελεῖ\FnParseFormGloss{g}{future active indicative 3rd singular}{ἀναγγέλλω}{to tell} ὑμῖν. 
\MainTextVerseMark{16}{14}ἐκεῖνος ἐμὲ δοξάσει,\FnParse{a}{future active indicative 3rd singular} ὅτι ἐκ τοῦ ἐμοῦ λήμψεται\FnParse{b}{future middle indicative 3rd singular} καὶ ἀναγγελεῖ\FnParseFormGloss{c}{future active indicative 3rd singular}{ἀναγγέλλω}{to tell} ὑμῖν. 
\MainTextVerseMark{16}{15}πάντα ὅσα ἔχει\FnParse{a}{present active indicative 3rd singular} ὁ πατὴρ ἐμά ἐστιν·\FnParse{b}{present active indicative 3rd singular} διὰ τοῦτο εἶπον\FnParse{c}{aorist active indicative 1st singular} ὅτι ἐκ τοῦ ἐμοῦ λαμβάνει\FnParse{d}{present active indicative 3rd singular} καὶ ἀναγγελεῖ\FnParseFormGloss{e}{future active indicative 3rd singular}{ἀναγγέλλω}{to tell} ὑμῖν. 
\MainTextVerseMark{16}{16}Μικρὸν καὶ ⸀οὐκέτι θεωρεῖτέ\FnParse{a}{present active indicative 2nd plural} με, καὶ πάλιν μικρὸν καὶ ὄψεσθέ\FnParse{b}{future middle indicative 2nd plural} ⸀με. 
\MainTextVerseMark{16}{17}εἶπαν\FnParse{a}{aorist active indicative 3rd plural} οὖν ἐκ τῶν μαθητῶν αὐτοῦ πρὸς ἀλλήλους· Τί ἐστιν\FnParse{b}{present active indicative 3rd singular} τοῦτο ὃ λέγει\FnParse{c}{present active indicative 3rd singular} ἡμῖν· Μικρὸν καὶ οὐ θεωρεῖτέ\FnParse{d}{present active indicative 2nd plural} με, καὶ πάλιν μικρὸν καὶ ὄψεσθέ\FnParse{e}{future middle indicative 2nd plural} με; καί· ⸀Ὅτι ὑπάγω\FnParse{f}{present active indicative 1st singular} πρὸς τὸν πατέρα; 
\MainTextVerseMark{16}{18}ἔλεγον\FnParse{a}{imperfect active indicative 3rd plural} οὖν· ⸂Τί ἐστιν\FnParse{b}{present active indicative 3rd singular} τοῦτο⸃ ὃ ⸀λέγει\FnParse{c}{present active indicative 3rd singular} μικρόν; οὐκ οἴδαμεν\FnParse{d}{perfect active indicative 1st plural} τί λαλεῖ.\FnParse{e}{present active indicative 3rd singular} 
\MainTextVerseMark{16}{19}⸀ἔγνω\FnParse{a}{aorist active indicative 3rd singular} Ἰησοῦς ὅτι ἤθελον\FnParse{b}{imperfect active indicative 3rd plural} αὐτὸν ἐρωτᾶν,\FnParse{c}{present active infinitive} καὶ εἶπεν\FnParse{d}{aorist active indicative 3rd singular} αὐτοῖς· Περὶ τούτου ζητεῖτε\FnParse{e}{present active indicative 2nd plural} μετ’ ἀλλήλων ὅτι εἶπον·\FnParse{f}{aorist active indicative 1st singular} Μικρὸν καὶ οὐ θεωρεῖτέ\FnParse{g}{present active indicative 2nd plural} με, καὶ πάλιν μικρὸν καὶ ὄψεσθέ\FnParse{h}{future middle indicative 2nd plural} με; 
\MainTextVerseMark{16}{20}ἀμὴν ἀμὴν λέγω\FnParse{a}{present active indicative 1st singular} ὑμῖν ὅτι κλαύσετε\FnParse{b}{future active indicative 2nd plural} καὶ θρηνήσετε\FnParseFormGloss{c}{future active indicative 2nd plural}{θρηνέω}{to sing a funeral dirge} ὑμεῖς, ὁ δὲ κόσμος χαρήσεται·\FnParse{d}{future passive indicative 3rd singular} ⸀ὑμεῖς λυπηθήσεσθε,\FnParseFormGloss{e}{future passive indicative 2nd plural}{λυπέω}{to cause sorrow} ἀλλ’ ἡ λύπη\FnParseFormGloss{f}{nominative feminine singular}{λύπη}{sorrow} ὑμῶν εἰς χαρὰν γενήσεται.\FnParse{g}{future middle indicative 3rd singular} 
\MainTextVerseMark{16}{21}ἡ γυνὴ ὅταν τίκτῃ\FnParseFormGloss{a}{present active subjunctive 3rd singular}{τίκτω}{to give birth to} λύπην\FnParseFormGloss{b}{accusative feminine singular}{λύπη}{sorrow} ἔχει,\FnParse{c}{present active indicative 3rd singular} ὅτι ἦλθεν\FnParse{d}{aorist active indicative 3rd singular} ἡ ὥρα αὐτῆς· ὅταν δὲ γεννήσῃ\FnParse{e}{aorist active subjunctive 3rd singular} τὸ παιδίον, οὐκέτι μνημονεύει\FnParseFormGloss{f}{present active indicative 3rd singular}{μνημονεύω}{to remember} τῆς θλίψεως διὰ τὴν χαρὰν ὅτι ἐγεννήθη\FnParse{g}{aorist passive indicative 3rd singular} ἄνθρωπος εἰς τὸν κόσμον. 
\MainTextVerseMark{16}{22}καὶ ὑμεῖς οὖν ⸂νῦν μὲν λύπην⸃\FnParseFormGloss{a}{accusative feminine singular}{λύπη}{sorrow} ἔχετε·\FnParse{b}{present active indicative 2nd plural} πάλιν δὲ ὄψομαι\FnParse{c}{future middle indicative 1st singular} ὑμᾶς, καὶ χαρήσεται\FnParse{d}{future passive indicative 3rd singular} ὑμῶν ἡ καρδία, καὶ τὴν χαρὰν ὑμῶν οὐδεὶς ⸀αἴρει\FnParse{e}{present active indicative 3rd singular} ἀφ’ ὑμῶν. 
\MainTextVerseMark{16}{23}καὶ ἐν ἐκείνῃ τῇ ἡμέρᾳ ἐμὲ οὐκ ἐρωτήσετε\FnParse{a}{future active indicative 2nd plural} οὐδέν· ἀμὴν ἀμὴν λέγω\FnParse{b}{present active indicative 1st singular} ὑμῖν, ⸂ἄν τι⸃ αἰτήσητε\FnParse{c}{aorist active subjunctive 2nd plural} τὸν πατέρα ⸂δώσει\FnParse{d}{future active indicative 3rd singular} ὑμῖν ἐν τῷ ὀνόματί μου⸃. 
\MainTextVerseMark{16}{24}ἕως ἄρτι οὐκ ᾐτήσατε\FnParse{a}{aorist active indicative 2nd plural} οὐδὲν ἐν τῷ ὀνόματί μου· αἰτεῖτε\FnParse{b}{present active imperative 2nd plural} καὶ λήμψεσθε,\FnParse{c}{future middle indicative 2nd plural} ἵνα ἡ χαρὰ ὑμῶν ᾖ\FnParse{d}{present active subjunctive 3rd singular} πεπληρωμένη.\FnParse{e}{perfect passive participle nominative feminine singular} 
\MainTextVerseMark{16}{25}Ταῦτα ἐν παροιμίαις\FnParseFormGloss{a}{dative feminine plural}{παροιμία}{figure of speech} λελάληκα\FnParse{b}{perfect active indicative 1st singular} ὑμῖν· ⸀ἔρχεται\FnParse{c}{present middle indicative 3rd singular} ὥρα ὅτε οὐκέτι ἐν παροιμίαις\FnParseFormGloss{d}{dative feminine plural}{παροιμία}{figure of speech} λαλήσω\FnParse{e}{future active indicative 1st singular} ὑμῖν ἀλλὰ παρρησίᾳ περὶ τοῦ πατρὸς ⸀ἀπαγγελῶ\FnParse{f}{future active indicative 1st singular} ὑμῖν. 
\MainTextVerseMark{16}{26}ἐν ἐκείνῃ τῇ ἡμέρᾳ ἐν τῷ ὀνόματί μου αἰτήσεσθε,\FnParse{a}{future middle indicative 2nd plural} καὶ οὐ λέγω\FnParse{b}{present active indicative 1st singular} ὑμῖν ὅτι ἐγὼ ἐρωτήσω\FnParse{c}{future active indicative 1st singular} τὸν πατέρα περὶ ὑμῶν· 
\MainTextVerseMark{16}{27}αὐτὸς γὰρ ὁ πατὴρ φιλεῖ\FnParseFormGloss{a}{present active indicative 3rd singular}{φιλέω}{to love} ὑμᾶς, ὅτι ὑμεῖς ἐμὲ πεφιλήκατε\FnParseFormGloss{b}{perfect active indicative 2nd plural}{φιλέω}{to love} καὶ πεπιστεύκατε\FnParse{c}{perfect active indicative 2nd plural} ὅτι ἐγὼ παρὰ τοῦ ⸀θεοῦ ἐξῆλθον.\FnParse{d}{aorist active indicative 1st singular} 
\MainTextVerseMark{16}{28}ἐξῆλθον\FnParse{a}{aorist active indicative 1st singular} ⸀ἐκ τοῦ πατρὸς καὶ ἐλήλυθα\FnParse{b}{perfect active indicative 1st singular} εἰς τὸν κόσμον· πάλιν ἀφίημι\FnParse{c}{present active indicative 1st singular} τὸν κόσμον καὶ πορεύομαι\FnParse{d}{present middle indicative 1st singular} πρὸς τὸν πατέρα. 
\MainTextVerseMark{16}{29}⸀Λέγουσιν\FnParse{a}{present active indicative 3rd plural} οἱ μαθηταὶ αὐτοῦ· Ἴδε νῦν ⸀ἐν παρρησίᾳ λαλεῖς,\FnParse{b}{present active indicative 2nd singular} καὶ παροιμίαν\FnParseFormGloss{c}{accusative feminine singular}{παροιμία}{figure of speech} οὐδεμίαν λέγεις.\FnParse{d}{present active indicative 2nd singular} 
\MainTextVerseMark{16}{30}νῦν οἴδαμεν\FnParse{a}{perfect active indicative 1st plural} ὅτι οἶδας\FnParse{b}{perfect active indicative 2nd singular} πάντα καὶ οὐ χρείαν ἔχεις\FnParse{c}{present active indicative 2nd singular} ἵνα τίς σε ἐρωτᾷ·\FnParse{d}{present active subjunctive 3rd singular} ἐν τούτῳ πιστεύομεν\FnParse{e}{present active indicative 1st plural} ὅτι ἀπὸ θεοῦ ἐξῆλθες.\FnParse{f}{aorist active indicative 2nd singular} 
\MainTextVerseMark{16}{31}ἀπεκρίθη\FnParse{a}{aorist passive indicative 3rd singular} ⸀αὐτοῖς Ἰησοῦς· Ἄρτι πιστεύετε;\FnParse{b}{present active indicative 2nd plural} 
\MainTextVerseMark{16}{32}ἰδοὺ ἔρχεται\FnParse{a}{present middle indicative 3rd singular} ὥρα ⸀καὶ ἐλήλυθεν\FnParse{b}{perfect active indicative 3rd singular} ἵνα σκορπισθῆτε\FnParseFormGloss{c}{aorist passive subjunctive 2nd plural}{σκορπίζω}{to scatter} ἕκαστος εἰς τὰ ἴδια κἀμὲ μόνον ἀφῆτε·\FnParse{d}{aorist active subjunctive 2nd plural} καὶ οὐκ εἰμὶ\FnParse{e}{present active indicative 1st singular} μόνος, ὅτι ὁ πατὴρ μετ’ ἐμοῦ ἐστιν.\FnParse{f}{present active indicative 3rd singular} 
\MainTextVerseMark{16}{33}ταῦτα λελάληκα\FnParse{a}{perfect active indicative 1st singular} ὑμῖν ἵνα ἐν ἐμοὶ εἰρήνην ἔχητε·\FnParse{b}{present active subjunctive 2nd plural} ἐν τῷ κόσμῳ θλῖψιν ἔχετε,\FnParse{c}{present active indicative 2nd plural} ἀλλὰ θαρσεῖτε,\FnParseFormGloss{d}{present active imperative 2nd plural}{θαρσέω}{take heart!} ἐγὼ νενίκηκα\FnParseFormGloss{e}{perfect active indicative 1st singular}{νικάω}{to overcome} τὸν κόσμον. 
\ChapterHeader{17}{Chapter 17}

\MainTextVerseMark{17}{1}Ταῦτα ⸀ἐλάλησεν\FnParse{a}{aorist active indicative 3rd singular} Ἰησοῦς, καὶ ⸂ἐπάρας\FnParseFormGloss{b}{aorist active participle nominative masculine singular}{ἐπαίρω}{to lift up} τοὺς ὀφθαλμοὺς αὐτοῦ εἰς τὸν οὐρανὸν εἶπεν⸃·\FnParse{c}{aorist active indicative 3rd singular} Πάτερ, ἐλήλυθεν\FnParse{d}{perfect active indicative 3rd singular} ἡ ὥρα· δόξασόν\FnParse{e}{aorist active imperative 2nd singular} σου τὸν υἱόν, ⸀ἵνα ὁ ⸀υἱὸς δοξάσῃ\FnParse{f}{aorist active subjunctive 3rd singular} σέ, 
\MainTextVerseMark{17}{2}καθὼς ἔδωκας\FnParse{a}{aorist active indicative 2nd singular} αὐτῷ ἐξουσίαν πάσης σαρκός, ἵνα πᾶν ὃ δέδωκας\FnParse{b}{perfect active indicative 2nd singular} αὐτῷ ⸀δώσῃ\FnParse{c}{aorist active subjunctive 3rd singular} αὐτοῖς ζωὴν αἰώνιον. 
\MainTextVerseMark{17}{3}αὕτη δέ ἐστιν\FnParse{a}{present active indicative 3rd singular} ἡ αἰώνιος ζωὴ ἵνα γινώσκωσι\FnParse{b}{present active subjunctive 3rd plural} σὲ τὸν μόνον ἀληθινὸν\FnParseFormGloss{c}{accusative masculine singular}{ἀληθινός}{true} θεὸν καὶ ὃν ἀπέστειλας\FnParse{d}{aorist active indicative 2nd singular} Ἰησοῦν Χριστόν. 
\MainTextVerseMark{17}{4}ἐγώ σε ἐδόξασα\FnParse{a}{aorist active indicative 1st singular} ἐπὶ τῆς γῆς, τὸ ἔργον ⸀τελειώσας\FnParseFormGloss{b}{aorist active participle nominative masculine singular}{τελειόω}{to perfect} ὃ δέδωκάς\FnParse{c}{perfect active indicative 2nd singular} μοι ἵνα ποιήσω·\FnParse{d}{aorist active subjunctive 1st singular} 
\MainTextVerseMark{17}{5}καὶ νῦν δόξασόν\FnParse{a}{aorist active imperative 2nd singular} με σύ, πάτερ, παρὰ σεαυτῷ τῇ δόξῃ ᾗ εἶχον\FnParse{b}{imperfect active indicative 1st singular} πρὸ τοῦ τὸν κόσμον εἶναι\FnParse{c}{present active infinitive} παρὰ σοί. 
\MainTextVerseMark{17}{6}Ἐφανέρωσά\FnParse{a}{aorist active indicative 1st singular} σου τὸ ὄνομα τοῖς ἀνθρώποις οὓς ⸀ἔδωκάς\FnParse{b}{aorist active indicative 2nd singular} μοι ἐκ τοῦ κόσμου. σοὶ ἦσαν\FnParse{c}{imperfect active indicative 3rd plural} κἀμοὶ αὐτοὺς ⸁ἔδωκας,\FnParse{d}{aorist active indicative 2nd singular} καὶ τὸν λόγον σου τετήρηκαν.\FnParse{e}{perfect active indicative 3rd plural} 
\MainTextVerseMark{17}{7}νῦν ἔγνωκαν\FnParse{a}{perfect active indicative 3rd plural} ὅτι πάντα ὅσα ⸀δέδωκάς\FnParse{b}{perfect active indicative 2nd singular} μοι παρὰ σοῦ ⸀εἰσιν·\FnParse{c}{present active indicative 3rd plural} 
\MainTextVerseMark{17}{8}ὅτι τὰ ῥήματα ἃ ⸀ἔδωκάς\FnParse{a}{aorist active indicative 2nd singular} μοι δέδωκα\FnParse{b}{perfect active indicative 1st singular} αὐτοῖς, καὶ αὐτοὶ ἔλαβον\FnParse{c}{aorist active indicative 3rd plural} καὶ ἔγνωσαν\FnParse{d}{aorist active indicative 3rd plural} ἀληθῶς\FnFormGloss{e}{ἀληθῶς}{truly} ὅτι παρὰ σοῦ ἐξῆλθον,\FnParse{f}{aorist active indicative 1st singular} καὶ ἐπίστευσαν\FnParse{g}{aorist active indicative 3rd plural} ὅτι σύ με ἀπέστειλας.\FnParse{h}{aorist active indicative 2nd singular} 
\MainTextVerseMark{17}{9}ἐγὼ περὶ αὐτῶν ἐρωτῶ·\FnParse{a}{present active indicative 1st singular} οὐ περὶ τοῦ κόσμου ἐρωτῶ\FnParse{b}{present active indicative 1st singular} ἀλλὰ περὶ ὧν δέδωκάς\FnParse{c}{perfect active indicative 2nd singular} μοι, ὅτι σοί εἰσιν,\FnParse{d}{present active indicative 3rd plural} 
\MainTextVerseMark{17}{10}καὶ τὰ ἐμὰ πάντα σά\FnParseFormGloss{a}{nominative neuter plural}{σός}{your} ἐστιν\FnParse{b}{present active indicative 3rd singular} καὶ τὰ σὰ\FnParseFormGloss{c}{nominative neuter plural}{σός}{your} ἐμά, καὶ δεδόξασμαι\FnParse{d}{perfect passive indicative 1st singular} ἐν αὐτοῖς. 
\MainTextVerseMark{17}{11}καὶ οὐκέτι εἰμὶ\FnParse{a}{present active indicative 1st singular} ἐν τῷ κόσμῳ, καὶ ⸀αὐτοὶ ἐν τῷ κόσμῳ εἰσίν,\FnParse{b}{present active indicative 3rd plural} κἀγὼ πρὸς σὲ ἔρχομαι.\FnParse{c}{present middle indicative 1st singular} πάτερ ἅγιε, τήρησον\FnParse{d}{aorist active imperative 2nd singular} αὐτοὺς ἐν τῷ ὀνόματί σου ᾧ δέδωκάς\FnParse{e}{perfect active indicative 2nd singular} μοι, ἵνα ὦσιν\FnParse{f}{present active subjunctive 3rd plural} ἓν ⸀καθὼς ἡμεῖς. 
\MainTextVerseMark{17}{12}ὅτε ἤμην\FnParse{a}{imperfect middle indicative 1st singular} μετ’ ⸀αὐτῶν ἐγὼ ἐτήρουν\FnParse{b}{imperfect active indicative 1st singular} αὐτοὺς ἐν τῷ ὀνόματί σου ⸀ᾧ δέδωκάς\FnParse{c}{perfect active indicative 2nd singular} μοι, ⸀καὶ ἐφύλαξα,\FnParse{d}{aorist active indicative 1st singular} καὶ οὐδεὶς ἐξ αὐτῶν ἀπώλετο\FnParse{e}{aorist middle indicative 3rd singular} εἰ μὴ ὁ υἱὸς τῆς ἀπωλείας,\FnParseFormGloss{f}{genitive feminine singular}{ἀπώλεια}{destruction} ἵνα ἡ γραφὴ πληρωθῇ.\FnParse{g}{aorist passive subjunctive 3rd singular} 
\MainTextVerseMark{17}{13}νῦν δὲ πρὸς σὲ ἔρχομαι,\FnParse{a}{present middle indicative 1st singular} καὶ ταῦτα λαλῶ\FnParse{b}{present active indicative 1st singular} ἐν τῷ κόσμῳ ἵνα ἔχωσιν\FnParse{c}{present active subjunctive 3rd plural} τὴν χαρὰν τὴν ἐμὴν πεπληρωμένην\FnParse{d}{perfect passive participle accusative feminine singular} ἐν ⸀ἑαυτοῖς. 
\MainTextVerseMark{17}{14}ἐγὼ δέδωκα\FnParse{a}{perfect active indicative 1st singular} αὐτοῖς τὸν λόγον σου, καὶ ὁ κόσμος ἐμίσησεν\FnParse{b}{aorist active indicative 3rd singular} αὐτούς, ὅτι οὐκ εἰσὶν\FnParse{c}{present active indicative 3rd plural} ἐκ τοῦ κόσμου καθὼς ἐγὼ οὐκ εἰμὶ\FnParse{d}{present active indicative 1st singular} ἐκ τοῦ κόσμου. 
\MainTextVerseMark{17}{15}οὐκ ἐρωτῶ\FnParse{a}{present active indicative 1st singular} ἵνα ἄρῃς\FnParse{b}{aorist active subjunctive 2nd singular} αὐτοὺς ἐκ τοῦ κόσμου ἀλλ’ ἵνα τηρήσῃς\FnParse{c}{aorist active subjunctive 2nd singular} αὐτοὺς ἐκ τοῦ πονηροῦ. 
\MainTextVerseMark{17}{16}ἐκ τοῦ κόσμου οὐκ εἰσὶν\FnParse{a}{present active indicative 3rd plural} καθὼς ἐγὼ ⸂οὐκ εἰμὶ\FnParse{b}{present active indicative 1st singular} ἐκ τοῦ κόσμου⸃. 
\MainTextVerseMark{17}{17}ἁγίασον\FnParseFormGloss{a}{aorist active imperative 2nd singular}{ἁγιάζω}{to sanctify} αὐτοὺς ἐν τῇ ⸀ἀληθείᾳ· ὁ λόγος ὁ σὸς\FnParseFormGloss{b}{nominative masculine singular}{σός}{your} ἀλήθειά ἐστιν.\FnParse{c}{present active indicative 3rd singular} 
\MainTextVerseMark{17}{18}καθὼς ἐμὲ ἀπέστειλας\FnParse{a}{aorist active indicative 2nd singular} εἰς τὸν κόσμον, κἀγὼ ἀπέστειλα\FnParse{b}{aorist active indicative 1st singular} αὐτοὺς εἰς τὸν κόσμον· 
\MainTextVerseMark{17}{19}καὶ ὑπὲρ αὐτῶν ἐγὼ ἁγιάζω\FnParseFormGloss{a}{present active indicative 1st singular}{ἁγιάζω}{to sanctify} ἐμαυτόν, ἵνα ⸂ὦσιν\FnParse{b}{present active subjunctive 3rd plural} καὶ αὐτοὶ⸃ ἡγιασμένοι\FnParseFormGloss{c}{perfect passive participle nominative masculine plural}{ἁγιάζω}{to sanctify} ἐν ἀληθείᾳ. 
\MainTextVerseMark{17}{20}Οὐ περὶ τούτων δὲ ἐρωτῶ\FnParse{a}{present active indicative 1st singular} μόνον, ἀλλὰ καὶ περὶ τῶν πιστευόντων\FnParse{b}{present active participle genitive masculine plural} διὰ τοῦ λόγου αὐτῶν εἰς ἐμέ, 
\MainTextVerseMark{17}{21}ἵνα πάντες ἓν ὦσιν,\FnParse{a}{present active subjunctive 3rd plural} καθὼς σύ, ⸀πάτερ, ἐν ἐμοὶ κἀγὼ ἐν σοί, ἵνα καὶ αὐτοὶ ἐν ⸀ἡμῖν ὦσιν,\FnParse{b}{present active subjunctive 3rd plural} ἵνα ὁ κόσμος ⸀πιστεύῃ\FnParse{c}{present active subjunctive 3rd singular} ὅτι σύ με ἀπέστειλας.\FnParse{d}{aorist active indicative 2nd singular} 
\MainTextVerseMark{17}{22}κἀγὼ τὴν δόξαν ἣν δέδωκάς\FnParse{a}{perfect active indicative 2nd singular} μοι δέδωκα\FnParse{b}{perfect active indicative 1st singular} αὐτοῖς, ἵνα ὦσιν\FnParse{c}{present active subjunctive 3rd plural} ἓν καθὼς ἡμεῖς ⸀ἕν, 
\MainTextVerseMark{17}{23}ἐγὼ ἐν αὐτοῖς καὶ σὺ ἐν ἐμοί, ἵνα ὦσιν\FnParse{a}{present active subjunctive 3rd plural} τετελειωμένοι\FnParseFormGloss{b}{perfect passive participle nominative masculine plural}{τελειόω}{to perfect} εἰς ἕν, ⸀ἵνα γινώσκῃ\FnParse{c}{present active subjunctive 3rd singular} ὁ κόσμος ὅτι σύ με ἀπέστειλας\FnParse{d}{aorist active indicative 2nd singular} καὶ ἠγάπησας\FnParse{e}{aorist active indicative 2nd singular} αὐτοὺς καθὼς ἐμὲ ἠγάπησας.\FnParse{f}{aorist active indicative 2nd singular} 
\MainTextVerseMark{17}{24}⸀πάτερ, ⸀ὃ δέδωκάς\FnParse{a}{perfect active indicative 2nd singular} μοι, θέλω\FnParse{b}{present active indicative 1st singular} ἵνα ὅπου εἰμὶ\FnParse{c}{present active indicative 1st singular} ἐγὼ κἀκεῖνοι\FnParseFormGloss{d}{nominative masculine plural}{κἀκεῖνος}{and that one} ὦσιν\FnParse{e}{present active subjunctive 3rd plural} μετ’ ἐμοῦ, ἵνα θεωρῶσιν\FnParse{f}{present active subjunctive 3rd plural} τὴν δόξαν τὴν ἐμὴν ἣν ⸀δέδωκάς\FnParse{g}{perfect active indicative 2nd singular} μοι, ὅτι ἠγάπησάς\FnParse{h}{aorist active indicative 2nd singular} με πρὸ καταβολῆς\FnParseFormGloss{i}{genitive feminine singular}{καταβολή}{creation} κόσμου. 
\MainTextVerseMark{17}{25}⸀Πάτερ δίκαιε, καὶ ὁ κόσμος σε οὐκ ἔγνω,\FnParse{a}{aorist active indicative 3rd singular} ἐγὼ δέ σε ἔγνων,\FnParse{b}{aorist active indicative 1st singular} καὶ οὗτοι ἔγνωσαν\FnParse{c}{aorist active indicative 3rd plural} ὅτι σύ με ἀπέστειλας,\FnParse{d}{aorist active indicative 2nd singular} 
\MainTextVerseMark{17}{26}καὶ ἐγνώρισα\FnParseFormGloss{a}{aorist active indicative 1st singular}{γνωρίζω}{to make known} αὐτοῖς τὸ ὄνομά σου καὶ γνωρίσω,\FnParseFormGloss{b}{future active indicative 1st singular}{γνωρίζω}{to make known} ἵνα ἡ ἀγάπη ἣν ἠγάπησάς\FnParse{c}{aorist active indicative 2nd singular} με ἐν αὐτοῖς ᾖ\FnParse{d}{present active subjunctive 3rd singular} κἀγὼ ἐν αὐτοῖς. 
\ChapterHeader{18}{Chapter 18}

\MainTextVerseMark{18}{1}Ταῦτα ⸀εἰπὼν\FnParse{a}{aorist active participle nominative masculine singular} Ἰησοῦς ἐξῆλθεν\FnParse{b}{aorist active indicative 3rd singular} σὺν τοῖς μαθηταῖς αὐτοῦ πέραν\FnFormGloss{c}{πέραν}{on the other side} τοῦ χειμάρρου\FnParseFormGloss{d}{genitive masculine singular}{χείμαρρος}{valley} ⸀τοῦ Κεδρὼν\FnParseFormGloss{e}{genitive masculine singular}{Κεδρών}{Kidron} ὅπου ἦν\FnParse{f}{imperfect active indicative 3rd singular} κῆπος,\FnParseFormGloss{g}{nominative masculine singular}{κῆπος}{garden} εἰς ὃν εἰσῆλθεν\FnParse{h}{aorist active indicative 3rd singular} αὐτὸς καὶ οἱ μαθηταὶ αὐτοῦ. 
\MainTextVerseMark{18}{2}ᾔδει\FnParse{a}{pluperfect active indicative 3rd singular} δὲ καὶ Ἰούδας ὁ παραδιδοὺς\FnParse{b}{present active participle nominative masculine singular} αὐτὸν τὸν τόπον, ὅτι πολλάκις\FnFormGloss{c}{πολλάκις}{many times} ⸀συνήχθη\FnParse{d}{aorist passive indicative 3rd singular} Ἰησοῦς ἐκεῖ μετὰ τῶν μαθητῶν αὐτοῦ. 
\MainTextVerseMark{18}{3}ὁ οὖν Ἰούδας λαβὼν\FnParse{a}{aorist active participle nominative masculine singular} τὴν σπεῖραν\FnParseFormGloss{b}{accusative feminine singular}{σπεῖρα}{company of soldiers} καὶ ἐκ τῶν ἀρχιερέων καὶ ⸂ἐκ τῶν⸃ Φαρισαίων ὑπηρέτας\FnParseFormGloss{c}{accusative masculine plural}{ὑπηρέτης}{servant} ἔρχεται\FnParse{d}{present middle indicative 3rd singular} ἐκεῖ μετὰ φανῶν\FnParseFormGloss{e}{genitive masculine plural}{φανός}{torch} καὶ λαμπάδων\FnParseFormGloss{f}{genitive feminine plural}{λαμπάς}{lamp} καὶ ὅπλων.\FnParseFormGloss{g}{genitive neuter plural}{ὅπλον}{instrument} 
\MainTextVerseMark{18}{4}Ἰησοῦς ⸀οὖν εἰδὼς\FnParse{a}{perfect active participle nominative masculine singular} πάντα τὰ ἐρχόμενα\FnParse{b}{present middle participle accusative neuter plural} ἐπ’ αὐτὸν ⸂ἐξῆλθεν,\FnParse{c}{aorist active indicative 3rd singular} καὶ λέγει⸃\FnParse{d}{present active indicative 3rd singular} αὐτοῖς· Τίνα ζητεῖτε;\FnParse{e}{present active indicative 2nd plural} 
\MainTextVerseMark{18}{5}ἀπεκρίθησαν\FnParse{a}{aorist passive indicative 3rd plural} αὐτῷ· Ἰησοῦν τὸν Ναζωραῖον.\FnParseFormGloss{b}{accusative masculine singular}{Ναζωραῖος}{Nazarene} λέγει\FnParse{c}{present active indicative 3rd singular} ⸀αὐτοῖς· Ἐγώ εἰμι.\FnParse{d}{present active indicative 1st singular} εἱστήκει\FnParse{e}{pluperfect active indicative 3rd singular} δὲ καὶ Ἰούδας ὁ παραδιδοὺς\FnParse{f}{present active participle nominative masculine singular} αὐτὸν μετ’ αὐτῶν. 
\MainTextVerseMark{18}{6}ὡς οὖν εἶπεν\FnParse{a}{aorist active indicative 3rd singular} ⸀αὐτοῖς· Ἐγώ εἰμι,\FnParse{b}{present active indicative 1st singular} ἀπῆλθον\FnParse{c}{aorist active indicative 3rd plural} εἰς τὰ ὀπίσω καὶ ἔπεσαν\FnParse{d}{aorist active indicative 3rd plural} χαμαί.\FnFormGloss{e}{χαμαί}{to the ground} 
\MainTextVerseMark{18}{7}πάλιν οὖν ⸂ἐπηρώτησεν\FnParse{a}{aorist active indicative 3rd singular} αὐτούς⸃· Τίνα ζητεῖτε;\FnParse{b}{present active indicative 2nd plural} οἱ δὲ εἶπαν·\FnParse{c}{aorist active indicative 3rd plural} Ἰησοῦν τὸν Ναζωραῖον.\FnParseFormGloss{d}{accusative masculine singular}{Ναζωραῖος}{Nazarene} 
\MainTextVerseMark{18}{8}ἀπεκρίθη\FnParse{a}{aorist passive indicative 3rd singular} Ἰησοῦς· Εἶπον\FnParse{b}{aorist active indicative 1st singular} ὑμῖν ὅτι ἐγώ εἰμι·\FnParse{c}{present active indicative 1st singular} εἰ οὖν ἐμὲ ζητεῖτε,\FnParse{d}{present active indicative 2nd plural} ἄφετε\FnParse{e}{aorist active imperative 2nd plural} τούτους ὑπάγειν·\FnParse{f}{present active infinitive} 
\MainTextVerseMark{18}{9}ἵνα πληρωθῇ\FnParse{a}{aorist passive subjunctive 3rd singular} ὁ λόγος ὃν εἶπεν\FnParse{b}{aorist active indicative 3rd singular} ὅτι Οὓς δέδωκάς\FnParse{c}{perfect active indicative 2nd singular} μοι οὐκ ἀπώλεσα\FnParse{d}{aorist active indicative 1st singular} ἐξ αὐτῶν οὐδένα. 
\MainTextVerseMark{18}{10}Σίμων οὖν Πέτρος ἔχων\FnParse{a}{present active participle nominative masculine singular} μάχαιραν\FnParseFormGloss{b}{accusative feminine singular}{μάχαιρα}{sword} εἵλκυσεν\FnParseFormGloss{c}{aorist active indicative 3rd singular}{ἕλκω}{to drag} αὐτὴν καὶ ἔπαισεν\FnParseFormGloss{d}{aorist active indicative 3rd singular}{παίω}{to strike} τὸν τοῦ ἀρχιερέως δοῦλον καὶ ἀπέκοψεν\FnParseFormGloss{e}{aorist active indicative 3rd singular}{ἀποκόπτω}{to cut off} αὐτοῦ τὸ ⸀ὠτάριον\FnParseFormGloss{f}{accusative neuter singular}{ὠτάριον}{ear} τὸ δεξιόν. ἦν\FnParse{g}{imperfect active indicative 3rd singular} δὲ ὄνομα τῷ δούλῳ Μάλχος.\FnParseFormGloss{h}{nominative masculine singular}{Μάλχος}{Malchus} 
\MainTextVerseMark{18}{11}εἶπεν\FnParse{a}{aorist active indicative 3rd singular} οὖν ὁ Ἰησοῦς τῷ Πέτρῳ· Βάλε\FnParse{b}{aorist active imperative 2nd singular} τὴν ⸀μάχαιραν\FnParseFormGloss{c}{accusative feminine singular}{μάχαιρα}{sword} εἰς τὴν θήκην·\FnParseFormGloss{d}{accusative feminine singular}{θήκη}{sheath} τὸ ποτήριον ὃ δέδωκέν\FnParse{e}{perfect active indicative 3rd singular} μοι ὁ πατὴρ οὐ μὴ πίω\FnParse{f}{aorist active subjunctive 1st singular} αὐτό; 
\MainTextVerseMark{18}{12}Ἡ οὖν σπεῖρα\FnParseFormGloss{a}{nominative feminine singular}{σπεῖρα}{company of soldiers} καὶ ὁ χιλίαρχος\FnParseFormGloss{b}{nominative masculine singular}{χιλίαρχος}{military officer} καὶ οἱ ὑπηρέται\FnParseFormGloss{c}{nominative masculine plural}{ὑπηρέτης}{servant} τῶν Ἰουδαίων συνέλαβον\FnParseFormGloss{d}{aorist active indicative 3rd plural}{συλλαμβάνω}{to seize} τὸν Ἰησοῦν καὶ ἔδησαν\FnParse{e}{aorist active indicative 3rd plural} αὐτὸν 
\MainTextVerseMark{18}{13}καὶ ⸀ἤγαγον\FnParse{a}{aorist active indicative 3rd plural} πρὸς Ἅνναν\FnParseFormGloss{b}{accusative masculine singular}{Ἅννας}{Annas} πρῶτον· ἦν\FnParse{c}{imperfect active indicative 3rd singular} γὰρ πενθερὸς\FnParseFormGloss{d}{nominative masculine singular}{πενθερός}{father-in-law} τοῦ Καϊάφα,\FnParseFormGloss{e}{genitive masculine singular}{Καϊάφας}{Caiaphas} ὃς ἦν\FnParse{f}{imperfect active indicative 3rd singular} ἀρχιερεὺς τοῦ ἐνιαυτοῦ\FnParseFormGloss{g}{genitive masculine singular}{ἐνιαυτός}{year} ἐκείνου· 
\MainTextVerseMark{18}{14}ἦν\FnParse{a}{imperfect active indicative 3rd singular} δὲ Καϊάφας\FnParseFormGloss{b}{nominative masculine singular}{Καϊάφας}{Caiaphas} ὁ συμβουλεύσας\FnParseFormGloss{c}{aorist active participle nominative masculine singular}{συμβουλεύω}{to advise} τοῖς Ἰουδαίοις ὅτι συμφέρει\FnParseFormGloss{d}{present active indicative 3rd singular}{συμφέρω}{to bring together} ἕνα ἄνθρωπον ⸀ἀποθανεῖν\FnParse{e}{aorist active infinitive} ὑπὲρ τοῦ λαοῦ. 
\MainTextVerseMark{18}{15}Ἠκολούθει\FnParse{a}{imperfect active indicative 3rd singular} δὲ τῷ Ἰησοῦ Σίμων Πέτρος ⸀καὶ ἄλλος μαθητής. ὁ δὲ μαθητὴς ἐκεῖνος ἦν\FnParse{b}{imperfect active indicative 3rd singular} γνωστὸς\FnParseFormGloss{c}{nominative masculine singular}{γνωστός}{known} τῷ ἀρχιερεῖ καὶ συνεισῆλθεν\FnParseFormGloss{d}{aorist active indicative 3rd singular}{συνεισέρχομαι}{to enter together with} τῷ Ἰησοῦ εἰς τὴν αὐλὴν\FnParseFormGloss{e}{accusative feminine singular}{αὐλή}{palace} τοῦ ἀρχιερέως, 
\MainTextVerseMark{18}{16}ὁ δὲ Πέτρος εἱστήκει\FnParse{a}{pluperfect active indicative 3rd singular} πρὸς τῇ θύρᾳ ἔξω. ἐξῆλθεν\FnParse{b}{aorist active indicative 3rd singular} οὖν ὁ μαθητὴς ὁ ἄλλος ⸀ὁ γνωστὸς\FnParseFormGloss{c}{nominative masculine singular}{γνωστός}{known} ⸂τοῦ ἀρχιερέως⸃ καὶ εἶπεν\FnParse{d}{aorist active indicative 3rd singular} τῇ θυρωρῷ\FnParseFormGloss{e}{dative feminine singular}{θυρωρός}{doorkeeper} καὶ εἰσήγαγεν\FnParseFormGloss{f}{aorist active indicative 3rd singular}{εἰσάγω}{to bring in} τὸν Πέτρον. 
\MainTextVerseMark{18}{17}λέγει\FnParse{a}{present active indicative 3rd singular} οὖν ⸂τῷ Πέτρῳ ἡ παιδίσκη\FnParseFormGloss{b}{nominative feminine singular}{παιδίσκη}{female servant} ἡ θυρωρός⸃·\FnParseFormGloss{c}{nominative feminine singular}{θυρωρός}{doorkeeper} Μὴ καὶ σὺ ἐκ τῶν μαθητῶν εἶ\FnParse{d}{present active indicative 2nd singular} τοῦ ἀνθρώπου τούτου; λέγει\FnParse{e}{present active indicative 3rd singular} ἐκεῖνος· Οὐκ εἰμί.\FnParse{f}{present active indicative 1st singular} 
\MainTextVerseMark{18}{18}εἱστήκεισαν\FnParse{a}{pluperfect active indicative 3rd plural} δὲ οἱ δοῦλοι καὶ οἱ ὑπηρέται\FnParseFormGloss{b}{nominative masculine plural}{ὑπηρέτης}{servant} ἀνθρακιὰν\FnParseFormGloss{c}{accusative feminine singular}{ἀνθρακιά}{charcoal fire} πεποιηκότες,\FnParse{d}{perfect active participle nominative masculine plural} ὅτι ψῦχος\FnParseFormGloss{e}{nominative neuter singular}{ψῦχος}{cold} ἦν,\FnParse{f}{imperfect active indicative 3rd singular} καὶ ἐθερμαίνοντο·\FnParseFormGloss{g}{imperfect middle indicative 3rd plural}{θερμαίνομαι}{to keep warm} ἦν\FnParse{h}{imperfect active indicative 3rd singular} δὲ ⸂καὶ ὁ Πέτρος μετ’ αὐτῶν⸃ ἑστὼς\FnParse{i}{perfect active participle nominative masculine singular} καὶ θερμαινόμενος.\FnParseFormGloss{j}{present middle participle nominative masculine singular}{θερμαίνομαι}{to keep warm} 
\MainTextVerseMark{18}{19}Ὁ οὖν ἀρχιερεὺς ἠρώτησεν\FnParse{a}{aorist active indicative 3rd singular} τὸν Ἰησοῦν περὶ τῶν μαθητῶν αὐτοῦ καὶ περὶ τῆς διδαχῆς αὐτοῦ. 
\MainTextVerseMark{18}{20}ἀπεκρίθη\FnParse{a}{aorist passive indicative 3rd singular} ⸀αὐτῷ Ἰησοῦς· Ἐγὼ παρρησίᾳ ⸀λελάληκα\FnParse{b}{perfect active indicative 1st singular} τῷ κόσμῳ· ἐγὼ πάντοτε ἐδίδαξα\FnParse{c}{aorist active indicative 1st singular} ἐν συναγωγῇ καὶ ἐν τῷ ἱερῷ, ὅπου ⸀πάντες οἱ Ἰουδαῖοι συνέρχονται,\FnParse{d}{present middle indicative 3rd plural} καὶ ἐν κρυπτῷ\FnParseFormGloss{e}{dative neuter singular}{κρυπτός}{hidden} ἐλάλησα\FnParse{f}{aorist active indicative 1st singular} οὐδέν· 
\MainTextVerseMark{18}{21}τί με ⸂ἐρωτᾷς;\FnParse{a}{present active indicative 2nd singular} ἐρώτησον⸃\FnParse{b}{aorist active imperative 2nd singular} τοὺς ἀκηκοότας\FnParse{c}{perfect active participle accusative masculine plural} τί ἐλάλησα\FnParse{d}{aorist active indicative 1st singular} αὐτοῖς· ἴδε οὗτοι οἴδασιν\FnParse{e}{perfect active indicative 3rd plural} ἃ εἶπον\FnParse{f}{aorist active indicative 1st singular} ἐγώ. 
\MainTextVerseMark{18}{22}ταῦτα δὲ αὐτοῦ εἰπόντος\FnParse{a}{aorist active participle genitive masculine singular} εἷς ⸂παρεστηκὼς\FnParse{b}{perfect active participle nominative masculine singular} τῶν ὑπηρετῶν⸃\FnParseFormGloss{c}{genitive masculine plural}{ὑπηρέτης}{servant} ἔδωκεν\FnParse{d}{aorist active indicative 3rd singular} ῥάπισμα\FnParseFormGloss{e}{accusative neuter singular}{ῥάπισμα}{slap} τῷ Ἰησοῦ εἰπών·\FnParse{f}{aorist active participle nominative masculine singular} Οὕτως ἀποκρίνῃ\FnParse{g}{present middle indicative 2nd singular} τῷ ἀρχιερεῖ; 
\MainTextVerseMark{18}{23}ἀπεκρίθη\FnParse{a}{aorist passive indicative 3rd singular} ⸀αὐτῷ Ἰησοῦς· Εἰ κακῶς\FnFormGloss{b}{κακῶς}{ill} ἐλάλησα,\FnParse{c}{aorist active indicative 1st singular} μαρτύρησον\FnParse{d}{aorist active imperative 2nd singular} περὶ τοῦ κακοῦ· εἰ δὲ καλῶς, τί με δέρεις;\FnParseFormGloss{e}{present active indicative 2nd singular}{δέρω}{to beat up} 
\MainTextVerseMark{18}{24}ἀπέστειλεν\FnParse{a}{aorist active indicative 3rd singular} ⸀οὖν αὐτὸν ὁ Ἅννας\FnParseFormGloss{b}{nominative masculine singular}{Ἅννας}{Annas} δεδεμένον\FnParse{c}{perfect passive participle accusative masculine singular} πρὸς Καϊάφαν\FnParseFormGloss{d}{accusative masculine singular}{Καϊάφας}{Caiaphas} τὸν ἀρχιερέα. 
\MainTextVerseMark{18}{25}Ἦν\FnParse{a}{imperfect active indicative 3rd singular} δὲ Σίμων Πέτρος ἑστὼς\FnParse{b}{perfect active participle nominative masculine singular} καὶ θερμαινόμενος.\FnParseFormGloss{c}{present middle participle nominative masculine singular}{θερμαίνομαι}{to keep warm} εἶπον\FnParse{d}{aorist active indicative 3rd plural} οὖν αὐτῷ· Μὴ καὶ σὺ ἐκ τῶν μαθητῶν αὐτοῦ εἶ;\FnParse{e}{present active indicative 2nd singular} ⸀ἠρνήσατο\FnParse{f}{aorist middle indicative 3rd singular} ἐκεῖνος καὶ εἶπεν·\FnParse{g}{aorist active indicative 3rd singular} Οὐκ εἰμί.\FnParse{h}{present active indicative 1st singular} 
\MainTextVerseMark{18}{26}λέγει\FnParse{a}{present active indicative 3rd singular} εἷς ἐκ τῶν δούλων τοῦ ἀρχιερέως, συγγενὴς\FnParseFormGloss{b}{nominative masculine singular}{συγγενής}{family} ὢν\FnParse{c}{present active participle nominative masculine singular} οὗ ἀπέκοψεν\FnParseFormGloss{d}{aorist active indicative 3rd singular}{ἀποκόπτω}{to cut off} Πέτρος τὸ ὠτίον·\FnParseFormGloss{e}{accusative neuter singular}{ὠτίον}{ear} Οὐκ ἐγώ σε εἶδον\FnParse{f}{aorist active indicative 1st singular} ἐν τῷ κήπῳ\FnParseFormGloss{g}{dative masculine singular}{κῆπος}{garden} μετ’ αὐτοῦ; 
\MainTextVerseMark{18}{27}πάλιν οὖν ⸀ἠρνήσατο\FnParse{a}{aorist middle indicative 3rd singular} Πέτρος· καὶ εὐθέως ἀλέκτωρ\FnParseFormGloss{b}{nominative masculine singular}{ἀλέκτωρ}{rooster} ἐφώνησεν.\FnParse{c}{aorist active indicative 3rd singular} 
\MainTextVerseMark{18}{28}Ἄγουσιν\FnParse{a}{present active indicative 3rd plural} οὖν τὸν Ἰησοῦν ἀπὸ τοῦ Καϊάφα\FnParseFormGloss{b}{genitive masculine singular}{Καϊάφας}{Caiaphas} εἰς τὸ πραιτώριον·\FnParseFormGloss{c}{accusative neuter singular}{πραιτώριον}{Praetorium} ἦν\FnParse{d}{imperfect active indicative 3rd singular} δὲ πρωΐ·\FnFormGloss{e}{πρωΐ}{early in the morning} καὶ αὐτοὶ οὐκ εἰσῆλθον\FnParse{f}{aorist active indicative 3rd plural} εἰς τὸ πραιτώριον,\FnParseFormGloss{g}{accusative neuter singular}{πραιτώριον}{Praetorium} ἵνα μὴ μιανθῶσιν\FnParseFormGloss{h}{aorist passive subjunctive 3rd plural}{μιαίνω}{to pollute} ⸀ἀλλὰ φάγωσιν\FnParse{i}{aorist active subjunctive 3rd plural} τὸ πάσχα.\FnParseFormGloss{j}{accusative neuter singular}{πάσχα}{Passover} 
\MainTextVerseMark{18}{29}ἐξῆλθεν\FnParse{a}{aorist active indicative 3rd singular} οὖν ὁ Πιλᾶτος ⸀ἔξω πρὸς αὐτοὺς καὶ ⸀φησίν·\FnParse{b}{present active indicative 3rd singular} Τίνα κατηγορίαν\FnParseFormGloss{c}{accusative feminine singular}{κατηγορία}{charge} φέρετε\FnParse{d}{present active indicative 2nd plural} ⸀κατὰ τοῦ ἀνθρώπου τούτου; 
\MainTextVerseMark{18}{30}ἀπεκρίθησαν\FnParse{a}{aorist passive indicative 3rd plural} καὶ εἶπαν\FnParse{b}{aorist active indicative 3rd plural} αὐτῷ· Εἰ μὴ ἦν\FnParse{c}{imperfect active indicative 3rd singular} οὗτος ⸂κακὸν ποιῶν⸃,\FnParse{d}{present active participle nominative masculine singular} οὐκ ἄν σοι παρεδώκαμεν\FnParse{e}{aorist active indicative 1st plural} αὐτόν. 
\MainTextVerseMark{18}{31}εἶπεν\FnParse{a}{aorist active indicative 3rd singular} οὖν αὐτοῖς ⸀ὁ Πιλᾶτος· Λάβετε\FnParse{b}{aorist active imperative 2nd plural} αὐτὸν ὑμεῖς, καὶ κατὰ τὸν νόμον ὑμῶν κρίνατε\FnParse{c}{aorist active imperative 2nd plural} αὐτόν. ⸀εἶπον\FnParse{d}{aorist active indicative 3rd plural} αὐτῷ οἱ Ἰουδαῖοι· Ἡμῖν οὐκ ἔξεστιν\FnParse{e}{present active indicative 3rd singular} ἀποκτεῖναι\FnParse{f}{aorist active infinitive} οὐδένα· 
\MainTextVerseMark{18}{32}ἵνα ὁ λόγος τοῦ Ἰησοῦ πληρωθῇ\FnParse{a}{aorist passive subjunctive 3rd singular} ὃν εἶπεν\FnParse{b}{aorist active indicative 3rd singular} σημαίνων\FnParseFormGloss{c}{present active participle nominative masculine singular}{σημαίνω}{to make known} ποίῳ θανάτῳ ἤμελλεν\FnParse{d}{imperfect active indicative 3rd singular} ἀποθνῄσκειν.\FnParse{e}{present active infinitive} 
\MainTextVerseMark{18}{33}Εἰσῆλθεν\FnParse{a}{aorist active indicative 3rd singular} οὖν ⸂πάλιν εἰς τὸ πραιτώριον⸃\FnParseFormGloss{b}{accusative neuter singular}{πραιτώριον}{Praetorium} ὁ Πιλᾶτος καὶ ἐφώνησεν\FnParse{c}{aorist active indicative 3rd singular} τὸν Ἰησοῦν καὶ εἶπεν\FnParse{d}{aorist active indicative 3rd singular} αὐτῷ· Σὺ εἶ\FnParse{e}{present active indicative 2nd singular} ὁ βασιλεὺς τῶν Ἰουδαίων; 
\MainTextVerseMark{18}{34}⸀ἀπεκρίθη\FnParse{a}{aorist passive indicative 3rd singular} Ἰησοῦς· ⸂Ἀπὸ σεαυτοῦ⸃ σὺ τοῦτο λέγεις\FnParse{b}{present active indicative 2nd singular} ἢ ἄλλοι ⸂εἶπόν\FnParse{c}{aorist active indicative 3rd plural} σοι⸃ περὶ ἐμοῦ; 
\MainTextVerseMark{18}{35}ἀπεκρίθη\FnParse{a}{aorist passive indicative 3rd singular} ὁ Πιλᾶτος· Μήτι\FnFormGloss{b}{μήτι}{[interrogative particle]} ἐγὼ Ἰουδαῖός εἰμι;\FnParse{c}{present active indicative 1st singular} τὸ ἔθνος τὸ σὸν\FnParseFormGloss{d}{nominative neuter singular}{σός}{your} καὶ οἱ ἀρχιερεῖς παρέδωκάν\FnParse{e}{aorist active indicative 3rd plural} σε ἐμοί· τί ἐποίησας;\FnParse{f}{aorist active indicative 2nd singular} 
\MainTextVerseMark{18}{36}ἀπεκρίθη\FnParse{a}{aorist passive indicative 3rd singular} Ἰησοῦς· Ἡ βασιλεία ἡ ἐμὴ οὐκ ἔστιν\FnParse{b}{present active indicative 3rd singular} ἐκ τοῦ κόσμου τούτου· εἰ ἐκ τοῦ κόσμου τούτου ἦν\FnParse{c}{imperfect active indicative 3rd singular} ἡ βασιλεία ἡ ἐμή, οἱ ὑπηρέται\FnParseFormGloss{d}{nominative masculine plural}{ὑπηρέτης}{servant} ⸂οἱ ἐμοὶ ἠγωνίζοντο\FnParseFormGloss{e}{imperfect middle indicative 3rd plural}{ἀγωνίζομαι}{to fight} ἄν⸃, ἵνα μὴ παραδοθῶ\FnParse{f}{aorist passive subjunctive 1st singular} τοῖς Ἰουδαίοις· νῦν δὲ ἡ βασιλεία ἡ ἐμὴ οὐκ ἔστιν\FnParse{g}{present active indicative 3rd singular} ἐντεῦθεν.\FnFormGloss{h}{ἐντεῦθεν}{from here} 
\MainTextVerseMark{18}{37}εἶπεν\FnParse{a}{aorist active indicative 3rd singular} οὖν αὐτῷ ὁ Πιλᾶτος· Οὐκοῦν\FnFormGloss{b}{οὐκοῦν}{so} βασιλεὺς εἶ\FnParse{c}{present active indicative 2nd singular} σύ; ἀπεκρίθη\FnParse{d}{aorist passive indicative 3rd singular} ⸀ὁ Ἰησοῦς· Σὺ λέγεις\FnParse{e}{present active indicative 2nd singular} ὅτι βασιλεύς ⸀εἰμι.\FnParse{f}{present active indicative 1st singular} ἐγὼ εἰς τοῦτο γεγέννημαι\FnParse{g}{perfect passive indicative 1st singular} καὶ εἰς τοῦτο ἐλήλυθα\FnParse{h}{perfect active indicative 1st singular} εἰς τὸν κόσμον ἵνα μαρτυρήσω\FnParse{i}{aorist active subjunctive 1st singular} τῇ ἀληθείᾳ· πᾶς ὁ ὢν\FnParse{j}{present active participle nominative masculine singular} ἐκ τῆς ἀληθείας ἀκούει\FnParse{k}{present active indicative 3rd singular} μου τῆς φωνῆς. 
\MainTextVerseMark{18}{38}λέγει\FnParse{a}{present active indicative 3rd singular} αὐτῷ ὁ Πιλᾶτος· Τί ἐστιν\FnParse{b}{present active indicative 3rd singular} ἀλήθεια; Καὶ τοῦτο εἰπὼν\FnParse{c}{aorist active participle nominative masculine singular} πάλιν ἐξῆλθεν\FnParse{d}{aorist active indicative 3rd singular} πρὸς τοὺς Ἰουδαίους, καὶ λέγει\FnParse{e}{present active indicative 3rd singular} αὐτοῖς· Ἐγὼ οὐδεμίαν ⸂εὑρίσκω\FnParse{f}{present active indicative 1st singular} ἐν αὐτῷ αἰτίαν⸃·\FnParseFormGloss{g}{accusative feminine singular}{αἰτία}{charge} 
\MainTextVerseMark{18}{39}ἔστιν\FnParse{a}{present active indicative 3rd singular} δὲ συνήθεια\FnParseFormGloss{b}{nominative feminine singular}{συνήθεια}{custom} ὑμῖν ἵνα ἕνα ⸂ἀπολύσω\FnParse{c}{aorist active subjunctive 1st singular} ὑμῖν⸃ ἐν τῷ πάσχα·\FnParseFormGloss{d}{dative neuter singular}{πάσχα}{Passover} βούλεσθε\FnParse{e}{present middle indicative 2nd plural} οὖν ⸄ἀπολύσω\FnParse{f}{aorist active subjunctive 1st singular} ὑμῖν⸅ τὸν βασιλέα τῶν Ἰουδαίων; 
\MainTextVerseMark{18}{40}ἐκραύγασαν\FnParseFormGloss{a}{aorist active indicative 3rd plural}{κραυγάζω}{to shout} οὖν ⸀πάλιν λέγοντες·\FnParse{b}{present active participle nominative masculine plural} Μὴ τοῦτον ἀλλὰ τὸν Βαραββᾶν.\FnParseFormGloss{c}{accusative masculine singular}{Βαραββᾶς}{Barabbas} ἦν\FnParse{d}{imperfect active indicative 3rd singular} δὲ ὁ Βαραββᾶς\FnParseFormGloss{e}{nominative masculine singular}{Βαραββᾶς}{Barabbas} λῃστής.\FnParseFormGloss{f}{nominative masculine singular}{λῃστής}{robber} 
\ChapterHeader{19}{Chapter 19}

\MainTextVerseMark{19}{1}Τότε οὖν ἔλαβεν\FnParse{a}{aorist active indicative 3rd singular} ὁ Πιλᾶτος τὸν Ἰησοῦν καὶ ἐμαστίγωσεν.\FnParseFormGloss{b}{aorist active indicative 3rd singular}{μαστιγόω}{to flog} 
\MainTextVerseMark{19}{2}καὶ οἱ στρατιῶται\FnParseFormGloss{a}{nominative masculine plural}{στρατιώτης}{soldier} πλέξαντες\FnParseFormGloss{b}{aorist active participle nominative masculine plural}{πλέκω}{to twist together} στέφανον\FnParseFormGloss{c}{accusative masculine singular}{στέφανος}{a crown} ἐξ ἀκανθῶν\FnParseFormGloss{d}{genitive feminine plural}{ἄκανθα}{thorn} ἐπέθηκαν\FnParse{e}{aorist active indicative 3rd plural} αὐτοῦ τῇ κεφαλῇ, καὶ ἱμάτιον πορφυροῦν\FnParseFormGloss{f}{accusative neuter singular}{πορφυροῦς}{purple} περιέβαλον\FnParseFormGloss{g}{aorist active indicative 3rd plural}{περιβάλλω}{to dress} αὐτόν, 
\MainTextVerseMark{19}{3}καὶ ⸂ἤρχοντο\FnParse{a}{imperfect middle indicative 3rd plural} πρὸς αὐτὸν καὶ⸃ ἔλεγον·\FnParse{b}{imperfect active indicative 3rd plural} Χαῖρε,\FnParse{c}{present active imperative 2nd singular} ὁ βασιλεὺς τῶν Ἰουδαίων· καὶ ἐδίδοσαν\FnParse{d}{imperfect active indicative 3rd plural} αὐτῷ ῥαπίσματα.\FnParseFormGloss{e}{accusative neuter plural}{ῥάπισμα}{slap} 
\MainTextVerseMark{19}{4}⸂καὶ ἐξῆλθεν⸃\FnParse{a}{aorist active indicative 3rd singular} πάλιν ἔξω ὁ Πιλᾶτος καὶ λέγει\FnParse{b}{present active indicative 3rd singular} αὐτοῖς· Ἴδε ἄγω\FnParse{c}{present active indicative 1st singular} ὑμῖν αὐτὸν ἔξω, ἵνα γνῶτε\FnParse{d}{aorist active subjunctive 2nd plural} ὅτι ⸂οὐδεμίαν αἰτίαν\FnParseFormGloss{e}{accusative feminine singular}{αἰτία}{charge} εὑρίσκω\FnParse{f}{present active indicative 1st singular} ἐν αὐτῷ⸃. 
\MainTextVerseMark{19}{5}ἐξῆλθεν\FnParse{a}{aorist active indicative 3rd singular} οὖν ὁ Ἰησοῦς ἔξω, φορῶν\FnParseFormGloss{b}{present active participle nominative masculine singular}{φορέω}{to wear} τὸν ἀκάνθινον\FnParseFormGloss{c}{accusative masculine singular}{ἀκάνθινος}{of thorns} στέφανον\FnParseFormGloss{d}{accusative masculine singular}{στέφανος}{a crown} καὶ τὸ πορφυροῦν\FnParseFormGloss{e}{accusative neuter singular}{πορφυροῦς}{purple} ἱμάτιον. καὶ λέγει\FnParse{f}{present active indicative 3rd singular} αὐτοῖς· ⸀Ἰδοὺ ὁ ἄνθρωπος. 
\MainTextVerseMark{19}{6}ὅτε οὖν εἶδον\FnParse{a}{aorist active indicative 3rd plural} αὐτὸν οἱ ἀρχιερεῖς καὶ οἱ ὑπηρέται\FnParseFormGloss{b}{nominative masculine plural}{ὑπηρέτης}{servant} ἐκραύγασαν\FnParseFormGloss{c}{aorist active indicative 3rd plural}{κραυγάζω}{to shout} λέγοντες·\FnParse{d}{present active participle nominative masculine plural} Σταύρωσον\FnParse{e}{aorist active imperative 2nd singular} ⸀σταύρωσον.\FnParse{f}{aorist active imperative 2nd singular} λέγει\FnParse{g}{present active indicative 3rd singular} αὐτοῖς ὁ Πιλᾶτος· Λάβετε\FnParse{h}{aorist active imperative 2nd plural} αὐτὸν ὑμεῖς καὶ σταυρώσατε,\FnParse{i}{aorist active imperative 2nd plural} ἐγὼ γὰρ οὐχ εὑρίσκω\FnParse{j}{present active indicative 1st singular} ἐν αὐτῷ αἰτίαν.\FnParseFormGloss{k}{accusative feminine singular}{αἰτία}{charge} 
\MainTextVerseMark{19}{7}ἀπεκρίθησαν\FnParse{a}{aorist passive indicative 3rd plural} αὐτῷ οἱ Ἰουδαῖοι· Ἡμεῖς νόμον ἔχομεν,\FnParse{b}{present active indicative 1st plural} καὶ κατὰ τὸν ⸀νόμον ὀφείλει\FnParse{c}{present active indicative 3rd singular} ἀποθανεῖν,\FnParse{d}{aorist active infinitive} ὅτι ⸂υἱὸν θεοῦ ἑαυτὸν⸃ ἐποίησεν.\FnParse{e}{aorist active indicative 3rd singular} 
\MainTextVerseMark{19}{8}Ὅτε οὖν ἤκουσεν\FnParse{a}{aorist active indicative 3rd singular} ὁ Πιλᾶτος τοῦτον τὸν λόγον, μᾶλλον ἐφοβήθη,\FnParse{b}{aorist passive indicative 3rd singular} 
\MainTextVerseMark{19}{9}καὶ εἰσῆλθεν\FnParse{a}{aorist active indicative 3rd singular} εἰς τὸ πραιτώριον\FnParseFormGloss{b}{accusative neuter singular}{πραιτώριον}{Praetorium} πάλιν καὶ λέγει\FnParse{c}{present active indicative 3rd singular} τῷ Ἰησοῦ· Πόθεν\FnFormGloss{d}{πόθεν}{from where} εἶ\FnParse{e}{present active indicative 2nd singular} σύ; ὁ δὲ Ἰησοῦς ἀπόκρισιν\FnParseFormGloss{f}{accusative feminine singular}{ἀπόκρισις}{answer} οὐκ ἔδωκεν\FnParse{g}{aorist active indicative 3rd singular} αὐτῷ. 
\MainTextVerseMark{19}{10}λέγει\FnParse{a}{present active indicative 3rd singular} οὖν αὐτῷ ὁ Πιλᾶτος· Ἐμοὶ οὐ λαλεῖς;\FnParse{b}{present active indicative 2nd singular} οὐκ οἶδας\FnParse{c}{perfect active indicative 2nd singular} ὅτι ἐξουσίαν ἔχω\FnParse{d}{present active indicative 1st singular} ⸂ἀπολῦσαί\FnParse{e}{aorist active infinitive} σε καὶ ἐξουσίαν ἔχω\FnParse{f}{present active indicative 1st singular} σταυρῶσαί⸃\FnParse{g}{aorist active infinitive} σε; 
\MainTextVerseMark{19}{11}ἀπεκρίθη\FnParse{a}{aorist passive indicative 3rd singular} ⸀αὐτῷ Ἰησοῦς· Οὐκ εἶχες\FnParse{b}{imperfect active indicative 2nd singular} ἐξουσίαν ⸂κατ’ ἐμοῦ οὐδεμίαν⸃ εἰ μὴ ἦν\FnParse{c}{imperfect active indicative 3rd singular} ⸂δεδομένον\FnParse{d}{perfect passive participle nominative neuter singular} σοι⸃ ἄνωθεν·\FnFormGloss{e}{ἄνωθεν}{from above} διὰ τοῦτο ὁ ⸀παραδούς\FnParse{f}{aorist active participle nominative masculine singular} μέ σοι μείζονα ἁμαρτίαν ἔχει.\FnParse{g}{present active indicative 3rd singular} 
\MainTextVerseMark{19}{12}ἐκ τούτου ⸂ὁ Πιλᾶτος ἐζήτει⸃\FnParse{a}{imperfect active indicative 3rd singular} ἀπολῦσαι\FnParse{b}{aorist active infinitive} αὐτόν· οἱ δὲ Ἰουδαῖοι ⸀ἐκραύγασαν\FnParseFormGloss{c}{aorist active indicative 3rd plural}{κραυγάζω}{to shout} λέγοντες·\FnParse{d}{present active participle nominative masculine plural} Ἐὰν τοῦτον ἀπολύσῃς,\FnParse{e}{aorist active subjunctive 2nd singular} οὐκ εἶ\FnParse{f}{present active indicative 2nd singular} φίλος\FnParseFormGloss{g}{nominative masculine singular}{φίλος}{friendly} τοῦ Καίσαρος·\FnParseFormGloss{h}{genitive masculine singular}{Καῖσαρ}{Caesar} πᾶς ὁ βασιλέα ἑαυτὸν ποιῶν\FnParse{i}{present active participle nominative masculine singular} ἀντιλέγει\FnParseFormGloss{j}{present active indicative 3rd singular}{ἀντιλέγω}{to speak against} τῷ Καίσαρι.\FnParseFormGloss{k}{dative masculine singular}{Καῖσαρ}{Caesar} 
\MainTextVerseMark{19}{13}Ὁ οὖν Πιλᾶτος ἀκούσας\FnParse{a}{aorist active participle nominative masculine singular} ⸂τῶν λόγων τούτων⸃ ἤγαγεν\FnParse{b}{aorist active indicative 3rd singular} ἔξω τὸν Ἰησοῦν, καὶ ἐκάθισεν\FnParse{c}{aorist active indicative 3rd singular} ⸀ἐπὶ βήματος\FnParseFormGloss{d}{genitive neuter singular}{βῆμα}{judicial court} εἰς τόπον λεγόμενον\FnParse{e}{present passive participle accusative masculine singular} Λιθόστρωτον,\FnParseFormGloss{f}{accusative masculine singular}{λιθόστρωτος}{stone pavement} Ἑβραϊστὶ\FnFormGloss{g}{Ἑβραϊστί}{in Aramaic} δὲ Γαββαθα.\FnParseFormGloss{h}{accusative neuter singular}{Γαββαθα}{Gabbatha} 
\MainTextVerseMark{19}{14}ἦν\FnParse{a}{imperfect active indicative 3rd singular} δὲ παρασκευὴ\FnParseFormGloss{b}{nominative feminine singular}{παρασκευή}{Preparation Day} τοῦ πάσχα,\FnParseFormGloss{c}{genitive neuter singular}{πάσχα}{Passover} ὥρα ⸂ἦν\FnParse{d}{imperfect active indicative 3rd singular} ὡς⸃ ἕκτη.\FnParseFormGloss{e}{nominative feminine singular}{ἕκτος}{sixth} καὶ λέγει\FnParse{f}{present active indicative 3rd singular} τοῖς Ἰουδαίοις· Ἴδε ὁ βασιλεὺς ὑμῶν. 
\MainTextVerseMark{19}{15}⸂ἐκραύγασαν\FnParseFormGloss{a}{aorist active indicative 3rd plural}{κραυγάζω}{to shout} οὖν ἐκεῖνοι⸃· Ἆρον\FnParse{b}{aorist active imperative 2nd singular} ἆρον,\FnParse{c}{aorist active imperative 2nd singular} σταύρωσον\FnParse{d}{aorist active imperative 2nd singular} αὐτόν. λέγει\FnParse{e}{present active indicative 3rd singular} αὐτοῖς ὁ Πιλᾶτος· Τὸν βασιλέα ὑμῶν σταυρώσω;\FnParse{f}{aorist active subjunctive 1st singular} ἀπεκρίθησαν\FnParse{g}{aorist passive indicative 3rd plural} οἱ ἀρχιερεῖς· Οὐκ ἔχομεν\FnParse{h}{present active indicative 1st plural} βασιλέα εἰ μὴ Καίσαρα.\FnParseFormGloss{i}{accusative masculine singular}{Καῖσαρ}{Caesar} 
\MainTextVerseMark{19}{16}τότε οὖν παρέδωκεν\FnParse{a}{aorist active indicative 3rd singular} αὐτὸν αὐτοῖς ἵνα σταυρωθῇ.\FnParse{b}{aorist passive subjunctive 3rd singular} Παρέλαβον\FnParse{c}{aorist active indicative 3rd plural} ⸀οὖν τὸν ⸀Ἰησοῦν· 
\MainTextVerseMark{19}{17}καὶ βαστάζων\FnParseFormGloss{a}{present active participle nominative masculine singular}{βαστάζω}{to carry} ⸂αὑτῷ τὸν σταυρὸν⸃\FnParseFormGloss{b}{accusative masculine singular}{σταυρός}{cross} ἐξῆλθεν\FnParse{c}{aorist active indicative 3rd singular} εἰς ⸀τὸν λεγόμενον\FnParse{d}{present passive participle accusative masculine singular} Κρανίου\FnParseFormGloss{e}{genitive neuter singular}{κρανίον}{skull} Τόπον, ⸀ὃ λέγεται\FnParse{f}{present passive indicative 3rd singular} Ἑβραϊστὶ\FnFormGloss{g}{Ἑβραϊστί}{in Aramaic} Γολγοθα,\FnParseFormGloss{h}{nominative feminine singular}{Γολγοθᾶ}{Golgotha} 
\MainTextVerseMark{19}{18}ὅπου αὐτὸν ἐσταύρωσαν,\FnParse{a}{aorist active indicative 3rd plural} καὶ μετ’ αὐτοῦ ἄλλους δύο ἐντεῦθεν\FnFormGloss{b}{ἐντεῦθεν}{from here} καὶ ἐντεῦθεν,\FnFormGloss{c}{ἐντεῦθεν}{from here} μέσον δὲ τὸν Ἰησοῦν. 
\MainTextVerseMark{19}{19}ἔγραψεν\FnParse{a}{aorist active indicative 3rd singular} δὲ καὶ τίτλον\FnParseFormGloss{b}{accusative masculine singular}{τίτλος}{sign} ὁ Πιλᾶτος καὶ ἔθηκεν\FnParse{c}{aorist active indicative 3rd singular} ἐπὶ τοῦ σταυροῦ·\FnParseFormGloss{d}{genitive masculine singular}{σταυρός}{cross} ἦν\FnParse{e}{imperfect active indicative 3rd singular} δὲ γεγραμμένον·\FnParse{f}{perfect passive participle nominative neuter singular} Ἰησοῦς ὁ Ναζωραῖος\FnParseFormGloss{g}{nominative masculine singular}{Ναζωραῖος}{Nazarene} ὁ βασιλεὺς τῶν Ἰουδαίων. 
\MainTextVerseMark{19}{20}τοῦτον οὖν τὸν τίτλον\FnParseFormGloss{a}{accusative masculine singular}{τίτλος}{sign} πολλοὶ ἀνέγνωσαν\FnParse{b}{aorist active indicative 3rd plural} τῶν Ἰουδαίων, ὅτι ἐγγὺς ἦν\FnParse{c}{imperfect active indicative 3rd singular} ὁ τόπος τῆς πόλεως ὅπου ἐσταυρώθη\FnParse{d}{aorist passive indicative 3rd singular} ὁ Ἰησοῦς· καὶ ἦν\FnParse{e}{imperfect active indicative 3rd singular} γεγραμμένον\FnParse{f}{perfect passive participle nominative neuter singular} Ἑβραϊστί,\FnFormGloss{g}{Ἑβραϊστί}{in Aramaic} ⸂Ῥωμαϊστί,\FnFormGloss{h}{Ῥωμαϊστί}{in Latin} Ἑλληνιστί⸃.\FnFormGloss{i}{Ἑλληνιστί}{in Greek} 
\MainTextVerseMark{19}{21}ἔλεγον\FnParse{a}{imperfect active indicative 3rd plural} οὖν τῷ Πιλάτῳ οἱ ἀρχιερεῖς τῶν Ἰουδαίων· Μὴ γράφε·\FnParse{b}{present active imperative 2nd singular} Ὁ βασιλεὺς τῶν Ἰουδαίων, ἀλλ’ ὅτι ἐκεῖνος εἶπεν\FnParse{c}{aorist active indicative 3rd singular} Βασιλεὺς ⸂τῶν Ἰουδαίων εἰμί⸃.\FnParse{d}{present active indicative 1st singular} 
\MainTextVerseMark{19}{22}ἀπεκρίθη\FnParse{a}{aorist passive indicative 3rd singular} ὁ Πιλᾶτος· Ὃ γέγραφα\FnParse{b}{perfect active indicative 1st singular} γέγραφα.\FnParse{c}{perfect active indicative 1st singular} 
\MainTextVerseMark{19}{23}οἱ οὖν στρατιῶται\FnParseFormGloss{a}{nominative masculine plural}{στρατιώτης}{soldier} ὅτε ἐσταύρωσαν\FnParse{b}{aorist active indicative 3rd plural} τὸν Ἰησοῦν ἔλαβον\FnParse{c}{aorist active indicative 3rd plural} τὰ ἱμάτια αὐτοῦ καὶ ἐποίησαν\FnParse{d}{aorist active indicative 3rd plural} τέσσαρα μέρη, ἑκάστῳ στρατιώτῃ\FnParseFormGloss{e}{dative masculine singular}{στρατιώτης}{soldier} μέρος, καὶ τὸν χιτῶνα.\FnParseFormGloss{f}{accusative masculine singular}{χιτών}{tunic} ἦν\FnParse{g}{imperfect active indicative 3rd singular} δὲ ὁ χιτὼν\FnParseFormGloss{h}{nominative masculine singular}{χιτών}{tunic} ἄραφος,\FnParseFormGloss{i}{nominative masculine singular}{ἄραφος}{seamless} ἐκ τῶν ἄνωθεν\FnFormGloss{j}{ἄνωθεν}{from above} ὑφαντὸς\FnParseFormGloss{k}{nominative masculine singular}{ὑφαντός}{woven} δι’ ὅλου· 
\MainTextVerseMark{19}{24}εἶπαν\FnParse{a}{aorist active indicative 3rd plural} οὖν πρὸς ἀλλήλους· Μὴ σχίσωμεν\FnParseFormGloss{b}{aorist active subjunctive 1st plural}{σχίζω}{to tear} αὐτόν, ἀλλὰ λάχωμεν\FnParseFormGloss{c}{aorist active subjunctive 1st plural}{λαγχάνω}{to choose by lot} περὶ αὐτοῦ τίνος ἔσται·\FnParse{d}{future middle indicative 3rd singular} ἵνα ἡ γραφὴ πληρωθῇ\FnParse{e}{aorist passive subjunctive 3rd singular} ⸂ἡ λέγουσα⸃·\FnParse{f}{present active participle nominative feminine singular} Διεμερίσαντο\FnParseFormGloss{g}{aorist middle indicative 3rd plural}{διαμερίζω}{to divide} τὰ ἱμάτιά μου ἑαυτοῖς καὶ ἐπὶ τὸν ἱματισμόν\FnParseFormGloss{h}{accusative masculine singular}{ἱματισμός}{clothing} μου ἔβαλον\FnParse{i}{aorist active indicative 3rd plural} κλῆρον.\FnParseFormGloss{j}{accusative masculine singular}{κλῆρος}{lots} Οἱ μὲν οὖν στρατιῶται\FnParseFormGloss{k}{nominative masculine plural}{στρατιώτης}{soldier} ταῦτα ἐποίησαν.\FnParse{l}{aorist active indicative 3rd plural} 
\MainTextVerseMark{19}{25}Εἱστήκεισαν\FnParse{a}{pluperfect active indicative 3rd plural} δὲ παρὰ τῷ σταυρῷ\FnParseFormGloss{b}{dative masculine singular}{σταυρός}{cross} τοῦ Ἰησοῦ ἡ μήτηρ αὐτοῦ καὶ ἡ ἀδελφὴ\FnParseFormGloss{c}{nominative feminine singular}{ἀδελφή}{sister} τῆς μητρὸς αὐτοῦ, Μαρία\FnParseFormGloss{d}{nominative feminine singular}{Μαρία}{Mary} ἡ τοῦ Κλωπᾶ\FnParseFormGloss{e}{genitive masculine singular}{Κλωπᾶς}{Clopas} καὶ Μαρία\FnParseFormGloss{f}{nominative feminine singular}{Μαρία}{Mary} ἡ Μαγδαληνή.\FnParseFormGloss{g}{nominative feminine singular}{Μαγδαληνή}{Magdalene} 
\MainTextVerseMark{19}{26}Ἰησοῦς οὖν ἰδὼν\FnParse{a}{aorist active participle nominative masculine singular} τὴν μητέρα καὶ τὸν μαθητὴν παρεστῶτα\FnParse{b}{perfect active participle accusative masculine singular} ὃν ἠγάπα\FnParse{c}{imperfect active indicative 3rd singular} λέγει\FnParse{d}{present active indicative 3rd singular} τῇ ⸀μητρί· Γύναι, ⸀ἴδε ὁ υἱός σου· 
\MainTextVerseMark{19}{27}εἶτα\FnFormGloss{a}{εἶτα}{then} λέγει\FnParse{b}{present active indicative 3rd singular} τῷ μαθητῇ· ⸀Ἴδε ἡ μήτηρ σου. καὶ ἀπ’ ἐκείνης τῆς ὥρας ἔλαβεν\FnParse{c}{aorist active indicative 3rd singular} ὁ μαθητὴς αὐτὴν εἰς τὰ ἴδια. 
\MainTextVerseMark{19}{28}Μετὰ τοῦτο ⸀εἰδὼς\FnParse{a}{perfect active participle nominative masculine singular} ὁ Ἰησοῦς ὅτι ⸂ἤδη πάντα⸃ τετέλεσται\FnParseFormGloss{b}{perfect passive indicative 3rd singular}{τελέω}{to finish} ἵνα τελειωθῇ\FnParseFormGloss{c}{aorist passive subjunctive 3rd singular}{τελειόω}{to perfect} ἡ γραφὴ λέγει·\FnParse{d}{present active indicative 3rd singular} Διψῶ.\FnParseFormGloss{e}{present active indicative 1st singular}{διψάω}{to be thirsty} 
\MainTextVerseMark{19}{29}⸀σκεῦος\FnParseFormGloss{a}{nominative neuter singular}{σκεῦος}{possession} ἔκειτο\FnParseFormGloss{b}{imperfect middle indicative 3rd singular}{κεῖμαι}{to lay} ὄξους\FnParseFormGloss{c}{genitive neuter singular}{ὄξος}{wine vinegar} μεστόν·\FnParseFormGloss{d}{nominative neuter singular}{μεστός}{full} ⸂σπόγγον\FnParseFormGloss{e}{accusative masculine singular}{σπόγγος}{sponge} οὖν μεστὸν\FnParseFormGloss{f}{accusative masculine singular}{μεστός}{full} τοῦ⸃ ὄξους\FnParseFormGloss{g}{genitive neuter singular}{ὄξος}{wine vinegar} ⸀ὑσσώπῳ\FnParseFormGloss{h}{dative feminine singular}{ὕσσωπος}{hyssop} περιθέντες\FnParseFormGloss{i}{aorist active participle nominative masculine plural}{περιτίθημι}{to put on} προσήνεγκαν\FnParse{j}{aorist active indicative 3rd plural} αὐτοῦ τῷ στόματι. 
\MainTextVerseMark{19}{30}ὅτε οὖν ἔλαβεν\FnParse{a}{aorist active indicative 3rd singular} τὸ ὄξος\FnParseFormGloss{b}{accusative neuter singular}{ὄξος}{wine vinegar} ὁ Ἰησοῦς εἶπεν·\FnParse{c}{aorist active indicative 3rd singular} Τετέλεσται,\FnParseFormGloss{d}{perfect passive indicative 3rd singular}{τελέω}{to finish} καὶ κλίνας\FnParseFormGloss{e}{aorist active participle nominative masculine singular}{κλίνω}{to bow down} τὴν κεφαλὴν παρέδωκεν\FnParse{f}{aorist active indicative 3rd singular} τὸ πνεῦμα. 
\MainTextVerseMark{19}{31}Οἱ οὖν Ἰουδαῖοι, ⸂ἐπεὶ\FnFormGloss{a}{ἐπεί}{since} παρασκευὴ\FnParseFormGloss{b}{nominative feminine singular}{παρασκευή}{Preparation Day} ἦν,\FnParse{c}{imperfect active indicative 3rd singular} ἵνα μὴ μείνῃ\FnParse{d}{aorist active subjunctive 3rd singular} ἐπὶ τοῦ σταυροῦ\FnParseFormGloss{e}{genitive masculine singular}{σταυρός}{cross} τὰ σώματα ἐν τῷ σαββάτῳ⸃, ἦν\FnParse{f}{imperfect active indicative 3rd singular} γὰρ μεγάλη ἡ ἡμέρα ἐκείνου τοῦ σαββάτου, ἠρώτησαν\FnParse{g}{aorist active indicative 3rd plural} τὸν Πιλᾶτον ἵνα κατεαγῶσιν\FnParseFormGloss{h}{aorist passive subjunctive 3rd plural}{κατάγνυμι}{to break} αὐτῶν τὰ σκέλη\FnParseFormGloss{i}{accusative neuter plural}{σκέλος}{leg} καὶ ἀρθῶσιν.\FnParse{j}{aorist passive subjunctive 3rd plural} 
\MainTextVerseMark{19}{32}ἦλθον\FnParse{a}{aorist active indicative 3rd plural} οὖν οἱ στρατιῶται,\FnParseFormGloss{b}{nominative masculine plural}{στρατιώτης}{soldier} καὶ τοῦ μὲν πρώτου κατέαξαν\FnParseFormGloss{c}{aorist active indicative 3rd plural}{κατάγνυμι}{to break} τὰ σκέλη\FnParseFormGloss{d}{accusative neuter plural}{σκέλος}{leg} καὶ τοῦ ἄλλου τοῦ συσταυρωθέντος\FnParseFormGloss{e}{aorist passive participle genitive masculine singular}{συσταυρόομαι}{to be crucified with} αὐτῷ· 
\MainTextVerseMark{19}{33}ἐπὶ δὲ τὸν Ἰησοῦν ἐλθόντες,\FnParse{a}{aorist active participle nominative masculine plural} ὡς εἶδον\FnParse{b}{aorist active indicative 3rd plural} ⸂ἤδη αὐτὸν⸃ τεθνηκότα,\FnParseFormGloss{c}{perfect active participle accusative masculine singular}{θνῄσκω}{to have died} οὐ κατέαξαν\FnParseFormGloss{d}{aorist active indicative 3rd plural}{κατάγνυμι}{to break} αὐτοῦ τὰ σκέλη,\FnParseFormGloss{e}{accusative neuter plural}{σκέλος}{leg} 
\MainTextVerseMark{19}{34}ἀλλ’ εἷς τῶν στρατιωτῶν\FnParseFormGloss{a}{genitive masculine plural}{στρατιώτης}{soldier} λόγχῃ\FnParseFormGloss{b}{dative feminine singular}{λόγχη}{spear} αὐτοῦ τὴν πλευρὰν\FnParseFormGloss{c}{accusative feminine singular}{πλευρά}{side} ἔνυξεν,\FnParseFormGloss{d}{aorist active indicative 3rd singular}{νύσσω}{to pierce} καὶ ⸂ἐξῆλθεν\FnParse{e}{aorist active indicative 3rd singular} εὐθὺς⸃ αἷμα καὶ ὕδωρ. 
\MainTextVerseMark{19}{35}καὶ ὁ ἑωρακὼς\FnParse{a}{perfect active participle nominative masculine singular} μεμαρτύρηκεν,\FnParse{b}{perfect active indicative 3rd singular} καὶ ἀληθινὴ\FnParseFormGloss{c}{nominative feminine singular}{ἀληθινός}{true} ⸂αὐτοῦ ἐστιν⸃\FnParse{d}{present active indicative 3rd singular} ἡ μαρτυρία, καὶ ἐκεῖνος οἶδεν\FnParse{e}{perfect active indicative 3rd singular} ὅτι ἀληθῆ\FnParseFormGloss{f}{accusative neuter plural}{ἀληθής}{true} λέγει,\FnParse{g}{present active indicative 3rd singular} ἵνα ⸀καὶ ὑμεῖς ⸀πιστεύητε.\FnParse{h}{present active subjunctive 2nd plural} 
\MainTextVerseMark{19}{36}ἐγένετο\FnParse{a}{aorist middle indicative 3rd singular} γὰρ ταῦτα ἵνα ἡ γραφὴ πληρωθῇ·\FnParse{b}{aorist passive subjunctive 3rd singular} Ὀστοῦν\FnParseFormGloss{c}{nominative neuter singular}{ὀστέον}{bone} οὐ συντριβήσεται\FnParseFormGloss{d}{future passive indicative 3rd singular}{συντρίβω}{to break} ⸀αὐτοῦ. 
\MainTextVerseMark{19}{37}καὶ πάλιν ἑτέρα γραφὴ λέγει·\FnParse{a}{present active indicative 3rd singular} Ὄψονται\FnParse{b}{future middle indicative 3rd plural} εἰς ὃν ἐξεκέντησαν.\FnParseFormGloss{c}{aorist active indicative 3rd plural}{ἐκκεντέω}{to pierce} 
\MainTextVerseMark{19}{38}Μετὰ ⸀δὲ ταῦτα ἠρώτησεν\FnParse{a}{aorist active indicative 3rd singular} τὸν Πιλᾶτον Ἰωσὴφ ⸀ἀπὸ Ἁριμαθαίας,\FnParseFormGloss{b}{genitive feminine singular}{Ἁριμαθαία}{Arimathea} ὢν\FnParse{c}{present active participle nominative masculine singular} μαθητὴς τοῦ Ἰησοῦ κεκρυμμένος\FnParseFormGloss{d}{perfect passive participle nominative masculine singular}{κρύπτω}{to hide} δὲ διὰ τὸν φόβον τῶν Ἰουδαίων, ἵνα ἄρῃ\FnParse{e}{aorist active subjunctive 3rd singular} τὸ σῶμα τοῦ Ἰησοῦ· καὶ ἐπέτρεψεν\FnParseFormGloss{f}{aorist active indicative 3rd singular}{ἐπιτρέπω}{to let} ὁ Πιλᾶτος. ἦλθεν\FnParse{g}{aorist active indicative 3rd singular} οὖν καὶ ἦρεν\FnParse{h}{aorist active indicative 3rd singular} τὸ σῶμα ⸀αὐτοῦ. 
\MainTextVerseMark{19}{39}ἦλθεν\FnParse{a}{aorist active indicative 3rd singular} δὲ καὶ Νικόδημος,\FnParseFormGloss{b}{nominative masculine singular}{Νικόδημος}{Nicodemus} ὁ ἐλθὼν\FnParse{c}{aorist active participle nominative masculine singular} πρὸς ⸀αὐτὸν νυκτὸς τὸ πρῶτον, φέρων\FnParse{d}{present active participle nominative masculine singular} ⸀μίγμα\FnParseFormGloss{e}{accusative neuter singular}{μίγμα}{mixture} σμύρνης\FnParseFormGloss{f}{genitive feminine singular}{σμύρνα}{myrrh} καὶ ἀλόης\FnParseFormGloss{g}{genitive feminine singular}{ἀλόη}{aloes} ὡς λίτρας\FnParseFormGloss{h}{accusative feminine plural}{λίτρα}{pound} ἑκατόν.\FnParseFormGloss{i}{accusative feminine plural}{ἑκατόν}{hundred} 
\MainTextVerseMark{19}{40}ἔλαβον\FnParse{a}{aorist active indicative 3rd plural} οὖν τὸ σῶμα τοῦ Ἰησοῦ καὶ ἔδησαν\FnParse{b}{aorist active indicative 3rd plural} ⸀αὐτὸ ὀθονίοις\FnParseFormGloss{c}{dative neuter plural}{ὀθόνιον}{strips of linen} μετὰ τῶν ἀρωμάτων,\FnParseFormGloss{d}{genitive neuter plural}{ἄρωμα}{spices} καθὼς ἔθος\FnParseFormGloss{e}{nominative neuter singular}{ἔθος}{custom} ἐστὶν\FnParse{f}{present active indicative 3rd singular} τοῖς Ἰουδαίοις ἐνταφιάζειν.\FnParseFormGloss{g}{present active infinitive}{ἐνταφιάζω}{to prepare for burial} 
\MainTextVerseMark{19}{41}ἦν\FnParse{a}{imperfect active indicative 3rd singular} δὲ ἐν τῷ τόπῳ ὅπου ἐσταυρώθη\FnParse{b}{aorist passive indicative 3rd singular} κῆπος,\FnParseFormGloss{c}{nominative masculine singular}{κῆπος}{garden} καὶ ἐν τῷ κήπῳ\FnParseFormGloss{d}{dative masculine singular}{κῆπος}{garden} μνημεῖον καινόν, ἐν ᾧ οὐδέπω\FnFormGloss{e}{οὐδέπω}{not yet} οὐδεὶς ⸂ἦν\FnParse{f}{imperfect active indicative 3rd singular} τεθειμένος⸃·\FnParse{g}{perfect passive participle nominative masculine singular} 
\MainTextVerseMark{19}{42}ἐκεῖ οὖν διὰ τὴν παρασκευὴν\FnParseFormGloss{a}{accusative feminine singular}{παρασκευή}{Preparation Day} τῶν Ἰουδαίων, ὅτι ἐγγὺς ἦν\FnParse{b}{imperfect active indicative 3rd singular} τὸ μνημεῖον, ἔθηκαν\FnParse{c}{aorist active indicative 3rd plural} τὸν Ἰησοῦν. 
\ChapterHeader{20}{Chapter 20}

\MainTextVerseMark{20}{1}Τῇ δὲ μιᾷ τῶν σαββάτων Μαρία\FnParseFormGloss{a}{nominative feminine singular}{Μαρία}{Mary} ἡ Μαγδαληνὴ\FnParseFormGloss{b}{nominative feminine singular}{Μαγδαληνή}{Magdalene} ἔρχεται\FnParse{c}{present middle indicative 3rd singular} πρωῒ\FnFormGloss{d}{πρωΐ}{early in the morning} σκοτίας\FnParseFormGloss{e}{genitive feminine singular}{σκοτία}{darkness} ἔτι οὔσης\FnParse{f}{present active participle genitive feminine singular} εἰς τὸ μνημεῖον, καὶ βλέπει\FnParse{g}{present active indicative 3rd singular} τὸν λίθον ἠρμένον\FnParse{h}{perfect passive participle accusative masculine singular} ἐκ τοῦ μνημείου. 
\MainTextVerseMark{20}{2}τρέχει\FnParseFormGloss{a}{present active indicative 3rd singular}{τρέχω}{to run} οὖν καὶ ἔρχεται\FnParse{b}{present middle indicative 3rd singular} πρὸς Σίμωνα Πέτρον καὶ πρὸς τὸν ἄλλον μαθητὴν ὃν ἐφίλει\FnParseFormGloss{c}{imperfect active indicative 3rd singular}{φιλέω}{to love} ὁ Ἰησοῦς, καὶ λέγει\FnParse{d}{present active indicative 3rd singular} αὐτοῖς· Ἦραν\FnParse{e}{aorist active indicative 3rd plural} τὸν κύριον ἐκ τοῦ μνημείου, καὶ οὐκ οἴδαμεν\FnParse{f}{perfect active indicative 1st plural} ποῦ ἔθηκαν\FnParse{g}{aorist active indicative 3rd plural} αὐτόν. 
\MainTextVerseMark{20}{3}ἐξῆλθεν\FnParse{a}{aorist active indicative 3rd singular} οὖν ὁ Πέτρος καὶ ὁ ἄλλος μαθητής, καὶ ἤρχοντο\FnParse{b}{imperfect middle indicative 3rd plural} εἰς τὸ μνημεῖον. 
\MainTextVerseMark{20}{4}ἔτρεχον\FnParseFormGloss{a}{imperfect active indicative 3rd plural}{τρέχω}{to run} δὲ οἱ δύο ὁμοῦ·\FnFormGloss{b}{ὁμοῦ}{together} καὶ ὁ ἄλλος μαθητὴς προέδραμεν\FnParseFormGloss{c}{aorist active indicative 3rd singular}{προτρέχω}{to run ahead} τάχιον\FnParseFormGloss{d}{accusative neuter singular}{τάχιον}{quickly} τοῦ Πέτρου καὶ ἦλθεν\FnParse{e}{aorist active indicative 3rd singular} πρῶτος εἰς τὸ μνημεῖον, 
\MainTextVerseMark{20}{5}καὶ παρακύψας\FnParseFormGloss{a}{aorist active participle nominative masculine singular}{παρακύπτω}{to bend over} βλέπει\FnParse{b}{present active indicative 3rd singular} κείμενα\FnParseFormGloss{c}{present middle participle accusative neuter plural}{κεῖμαι}{to lay} τὰ ὀθόνια,\FnParseFormGloss{d}{accusative neuter plural}{ὀθόνιον}{strips of linen} οὐ μέντοι\FnFormGloss{e}{μέντοι}{but} εἰσῆλθεν.\FnParse{f}{aorist active indicative 3rd singular} 
\MainTextVerseMark{20}{6}ἔρχεται\FnParse{a}{present middle indicative 3rd singular} οὖν ⸀καὶ Σίμων Πέτρος ἀκολουθῶν\FnParse{b}{present active participle nominative masculine singular} αὐτῷ, καὶ εἰσῆλθεν\FnParse{c}{aorist active indicative 3rd singular} εἰς τὸ μνημεῖον· καὶ θεωρεῖ\FnParse{d}{present active indicative 3rd singular} τὰ ὀθόνια\FnParseFormGloss{e}{accusative neuter plural}{ὀθόνιον}{strips of linen} κείμενα,\FnParseFormGloss{f}{present middle participle accusative neuter plural}{κεῖμαι}{to lay} 
\MainTextVerseMark{20}{7}καὶ τὸ σουδάριον,\FnParseFormGloss{a}{accusative neuter singular}{σουδάριον}{piece of cloth} ὃ ἦν\FnParse{b}{imperfect active indicative 3rd singular} ἐπὶ τῆς κεφαλῆς αὐτοῦ, οὐ μετὰ τῶν ὀθονίων\FnParseFormGloss{c}{genitive neuter plural}{ὀθόνιον}{strips of linen} κείμενον\FnParseFormGloss{d}{present middle participle accusative neuter singular}{κεῖμαι}{to lay} ἀλλὰ χωρὶς ἐντετυλιγμένον\FnParseFormGloss{e}{perfect passive participle accusative neuter singular}{ἐντυλίσσω}{to wrap up} εἰς ἕνα τόπον· 
\MainTextVerseMark{20}{8}τότε οὖν εἰσῆλθεν\FnParse{a}{aorist active indicative 3rd singular} καὶ ὁ ἄλλος μαθητὴς ὁ ἐλθὼν\FnParse{b}{aorist active participle nominative masculine singular} πρῶτος εἰς τὸ μνημεῖον, καὶ εἶδεν\FnParse{c}{aorist active indicative 3rd singular} καὶ ἐπίστευσεν·\FnParse{d}{aorist active indicative 3rd singular} 
\MainTextVerseMark{20}{9}οὐδέπω\FnFormGloss{a}{οὐδέπω}{not yet} γὰρ ᾔδεισαν\FnParse{b}{pluperfect active indicative 3rd plural} τὴν γραφὴν ὅτι δεῖ\FnParse{c}{present active indicative 3rd singular} αὐτὸν ἐκ νεκρῶν ἀναστῆναι.\FnParse{d}{aorist active infinitive} 
\MainTextVerseMark{20}{10}ἀπῆλθον\FnParse{a}{aorist active indicative 3rd plural} οὖν πάλιν πρὸς ⸀αὑτοὺς οἱ μαθηταί. 
\MainTextVerseMark{20}{11}Μαρία\FnParseFormGloss{a}{nominative feminine singular}{Μαρία}{Mary} δὲ εἱστήκει\FnParse{b}{pluperfect active indicative 3rd singular} πρὸς ⸂τῷ μνημείῳ ἔξω κλαίουσα⸃.\FnParse{c}{present active participle nominative feminine singular} ὡς οὖν ἔκλαιεν\FnParse{d}{imperfect active indicative 3rd singular} παρέκυψεν\FnParseFormGloss{e}{aorist active indicative 3rd singular}{παρακύπτω}{to bend over} εἰς τὸ μνημεῖον, 
\MainTextVerseMark{20}{12}καὶ θεωρεῖ\FnParse{a}{present active indicative 3rd singular} δύο ἀγγέλους ἐν λευκοῖς\FnParseFormGloss{b}{dative neuter plural}{λευκός}{white} καθεζομένους,\FnParseFormGloss{c}{present middle participle accusative masculine plural}{καθέζομαι}{to sit down} ἕνα πρὸς τῇ κεφαλῇ καὶ ἕνα πρὸς τοῖς ποσίν, ὅπου ἔκειτο\FnParseFormGloss{d}{imperfect middle indicative 3rd singular}{κεῖμαι}{to lay} τὸ σῶμα τοῦ Ἰησοῦ. 
\MainTextVerseMark{20}{13}καὶ λέγουσιν\FnParse{a}{present active indicative 3rd plural} αὐτῇ ἐκεῖνοι· Γύναι, τί κλαίεις;\FnParse{b}{present active indicative 2nd singular} λέγει\FnParse{c}{present active indicative 3rd singular} αὐτοῖς ὅτι Ἦραν\FnParse{d}{aorist active indicative 3rd plural} τὸν κύριόν μου, καὶ οὐκ οἶδα\FnParse{e}{perfect active indicative 1st singular} ποῦ ἔθηκαν\FnParse{f}{aorist active indicative 3rd plural} αὐτόν. 
\MainTextVerseMark{20}{14}⸀ταῦτα εἰποῦσα\FnParse{a}{aorist active participle nominative feminine singular} ἐστράφη\FnParseFormGloss{b}{aorist passive indicative 3rd singular}{στρέφω}{to turn} εἰς τὰ ὀπίσω, καὶ θεωρεῖ\FnParse{c}{present active indicative 3rd singular} τὸν Ἰησοῦν ἑστῶτα,\FnParse{d}{perfect active participle accusative masculine singular} καὶ οὐκ ᾔδει\FnParse{e}{pluperfect active indicative 3rd singular} ὅτι Ἰησοῦς ἐστιν.\FnParse{f}{present active indicative 3rd singular} 
\MainTextVerseMark{20}{15}λέγει\FnParse{a}{present active indicative 3rd singular} ⸀αὐτῇ Ἰησοῦς· Γύναι, τί κλαίεις;\FnParse{b}{present active indicative 2nd singular} τίνα ζητεῖς;\FnParse{c}{present active indicative 2nd singular} ἐκείνη δοκοῦσα\FnParse{d}{present active participle nominative feminine singular} ὅτι ὁ κηπουρός\FnParseFormGloss{e}{nominative masculine singular}{κηπουρός}{gardener} ἐστιν\FnParse{f}{present active indicative 3rd singular} λέγει\FnParse{g}{present active indicative 3rd singular} αὐτῷ· Κύριε, εἰ σὺ ἐβάστασας\FnParseFormGloss{h}{aorist active indicative 2nd singular}{βαστάζω}{to carry} αὐτόν, εἰπέ\FnParse{i}{aorist active imperative 2nd singular} μοι ποῦ ἔθηκας\FnParse{j}{aorist active indicative 2nd singular} αὐτόν, κἀγὼ αὐτὸν ἀρῶ.\FnParse{k}{future active indicative 1st singular} 
\MainTextVerseMark{20}{16}λέγει\FnParse{a}{present active indicative 3rd singular} ⸀αὐτῇ Ἰησοῦς· ⸀Μαριάμ.\FnParseFormGloss{b}{feminine singular}{Μαριάμ}{Mary} στραφεῖσα\FnParseFormGloss{c}{aorist passive participle nominative feminine singular}{στρέφω}{to turn} ἐκείνη λέγει\FnParse{d}{present active indicative 3rd singular} αὐτῷ ⸀Ἑβραϊστί·\FnFormGloss{e}{Ἑβραϊστί}{in Aramaic} Ραββουνι\FnParseFormGloss{f}{masculine singular}{ραββουνι}{Rabboni} (ὃ λέγεται\FnParse{g}{present passive indicative 3rd singular} Διδάσκαλε). 
\MainTextVerseMark{20}{17}λέγει\FnParse{a}{present active indicative 3rd singular} ⸀αὐτῇ Ἰησοῦς· Μή μου ἅπτου,\FnParse{b}{present middle imperative 2nd singular} οὔπω\FnFormGloss{c}{οὔπω}{not yet} γὰρ ἀναβέβηκα\FnParse{d}{perfect active indicative 1st singular} πρὸς τὸν ⸀πατέρα· πορεύου\FnParse{e}{present middle imperative 2nd singular} δὲ πρὸς τοὺς ἀδελφούς μου καὶ εἰπὲ\FnParse{f}{aorist active imperative 2nd singular} αὐτοῖς· Ἀναβαίνω\FnParse{g}{present active indicative 1st singular} πρὸς τὸν πατέρα μου καὶ πατέρα ὑμῶν καὶ θεόν μου καὶ θεὸν ὑμῶν. 
\MainTextVerseMark{20}{18}ἔρχεται\FnParse{a}{present middle indicative 3rd singular} ⸀Μαριὰμ\FnParseFormGloss{b}{nominative feminine singular}{Μαριάμ}{Mary} ἡ Μαγδαληνὴ\FnParseFormGloss{c}{nominative feminine singular}{Μαγδαληνή}{Magdalene} ⸀ἀγγέλλουσα\FnParseFormGloss{d}{present active participle nominative feminine singular}{ἀγγέλλω}{to bring news} τοῖς μαθηταῖς ὅτι ⸀Ἑώρακα\FnParse{e}{perfect active indicative 1st singular} τὸν κύριον καὶ ταῦτα εἶπεν\FnParse{f}{aorist active indicative 3rd singular} αὐτῇ. 
\MainTextVerseMark{20}{19}Οὔσης\FnParse{a}{present active participle genitive feminine singular} οὖν ὀψίας\FnParseFormGloss{b}{genitive feminine singular}{ὀψία}{evening} τῇ ἡμέρᾳ ἐκείνῃ τῇ ⸀μιᾷ σαββάτων, καὶ τῶν θυρῶν κεκλεισμένων\FnParseFormGloss{c}{perfect passive participle genitive feminine plural}{κλείω}{to close} ὅπου ἦσαν\FnParse{d}{imperfect active indicative 3rd plural} οἱ ⸀μαθηταὶ διὰ τὸν φόβον τῶν Ἰουδαίων, ἦλθεν\FnParse{e}{aorist active indicative 3rd singular} ὁ Ἰησοῦς καὶ ἔστη\FnParse{f}{aorist active indicative 3rd singular} εἰς τὸ μέσον, καὶ λέγει\FnParse{g}{present active indicative 3rd singular} αὐτοῖς· Εἰρήνη ὑμῖν. 
\MainTextVerseMark{20}{20}καὶ τοῦτο εἰπὼν\FnParse{a}{aorist active participle nominative masculine singular} ⸀ἔδειξεν\FnParse{b}{aorist active indicative 3rd singular} τὰς χεῖρας καὶ τὴν πλευρὰν\FnParseFormGloss{c}{accusative feminine singular}{πλευρά}{side} ⸀αὐτοῖς. ἐχάρησαν\FnParse{d}{aorist passive indicative 3rd plural} οὖν οἱ μαθηταὶ ἰδόντες\FnParse{e}{aorist active participle nominative masculine plural} τὸν κύριον. 
\MainTextVerseMark{20}{21}εἶπεν\FnParse{a}{aorist active indicative 3rd singular} οὖν αὐτοῖς ⸂ὁ Ἰησοῦς⸃ πάλιν· Εἰρήνη ὑμῖν· καθὼς ἀπέσταλκέν\FnParse{b}{perfect active indicative 3rd singular} με ὁ πατήρ, κἀγὼ πέμπω\FnParse{c}{present active indicative 1st singular} ὑμᾶς. 
\MainTextVerseMark{20}{22}καὶ τοῦτο εἰπὼν\FnParse{a}{aorist active participle nominative masculine singular} ἐνεφύσησεν\FnParseFormGloss{b}{aorist active indicative 3rd singular}{ἐμφυσάω}{to breathe on} καὶ λέγει\FnParse{c}{present active indicative 3rd singular} αὐτοῖς· Λάβετε\FnParse{d}{aorist active imperative 2nd plural} πνεῦμα ἅγιον· 
\MainTextVerseMark{20}{23}ἄν τινων ἀφῆτε\FnParse{a}{aorist active subjunctive 2nd plural} τὰς ἁμαρτίας ⸀ἀφέωνται\FnParse{b}{perfect passive indicative 3rd plural} αὐτοῖς· ἄν τινων κρατῆτε\FnParse{c}{present active subjunctive 2nd plural} κεκράτηνται.\FnParse{d}{perfect passive indicative 3rd plural} 
\MainTextVerseMark{20}{24}Θωμᾶς\FnParseFormGloss{a}{nominative masculine singular}{Θωμᾶς}{Thomas} δὲ εἷς ἐκ τῶν δώδεκα, ὁ λεγόμενος\FnParse{b}{present passive participle nominative masculine singular} Δίδυμος,\FnParseFormGloss{c}{nominative masculine singular}{Δίδυμος}{Didymus} οὐκ ἦν\FnParse{d}{imperfect active indicative 3rd singular} μετ’ αὐτῶν ὅτε ⸀ἦλθεν\FnParse{e}{aorist active indicative 3rd singular} Ἰησοῦς. 
\MainTextVerseMark{20}{25}ἔλεγον\FnParse{a}{imperfect active indicative 3rd plural} οὖν αὐτῷ οἱ ἄλλοι μαθηταί· Ἑωράκαμεν\FnParse{b}{perfect active indicative 1st plural} τὸν κύριον. ὁ δὲ εἶπεν\FnParse{c}{aorist active indicative 3rd singular} αὐτοῖς· Ἐὰν μὴ ἴδω\FnParse{d}{aorist active subjunctive 1st singular} ἐν ταῖς χερσὶν αὐτοῦ τὸν τύπον\FnParseFormGloss{e}{accusative masculine singular}{τύπος}{pattern} τῶν ἥλων\FnParseFormGloss{f}{genitive masculine plural}{ἧλος}{nail} καὶ βάλω\FnParse{g}{aorist active subjunctive 1st singular} τὸν δάκτυλόν\FnParseFormGloss{h}{accusative masculine singular}{δάκτυλος}{finger} μου εἰς τὸν τύπον\FnParseFormGloss{i}{accusative masculine singular}{τύπος}{pattern} τῶν ἥλων\FnParseFormGloss{j}{genitive masculine plural}{ἧλος}{nail} καὶ βάλω\FnParse{k}{aorist active subjunctive 1st singular} ⸂μου τὴν χεῖρα⸃ εἰς τὴν πλευρὰν\FnParseFormGloss{l}{accusative feminine singular}{πλευρά}{side} αὐτοῦ, οὐ μὴ πιστεύσω.\FnParse{m}{aorist active subjunctive 1st singular} 
\MainTextVerseMark{20}{26}Καὶ μεθ’ ἡμέρας ὀκτὼ\FnParseFormGloss{a}{accusative feminine plural}{ὀκτώ}{eight} πάλιν ἦσαν\FnParse{b}{imperfect active indicative 3rd plural} ἔσω\FnFormGloss{c}{ἔσω}{in} οἱ μαθηταὶ αὐτοῦ καὶ Θωμᾶς\FnParseFormGloss{d}{nominative masculine singular}{Θωμᾶς}{Thomas} μετ’ αὐτῶν. ἔρχεται\FnParse{e}{present middle indicative 3rd singular} ὁ Ἰησοῦς τῶν θυρῶν κεκλεισμένων,\FnParseFormGloss{f}{perfect passive participle genitive feminine plural}{κλείω}{to close} καὶ ἔστη\FnParse{g}{aorist active indicative 3rd singular} εἰς τὸ μέσον καὶ εἶπεν·\FnParse{h}{aorist active indicative 3rd singular} Εἰρήνη ὑμῖν. 
\MainTextVerseMark{20}{27}εἶτα\FnFormGloss{a}{εἶτα}{then} λέγει\FnParse{b}{present active indicative 3rd singular} τῷ Θωμᾷ·\FnParseFormGloss{c}{dative masculine singular}{Θωμᾶς}{Thomas} Φέρε\FnParse{d}{present active imperative 2nd singular} τὸν δάκτυλόν\FnParseFormGloss{e}{accusative masculine singular}{δάκτυλος}{finger} σου ὧδε καὶ ἴδε\FnParse{f}{aorist active imperative 2nd singular} τὰς χεῖράς μου, καὶ φέρε\FnParse{g}{present active imperative 2nd singular} τὴν χεῖρά σου καὶ βάλε\FnParse{h}{aorist active imperative 2nd singular} εἰς τὴν πλευράν\FnParseFormGloss{i}{accusative feminine singular}{πλευρά}{side} μου, καὶ μὴ γίνου\FnParse{j}{present middle imperative 2nd singular} ἄπιστος\FnParseFormGloss{k}{nominative masculine singular}{ἄπιστος}{unbelieving} ἀλλὰ πιστός. 
\MainTextVerseMark{20}{28}⸀ἀπεκρίθη\FnParse{a}{aorist passive indicative 3rd singular} Θωμᾶς\FnParseFormGloss{b}{nominative masculine singular}{Θωμᾶς}{Thomas} καὶ εἶπεν\FnParse{c}{aorist active indicative 3rd singular} αὐτῷ· Ὁ κύριός μου καὶ ὁ θεός μου. 
\MainTextVerseMark{20}{29}λέγει\FnParse{a}{present active indicative 3rd singular} αὐτῷ ὁ Ἰησοῦς· Ὅτι ἑώρακάς\FnParse{b}{perfect active indicative 2nd singular} με πεπίστευκας;\FnParse{c}{perfect active indicative 2nd singular} μακάριοι οἱ μὴ ἰδόντες\FnParse{d}{aorist active participle nominative masculine plural} καὶ πιστεύσαντες.\FnParse{e}{aorist active participle nominative masculine plural} 
\MainTextVerseMark{20}{30}Πολλὰ μὲν οὖν καὶ ἄλλα σημεῖα ἐποίησεν\FnParse{a}{aorist active indicative 3rd singular} ὁ Ἰησοῦς ἐνώπιον τῶν ⸀μαθητῶν, ἃ οὐκ ἔστιν\FnParse{b}{present active indicative 3rd singular} γεγραμμένα\FnParse{c}{perfect passive participle nominative neuter plural} ἐν τῷ βιβλίῳ τούτῳ· 
\MainTextVerseMark{20}{31}ταῦτα δὲ γέγραπται\FnParse{a}{perfect passive indicative 3rd singular} ἵνα ⸀πιστεύητε\FnParse{b}{present active subjunctive 2nd plural} ὅτι Ἰησοῦς ἐστιν\FnParse{c}{present active indicative 3rd singular} ὁ χριστὸς ὁ υἱὸς τοῦ θεοῦ, καὶ ἵνα πιστεύοντες\FnParse{d}{present active participle nominative masculine plural} ζωὴν ἔχητε\FnParse{e}{present active subjunctive 2nd plural} ἐν τῷ ὀνόματι αὐτοῦ. 
\ChapterHeader{21}{Chapter 21}

\MainTextVerseMark{21}{1}Μετὰ ταῦτα ἐφανέρωσεν\FnParse{a}{aorist active indicative 3rd singular} ἑαυτὸν πάλιν ⸀ὁ Ἰησοῦς τοῖς μαθηταῖς ἐπὶ τῆς θαλάσσης τῆς Τιβεριάδος·\FnParseFormGloss{b}{genitive feminine singular}{Τιβεριάς}{Tiberias} ἐφανέρωσεν\FnParse{c}{aorist active indicative 3rd singular} δὲ οὕτως. 
\MainTextVerseMark{21}{2}ἦσαν\FnParse{a}{imperfect active indicative 3rd plural} ὁμοῦ\FnFormGloss{b}{ὁμοῦ}{together} Σίμων Πέτρος καὶ Θωμᾶς\FnParseFormGloss{c}{nominative masculine singular}{Θωμᾶς}{Thomas} ὁ λεγόμενος\FnParse{d}{present passive participle nominative masculine singular} Δίδυμος\FnParseFormGloss{e}{nominative masculine singular}{Δίδυμος}{Didymus} καὶ Ναθαναὴλ\FnParseFormGloss{f}{nominative masculine singular}{Ναθαναήλ}{Nathanael} ὁ ἀπὸ Κανὰ\FnParseFormGloss{g}{genitive feminine singular}{Κανά}{Cana} τῆς Γαλιλαίας καὶ οἱ τοῦ Ζεβεδαίου\FnParseFormGloss{h}{genitive masculine singular}{Ζεβεδαῖος}{Zebedee} καὶ ἄλλοι ἐκ τῶν μαθητῶν αὐτοῦ δύο. 
\MainTextVerseMark{21}{3}λέγει\FnParse{a}{present active indicative 3rd singular} αὐτοῖς Σίμων Πέτρος· Ὑπάγω\FnParse{b}{present active indicative 1st singular} ἁλιεύειν·\FnParseFormGloss{c}{present active infinitive}{ἁλιεύω}{to catch fish} λέγουσιν\FnParse{d}{present active indicative 3rd plural} αὐτῷ· Ἐρχόμεθα\FnParse{e}{present middle indicative 1st plural} καὶ ἡμεῖς σὺν σοί. ἐξῆλθον\FnParse{f}{aorist active indicative 3rd plural} καὶ ἐνέβησαν\FnParseFormGloss{g}{aorist active indicative 3rd plural}{ἐμβαίνω}{to get into} εἰς τὸ ⸀πλοῖον, καὶ ἐν ἐκείνῃ τῇ νυκτὶ ἐπίασαν\FnParseFormGloss{h}{aorist active indicative 3rd plural}{πιάζω}{to seize} οὐδέν. 
\MainTextVerseMark{21}{4}Πρωΐας\FnParseFormGloss{a}{genitive feminine singular}{πρωΐα}{early morning} δὲ ἤδη ⸀γενομένης\FnParse{b}{aorist middle participle genitive feminine singular} ⸀ἔστη\FnParse{c}{aorist active indicative 3rd singular} Ἰησοῦς εἰς τὸν αἰγιαλόν·\FnParseFormGloss{d}{accusative masculine singular}{αἰγιαλός}{shore} οὐ μέντοι\FnFormGloss{e}{μέντοι}{but} ᾔδεισαν\FnParse{f}{pluperfect active indicative 3rd plural} οἱ μαθηταὶ ὅτι Ἰησοῦς ἐστιν.\FnParse{g}{present active indicative 3rd singular} 
\MainTextVerseMark{21}{5}λέγει\FnParse{a}{present active indicative 3rd singular} οὖν αὐτοῖς ⸀ὁ Ἰησοῦς· Παιδία, μή τι προσφάγιον\FnParseFormGloss{b}{accusative neuter singular}{προσφάγιον}{fish} ἔχετε;\FnParse{c}{present active indicative 2nd plural} ἀπεκρίθησαν\FnParse{d}{aorist passive indicative 3rd plural} αὐτῷ· Οὔ.\FnFormGloss{e}{οὔ}{no} 
\MainTextVerseMark{21}{6}ὁ δὲ εἶπεν\FnParse{a}{aorist active indicative 3rd singular} αὐτοῖς· Βάλετε\FnParse{b}{aorist active imperative 2nd plural} εἰς τὰ δεξιὰ μέρη τοῦ πλοίου τὸ δίκτυον,\FnParseFormGloss{c}{accusative neuter singular}{δίκτυον}{net} καὶ εὑρήσετε.\FnParse{d}{future active indicative 2nd plural} ἔβαλον\FnParse{e}{aorist active indicative 3rd plural} οὖν, καὶ οὐκέτι αὐτὸ ἑλκύσαι\FnParseFormGloss{f}{aorist active infinitive}{ἕλκω}{to drag} ⸀ἴσχυον\FnParseFormGloss{g}{imperfect active indicative 3rd plural}{ἰσχύω}{to be strong} ἀπὸ τοῦ πλήθους τῶν ἰχθύων.\FnParseFormGloss{h}{genitive masculine plural}{ἰχθύς}{fish} 
\MainTextVerseMark{21}{7}λέγει\FnParse{a}{present active indicative 3rd singular} οὖν ὁ μαθητὴς ἐκεῖνος ὃν ἠγάπα\FnParse{b}{imperfect active indicative 3rd singular} ὁ Ἰησοῦς τῷ Πέτρῳ· Ὁ κύριός ἐστιν.\FnParse{c}{present active indicative 3rd singular} Σίμων οὖν Πέτρος, ἀκούσας\FnParse{d}{aorist active participle nominative masculine singular} ὅτι ὁ κύριός ἐστιν,\FnParse{e}{present active indicative 3rd singular} τὸν ἐπενδύτην\FnParseFormGloss{f}{accusative masculine singular}{ἐπενδύτης}{outer garment} διεζώσατο,\FnParseFormGloss{g}{aorist middle indicative 3rd singular}{διαζώννυμι}{to wrap around} ἦν\FnParse{h}{imperfect active indicative 3rd singular} γὰρ γυμνός,\FnParseFormGloss{i}{nominative masculine singular}{γυμνός}{naked} καὶ ἔβαλεν\FnParse{j}{aorist active indicative 3rd singular} ἑαυτὸν εἰς τὴν θάλασσαν· 
\MainTextVerseMark{21}{8}οἱ δὲ ἄλλοι μαθηταὶ τῷ πλοιαρίῳ\FnParseFormGloss{a}{dative neuter singular}{πλοιάριον}{boat} ἦλθον,\FnParse{b}{aorist active indicative 3rd plural} οὐ γὰρ ἦσαν\FnParse{c}{imperfect active indicative 3rd plural} μακρὰν\FnParseFormGloss{d}{accusative feminine singular}{μακρός}{lengthy} ἀπὸ τῆς γῆς ἀλλὰ ὡς ἀπὸ πηχῶν\FnParseFormGloss{e}{genitive masculine plural}{πῆχυς}{cubit} διακοσίων,\FnParseFormGloss{f}{genitive masculine plural}{διακόσιοι}{two hundred} σύροντες\FnParseFormGloss{g}{present active participle nominative masculine plural}{σύρω}{to drag} τὸ δίκτυον\FnParseFormGloss{h}{accusative neuter singular}{δίκτυον}{net} τῶν ἰχθύων.\FnParseFormGloss{i}{genitive masculine plural}{ἰχθύς}{fish} 
\MainTextVerseMark{21}{9}Ὡς οὖν ἀπέβησαν\FnParseFormGloss{a}{aorist active indicative 3rd plural}{ἀποβαίνω}{to leave} εἰς τὴν γῆν βλέπουσιν\FnParse{b}{present active indicative 3rd plural} ἀνθρακιὰν\FnParseFormGloss{c}{accusative feminine singular}{ἀνθρακιά}{charcoal fire} κειμένην\FnParseFormGloss{d}{present middle participle accusative feminine singular}{κεῖμαι}{to lay} καὶ ὀψάριον\FnParseFormGloss{e}{accusative neuter singular}{ὀψάριον}{fish} ἐπικείμενον\FnParseFormGloss{f}{present middle participle accusative neuter singular}{ἐπίκειμαι}{to lay upon} καὶ ἄρτον. 
\MainTextVerseMark{21}{10}λέγει\FnParse{a}{present active indicative 3rd singular} αὐτοῖς ὁ Ἰησοῦς· Ἐνέγκατε\FnParse{b}{aorist active imperative 2nd plural} ἀπὸ τῶν ὀψαρίων\FnParseFormGloss{c}{genitive neuter plural}{ὀψάριον}{fish} ὧν ἐπιάσατε\FnParseFormGloss{d}{aorist active indicative 2nd plural}{πιάζω}{to seize} νῦν. 
\MainTextVerseMark{21}{11}ἀνέβη\FnParse{a}{aorist active indicative 3rd singular} ⸀οὖν Σίμων Πέτρος καὶ εἵλκυσεν\FnParseFormGloss{b}{aorist active indicative 3rd singular}{ἕλκω}{to drag} τὸ δίκτυον\FnParseFormGloss{c}{accusative neuter singular}{δίκτυον}{net} ⸂εἰς τὴν γῆν⸃ μεστὸν\FnParseFormGloss{d}{accusative neuter singular}{μεστός}{full} ἰχθύων\FnParseFormGloss{e}{genitive masculine plural}{ἰχθύς}{fish} μεγάλων ἑκατὸν\FnParseFormGloss{f}{genitive masculine plural}{ἑκατόν}{hundred} πεντήκοντα\FnParseFormGloss{g}{genitive masculine plural}{πεντήκοντα}{fifty} τριῶν· καὶ τοσούτων\FnParseFormGloss{h}{genitive masculine plural}{τοσοῦτος}{so great} ὄντων\FnParse{i}{present active participle genitive masculine plural} οὐκ ἐσχίσθη\FnParseFormGloss{j}{aorist passive indicative 3rd singular}{σχίζω}{to tear} τὸ δίκτυον.\FnParseFormGloss{k}{nominative neuter singular}{δίκτυον}{net} 
\MainTextVerseMark{21}{12}λέγει\FnParse{a}{present active indicative 3rd singular} αὐτοῖς ὁ Ἰησοῦς· Δεῦτε\FnFormGloss{b}{δεῦτε}{come} ἀριστήσατε.\FnParseFormGloss{c}{aorist active imperative 2nd plural}{ἀριστάω}{to eat} οὐδεὶς ⸀δὲ ἐτόλμα\FnParseFormGloss{d}{imperfect active indicative 3rd singular}{τολμάω}{to dare} τῶν μαθητῶν ἐξετάσαι\FnParseFormGloss{e}{aorist active infinitive}{ἐξετάζω}{to make a search} αὐτόν· Σὺ τίς εἶ;\FnParse{f}{present active indicative 2nd singular} εἰδότες\FnParse{g}{perfect active participle nominative masculine plural} ὅτι ὁ κύριός ἐστιν.\FnParse{h}{present active indicative 3rd singular} 
\MainTextVerseMark{21}{13}⸀ἔρχεται\FnParse{a}{present middle indicative 3rd singular} ⸀ὁ Ἰησοῦς καὶ λαμβάνει\FnParse{b}{present active indicative 3rd singular} τὸν ἄρτον καὶ δίδωσιν\FnParse{c}{present active indicative 3rd singular} αὐτοῖς, καὶ τὸ ὀψάριον\FnParseFormGloss{d}{accusative neuter singular}{ὀψάριον}{fish} ὁμοίως. 
\MainTextVerseMark{21}{14}τοῦτο ἤδη τρίτον ἐφανερώθη\FnParse{a}{aorist passive indicative 3rd singular} ⸀ὁ Ἰησοῦς τοῖς ⸀μαθηταῖς ἐγερθεὶς\FnParse{b}{aorist passive participle nominative masculine singular} ἐκ νεκρῶν. 
\MainTextVerseMark{21}{15}Ὅτε οὖν ἠρίστησαν\FnParseFormGloss{a}{aorist active indicative 3rd plural}{ἀριστάω}{to eat} λέγει\FnParse{b}{present active indicative 3rd singular} τῷ Σίμωνι Πέτρῳ ὁ Ἰησοῦς· Σίμων ⸀Ἰωάννου, ἀγαπᾷς\FnParse{c}{present active indicative 2nd singular} με πλέον τούτων; λέγει\FnParse{d}{present active indicative 3rd singular} αὐτῷ· Ναί, κύριε, σὺ οἶδας\FnParse{e}{perfect active indicative 2nd singular} ὅτι φιλῶ\FnParseFormGloss{f}{present active indicative 1st singular}{φιλέω}{to love} σε. λέγει\FnParse{g}{present active indicative 3rd singular} αὐτῷ· Βόσκε\FnParseFormGloss{h}{present active imperative 2nd singular}{βόσκω}{to feed} τὰ ἀρνία μου. 
\MainTextVerseMark{21}{16}λέγει\FnParse{a}{present active indicative 3rd singular} αὐτῷ πάλιν δεύτερον· Σίμων ⸀Ἰωάννου, ἀγαπᾷς\FnParse{b}{present active indicative 2nd singular} με; λέγει\FnParse{c}{present active indicative 3rd singular} αὐτῷ· Ναί, κύριε, σὺ οἶδας\FnParse{d}{perfect active indicative 2nd singular} ὅτι φιλῶ\FnParseFormGloss{e}{present active indicative 1st singular}{φιλέω}{to love} σε. λέγει\FnParse{f}{present active indicative 3rd singular} αὐτῷ· Ποίμαινε\FnParseFormGloss{g}{present active imperative 2nd singular}{ποιμαίνω}{to shepherd} τὰ πρόβατά μου. 
\MainTextVerseMark{21}{17}λέγει\FnParse{a}{present active indicative 3rd singular} αὐτῷ τὸ τρίτον· Σίμων ⸀Ἰωάννου, φιλεῖς\FnParseFormGloss{b}{present active indicative 2nd singular}{φιλέω}{to love} με; ἐλυπήθη\FnParseFormGloss{c}{aorist passive indicative 3rd singular}{λυπέω}{to cause sorrow} ὁ Πέτρος ὅτι εἶπεν\FnParse{d}{aorist active indicative 3rd singular} αὐτῷ τὸ τρίτον· Φιλεῖς\FnParseFormGloss{e}{present active indicative 2nd singular}{φιλέω}{to love} με; καὶ ⸀εἶπεν\FnParse{f}{aorist active indicative 3rd singular} αὐτῷ· Κύριε, ⸂πάντα σὺ⸃ οἶδας,\FnParse{g}{perfect active indicative 2nd singular} σὺ γινώσκεις\FnParse{h}{present active indicative 2nd singular} ὅτι φιλῶ\FnParseFormGloss{i}{present active indicative 1st singular}{φιλέω}{to love} σε. λέγει\FnParse{j}{present active indicative 3rd singular} αὐτῷ ⸀ὁ Ἰησοῦς· Βόσκε\FnParseFormGloss{k}{present active imperative 2nd singular}{βόσκω}{to feed} τὰ ⸀πρόβατά μου. 
\MainTextVerseMark{21}{18}ἀμὴν ἀμὴν λέγω\FnParse{a}{present active indicative 1st singular} σοι, ὅτε ἦς\FnParse{b}{imperfect active indicative 2nd singular} νεώτερος,\FnParseFormGloss{c}{nominative masculine singular}{νέος}{new} ἐζώννυες\FnParseFormGloss{d}{imperfect active indicative 2nd singular}{ζώννυμι}{to dress} σεαυτὸν καὶ περιεπάτεις\FnParse{e}{imperfect active indicative 2nd singular} ὅπου ἤθελες·\FnParse{f}{imperfect active indicative 2nd singular} ὅταν δὲ γηράσῃς,\FnParseFormGloss{g}{aorist active subjunctive 2nd singular}{γηράσκω}{to grow old} ἐκτενεῖς\FnParseFormGloss{h}{future active indicative 2nd singular}{ἐκτείνω}{to stretch out} τὰς χεῖράς σου, καὶ ἄλλος ⸂σε ζώσει⸃\FnParseFormGloss{i}{future active indicative 3rd singular}{ζώννυμι}{to dress} καὶ οἴσει\FnParse{j}{future active indicative 3rd singular} ὅπου οὐ θέλεις.\FnParse{k}{present active indicative 2nd singular} 
\MainTextVerseMark{21}{19}τοῦτο δὲ εἶπεν\FnParse{a}{aorist active indicative 3rd singular} σημαίνων\FnParseFormGloss{b}{present active participle nominative masculine singular}{σημαίνω}{to make known} ποίῳ θανάτῳ δοξάσει\FnParse{c}{future active indicative 3rd singular} τὸν θεόν. καὶ τοῦτο εἰπὼν\FnParse{d}{aorist active participle nominative masculine singular} λέγει\FnParse{e}{present active indicative 3rd singular} αὐτῷ· Ἀκολούθει\FnParse{f}{present active imperative 2nd singular} μοι. 
\MainTextVerseMark{21}{20}⸀Ἐπιστραφεὶς\FnParse{a}{aorist passive participle nominative masculine singular} ὁ Πέτρος βλέπει\FnParse{b}{present active indicative 3rd singular} τὸν μαθητὴν ὃν ἠγάπα\FnParse{c}{imperfect active indicative 3rd singular} ὁ Ἰησοῦς ἀκολουθοῦντα,\FnParse{d}{present active participle accusative masculine singular} ὃς καὶ ἀνέπεσεν\FnParseFormGloss{e}{aorist active indicative 3rd singular}{ἀναπίπτω}{recline} ἐν τῷ δείπνῳ\FnParseFormGloss{f}{dative neuter singular}{δεῖπνον}{banquet} ἐπὶ τὸ στῆθος\FnParseFormGloss{g}{accusative neuter singular}{στῆθος}{chest} αὐτοῦ καὶ εἶπεν·\FnParse{h}{aorist active indicative 3rd singular} Κύριε, τίς ἐστιν\FnParse{i}{present active indicative 3rd singular} ὁ παραδιδούς\FnParse{j}{present active participle nominative masculine singular} σε; 
\MainTextVerseMark{21}{21}τοῦτον ⸀οὖν ἰδὼν\FnParse{a}{aorist active participle nominative masculine singular} ὁ Πέτρος λέγει\FnParse{b}{present active indicative 3rd singular} τῷ Ἰησοῦ· Κύριε, οὗτος δὲ τί; 
\MainTextVerseMark{21}{22}λέγει\FnParse{a}{present active indicative 3rd singular} αὐτῷ ὁ Ἰησοῦς· Ἐὰν αὐτὸν θέλω\FnParse{b}{present active subjunctive 1st singular} μένειν\FnParse{c}{present active infinitive} ἕως ἔρχομαι,\FnParse{d}{present middle indicative 1st singular} τί πρὸς σέ; σύ ⸂μοι ἀκολούθει⸃.\FnParse{e}{present active imperative 2nd singular} 
\MainTextVerseMark{21}{23}ἐξῆλθεν\FnParse{a}{aorist active indicative 3rd singular} οὖν ⸂οὗτος ὁ λόγος⸃ εἰς τοὺς ἀδελφοὺς ὅτι ὁ μαθητὴς ἐκεῖνος οὐκ ἀποθνῄσκει.\FnParse{b}{present active indicative 3rd singular} ⸂οὐκ εἶπεν\FnParse{c}{aorist active indicative 3rd singular} δὲ⸃ αὐτῷ ὁ Ἰησοῦς ὅτι οὐκ ἀποθνῄσκει\FnParse{d}{present active indicative 3rd singular} ἀλλ’· Ἐὰν αὐτὸν θέλω\FnParse{e}{present active subjunctive 1st singular} μένειν\FnParse{f}{present active infinitive} ἕως ἔρχομαι,\FnParse{g}{present middle indicative 1st singular} τί πρὸς σέ; 
\MainTextVerseMark{21}{24}Οὗτός ἐστιν\FnParse{a}{present active indicative 3rd singular} ὁ μαθητὴς ὁ μαρτυρῶν\FnParse{b}{present active participle nominative masculine singular} περὶ τούτων καὶ ⸀ὁ γράψας\FnParse{c}{aorist active participle nominative masculine singular} ταῦτα, καὶ οἴδαμεν\FnParse{d}{perfect active indicative 1st plural} ὅτι ἀληθὴς\FnParseFormGloss{e}{nominative feminine singular}{ἀληθής}{true} ⸂αὐτοῦ ἡ μαρτυρία ἐστίν⸃.\FnParse{f}{present active indicative 3rd singular} 
\MainTextVerseMark{21}{25}ἔστιν\FnParse{a}{present active indicative 3rd singular} δὲ καὶ ἄλλα πολλὰ ⸀ἃ ἐποίησεν\FnParse{b}{aorist active indicative 3rd singular} ὁ Ἰησοῦς, ἅτινα ἐὰν γράφηται\FnParse{c}{present passive subjunctive 3rd singular} καθ’ ἕν, οὐδ’ αὐτὸν οἶμαι\FnParseFormGloss{d}{present middle indicative 1st singular}{οἶμαι}{to think} τὸν κόσμον ⸀χωρήσειν\FnParseFormGloss{e}{future active infinitive}{χωρέω}{to go} τὰ γραφόμενα\FnParse{f}{present passive participle accusative neuter plural} ⸀βιβλία. 
\end{document}